%% This is a LaTeX document.  Hey, Emacs, -*- latex -*- , get it?
\documentclass[12pt,a4paper]{article}
\usepackage[a4paper,margin=2.5cm]{geometry}
\usepackage[french]{babel}
\usepackage[utf8]{inputenc}
\usepackage[T1]{fontenc}
\usepackage{lmodern}
\usepackage{newtxtext}
\usepackage{amsmath}
\usepackage{amsfonts}
\usepackage{amssymb}
\usepackage{amsthm}
%
\usepackage{mathrsfs}
\usepackage{wasysym}
\usepackage{url}
%
\usepackage{makeidx}
%% Self-note: compile index with:
%% texindy -C utf8 -L french notes-mitro206.idx
%
\usepackage{graphics}
\usepackage[usenames,dvipsnames]{xcolor}
\usepackage{tikz}
\usetikzlibrary{matrix,calc}
\usepackage[hyperindex=false]{hyperref}
%
\theoremstyle{definition}
\newtheorem{comcnt}{Tout}[subsection]
\newcommand\thingy{%
\refstepcounter{comcnt}\smallskip\noindent\textbf{\thecomcnt.} }
\newcommand\exercice{%
\refstepcounter{comcnt}\bigskip\noindent\textbf{Exercice~\thecomcnt.}\par\nobreak}
\newtheorem{defn}[comcnt]{Définition}
\newtheorem{prop}[comcnt]{Proposition}
\newtheorem{lem}[comcnt]{Lemme}
\newtheorem{thm}[comcnt]{Théorème}
\newtheorem{cor}[comcnt]{Corollaire}
\newtheorem{scho}[comcnt]{Scholie}
\newtheorem{algo}[comcnt]{Algorithme}
\renewcommand{\qedsymbol}{\smiley}
%
\newcommand{\outnb}{\operatorname{outnb}}
\newcommand{\downstr}{\operatorname{downstr}}
\newcommand{\precs}{\operatorname{precs}}
\newcommand{\mex}{\operatorname{mex}}
\newcommand{\id}{\operatorname{id}}
\newcommand{\limp}{\Longrightarrow}
\newcommand{\gr}{\operatorname{gr}}
\newcommand{\rk}{\operatorname{rk}}
\newcommand{\fuzzy}{\mathrel{\|}}
%
\newcommand{\dblunderline}[1]{\underline{\underline{#1}}}
%
\DeclareUnicodeCharacter{00A0}{~}
%
\DeclareMathSymbol{\tiret}{\mathord}{operators}{"7C}
\DeclareMathSymbol{\traitdunion}{\mathord}{operators}{"2D}
%
\DeclareFontFamily{U}{manual}{} 
\DeclareFontShape{U}{manual}{m}{n}{ <->  manfnt }{}
\newcommand{\manfntsymbol}[1]{%
    {\fontencoding{U}\fontfamily{manual}\selectfont\symbol{#1}}}
\newcommand{\dbend}{\manfntsymbol{127}}% Z-shaped
\newcommand{\danger}{\noindent\hangindent\parindent\hangafter=-2%
  \hbox to0pt{\hskip-\hangindent\dbend\hfill}}
%
\newcommand{\spaceout}{\hskip1emplus2emminus.5em}
\newif\ifcorrige
\corrigetrue
\newenvironment{corrige}%
{\ifcorrige\relax\else\setbox0=\vbox\bgroup\fi%
\smallbreak\noindent{\underbar{\textit{Corrigé.}}\quad}}
{{\hbox{}\nobreak\hfill\checkmark}%
\ifcorrige\relax\else\egroup\fi\par}
%
\newcommand{\defin}[2][]{\def\latexsucks{#1}\ifx\latexsucks\empty\index{#2}\else\index{\latexsucks}\fi\textbf{#2}}
%
%
%
\makeindex
\begin{document}
\title{Théories des jeux\\(notes de cours)}
\author{David A. Madore}
\maketitle

\centerline{\textbf{CSC-4MI06-TP / MITRO206}}

{\footnotesize
\immediate\write18{sh ./vc > vcline.tex}
\begin{center}
Git: \input{vcline.tex}
\end{center}
\immediate\write18{echo ' (stale)' >> vcline.tex}
\par}

\pretolerance=8000
\tolerance=50000



%
%
%

{\footnotesize
\tableofcontents
\par}

\bigbreak

\section{Introduction et typologie}

\subsection{La notion de jeu mathématique : généralités}

\thingy Il n'existe pas une théorie des jeux mais \emph{des}
théorie\emph{s} des jeux.

Il n'est pas possible de donner une définition générale
précise de la notion de « jeu mathématique ».  On verra plus loin des
définitions précises de certains types de jeux (p. ex., les jeux
impartiaux à information parfaite), mais il n'existe pas de définition
générale utile qui s'applique à tous ces types, et à partir de
laquelle on pourrait développer une théorie intéressante.

Pire, différentes disciplines se sont développées sous le nom de
« théorie des jeux », chacune donnant une définition différente de ce
qu'est un « jeu ».  Par exemple, l'étude des jeux « en forme normale »
(=jeux définis par des matrices de gains), la théorie combinatoire des
jeux (jeux à information parfaite), la théorie des jeux logiques, la
théorie des jeux différentiels, etc.  Il n'existe donc pas une mais
plusieurs théories des jeux.

Ces différentes théories des jeux intersectent différentes branches
des mathématiques ou d'autres sciences : probabilités,
optimisation/contrôle, combinatoire, logique, calculabilité,
complexité, analyse/EDP ou encore (en-dehors ou en marge des
mathématiques), économie, cryptographie, physique quantique,
cybernétique, biologie, sociologie, linguistique, philosophie.

Il va de soi qu'on ne pourra dans ce cours donner qu'un aperçu de
quelques unes de ces théories des jeux.


\thingy Une tentative pour approcher la notion de jeu mathématique :
le jeu possède un \defin{état}, qui évolue dans un ensemble (fini ou
infini) d'états ou \defin{positions} possibles ; un certain nombre de
\defin{joueurs} choisissent, simultanément ou consécutivement, un
\defin{coup} à jouer parmi différentes \defin[option]{options}, en fonction
de l'état courant, ou peut-être seulement d'une fonction de l'état
courant ; ce coup peut éventuellement faire intervenir un aléa (hasard
voulu par le joueur) ; l'état du jeu évolue en fonction des coups des
joueurs et éventuellement d'un autre aléa (hasard intrinsèque au
jeu) ; au bout d'un certain nombre de coups (fini ou infini), la règle
du jeu attribue, en fonction de l'état final, ou de son évolution
complète, un \defin{gain} à chaque joueur, ce gain pouvant être un
réel (gain numérique), l'étiquette « gagné » / « perdu », ou encore
autre chose, et chaque joueur cherche en priorité à maximiser son gain
(i.e., à gagner le plus possible, ou à gagner tout court), ou dans le
cas probabiliste, son espérance de gain.

Mais même cette définition très vague est incomplète !, par exemple
dans le cas des jeux différentiels, les coups n'ont pas lieu tour à
tour mais continûment.

Une \defin{stratégie} d'un joueur est la fonction par laquelle il
choisit son coup à jouer en fonction de l'état du jeu (ou de la
fonction de l'état qui lui est présentée), et d'aléa éventuel.  On
peut ainsi résumer le jeu en : chaque joueur choisit une stratégie, et
la règle du jeu définit alors un gain pour chaque joueur.  Les
stratégies peuvent être contraintes de différentes manières (par
exemple : être calculables par une machine de Turing).  Une stratégie
est dite \defin{gagnante} si le joueur qui l'utilise gagne le jeu
(supposé avoir une notion de « joueur gagnant ») quels que soient les
coups choisis par l'autre joueur.

Il faut aussi se poser la question de si les joueurs peuvent
communiquer entre eux (et si oui, s'ils peuvent prouver leur honnêteté
ou s'engager irrévocablement quant au coup qu'ils vont jouer, etc.).
Dans certains cas, on peut aussi être amené à supposer que les joueurs
ne connaissent pas toute la règle du jeu (voir « information
  complète » ci-dessous).


\subsection{Quelques types de jeux}

\thingy Le \defin{nombre de joueurs} est généralement $2$.  On peut
néanmoins étudier des jeux multi-joueurs, ce qui pose des questions
d'alliances et compliquer la question des buts (un joueur peut être
incapable de gagner lui-même mais être en situation de décider quel
autre joueur gagnera : on parle de « kingmaker »).  On peut aussi
étudier des jeux à un seul joueur (jouant contre le hasard), voire à
zéro joueurs (systèmes dynamiques), mais ceux-ci relèvent plutôt
d'autres domaines.  Dans ce cours, on s'intéressera (presque
uniquement) aux jeux à deux joueurs.

\thingy Les joueurs peuvent avoir \textbf{des intérêts communs,
  opposés, ou toute situation intermédiaire}.

Le cas d'intérêts communs est celui où tous les joueurs ont le même
gain.  Si les joueurs peuvent parfaitement communiquer, on est alors
essentiellement ramené à un jeu à un seul joueur : on s'intéresse donc
ici surtout aux situations où la communication est imparfaite.

Le cas de deux joueurs d'intérêts opposés est le plus courant : dans
le cas de gains numériques, on le modélise en faisant des gains d'un
joueur l'opposé des gains de l'autre — on parle alors de \defin[somme nulle (jeu à)]{jeu à
  somme nulle} ; ou bien la règle fera qu'un et un seul joueur aura
gagné et l'autre perdu (mais parfois, elle peut aussi admettre le
match nul).

Toute autre situation intermédiaire est possible.  Mais on conviendra
bien que le but de chaque joueur est de maximiser son propre gain,
sans considération des gains des autres joueurs.

\thingy Le jeu peut être \index{partial (jeu)|see{partisan}}\index{partisan (jeu)}\defin[impartial (jeu)]{partial/partisan ou impartial}.  Un
jeu impartial est un jeu où tous les joueurs sont traités de façon
équivalente par la règle (le sens de « équivalent » étant à définir
plus précisément selon le type de jeu).

\thingy\label{intro-simultaneous-or-sequential} Les coups des joueurs
peuvent avoir lieu \textbf{simultanément ou séquentiellement}.

Formellement, il s'agit seulement d'une différence de présentation.
On peut toujours ramener des coups séquentiels à plusieurs coups
simultanés en n'offrant qu'une seule option à tous les joueurs sauf
l'un, et réciproquement, on peut ramener des coups simultanés à des
coups séquentiels en cachant à chaque joueur l'information de ce que
l'autre a joué.  La question \ref{question-preposing-moves} est
cependant plus intéressante.

\thingy Le jeu peut être à \defin[information parfaite (jeu à)]{information parfaite} ou non.  Un
jeu à information parfaite est un jeu dont la règle ne fait pas
intervenir le hasard et où chaque joueur joue séquentiellement en
ayant la connaissance complète de l'état du jeu et de tous les coups
effectués antérieurement par tous les autres joueurs.

(Cette notion est parfois distinguée de la notion plus faible
d'\defin[information complète (jeu à)]{information complète}, qui souligne que les joueurs ont
connaissance complète de la \emph{règle} du jeu, i.e., des gains
finaux et des options disponibles à chaque joueur.  Néanmoins, on peut
formellement ramener un jeu à information incomplète en jeu à
information complète en regroupant toute l'inconnue sur les règles du
jeu dans des coups d'un joueur appelé « la nature ».  Dans ce cours,
on ne considérera que des jeux à information parfaite [et toute
  occurrence des mots « information complète » sera probablement un
  lapsus pour « information parfaite »].)

\thingy Le nombre de positions (= états possibles), comme le nombre
d'options dans une position donnée, ou comme le nombre de coups, peut
être \textbf{fini ou infini}.  Même si l'étude des jeux finis (de
différentes manières) est la plus intéressante pour des raisons
pratiques, toutes sortes de jeux infinis peuvent être considérés, par
exemple en logique (voir plus loin sur l'axiome de détermination).
Pour un jeu à durée infinie, le gagnant pourra être déterminé, par
exemple, par toute la suite des coups effectués par les deux joueurs ;
on peut même introduire des coups après un nombre infini de coups,
etc.

De même, l'ensemble des positions, des options ou des temps peut être
\textbf{discret ou continu}.  Dans ce cours, on s'intéressera presque
exclusivement au cas discret (on écartera, par exemple, la théorie des
jeux différentiels).


\subsection{Quelques exemples en vrac}

\thingy Le jeu de \defin{pile ou face} entre Pauline et Florian.  On
tire une pièce non-truquée : si elle tombe sur pile, Pauline gagne, si
c'est face, c'est Florian.  Aucun des joueurs n'a de choix à faire.
Chacun a une probabilité $\frac{1}{2}$ de gagner, ou une espérance de
$0$ si les gains sont $+1$ au gagnant et $-1$ au perdant (il s'agit
donc d'un jeu à somme nulle).

Variante entre Alice et Bob : maintenant, Alice choisit « pile » ou
« face » avant qu'on (Bob) tire la pièce.  Si Alice a bien prévu, elle
gagne, sinon c'est Bob.  Ici, seule Alice a un choix à faire.
Néanmoins, il n'y a pas de stratégie intéressante : la stratégie
consistant à choisir « pile » offre la même espérance que celle
consistant à choisir « face », et il n'existe pas de stratégie
(c'est-à-dire, de stratégie mesurable par rapport à l'information dont
dispose Alice) offrant une meilleure espérance.

\thingy Variante : Alice choisit « pile » ou « face », l'écrit dans
une enveloppe scellée sans la montrer à Bob (elle s'\emph{engage} sur
son choix), et Bob, plutôt que tirer une pièce, choisit le côté qu'il
montre.  Si Alice a bien deviné le choix de Bob, Alice gagne, sinon
c'est Bob.  Variante : Bob choisit une carte dans un jeu de 52 cartes
sans la montrer à Alice, et Alice doit deviner si la carte est noire ou
rouge.

Variante équivalente : Alice choisit « Alice » ou « Bob » et Bob
choisit simultanément « gagne » ou « perd ».  Si la phrase obtenue en
combinant ces deux mots est « Alice gagne » ou « Bob perd », alors
Alice gagne, si c'est « Alice perd » ou « Bob gagne », alors Bob
gagne.  Encore une variante : Alice et Bob choisissent simultanément
un bit (élément de $\{\mathtt{0},\mathtt{1}\}$), si le XOR de ces deux
bits vaut $\mathtt{0}$ alors Alice gagne, s'il vaut $\mathtt{1}$ c'est
Bob.  Ce jeu est impartial (même s'il n'est pas parfaitement
symétrique entre les joueurs) : Alice n'a pas d'avantage particulier
sur Bob (ce qui est assez évident sur ces dernières variantes).

\begin{center}
\begin{tabular}{r|cc}
$\downarrow$Alice, Bob$\rightarrow$&$\mathtt{0}$/« gagne »&$\mathtt{1}$/« perd »\\\hline
$\mathtt{0}$/« Alice »&$+1,-1$&$-1,+1$\\
$\mathtt{1}$/« Bob »&$-1,+1$&$+1,-1$\\
\end{tabular}
\end{center}

La notion de coups simultanés peut se convertir en coups engagés dans
une enveloppe scellée (cf. \ref{intro-simultaneous-or-sequential}).

On verra, et il est assez facile de comprendre intuitivement, que la
meilleure stratégie possible pour un joueur comme pour l'autre,
consiste à choisir l'une ou l'autre des deux options offertes avec
probabilité $\frac{1}{2}$ (ceci assure une espérance de gain nul quoi
que fasse l'autre joueur).

(En pratique, si on joue de façon répétée à ce jeu, il peut être
intéressant d'essayer d'exploiter le fait que les humains ont des
générateurs aléatoires assez mauvais, et d'arriver à prédire leurs
coups pour gagner.  Ceci est particulièrement amusant avec des petits
enfants.  Voir aussi la « battle of wits » du film \textit{Princess
  Bride} à ce sujet.)

\thingy\label{rock-paper-scissors} Le jeu de
\defin{pierre-papier-ciseaux} : Alice et Bob choisissent
simultanément un élément de l'ensemble $\{\textrm{pierre},\penalty0
\textrm{papier},\penalty0 \textrm{ciseaux}\}$.  S'ils ont choisi le
même élément, le jeu est nul ; sinon, papier gagne sur pierre, ciseaux
gagne sur papier et pierre gagne sur ciseaux (l'intérêt étant qu'il
s'agit d'un « ordre » cyclique, totalement symétrique entre les
options).  Il s'agit toujours d'un jeu à somme nulle (disons que
gagner vaut $+1$ et perdre vaut $-1$), et cette fois les deux joueurs
sont en situation complètement symétrique.

\begin{center}
\begin{tabular}{r|ccc}
$\downarrow$Alice, Bob$\rightarrow$&Pierre&Papier&Ciseaux\\\hline
Pierre&$0,0$&$-1,+1$&$+1,-1$\\
Papier&$+1,-1$&$0,0$&$-1,+1$\\
Ciseaux&$-1,+1$&$+1,-1$&$0,0$\\
\end{tabular}
\end{center}

On verra que la meilleure stratégie possible consiste à choisir
chacune des options avec probabilité $\frac{1}{3}$ (ceci assure une
espérance de gain nul quoi que fasse l'autre joueur).

Ce jeu s'appelle aussi papier-ciseaux-puits, qui est exactement le
même si ce n'est que « pierre » s'appelle maintenant « puits » (donc
ciseaux gagne sur papier, puits gagne sur ciseaux et papier gagne sur
puits) : la stratégie optimale est évidemment la même.

Certains enfants, embrouillés par l'existence des deux variantes,
jouent à pierre-papier-ciseaux-puits, qui permet les quatre options,
et où on convient que la pierre tombe dans le puits : quelle est alors
la stratégie optimale ? il est facile de se convaincre qu'elle
consiste à ne jamais jouer pierre (qui est strictement « dominée » par
puits), et jouer papier, ciseaux ou puits avec probabilité
$\frac{1}{3}$ chacun (cette stratégie garantit un gain au moins nul
quoi que fasse l'autre adversaire, et même strictement positif s'il
joue pierre avec probabilité strictement positive).

\begin{center}
\begin{tabular}{r|cccc}
$\downarrow$Alice, Bob$\rightarrow$&Pierre&Papier&Ciseaux&Puits\\\hline
Pierre&$0,0$&$-1,+1$&$+1,-1$&$-1,+1$\\
Papier&$+1,-1$&$0,0$&$-1,+1$&$+1,-1$\\
Ciseaux&$-1,+1$&$+1,-1$&$0,0$&$-1,+1$\\
Puits&$+1,-1$&$-1,+1$&$+1,-1$&$0,0$\\
\end{tabular}
\end{center}

\thingy\label{prisonners-dilemma} Le \defin{dilemme du prisonnier} :
Alice et Bob choisissent simultanément une option parmi « coopérer »
ou « faire défaut ».  Les gains sont déterminés par la matrice
suivante :

\begin{center}
\begin{tabular}{r|cc}
$\downarrow$Alice, Bob$\rightarrow$&Coopère&Défaut\\\hline
Coopère&$2,2$&$0,4$\\
Défaut&$4,0$&$1,1$\\
\end{tabular}
\end{center}

Ou plus généralement, en remplaçant $4,2,1,0$ par quatre nombres
$T$ (tentation), $R$ (récompense), $P$ (punition) et
$S$ (\textit{sucker}) tels que $T>R>P>S$.  Ces inégalités font que
chaque joueur a intérêt à faire défaut, quelle que soit l'option
choisie par l'autre joueur : on se convaincra facilement que le seul
équilibre de Nash
(cf. \ref{definition-best-response-and-nash-equilibrium}) pour ce jeu
est celui où Alice et Bob font tous deux défaut ; pourtant, tous les
deux reçoivent moins dans cette situation que s'ils coopèrent
mutuellement.

Ce jeu a été énormément étudié du point de vue économique,
psychologique, politique, philosophique, etc., pour trouver des cadres
d'étude justifiant que la coopération est rationnelle, pour expliquer
en quoi le jeu itéré (=répété) diffère du jeu simple, ou pour montrer
que la notion d'équilibre de Nash est perfectible.

\thingy\label{dove-or-hawk} Le jeu du \defin[trouillard (jeu du)]{trouillard}, ou de la
\defin[colombe et faucon]{colombe et du faucon}, obtenu en modifiant les gains du
dilemme du prisonnier pour pénaliser le double défaut (maintenant
appelé rencontre faucon-faucon) plus lourdement que la coopération
(colombe) face au défaut.  Autrement dit :

\begin{center}
\begin{tabular}{r|cc}
$\downarrow$Alice, Bob$\rightarrow$&Colombe&Faucon\\\hline
Colombe&$2,2$&$0,4$\\
Faucon&$4,0$&$-4,-4$\\
\end{tabular}
\end{center}

Ou plus généralement, en remplaçant $4,2,0,-4$ par quatre nombres
$W$ (\textit{win}), $T$ (\textit{truce}), $L$ (\textit{loss}) et
$X$ (\textit{crash}) tels que $W>T>L>X$.  Ces inégalités font que
chaque joueur a intérêt à faire le contraire de ce que fait l'autre
(si Bob joue faucon, Alice a intérêt à jouer colombe, et si Bob joue
colombe, Alice a intérêt à jouer faucon).

(Pour justifier le nom de « jeu du trouillard », on peut évoquer le
scénario d'une course de voitures vers une falaise, à la façon du film
\textit{La Fureur de vivre} : jouer colombe, c'est arrêter sa voiture
avant d'arriver à la falaise, et jouer faucon, c'est ne pas s'arrêter
sauf si l'autre s'est arrêté : celui qui s'arrête passe pour un
trouillard et perd le jeu, mais si aucun ne s'arrête, les deux
voitures tombent dans la falaise, ce qui est pire que de passer pour
un trouillard.)

Ce jeu présente par exemple un intérêt en biologie, notamment pour ce
qui est de l'évolution des comportements.

On pourra se convaincre que ce jeu a trois équilibres de Nash
(cf. \ref{definition-best-response-and-nash-equilibrium} ; en gros, il
s'agit d'une situation dans laquelle aucun des joueurs n'améliorerait
son gain en changeant \emph{unilatéralement} la stratégie employée) :
l'un où Alice joue colombe et Bob joue faucon, un deuxième où c'est le
contraire, et un troisième où chacun joue colombe ou faucon avec les
probabilités respectives $\frac{L-X}{W-T + L-X}$ et $\frac{W-T}{W-T +
  L-X}$ (avec les valeurs ci-dessus : $\frac{2}{3}$ et
$\frac{1}{3}$), pour un gain espéré de $\frac{LW - TX}{W-T + L-X}$
(avec les valeurs ci-dessus : $\frac{4}{3}$).

\thingy\label{battle-of-sexes} La \defin{guerre des sexes}.  Alice et
Bob veulent faire du sport ensemble : Alice préfère l'alpinisme, Bob
préfère la boxe, mais tous les deux préfèrent faire quelque chose avec
l'autre que séparément.  D'où les gains suivants :

\begin{center}
\begin{tabular}{r|cc}
$\downarrow$Alice, Bob$\rightarrow$&Alpinisme&Boxe\\\hline
Alpinisme&$2,1$&$0,0$\\
Boxe&$0,0$&$1,2$\\
\end{tabular}
\end{center}

Ou plus généralement, en remplaçant $2,1,0$ par trois nombres
$P$ (préféré), $Q$ (autre), $N$ (nul) tels que $P>Q>N$.

Ce jeu présente par exemple un intérêt en sociologie, notamment pour
ce qui est de la synchronisation autour d'une ressource commune (par
exemple l'adoption d'un standard).

On pourra se convaincre que ce jeu a trois équilibres de Nash
(cf. \ref{definition-best-response-and-nash-equilibrium}) : l'un où
les deux joueurs vont à l'alpinisme, un deuxième où les deux vont à la
boxe, et un troisième où chacun va à son activité préférée avec
probabilité $\frac{P-N}{P+Q-2N}$ et à l'autre avec probabilité
$\frac{Q-N}{P+Q-2N}$ (avec les valeurs ci-dessus : $\frac{2}{3}$ et
$\frac{1}{3}$), pour un gain espéré de $\frac{PQ-N^2}{P+Q-2N}$ (avec
les valeurs ci-dessus : $\frac{2}{3}$).  Remarquablement, ce gain
espéré est inférieur à $Q$.

\thingy Le \defin[partage (jeu du)]{jeu du partage} ou \defin[ultimatum (jeu de l')]{de l'ultimatum} : Alice
et Bob ont $10$ points à se partager : Alice choisit un $k$ entre $0$
et $10$ entier (disons), la part qu'elle se propose de garder pour
elle, \emph{puis} Bob choisit, en fonction du $k$ proposé par Alice,
d'accepter ou de refuser le partage : s'il accepte, Alice reçoit le
gain $k$ et Bob reçoit le gain $10-k$, tandis que si Bob refuse, les
deux reçoivent $0$.  Cette fois, il ne s'agit pas d'un jeu à somme
nulle !

Variante : Alice choisit $k$ et \emph{simultanément} Bob choisit
$\varphi \colon \{0,\ldots,10\} \to \{\textrm{accepte},\penalty0
\textrm{refuse}\}$.  Si $\varphi(k) = \textrm{accepte}$ alors Alice
reçoit $k$ et Bob reçoit $10-k$, tandis que si $\varphi(k) =
\textrm{refuse}$ alors Alice et Bob reçoivent tous les deux $0$.  Ceci
revient (cf. \ref{question-preposing-moves}) à demander à Bob de
préparer sa réponse $\varphi(k)$ à tous les coups possibles d'Alice
(notons qu'Alice n'a pas connaissance de $\varphi$ quand elle
choisit $k$, les deux sont choisis simultanément).  On se convainc
facilement que si Bob accepte $k$, il devrait aussi accepter tous
les $k'\leq k$, d'où la nouvelle :

Variante : Alice choisit $k$ entre $0$ et $10$ (la somme qu'elle
propose de se garder) et \emph{simultanément} Bob choisit $\ell$ entre
$0$ et $10$ (le maximum qu'il accepte qu'Alice garde pour elle) : si
$k\leq \ell$ alors Alice reçoit $k$ et Bob reçoit $10-k$, tandis que
si $k>\ell$ alors Alice et Bob reçoivent tous les deux $0$.

Ce jeu peut sembler paradoxal pour la raison suivante : dans la
première forme proposée, une fois $k$ choisi, on il semble que Bob ait
toujours intérêt à accepter le partage dès que $k<10$ (il gagnera
quelque chose, alors que s'il refuse il ne gagne rien) ; pourtant, on
a aussi l'impression que refuser un partage pour $k>5$ correspond à
refuser un chantage (Alice dit en quelque sorte à Bob « si tu
n'acceptes pas la petite part que je te laisse, tu n'auras rien du
tout »).

Dans la troisième forme, qui est censée être équivalente, on verra
qu'il existe plusieurs équilibres de Nash, ceux où $\ell=k$ (les deux
joueurs sont d'accord sur le partage) et celui où $k=10$ et $\ell=0$
(les deux joueurs demandent tous les deux la totalité du butin, et
n'obtiennent rien).

\thingy Un jeu idiot : Alice et Bob choisissent simultanément chacun
un entier naturel.  Celui qui a choisi le plus grand gagne (en cas
d'égalité, on peut déclarer le nul, ou décider arbitrairement qu'Alice
gagne — ceci ne changera rien au problème).  Ce jeu résiste à toute
forme d'analyse intelligente, il n'existe pas de stratégie gagnante
(ni d'équilibre de Nash, cf. plus haut), on ne peut rien en dire
d'utile.

Cet exemple sert à illustrer le fait que dans l'étude des jeux sous
forme normale, l'hypothèse de finitude des choix sera généralement
essentielle.

\thingy\label{introduction-graph-game} Le jeu d'un graphe : soit $G$
un graphe orienté (cf. \ref{definitions-graphs} ci-dessous pour la
définition) et $x_0$ un sommet de $G$.  Partant de $x_0$, Alice et Bob
choisissent tour à tour une arête à emprunter pour arriver dans un
nouveau sommet (c'est-à-dire : Alice choisit un voisin sortant $x_1$
de $x_0$, puis Bob un voisin sortant $x_2$ de $x_1$, puis Alice $x_3$
de $x_2$ et ainsi de suite).  \emph{Le perdant est celui qui ne peut
  plus jouer}, et si ceci ne se produit jamais (si on définit un
chemin infini $x_0, x_1, x_2, x_3,\ldots$) alors la partie est
déclarée nulle (ceci ne peut pas se produire lorsque le graphe $G$ est
« bien-fondé »).  On verra qu'il s'agit là du cadre général dans
lequel on étudie la théorie combinatoire des jeux impartiaux à
information parfaite
(cf. \ref{definition-impartial-combinatorial-game}), et qu'un des
joueurs a forcément une stratégie gagnante ou bien les deux joueurs
une stratégie assurant le nul (si le nul est possible)
(cf. \ref{determinacy-of-perfect-information-games}).

Dans une variante du jeu, celui qui ne peut plus jouer gagne au lieu
de perdre : on parle alors de la variante « misère » du jeu.

On peut aussi considérer un graphe dont les arêtes peuvent être
coloriées de trois couleurs possibles : des arêtes rouges, qui ne
peuvent être suivies que par Alice, des arêtes bleues, qui ne peuvent
être suivies que par Bob, et des arêtes vertes (équivalentes à une
arête rouge \emph{et} une arête bleue entre les mêmes deux sommets),
qui peuvent être suivies par l'un ou l'autre joueur (le cas précédent
est donc équivalent à celui d'un graphe entièrement vert).  Il s'agira
là du cadre général dans lequel on étudie la théorie combinatoire des
jeux \emph{partiaux} à information parfaite : on verra que, si le nul
est rendu impossible, quatre cas sont possibles (Alice a une stratégie
gagnante qui que soit le joueur qui commence, ou Bob en a une, ou le
premier joueur a une stratégie gagnante, ou le second en a une).

\thingy\label{introduction-nim-game} Le \defin[nim (jeu de)]{jeu de nim} : un
certain nombre d'allumettes sont
arrangées en plusieurs lignes ; chacun leur tour, Alice et Bob
retirent des allumettes, au moins une à chaque fois, et autant qu'ils
veulent, mais \emph{d'une ligne seulement} ; le gagnant est celui qui
retire la dernière allumette (de façon équivalente, le perdant est
celui qui ne peut pas jouer).  Autrement dit, une position du jeu de
nim est une suite finie $(n_1,\ldots,n_r)$ d'entiers naturels
(représentant le nombre d'allumettes de chaque ligne), et un coup
possible à partir de cette position consiste à aller vers l'état
$(n'_1,\ldots,n'_r)$ où $n'_i = n_i$ pour tout $i$ sauf exactement un
pour lequel $n'_i < n_i$.  Il s'agit ici d'un jeu à deux joueurs
impartial à connaissance parfaite (un cas particulier du jeu général
défini en \ref{introduction-graph-game}).  On verra en \ref{subsection-nim-sum} que la théorie de
Grundy permet de décrire exactement la stratégie gagnante : en
anticipant sur la suite, il s'agit de calculer le XOR (= « ou
  exclusif », appelé aussi \index{nim (somme de)|see{somme de nim}}\index{somme de nim}\textit{somme de nim} dans ce contexte des
nombres $n_i$ d'allumettes des différentes lignes (écrits en
binaire) : ce XOR s'appelle la \index{Grundy (fonction de)}\textit{fonction de Grundy} de la
position, et le jeu est gagnable par le second joueur (c'est-à-dire,
celui qui \emph{vient de} jouer) si et seulement cette fonction de
Grundy vaut $0$.  (À titre d'exemple, la position de départ la plus
courante du jeu de nim est $(1,3,5,7)$, et comme $\mathtt{001} \oplus
\mathtt{011} \oplus \mathtt{101} \oplus \mathtt{111} = \mathtt{000}$
en binaire, en notant $\oplus$ pour le XOR, le second joueur a une
stratégie gagnante.)

On peut aussi jouer à la variante « misère » du jeu : celui qui prend
la dernière allumette a perdu (cf. le film \textit{L'année dernière à
  Marienbad}) ; néanmoins, elle se ramène assez facilement à la
variante « normale » (où celui qui prend la dernière allumette a
gagné), cette dernière ayant plus d'intérêt mathématique.

Le jeu de nim apparaît sous différents déguisements.  On peut par
exemple évoquer le suivant, complètement équivalent à ce qu'on vient
de dire : on place $r$ jetons sur un plateau formé d'une seule ligne
dont les cases sont numérotées $0,1,2,3,\ldots$ (de la gauche vers la
droite, pour fixer les idées).  Chacun tour à tour déplace un jeton
vers la gauche ; plusieurs jetons ont le droit de se trouver sur la
même case, et ils peuvent passer par-dessus l'un l'autre.  Le perdant
est celui qui ne peut plus jouer (parce que tous les jetons sont sur
la case la plus à gauche, $0$).  Il s'agit exactement du jeu de nim,
en considérant que la position où les jetons sont sur les cases
$n_1,\ldots,n_r$ correspond à celle du jeu de nim où il y a
$n_1,\ldots,n_r$ allumettes sur les différentes lignes.  C'est ce
point de vue qui suggère le type de jeux suivant :

\thingy Jeux de \defin{retournement de pièces}.  Ici une position est
une rangée de pièces (qui pourront être numérotées, de la gauche vers
la droite, de $0$ à $N-1$ ou de $1$ à $N$, selon la commodité du jeu),
chacune en position « pile vers le haut » (qu'on notera $\mathtt{0}$)
ou « face vers le haut » ($\mathtt{1}$).  Chaque joueur tour à tour va
retourner certaines pièces selon des règles propres au jeu, avec
toujours la règle générale que \textit{au moins une pièce est
  retournée, et la plus à droite à être retournée doit passer de face
  à pile} (d'autres pièces peuvent passer de pile à face, et d'autres
pièces plus à droite peuvent rester sur pile ou rester sur face, mais
la plus à droite parmi les pièces qui se font retourner devait être
face avant le mouvement et devient du coup pile).  Cette règle
générale assure que le nombre binaire formé de l'ensemble des pièces,
lues de la droite vers la gauche, diminue strictement à chaque coup,
et donc que le jeu termine forcément en temps fini.  Le joueur qui ne
peut plus jouer a perdu.

Il faut bien sûr mettre des règles supplémentaires restreignant les
retournements possibles, sinon le jeu n'a aucun intérêt (le premier
joueur met toutes les pièces à montrer pile et gagne immédiatement).
Quelques exemples de telles règles peuvent être :
\begin{itemize}
\item On ne peut retourner qu'une pièce à chaque coup.  Dans ce cas,
  seul importe le nombre de pièces montrant face, et les joueurs n'ont
  essentiellement aucun choix dans le coup à jouer : peu importe la
  pièce retournée (qui passe forcément de face à pile) ; si le nombre
  de pièces montrant face est pair, le second joueur gagne, tandis que
  s'il est impair, c'est le premier qui gagne.  Ce jeu est très peu
  intéressant.
\item On retourne exactement deux pièces à chaque coup (toujours avec
  la règle générale que la plus à droite des deux passe de face à
  pile).  Il s'agit de nouveau du jeu de nim déguisé (mais un peu
  mieux) : si les pièces sont numérotées à partir de $0$ (la plus à
  gauche), retourner les pièces $n$ et $n'<n$ peut se comprendre comme
  faire passer une ligne d'allumettes de $n$ à $n'$ allumettes, avec
  la différence que deux lignes identiques disparaissent mais on peut
  montrer que cette différence n'a aucun impact sur le jeu de nim
  (essentiellement parce que deux lignes identiques s'« annulent » :
  si un joueur prend des allumettes de l'une, l'autre peut faire le
  même coup sur l'autre).
\item On retourne \emph{une ou deux} pièces (toujours avec la règle
  générale que la plus à droite des deux passe de face à pile).  Il
  s'agit encore une fois de nim déguisé, mais cette fois en numérotant
  les pièces à partir de $1$ (retourner une seule pièce revient à
  vider une ligne de nim, en retourner deux revient à diminuer le
  nombre de pièces d'une ligne).
\item On retourne \emph{au plus trois} pièces (toujours avec la règle
  générale).  On peut décrire la stratégie gagnante dans ce jeu en
  rapport avec le code de parité binaire.  Plus généralement, les jeux
  où on retourne au plus $s$ pièces peuvent, pour les petites valeurs
  de $s$, être reliés à des codes correcteurs remarquables.
\item On retourne n'importe quel nombre de pièces, mais elles doivent
  être consécutives (et toujours avec la règle générale que la pièce
  retournée la plus à droite passe de face à pile).  Il est assez
  facile de décrire la stratégie gagnante de ce jeu.
\end{itemize}

On peut aussi considérer des jeux de retournement de pièces
bidimensionnels : une position est alors un damier, par exemple avec
$M$ lignes et $N$ colonnes (qu'on peut donc repérer comme
$\{0,\ldots,M-1\} \times \{0,\ldots,N-1\}$) avec une pièce à chaque
case, qui peut montrer pile ou face.  Donnons juste un exemple de tel
jeu : chaque joueur peut retourner soit une seule pièce, soit
exactement deux pièces de la même ligne, soit exactement deux pièces
de la même colonne, soit exactement quatre pièces formant les quatre
sommets d'un rectangle (i.e., définies par l'intersection de deux
lignes et de deux colonnes), avec la contrainte supplémentaire que
dans chaque cas la pièce la plus en bas à droite de celles retournées
doit passer de face à pile.

\thingy Le jeu de \defin{chomp} ou de la tablette de chocolat (ou
gaufre) empoisonnée.

On part d'une « tablette de chocolat » de taille $m\times n$,
c'est-à-dire le produit $\{0,\ldots,m-1\} \times \{0,\ldots,n-1\}$
dont les éléments (les couples $(i,j)$ avec $0\leq i<m$ et $0\leq
j<n$) sont appelés les « carrés » de la tablette ; un état général du
jeu sera un sous-ensemble de ce produit (l'ensemble des carrés restant
à manger).  Le carré $(0,0)$ est empoisonné et le but est de ne pas le
manger.  Un coup consiste à choisir un carré $(i,j)$ où mordre dans la
tablette, ce qui fait disparaître tous les carrés $(i',j')$ avec
$i'\geq i$ et $j'\geq j$.  Chaque joueur, tour à tour, effectue un
coup de la sorte, et le premier à mordre dans la case empoisonnée
$(0,0)$ a perdu (de façon équivalente, on ne peut pas mordre dedans,
ce qui se ramène au formalisme général où le premier qui ne peut pas
jouer a perdu).

On ne sait pas décrire la stratégie gagnante en général, mais on peut
montrer que, partant d'une tablette rectangulaire (ou carrée) $m\times
n$ (par opposition à une forme irrégulière quelconque), le
\emph{premier joueur} a forcément une stratégie gagnante.  En effet,
en admettant provisoirement
(cf. \ref{determinacy-of-perfect-information-games}) qu'un des deux
joueurs a une stratégie gagnante, montrons qu'il s'agit forcément du
premier ; pour cela, supposons par l'absurde que le second joueur ait
une stratégie gagnante, et considérons la réponse $(i,j)$ préconisée
par cette stratégie si le premier joueur joue en mordant la case
$(m-1,n-1)$ opposée à la case empoisonnée : à partir de l'état obtenu
en jouant cette réponse (i.e., toutes les cases $(i',j')$ avec $i'\geq
i$ et $j'\geq j$ ont été mangées), le joueur qui vient de jouer est
censé avoir une stratégie gagnante ; mais si le premier joueur jouait
directement en mordant en $(i,j)$, il se ramènerait à cet état, les
rôles des joueurs étant inversés, donc il aurait une stratégie
gagnante, et cela signifie qu'il en a une dès le premier tour.

\thingy\label{introduction-hackenbush} Le jeu de \defin{Hackenbush}
impartial, bicolore, ou tricolore.

Dans ce jeu, l'état est défini par un dessin, plus précisément un
graphe non orienté, pouvant avoir des arêtes multiples et des arêtes
reliant un sommet à lui-même, dont certains sommets sont « au sol »
(graphiquement représentés en les plaçant sur une droite horizontale
en bas du dessin).  Chaque sommet et chaque arête doit être « relié au
  sol », c'est-à-dire atteignable depuis un sommet au sol par une
succession d'arêtes.  De plus, dans le cas de Hackenbush bicolore,
chaque arête est coloriée rouge ou bleue, dans le cas de Hackenbush
tricolore elle peut aussi être verte, et dans le cas de Hackenbush
impartial il n'y a pas de couleur, ou, si on préfère, toutes les
arêtes sont vertes.

\begin{center}
\begin{tikzpicture}
\draw[very thin] (-0.5,0) -- (3.5,0);
\begin{scope}[every node/.style={circle,fill,inner sep=0.5mm}]
\node (P0) at (0,0) {};
\node (P1) at (1.5,0) {};
\node (P2) at (0.75,1) {};
\node (P3) at (0.75,2.5) {};
\node (P4) at (0,2) {};
\node (P5) at (1.5,2) {};
\node (Q0) at (3,0) {};
\node (Q1) at (3,1.5) {};
\node (Q2) at (3,3) {};
\end{scope}
\begin{scope}[line width=1.5pt]
\draw[color=green] (P0) -- (P2);
\draw[color=green] (P1) -- (P2);
\draw[color=green] (P2) -- (P3);
\draw[color=green] (P3) -- (P4);
\draw[color=green] (P3) -- (P5);
\draw[color=green] (P3) .. controls (1.0,2.75) and (1.0,3.0) .. (0.75,3.0) .. controls (0.5,3.0) and (0.5,2.75) .. (P3);
\draw[color=red] (Q0) -- (Q1);
\draw[color=blue] (Q1) -- (Q2);
\end{scope}
\end{tikzpicture}
\\{\footnotesize (Un état possible de Hackenbush.)}
\end{center}

Alice et Bob jouent tour à tour, chacun efface une arête du dessin, ce
qui fait disparaître du même coup toutes les arêtes et tous les
sommets qui ne sont plus reliés au sol.  (Par exemple, dans le dessin
représenté ci-dessus, si on efface l'arête rouge, l'arête bleue
au-dessus disparaît immédiatement ; si on efface l'une des deux arêtes
vertes reliées au sol, les « jambes » du « bonhomme » vert, rien de
particulier ne se passe mais si on efface la deuxième, toutes les
arêtes vertes disparaissent.)  Dans le jeu de Hackenbush impartial,
n'importe quel joueur peut effacer n'importe quelle arête ; dans le
jeu bicolore, seule Alice peut effacer les arêtes rouges et seul Bob
peut effacer les arêtes bleues ; dans le jeu tricolore, les arêtes
vertes sont effaçables par l'un ou l'autre joueur (mais dans tous les
cas, la disparition des arêtes non reliées au sol est automatique).
Le jeu se termine quand un joueur ne peut plus jouer, auquel cas il a
perdu (au Hackenbush impartial, cela signifie que le jeu se termine
quand un joueur finit de faire disparaître le dessin, auquel cas il a
gagné ; au Hackenbush bicolore ou tricolore, il se peut bien sûr qu'il
reste des arêtes de la couleur du joueur qui vient de jouer).

Le jeu de Hackenbush impartial possède une stratégie gagnante soit par
le premier soit par le second joueur (le dessin formé uniquement des
arêtes vertes ci-dessus, par exemple, est gagnable par le premier
joueur, le seul coup gagnant consistant à effacer le « corps » du
« bonhomme » pour ne laisser que ses jambes).  Le jeu de Hackenbush
bicolore possède une stratégie gagnante soit pour Alice, soit pour
Bob, soit pour le second joueur, mais jamais pour le premier (le
dessin formé par les arêtes rouge et bleue ci-dessus, par exemple, est
gagnable par Alice).  Le jeu de Hackenbush tricolore possède une
stratégie gagnante soit pour Alice, soit pour Bob, soit pour le
premier joueur, soit pour le second (l'ensemble du dessin ci-dessus,
par exemple, est gagnable par Alice).

\thingy\label{introduction-hydra-game}
Le \defin[hydre (jeu de l')]{jeu de l'hydre} : Hercule essaie de terrasser
l'hydre.  Le joueur qui joue l'hydre commence par dessiner (i.e.,
choisir) un arbre (fini, enraciné), la forme initiale de l'hydre.
Puis Hercule choisit une \emph{tête} de l'hydre, c'est-à-dire une
feuille $x$ de l'arbre, et la décapite en la supprimant de l'arbre.
L'hydre se reproduit alors de la façon suivante : soit $y$ le nœud
parent de $x$ dans l'arbre, et $z$ le nœud parent de $y$ (grand-parent
de $x$, donc) : si l'un ou l'autre n'existe pas, rien ne se passe
(l'hydre passe son tour) ; sinon, l'hydre choisit un entier naturel
$n$ (aussi grand qu'elle veut) et attache à $z$ autant de nouvelles
copies de $y$ (mais sans la tête $x$ qui a été décapitée) qu'elle le
souhaite.  Hercule gagne s'il réussit à décapiter le dernier nœud de
l'hydre ; l'hydre gagnerait si elle réussissait à survivre
indéfiniment.

\begin{center}
\begin{tikzpicture}[baseline=0]
\draw[very thin] (-1.5,0) -- (1.5,0);
\begin{scope}[every node/.style={circle,fill,inner sep=0.5mm}]
\node (P0) at (0,0) {};
\node (P1) at (0,1) {};
\node (P2) at (-1,2) {};
\node (P3) at (0,2) {};
\node (P4) at (1,2) {};
\node (P5) at (0.5,3) {};
\node (P6) at (1.5,3) {};
\end{scope}
\begin{scope}[line width=1.5pt]
\draw (P0) -- (P1);
\draw (P1) -- (P2);
\draw (P1) -- (P3);
\draw (P1) -- (P4);
\draw (P4) -- (P5);
\draw (P4) -- (P6);
\end{scope}
\begin{scope}[line width=3pt,red]
\draw ($(P6) + (-0.2,-0.2)$) -- ($(P6) + (0.2,0.2)$);
\draw ($(P6) + (-0.2,0.2)$) -- ($(P6) + (0.2,-0.2)$);
\end{scope}
\node[anchor=west] at (P6) {$x$};
\node[anchor=west] at (P4) {$y$};
\node[anchor=west] at (P1) {$z$};
\end{tikzpicture}
devient
\begin{tikzpicture}[baseline=0]
\draw[very thin] (-1.5,0) -- (1.5,0);
\begin{scope}[every node/.style={circle,fill,inner sep=0.5mm}]
\node (P0) at (0,0) {};
\node (P1) at (0,1) {};
\node (P2) at (-1,2) {};
\node (P3) at (0,2) {};
\node (P4) at (0.8,2) {};
\node (P5) at (0.6,3) {};
\node (P4b) at (1.2,2) {};
\node (P5b) at (1.2,3) {};
\node (P4c) at (1.7,2) {};
\node (P5c) at (1.7,3) {};
\node (P4d) at (2.2,2) {};
\node (P5d) at (2.2,3) {};
\end{scope}
\begin{scope}[line width=1.5pt]
\draw (P0) -- (P1);
\draw (P1) -- (P2);
\draw (P1) -- (P3);
\draw (P1) -- (P4);
\draw (P4) -- (P5);
\draw[blue] (P1) -- (P4b);
\draw[blue] (P4b) -- (P5b);
\draw[blue] (P1) -- (P4c);
\draw[blue] (P4c) -- (P5c);
\draw[blue] (P1) -- (P4d);
\draw[blue] (P4d) -- (P5d);
\end{scope}
\end{tikzpicture}
\end{center}

Ce jeu est particulier en ce que, mathématiquement, non seulement
Hercule possède une stratégie gagnante, mais en fait Hercule gagne
\emph{toujours}, quoi qu'il fasse et quoi que fasse l'hydre
(cf. \ref{subsection-hydra-game-again}).  Pourtant, en pratique,
l'hydre peut facilement s'arranger pour survivre un temps
inimaginablement long.

\thingy Le \defin[Choquet (jeu topologique de)]{jeu topologique de Choquet} : soit $X$ un espace
métrique (ou topologique) fixé à l'avance.  Uriel et Vania choisissent
tour à tour un ouvert non vide de ($X$ contenu dans) l'ouvert
précédemment choisi : i.e., Uriel choisit $\varnothing \neq U_0
\subseteq X$, puis Vania choisit $\varnothing \neq V_0 \subseteq U_0$,
puis Uriel choisit $\varnothing \neq U_1 \subseteq V_0$ et ainsi de
suite.  Le jeu continue pendant un nombre infini de tours indicés par
les entiers naturels.  À la fin, on a bien sûr $\bigcap_{n=0}^{\infty}
U_n = \bigcap_{n=0}^{\infty} V_n$ : on dit qu'Uriel gagne le jeu si
cette intersection est vide, Vania le gagne si elle est non-vide.  On
peut se convaincre que si $X = \mathbb{Q}$, alors Uriel possède une
stratégie gagnante, tandis que si $X = \mathbb{R}$ c'est Vania qui en
a une.

\thingy\label{introduction-gale-stewart-games} Les \defin[Gale-Stewart (jeu de)]{jeux de
  Gale-Stewart} (cf. partie \ref{section-gale-stewart-games}) : soit $A$ un sous-ensemble de
$\mathbb{N}^{\mathbb{N}}$ ou de $\{0,1\}^{\mathbb{N}}$ ou de $[0,1]$.
Alice et Bob choisissent tour à tour un élément de $\mathbb{N}$ (dans
le premier cas) ou de $\{0,1\}$ (dans les deux suivants).  Ils jouent
un nombre infini de tours, « à la fin » desquels la suite de leurs
coups définit un élément de $\mathbb{N}^{\mathbb{N}}$ ou de
$\{0,1\}^{\mathbb{N}}$ ou, en les considérant comme la suite des
chiffres binaires d'un réel commençant par $0.$, de $[0,1]$ : si cet
élément appartient à $A$, Alice gagne, sinon c'est Bob (la partie
n'est jamais nulle).  \emph{Il n'est pas vrai} qu'un des deux joueurs
possède forcément une stratégie gagnante.

\thingy Considérons le jeu suivant : Turing choisit publiquement une
machine de Turing (i.e., un programme sur ordinateur) et Blanche (son
adversaire) doit répondre soit « elle termine en $n$ étapes » où $n$
est un entier naturel (explicite), soit « elle ne termine pas ».  Dans
le premier cas, on lance l'exécution de la machine de Turing sur $n$
étapes, et si elle termine bien dans le temps annoncé, Blanche a
gagné, sinon c'est Turing qui a gagné.  Dans le second cas (i.e., si
Blanche a annoncé « elle ne termine pas »), c'est à Turing d'annoncer
soit « si, elle termine en $m$ étapes » où $m$ est un entier naturel
(explicite), soit « en effet, elle ne termine pas ».  Dans le premier
sous-cas, on lance l'exécution de la machine de Turing sur $m$ étapes,
et si elle termine bien dans le temps annoncé, Turing a gagné, sinon
c'est Blanche qui a gagné.  Dans le second sous-cas (i.e., si Turing a
confirmé « en effet, elle ne termine pas »), Blanche a gagné.

Dit de façon plus simple : Turing propose à Blanche de décider l'arrêt
d'une machine de Turing ; si Blanche prédit l'arrêt, elle doit donner
le nombre d'étapes et on peut vérifier cette affirmation ; si elle
prédit le contraire, c'est à Turing de la contredire le cas échéant
par une affirmation d'arrêt, qui sera elle aussi vérifiée.

La règle du jeu peut être implémentée algorithmiquement : i.e., on
peut vérifier (sur une machine de Turing !) qui gagne ou qui perd en
fonction des coups joués (puisque à chaque fois on fait des
vérifications finies).  Néanmoins, aucun des joueurs n'a de stratégie
gagnante \emph{algorithmique} (i.e., choisissant un coup
algorithmiquement en fonction des coups antérieurs).  En fait, Turing
n'a pas de stratégie gagnante du tout (quelle que soit la machine
qu'il choisit au premier coup, Blanche \emph{pourrait} répondre
correctement auquel cas Turing ne gagne pas).  Mais Blanche n'a pas de
stratégie gagnante algorithmique, car cela reviendrait à résoudre le
problème de l'arrêt.

Cet exemple illustre le fait qu'on ne peut pas espérer avoir un
algorithme qui calcule un coup gagnant dans n'importe quel jeu même si
on se limite aux jeux dont le gain est calculable algorithmiquement.

(On peut remplacer le problème de l'arrêt par n'importe quel problème
semi-décidable : si $f \colon \mathbb{N} \to \mathbb{N}$ est une
fonction algorithmiquement calculable dont l'image n'est pas
décidable, Turing choisit un élément $y$ de $\mathbb{N}$, Blanche doit
soit répondre « $y=f(n)$ » pour un $n$ explicite soit « $y$ n'est pas
  dans l'image », auquel cas Turing peut soit rétorquer « si, $y =
  f(m)$ » soit concéder que $y$ n'est pas dans l'image.  Autre
exemple : Turing choisit un énoncé mathématique, Blanche doit soit le
démontrer soit dire que ce n'est pas un théorème, et dans le second
cas c'est à Turing de le démontrer.)


\subsection{Remarques}

\thingy\label{question-preposing-moves} La question suivante mérite
l'attention : supposons que, dans un jeu, deux joueurs aient à jouer
deux coups successifs, disons que le joueur $A$ choisit une option $x$
parmi un certain ensemble $E$ (typiquement fini), \emph{puis} le
joueur $B$ choisit, en connaissant le $x$ choisi par $A$, une option
$y$ parmi un certain ensemble $F$ (typiquement fini).  Revient-il au
même de demander de choisir \emph{simultanément} pour $A$ un élément
de $E$ et pour $B$ un élément de l'ensemble $F^E$ des fonctions de $E$
dans $F$ ?  L'idée étant que $B$ choisit la fonction $\varphi$ qui,
selon le coup $x \in E$ joué par $A$, déterminera le coup $y :=
\varphi(x) \in F$ qu'il joue en réponse.  Au moins si $E$ est fini, on
peut imaginer que $B$ considère mentalement tous les coups que $A$
pourra jouer et choisit la réponse qu'il y apporterait, déterminant
ainsi la fonction $\varphi$ (si on préfère, $\varphi$ est une
stratégie locale pour le prochain coup de $B$).

En principe, les jeux ainsi considérés (le jeu initial, et celui où on
a demandé à $B$ d'anticiper son choix en le remplaçant par une
fonction du choix de $A$) devraient être équivalents.  En pratique, il
se peut qu'on les analyse différemment pour différentes raisons.

Notons que si on permet ou oblige $B$ à communiquer à $A$ la fonction
$\varphi$ qu'il a choisie, i.e., à s'\emph{engager} irrévocablement
sur le coup $y$ qu'il jouerait selon le coup $x$ de $A$, on peut
véritablement changer le jeu.


\subsection{Plan}

La partie \ref{section-games-in-normal-form} concerne les jeux en
forme normale et la notion d'équilibre de Nash : on gardera donc à
l'esprit les exemples tels que le dilemme du
prisonnier (\ref{prisonners-dilemma}), le
trouillard (\ref{dove-or-hawk}) et la bataille des
sexes (\ref{battle-of-sexes}).  On évoque plus particulièrement les
jeux à somme nulle en \ref{zero-sum-games} : on pensera alors à des
jeux comme pierre-papier-ciseaux (cf. \ref{rock-paper-scissors}).

La partie \ref{section-gale-stewart-games} introduit la notion de jeux de
Gale-Stewart et prouve un théorème fondamental de détermination (la
détermination des jeux \emph{ouverts}).

La partie \ref{section-well-founded-induction} introduit la notion de
graphe bien-fondée et d'induction bien-fondée qui est essentielle pour
la suite.  La partie \ref{section-ordinals} introduit la notion
d'ordinaux qui permet de généraliser beaucoup de résultats du fini à
l'infini.

La partie \ref{section-combinatorial-impartial-games} concerne la
théorie, dite « combinatoire », des jeux impartiaux à information
parfaite, dont le modèle est décrit en \ref{introduction-graph-game}
(sans coloriage) et dont l'archétype est le jeu de nim
(cf. \ref{introduction-nim-game}) ou le Hackenbush impartial (=vert)
(cf. \ref{introduction-hackenbush}).

On parlera ensuite rapidement dans la
partie \ref{section-combinatorial-partizan-games} des jeux
\emph{partisans} à information parfaite, dont l'archétype est le
Hackenbush bicolore ou tricolore, et de la théorie des nombres
surréels de Conway.


%
%
%

\section{Jeux en forme normale}\label{section-games-in-normal-form}

\setcounter{subsection}{-1}

\subsection{Rappels sur les convexes}

\thingy Pour cette section, on rappelle les définitions suivants : une
\index{affine (combinaison)|see{combinaison affine}}\defin{combinaison
  affine} ou \index{barycentrique (combinaison)|see{combinaison
    barycentrique}}\defin{combinaison barycentrique} d'éléments
$x_1,\ldots,x_m \in \mathbb{R}^n$ est une expression de la forme
$\sum_{i=1}^m \lambda_i x_i \in \mathbb{R}^n$ où
$\lambda_1,\ldots,\lambda_m$ vérifient $\sum_{i=1}^m \lambda_i = 1$
(on peut donc, si on préfère, la définir comme une expression de la
forme $\frac{\sum_{i=1}^m \lambda_i x_i}{\sum_{i=1}^m \lambda_i}$ où
$\lambda_1,\ldots,\lambda_m$ vérifient $\sum_{i=1}^m \lambda_i \neq
0$ : on parle alors de \defin{barycentre} de $x_1,\ldots,x_m$ affecté
des poids $\lambda_1,\ldots,\lambda_m$).

Lorsque $x_1,\ldots,x_m$ sont tels que la combinaison affine est
déterminée par ses coefficients, autrement dit lorsque l'application
$(\lambda_1,\ldots,\lambda_m) \mapsto \sum_{i=1}^m \lambda_i x_i$
définie sur l'ensemble des $(\lambda_1,\ldots,\lambda_m)$ tels que
$\sum_{i=1}^m \lambda_i = 1$, est injective, on dit que
$x_1,\ldots,x_m$ sont \index{affinement indépendants
  (points)}\defin{affinement indépendants} (ou affinement libres).  Il
revient au même de dire que $x_2-x_1,\ldots,x_m-x_1$ sont linéairement
indépendants.

\thingy Une \index{convexe (combinaison)|see{combinaison
    convexe}}\defin{combinaison convexe} est une combinaison affine
$\sum_{i=1}^m \lambda_i x_i$ où $\lambda_1,\ldots,\lambda_m$ sont
\emph{positifs} ($\lambda_1,\ldots,\lambda_m \geq 0$) et vérifient
$\sum_{i=1}^m \lambda_i = 1$, autrement dit, c'est un barycentre
affecté de coefficients positifs (non tous nuls).

Un \defin{convexe} de $\mathbb{R}^m$ est une partie stable par
combinaisons convexes (c'est-à-dire que $C$ est dit convexe lorsque si
$x_1,\ldots,x_m \in C$ et $\lambda_1,\ldots\lambda_m\geq 0$ vérifient
$\sum_{i=1}^m \lambda_i = 1$ alors $\sum_{i=1}^m \lambda_i x_i \in
C$).  Par associativité du barycentre, il revient au même d'imposer
cette condition pour $m=2$, c'est-à-dire que $C \subseteq
\mathbb{R}^m$ est convexe ssi on a $[x,y] \subseteq C$ pour tous
$x,y\in C$, où $[x,y] := \{\lambda x + (1-\lambda) y : \lambda \in
[0;1]\}$.

Une intersection quelconque de convexes étant manifestement un
convexe, on peut parler du plus petit convexe contenant une partie $A
\subseteq \mathbb{R}^m$, appelé \index{convexe
  (enveloppe)|see{enveloppe convexe}}\defin{enveloppe convexe}
de $A$ : il s'agit simplement de l'ensemble de toutes les combinaisons
convexes des points de $A$ (i.e., des ensembles \emph{finis} de points
de $A$).  L'enveloppe convexe d'une partie finie s'appelle un
\defin{polytope} (convexe).  L'enveloppe convexe d'une partie $A$
affinement libre et finie\footnote{Une partie affinement libre de
  $\mathbb{R}^m$ a au plus $m+1$ points, donc le mot « fini » est
  redondant dans ce cadre, mais la définition peut s'étendre en
  dimension infinie auquel cas il n'est plus inutile.} s'appelle
\defin{simplexe} de sommets $A$.

\thingy Une \index{affine (application)}\textbf{application affine}
$u\colon \mathbb{R}^p \to \mathbb{R}^q$ est une fonction qui préserve
les combinaisons affines (autrement dit, si $\sum_{i=1}^m \lambda_i =
1$ alors $u\Big(\sum_{i=1}^m \lambda_i x_i\Big) = \sum_{i=1}^m
\lambda_i\, u(x_i)$).  Il revient au même de dire que $u$ est la somme
d'une constante (dans $\mathbb{R}^q$) et d'une application linéaire
$\mathbb{R}^p \to \mathbb{R}^q$.

On aura fréquemment besoin du fait évident suivant :
\begin{prop}\label{affine-functions-take-extrema-at-boundary}
Soit $A \subseteq \mathbb{R}^m$ un ensemble fini et $C$ son enveloppe
convexe, et soit $u \colon \mathbb{R}^m \to \mathbb{R}$ une
application affine.  Alors $u(a) \leq M$ pour tout $a \in A$ implique
$u(s) \leq M$ pour tout $s \in C$.  Ou, si on préfère, $\max\{u(s) :
s\in C\}$ existe et vaut $\max\{u(a) : a\in A\}$ (et la même
affirmation en remplaçant $\max$ par $\min$ vaut aussi).
\end{prop}

En français : « une fonction affine sur un polytope atteint ses bornes
sur un sommet de ce dernier ».

\begin{proof}
Tout point $s$ de $C$ est combinaison convexe de points de $A$,
c'est-à-dire $s = \sum_{i=1}^n \lambda_i x_i$ avec $x_i \in A$ et
$\lambda_i \geq 0$ vérifiant $\sum_{i=1}^n \lambda_i = 1$.  Comme $u$
est affine, on a alors $u(s) = \sum_{i=1}^n \lambda_i\, u(x_i)$.  Si
$u(a) \leq M$ pour tout $a\in A$, ceci implique $u(s) \leq
\sum_{i=1}^n \lambda_i\, M = M$, comme annoncé.

Pour en déduire l'affirmation sur le $\max$, il suffit d'appliquer ce
qu'on vient de démontrer pour $M = \max\{u(a) : a\in A\}$ (qui existe
puisque $A$ est fini) : on a $u(s) \leq M$ pour tout $s \in C$, et
comme les éléments de $A$ sont dans $C$, ceci montre bien $\max\{u(s)
: s\in C\} = \max\{u(a) : a\in A\}$.
\end{proof}

\begin{prop}\label{affine-functions-take-no-strict-extrema-inside}
Reprenant le contexte de la proposition précédente ($A \subseteq
\mathbb{R}^m$ un ensemble fini, $C$ son enveloppe convexe, et $u
\colon \mathbb{R}^m \to \mathbb{R}$ affine), soit $M := \max_{a \in A}
u$ (dont on vient de voir que c'est aussi $\max_{s\in C} u$) : alors
une combinaison convexe $s$ à coefficients \emph{strictement positifs}
de points de $A$ (i.e., $s = \sum_{i=1}^n \lambda_i x_i$ avec $x_i \in
A$ et $\lambda_i > 0$ vérifiant $\sum_{i=1}^n \lambda_i = 1$) vérifie
$u(s) = M$, si et seulement si chacun des points en question vérifie
aussi cette propriété (i.e., on a $u(x_i) = M$ pour tout $i$).
\end{prop}

\begin{proof}
On vient de voir que $u(x_i) \leq M$ pour tout $i$ implique $u(s) \leq
M$.  Si on a $u(x_j) < M$ pour un certain $j$, alors on a $\lambda_j\,
u(x_j) < \lambda_j\, M$, donc en sommant avec les autres $\lambda_i\,
u(x_i) \leq \lambda_i M$ on trouve $u(s) = \sum_{i=1}^n \lambda_i\,
u(x_i) < \sum_{i=1}^n \lambda_i\, M = M$.  Réciproquement, si $u(x_i)
= M$ pour tout $i$, il est évident que $u(s) = M$.
\end{proof}

On peut aussi reformuler ce résultat en affirmant que la partie de $C$
où $u$ prend sa valeur maximale $M$ est l'enveloppe convexe de la
partie de $A$ où elle prend cette valeur.


\subsection{Généralités}

\begin{defn}\label{definition-game-in-normal-form}
Un \index{forme normale (jeu en)}\defin[normale (jeu en forme)]{jeu en forme normale} à $N$ joueurs est la donnée de $N$
ensembles finis $A_1,\ldots,A_N$ et de $N$ fonctions
$u_1,\ldots,u_N\colon A \to \mathbb{R}$ où $A := A_1 \times \cdots
\times A_N$.

Un élément de $A_i$ s'appelle une \defin{option} ou \index{pure (stratégie)}\defin{stratégie
  pure} pour le joueur $i$.  Un élément de $A := A_1 \times \cdots
\times A_N$ s'appelle un \defin{profil de stratégies pures}.  La
valeur $u_i(a)$ de la fonction $u_i$ sur un $a\in A$ s'appelle le
\defin{gain} du joueur $i$ selon le profil $a$.
\end{defn}

Le jeu doit se comprendre de la manière suivante : chaque joueur
choisit une option $a_i \in A_i$ indépendamment des autres, et chaque
joueur reçoit un gain égal à la valeur $u_i(a_1,\ldots,a_n)$ définie
par le profil $(a_1,\ldots,a_n)$ des choix effectués par tous les
joueurs.  Le but de chaque joueur est de maximiser son propre gain.

On utilisera le terme « option » ou « stratégie pure » selon qu'on
veut souligner que le joueur $i$ choisit effectivement $a_i$ ou décide
a priori de faire forcément ce choix-là.  Cette différence vient du
fait que les joueurs peuvent également jouer de façon probabiliste, ce
qui amène à introduire la notion de stratégie mixte :

\begin{defn}\label{definition-mixed-strategy-abst}
Donné un ensemble $B$ fini d'« options », on appelle \index{mixte
  (stratégie)}\defin{stratégie mixte} sur $B$ une fonction $s\colon
B\to\mathbb{R}$ telle que $s(b)\geq 0$ pour tout $b\in B$ et
$\sum_{b\in B} s(b) = 1$ : autrement dit, il s'agit d'une distribution
de probabilités sur $B$.

Le \defin[support (d'une stratégie mixte)]{support} de $s$ est
l'ensemble des options $b\in B$ pour lesquelles $s(b) > 0$.
\end{defn}

On identifiera un élément $b$ de $B$ à la fonction (stratégie pure)
$\delta_b \colon B\to\mathbb{R}$ qui prend la valeur $1$ en $b$ et $0$
ailleurs (distribution de probabilités de Dirac en $b$).  Ceci permet
de considérer $B$ comme une partie de l'ensemble $S_B$ des stratégies
mixtes, et d'identifier une stratégie mixte $s\colon B\to\mathbb{R}$
quelconque sur $B$ à la combinaison convexe $\sum_{b\in B} s(b)\cdot
b$ des stratégies pures : i.e., on peut considérer les stratégies
mixtes comme des combinaison convexes formelles\footnote{Le mot
  « formel » signifie ici que la combinaison n'a pas d'autre sens que
  comme l'expression par laquelle elle est définie, par exemple
  $\frac{1}{3}\text{Pierre} + \frac{1}{3}\text{Papier} +
  \frac{1}{3}\text{Ciseaux}$ est une combinaison convexe formelle de
  $\{\text{Pierre}, \text{Papier}, \text{Ciseaux}\}$ représentant la
  probabilité uniforme sur cet ensemble.} des éléments de $B$.

On passera indifféremment entre ces trois points de vue sur les
stratégies mixtes : fonctions $s\colon B\to\mathbb{R}$ (positives de
somme $1$), mesures de probabilités sur $B$, ou bien combinaisons
convexes formelles d'éléments de $B$.

En tout état de cause, l'ensemble $S_B$ des stratégies mixtes sur $B$
sera vu comme la partie de $\mathbb{R}^B$ définie par l'intersection
des demi-espaces de coordonnées positives et de l'hyperplan défini par
la somme des coordonnées égale à $1$ : il s'agit d'un \emph{convexe
  fermé}, qui est l'enveloppe convexe des points de $B$ identifiés
ainsi qu'on vient de le dire aux $\delta_b =
(0,\ldots,0,1,0,\ldots,0)$, c'est-à-dire que $S_B$ est le
\defin{simplexe} d'ensemble de sommets $B$ (identifié à l'ensemble
des $\delta_b$).

\begin{defn}\label{definition-mixed-strategy-game}
Pour un jeu comme défini en \ref{definition-game-in-normal-form}, une
stratégie mixte pour le joueur $i$ est donc une fonction $s\colon A_i
\to\mathbb{R}$ comme on vient de le dire
en \ref{definition-mixed-strategy-abst}.  On notera parfois $S_i$
l'ensemble des stratégies mixtes du joueur $i$.  Un \defin{profil de
  stratégies mixtes} est un élément du produit cartésien $S := S_1
\times \cdots \times S_N$.

Plus généralement, si $I \subseteq \{1,\ldots,N\}$ est un ensemble de
joueurs, un élément du produit $S_I := \prod_{j\in I} S_j$ s'appellera
un profil de stratégies mixtes pour l'ensemble $I$ de joueurs ; ceci
sera notamment utilisé si $I = \{1,\ldots,N\}\setminus\{i\}$ est
l'ensemble de tous les joueurs sauf le joueur $i$, auquel cas on
notera $S_{?i} := \prod_{j\neq i} S_j$ l'ensemble des profils.
Naturellement, si chaque composante est une stratégie pure, on pourra
parler de profil de stratégies pures.
\end{defn}

\thingy\label{remark-mixed-stragy-profile-versus-correlated-profile}
Il va de soi qu'un profil de stratégies mixtes, i.e., un élément de $S
:= S_1 \times \cdots \times S_N$, i.e., la donnée d'une distribution
de probabilité sur chaque $A_i$, \emph{n'est pas la même chose} qu'une
distribution de probabilités sur $A := A_1 \times \cdots \times A_N$.

Néanmoins, si $(s_1,\ldots,s_N) \in S$, on définit une distribution de
probabilités $s\colon A\to\mathbb{R}$ (un élément de $S_A$, si on
veut) de la façon suivante :
\[
s(a_1,\ldots,a_N) = s_1(a_1)\cdots s_N(a_N)
\tag{$*$}
\]
(produit des $s_i(a_i)$).

Probabilistiquement, cette formule revient à dire qu'on tire un
élément $a_i$ de $A_i$ selon la distribution $s_i$ \emph{de façon
  indépendante} pour chaque $i$ de manière à former un $N$-uplet
$(a_1,\ldots,a_N)$.  La formule ($*$) servira donc à représenter
l'hypothèse que les joueurs ont accès à des sources de hasard qui sont
indépendantes (si Alice tire son coup à pile ou face et que Bob fait
de même, les pièces tirées par Alice et Bob sont des variables
aléatoires indépendantes).  Les distributions de probabilités $s$
sur $A$ définies par la formule ($*$) sont précisément celles dont
composantes sont indépendantes, et qui sont alors complètement
déterminées par leurs \emph{distributions
  marginales}\footnote{\label{footnote-marginals}La $i$-ième
  \defin{marginale} d'une variable aléatoire sur $A_1\times \cdots
  \times A_N$ est simplement sa $i$-ième composante (= projection
  sur $A_i$).  La $i$-ième \textbf{distribution marginale} de la
  distribution de probabilités $s \colon A \to \mathbb{R}$ est donc la
  distribution de probabilités $s_i \colon A_i \to \mathbb{R}$ qui à
  $b\in A_i$ associe la somme des $s(a_1,\ldots,a_N)$ prise sur tous
  les $N$-uplets $(a_1,\ldots,a_N)$ tels que $a_i =
  b$.} $s_1,\ldots,s_N$.

On identifiera parfois abusivement l'élément $(s_1,\ldots,s_N) \in S$
à la distribution $s\colon A\to\mathbb{R}$ qu'on vient de décrire (ce
n'est pas un problème car $s_i$ se déduit de $s$ : précisément,
$s_i(b) = \sum_{a\,:\, a_i = b} s(a)$ où la somme est prise sur les $a
\in A$ tels que $a_i = b$, cf. note \ref{footnote-marginals}).

\danger Il faudra prendre garde au fait que, du coup, $S$ peut se voir
\emph{soit} comme la partie $S_1\times\cdots\times S_N$ de
$\mathbb{R}^{A_1}\times\cdots\times \mathbb{R}^{A_N}$ formé des
$N$-uplets $(s_1,\ldots,s_N)$, \emph{soit} comme la partie de
$\mathbb{R}^A = \mathbb{R}^{A_1\times\cdots\times A_N}$ formé des
fonctions de la forme $s\colon (a_1,\ldots,a_N) \mapsto s_1(a_1)\cdots
s_N(a_N)$ comme on l'a expliqué aux paragraphes précédents.  Ces deux
points de vue ont un sens et ont parfois un intérêt, mais ils ne
partagent pas les mêmes propriétés.  Par exemple, $S$ est convexe en
tant que partie $S_1\times\cdots\times S_N$ de
$\mathbb{R}^{A_1}\times\cdots\times \mathbb{R}^{A_N}$, mais pas en
tant que partie de $\mathbb{R}^A$.  Néanmoins, ce double point de vue
ne devrait pas causer de confusion.

Ceci conduit à faire la définition suivante :

\begin{defn}\label{definition-mixed-strategy-gain}
Donné un jeu en forme normale comme
en \ref{definition-game-in-normal-form}, si $s := (s_1,\ldots,s_N) \in
S_1 \times \cdots \times S_N$ est un profil de stratégies mixtes, on
appelle \defin[gain espéré]{gain [espéré]} du joueur $i$ selon ce
profil la quantité
\[
u_i(s) := \sum_{a\in A} s_1(a_1)\cdots s_N(a_N)\,u_i(a)
\]
c'est-à-dire selon la distribution de probabilités définie par la
formule ($*$)
de \ref{remark-mixed-stragy-profile-versus-correlated-profile}.  Ceci
étend $u_i$ en une function de $S := S_1\times\cdots \times S_N$
vers $\mathbb{R}$, qu'on dénotera par le même symbole car il n'en
résultera pas de confusion.
\end{defn}

Selon l'approche qu'on veut avoir, on peut dire qu'on a défini
$u_i(s)$ comme l'espérance de $u_i(a)$ si chaque $a_j$ est tiré
indépendamment selon la distribution de probabilité $s_i$ ; ou bien
qu'on a utilisé l'unique prolongement de $u_i$ au produit des
simplexes $S_i$ qui soit affine en chaque variable $s_i$
(« multi-affine »).



\subsection{Équilibres de Nash}

\begin{defn}\label{definition-best-response-and-nash-equilibrium}
Donné un jeu en forme normale comme
en \ref{definition-game-in-normal-form}, si $1 \leq i \leq N$ et si
$s_? := (s_1,\ldots,s_{i-1},s_{i+1},\ldots,s_N) \in S_{?i} := S_1
\times \cdots \times S_{i-1} \times S_{i+1} \times \cdots \times S_N$
est un profil de stratégies mixtes pour tous les joueurs autres que le
joueur $i$, on dit que la stratégie mixte $s_! \in S_i$ est une
\defin{meilleure réponse} (resp. la meilleure réponse
\textbf{stricte}) contre $s_?$ lorsque pour tout $t \in S_i$ on a
$u_i(s_?,s_!) \geq u_i(s_?,t)$ (resp. lorsque pour tout $t \in S_i$
différent de $s_!$ on a $u_i(s_?,s_!) > u_i(s_?,t)$), où $(s_?,t)$
désigne l'élément de $S_1\times \cdots \times S_N$ obtenu en insérant
$t \in S_i$ comme $i$-ième composante entre $s_{i-1}$ et $s_{i+1}$,
c'est-à-dire le gain [espéré] obtenu en jouant $t$ contre $s_?$.

Un profil de stratégies mixtes $s = (s_1,\ldots,s_N)$ (pour l'ensemble
des joueurs) est dit être un \index{Nash (équilibre
  de)}\defin{équilibre de Nash} (resp., un équilibre de Nash
\defin[strict (équilibre de Nash)]{strict}) lorsque pour tout $1\leq i
\leq N$, la stratégie $s_i$ pour le joueur $i$ est une meilleure
réponse (resp. la meilleure réponse stricte) contre le profil $s_{?i}
:= (s_1,\ldots,s_{i-1},s_{i+1},\ldots,s_N)$ pour les autres joueurs
obtenu en supprimant la composante $s_i$ de $s$.
\end{defn}

Concrètement, donc, un équilibre de Nash est donc un profil de
stratégies mixtes de l'ensemble des joueurs dans lequel \emph{aucun
  joueur n'a intérêt à changer unilatéralement sa stratégie} (au sens
où faire un tel changement lui apporterait une espérance de gain
strictement supérieure).  Un équilibre de Nash \emph{strict}
correspond à la situation où tout changement unilatéral de stratégie
d'un joueur lui apporterait une espérance de gain strictement
inférieure.

\begin{prop}\label{stupid-remark-best-mixed-strategies}
Donné un jeu en forme normale comme
en \ref{definition-game-in-normal-form}, si $1 \leq i \leq N$ et si
$s_?$ est un profil de stratégies mixtes pour tous les joueurs autres
que le joueur $i$, il existe une meilleure réponse pour le joueur $i$
qui est une stratégie pure.  De plus, si $s_!$ (stratégie mixte) est
une meilleure réponse contre $s_?$ si et seulement si \emph{chaque}
stratégie pure appartenant au support de $s_!$ est une meilleure
réponse possible contre $s_?$ ; et elles apportent toutes le même
gain.

En particulier, une meilleure réponse stricte est nécessairement une
stratégie pure.
\end{prop}
\begin{proof}
Tout ceci résulte essentiellement des propositions
\ref{affine-functions-take-extrema-at-boundary} et \ref{affine-functions-take-no-strict-extrema-inside} :
le gain espéré $u_i(s_?,t)$ du joueur $i$ (une fois fixé le
profil $s_?$ de tous les autres joueurs) est une fonction affine de la
stratégie $t \in S_i$ du joueur $i$, et une fonction affine sur un
simplexe prend son maximum sur un sommet du simplexe, et ne peut la
prendre à l'intérieur d'une face (= enveloppe convexe d'un
sous-ensemble des sommets) que si elle la prend également en chaque
sommet de cette face.

Plus précisément, soit $v = \max_{t \in S_i} u_i(s_?,t)$, dont la
proposition \ref{affine-functions-take-extrema-at-boundary} garantit
l'existence ; une meilleure réponse possible contre $s_?$ est
précisément un $t$ réalisant $u_i(s_?,t) = v$ : il résulte de la
proposition citée que $v$ est aussi $\max_{a \in A_i} u_i(s_?,a)$,
c'est-à-dire qu'il existe une stratégie \emph{pure} $a \in A_i$ qui
est une meilleure réponse possible contre $s_?$.  Par ailleurs, si
$s_!$ est une meilleure réponse, i.e., $u_i(s_?,s_!) = v$, alors la
proposition \ref{affine-functions-take-no-strict-extrema-inside}
(appliquée avec pour $x_1,\ldots,x_n$ les points du support de $s_!$)
montre que chaque stratégie pure $a$ appartenant au support de $s_!$
vérifie aussi $u_i(s_?,a) = v$.

Comme une meilleure réponse stricte est unique, elle est forcément
pure d'après ce qu'on vient de voir.
\end{proof}

\textbf{Reformulation.}  Un profil de stratégies mixtes
$(s_1,\ldots,s_N) \in S_1\times \cdots\times S_N$ est un équilibre de
Nash si et seulement si, pour chaque $i$,
\begin{itemize}
\item pour tout $a$ dans le support de $s_i$, le gain espéré
  $u_i(s_{?i},a)$ (rappelons que ceci est une notation pour
  $u_i(s_1,\ldots,s_{i-1},a,s_{i+1},\ldots,s_N$) prend la \emph{même}
  valeur, et
\item cette valeur est supérieure ou égale à $u_i(s_{?i},b)$ pour
  tout $b\in A_i$.
\end{itemize}

\smallskip

On vient de voir que \emph{lorsque les stratégies de tous les autres
  joueurs sont fixées} (à $s_?$ dans la proposition ci-dessus), le
joueur $i$ a une meilleure réponse en stratégie pure : on pourrait
être tenté d'en conclure, mais ce serait une erreur, que les
stratégies mixtes n'apportent donc rien à l'histoire, et qu'un
équilibre de Nash existe forcément en stratégies pures : \emph{ce
  n'est pas le cas}.  Néanmoins, les équilibres de Nash existent bien
en stratégies mixtes, et c'est le résultat central du domaine (ayant
essentiellement valu à son auteur le prix d'économie de la banque de
Suède à la mémoire d'Alfred Nobel en 1994) :

\begin{thm}[John Nash, 1951]\label{theorem-nash-equilibria}
Pour un jeu en forme normale comme
en \ref{definition-game-in-normal-form}, il existe un équilibre de
Nash.
\end{thm}

Pour démontrer le théorème en question, on utilise (et on admet) le
théorème du point fixe de Brouwer, qui affirme que :

\begin{thm}[L. E. J. Brouwer, 1910]\label{brouwer-fixed-point-theorem}
Si $K$ est un convexe compact de $\mathbb{R}^m$, et que $T \colon K
\to K$ est continue, alors il existe $x\in K$ tel que $T(x) = x$ (un
\emph{point fixe} de $T$, donc).
\end{thm}

L'idée intuitive de la démonstration suivante est : partant d'un
profil $s$ de stratégies, on peut définir continûment un nouveau
profil $s^\sharp$ en donnant plus de poids aux options qui donnent un
meilleur gain au joueur correspondant — si bien que $s^\sharp$ sera
différent de $s$ dès que $s^\sharp$ n'est pas un équilibre de Nash.  Comme
la fonction $T \colon s \mapsto s^\sharp$ doit avoir un point fixe, ce point
fixe sera un équilibre de Nash.

\begin{proof}[Démonstration de \ref{theorem-nash-equilibria}]
Si $s \in S$ et $1\leq i\leq N$, convenons de noter $s_{?i}$
l'effacement de la composante $s_i$ (c'est-à-dire le profil pour les
joueurs autres que $i$).  Si de plus $b \in A_i$, notons
$\varphi_{i,b}(s) = \max(0,\; u_i(s_{?i},b) - u_i(s))$ l'augmentation
du gain du joueur $i$ si on remplace sa stratégie $s_i$ par la
stratégie pure $b$ en laissant le profil $s_{?i}$ des autres joueurs
inchangé (ou bien $0$ s'il n'y a pas d'augmentation).  On remarquera
que $s$ est un équilibre de Nash si et seulement si les
$\varphi_{i,b}(s)$ sont nuls pour tout $1\leq i\leq N$ et tout $b\in
A_i$ (faire appel à la proposition précédente pour le « si »).  On
remarquera aussi que chaque $\varphi_{i,b}$ est une fonction continue
sur $S$.

Définissons maintenant $T\colon S\to S$ de la façon suivante : si $s
\in S$, on pose $T(s) = s^\sharp$, où $s^\sharp =
(s^\sharp_1,\ldots,s^\sharp_N)$ avec $s^\sharp_i$ le barycentre de
$s_i$ avec coefficient $1$ et des $a \in A_i$ avec les coefficients
$\varphi_{i,a}(s)$, autrement dit :
\[
\begin{aligned}
s^\sharp_i(a) &= \frac{s_i(a) + \varphi_{i,a}(s)}{\sum_{b\in A_i}(s_i(b) + \varphi_{i,b}(s))}\\
&= \frac{s_i(a) + \varphi_{i,a}(s)}{1 + \sum_{b\in A_i}\varphi_{i,b}(s)}
\end{aligned}
\]
(L'important est que $s^\sharp_i$ augmente strictement le poids des options
$a\in A_i$ telles que $u_i(s_{?i},a) > u_i(s)$ ; en fait, on pourrait
composer $\varphi$ à gauche par n'importe quelle fonction $\mathbb{R}
\to \mathbb{R}$ croissante, continue, nulle sur les négatifs et
strictement positive sur les réels strictement positifs, on a choisi
l'identité ci-dessus pour rendre l'expression plus simple à écrire,
mais elle peut donner l'impression qu'on commet une « erreur
  d'homogénéité » en ajoutant un gain à une probabilité.)

D'après la première expression donnée, il est clair qu'on a bien
$s^\sharp_i \in S_i$, et qu'on a donc bien défini une fonction
$T\colon S\to S$.  Cette fonction est continue (comme composée de
fonctions continues), donc admet un point fixe $s$
d'après \ref{brouwer-fixed-point-theorem} (notons qu'ici, $S$ est vu
comme le convexe $S_1\times\cdots\times S_N$ dans
$\mathbb{R}^{A_1}\times\cdots\times \mathbb{R}^{A_N}$).  On va montrer
que $s$ est un équilibre de Nash.

Si $1\leq i\leq N$, il existe $a$ dans le support de $s_i$
tel que $u_i(s_{?i},a) \leq u_i(s)$ (car, comme dans la preuve
de \ref{stupid-remark-best-mixed-strategies}, $u_i(s)$ est combinaison
convexe des $u_i(s_{?i},a)$ dont est supérieur au plus petit d'entre
eux) : c'est-à-dire $\varphi_{i,a}(s) = 0$.  Pour un tel $a$, la
seconde expression $s^\sharp_i(a) = s_i(a) / \big(1 + \sum_{b\in
  A_i}\varphi_{i,b}(s)\big)$ montre, en tenant compte du fait que
$s^\sharp_i = s_i$ puisque $s$ est un point fixe, que $\sum_{b\in A_i}
\varphi_{i,b}(s) = 0$, donc $\varphi_{i,b}(s) = 0$ pour tout $b$.  On
vient de voir que les $\varphi_{i,b}(s)$ sont nuls pour tout $i$ et
tout $b$, et on a expliqué que ceci signifie que $s$ est un équilibre
de Nash.
\end{proof}

\thingy La question du calcul \emph{algorithmique} des équilibres de
Nash, ou simplement d'\emph{un} équilibre de Nash, est délicate en
général, ou même si on se restreint à $N=2$ joueurs (on verra
en \ref{zero-sum-games-by-linear-programming-algorithm} ci-dessous que
dans le cas particulier des jeux à somme nulle elle se simplifie
considérablement).  Il existe un algorithme pour $N=2$, appelé
\index{Lemke-Howson (algorithme de)}algorithme de Lemke-Howson, qui
généralise l'algorithme du simplexe pour la programmation linéaire
pour trouver un équilibre de Nash d'un jeu en forme normale : comme
l'algorithme du simplexe, il est exponentiel dans le pire cas, mais se
comporte souvent bien « en pratique ».

Pour $N=2$, une méthode qui peut fonctionner dans des cas suffisamment
petits consiste à énumérer tous les supports
(cf. \ref{definition-mixed-strategy-abst}) possibles des stratégies
mixtes des joueurs dans un équilibre de Nash, c'est-à-dire toutes les
$(2^{\#A_1}-1)\penalty0\times\penalty0(2^{\#A_2}-1)$ données de
parties non vides de $A_1$ et $A_2$, et, pour chacune, appliquer le
raisonnement suivant : si $s_i$ est une meilleure réponse possible
pour le joueur $i$ (contre la stratégie $s_{?i}$ de l'autre joueur)
alors \emph{toutes les options du support de $s_i$ ont la même
  espérance de gain} (contre $s_{?i}$ ;
cf. \ref{stupid-remark-best-mixed-strategies}), ce qui fournit un jeu
d'égalités linéaires sur les valeurs de $s_{?i}$.  En rassemblant ces
inégalités (ainsi que celles qui affirment que la somme des valeurs de
$s_i$ et de $s_{?i}$ valent $1$), on arrive « normalement » à trouver
tous les équilibres de Nash possibles : voir \ref{dove-or-hawk}
ci-dessous, ainsi que les exercices
\ref{normal-form-game-exercise-two-by-two} et \ref{normal-form-game-exercise-three-by-three}
(dernières questions) pour des exemples.

(Pour $N\geq 3$, on sera amené en général à résoudre des systèmes
d'équations algébriques pour chaque combinaison de supports : ceci est
algorithmiquement possible en théorie en vertu d'un théorème de Tarski
et Seidenberg sur la décidabilité des systèmes d'équations algébriques
réels, mais possiblement inextricable dans la pratique.)

\thingy\label{dove-or-hawk-redux} Pour donner un exemple, revenons sur
le \index{trouillard (jeu du)}jeu du trouillard défini
en \ref{dove-or-hawk} et dressons un tableau de l'espérance de gain
pour quelques stratégies mixtes ainsi que la stratégie mixte
« générique » :

{\footnotesize
\begin{center}
\begin{tabular}{r|c|c|c|c|c|}
$\downarrow$Alice, Bob$\rightarrow$&Colombe&Faucon&$\frac{2}{3}C+\frac{1}{3}F$&$\frac{1}{2}C+\frac{1}{2}F$&$qC+(1-q)F$\\\hline
Colombe&$2,2$&\textcolor{blue}{$0,4$}&$\frac{4}{3},\frac{8}{3}$&$1,3$&{\tiny $2q, 4-2q$}\\\hline
Faucon&\textcolor{blue}{$4,0$}&$-4,-4$&$\frac{4}{3},-\frac{4}{3}$&$0,-2$&{\tiny $-4+8q, -4+4q$}\\\hline
$\frac{2}{3}C+\frac{1}{3}F$&$\frac{8}{3},\frac{4}{3}$&$-\frac{4}{3},\frac{4}{3}$&\textcolor{blue}{$\frac{4}{3},\frac{4}{3}$}&$\frac{2}{3},\frac{4}{3}$&{\tiny $-\frac{4}{3}+4q, \frac{4}{3}$}\\\hline
$\frac{1}{2}C+\frac{1}{2}F$&$3,1$&$-2,0$&$\frac{4}{3},\frac{2}{3}$&$\frac{1}{2},\frac{1}{2}$&$-2+5q,q$\\\hline
$pC+(1-p)F$&{\tiny $4-2p,2p$}&{\tiny $-4+2p,-4+8p$}&{\tiny $\frac{4}{3},-\frac{4}{3}+4p$}&{\tiny $p,-2+5p$}&$(\dagger)$\\\hline
\end{tabular}
\end{center}
\par\noindent
où $(\dagger) = (-4 + 4p + 8q - 6pq, - 4 + 8p + 4q - 6pq)$.
\par}

\smallskip

Les trois profils $(C,F)$ (i.e., Alice joue Colombe et Bob joue
Faucon), $(F,C)$ et $(\frac{2}{3}C+\frac{1}{3}F,
\frac{2}{3}C+\frac{1}{3}F)$ sont des équilibres de Nash : cela résulte
du fait que, quand on regarde la case correspondante du tableau, le
premier nombre est maximum sur la colonne (c'est-à-dire qu'Alice n'a
pas intérêt à changer unilatéralement sa stratégie) et le second est
maximum sur la ligne (c'est-à-dire que Bob n'a pas intérêt à changer
unilatéralement sa stratégie) : chacun de ces maxima peut se tester
simplement contre les stratégies pures de l'adversaire (i.e., les deux
premières lignes ou colonnes).  Pour trouver ces équilibres et
vérifier qu'il n'y en a pas d'autre, on peut :
\begin{itemize}
\item commencer par rechercher les équilibres de Nash où chacun des
  deux joueurs joue une stratégie pure (ceci revient à regarder
  chacune des quatre cases du tableau initial et chercher si le
  premier nombre est maximal sur la colonne et le second maximal sur
  la ligne) ;
\item considérer la possibilité d'un équilibre de Nash où l'un des
  joueurs joue une stratégie pure et l'autre une stratégie mixte
  supportée par les deux options : si par exemple Alice joue
  $pC+(1-p)F$ avec $p>0$ et Bob joue $a \in \{C,F\}$,
  d'après \ref{stupid-remark-best-mixed-strategies}, le gain d'Alice
  dans les cases $(C,a)$ et $(F,a)$ doit être le même, ce qui n'est
  pas le cas, donc il n'y a pas de tel équilibre (pas plus que dans le
  cas correspondant pour Bob) ;
\item enfin, rechercher les équilibres de Nash où chacun des deux
  joueurs joue une stratégie supportée par les deux options, soit ici
  $pC+(1-p)F$ pour Alice et $qC+(1-q)F$ pour Bob avec $0<p<1$ et
  $0<q<1$ : d'après \ref{stupid-remark-best-mixed-strategies} le gain
  espéré d'Alice doit être le même contre $qC+(1-q)F$ indifféremment
  qu'elle joue Colombe ou Faucon, ce qui implique $2q = -4+8q$ donc
  $q=\frac{2}{3}$, et symétriquement, le gain espéré de Bob doit être
  le même contre $pC+(1-p)F$ indifféremment qu'il joue Colombe ou
  Faucon, ce qui implique $2p = -4+8p$ donc $p=\frac{2}{3}$, si bien
  que le seul équilibre de Nash de cette nature est celui qu'on a
  décrit (et il faut vérifier qu'il en est bien un).
\end{itemize}

\thingy Mentionnons en complément une notion plus générale que celle
d'équilibre de Nash : si $s\colon A \to \mathbb{R}$ (où $A :=
A_1\times\cdots\times A_N$) est cette fois une distribution de
probabilités sur l'ensemble $A$ des profils de stratégies pures (le
rapport avec l'ensemble des profils de stratégies mixtes est explicité
en \ref{remark-mixed-stragy-profile-versus-correlated-profile}), on
dit que $s$ est un \index{corrélé (équilibre)}\defin{équilibre
  corrélé} lorsque pour tout $1 \leq i \leq N$ et pour tous $b,b' \in
A_i$ on a
\[
\sum_{a\,:\, a_i = b} s(a)\,(u_i(a) - u_i(a_{?i},b'))\geq 0
\]
où la somme est prise sur les $a \in A$ tels que $a_i = b$ et où
$u_i(a_{?i},b')$ désigne bien sûr la valeur de $u_i$ en l'élément de
$A$ égal à $a$ sauf que la $i$-ième coordonnée (qui vaut $b$) a été
remplacée par $b'$.

De façon plus intuitive, il faut imaginer qu'un « corrélateur » tire
au hasard un profil $a$ de stratégies pures selon la distribution $s$,
et la condition d'équilibre indiquée ci-dessus signifie que si chaque
joueur $i$ reçoit l'information ($b = a_i$) de l'option qui a été
tirée pour lui, tant que les autres joueurs suivent les instructions
($a_{?i}$) du corrélateur, il n'a pas intérêt à choisir une autre
option ($b'$) que celle qui lui est proposée.

Pour dire les choses autrement, faisons les définitions suivantes.
Lorsque $s_{?}$ une distribution de probabilités sur $A_{?i} :=
A_1\times \cdots \penalty500 \times A_{i-1} \penalty0 \times \penalty0
A_{i+1} \times\cdots \penalty500 \times A_N$ et $b \in A_i$, notons
$u_i(s_?, b) := \sum_{a_? \in A_{?i}} s_?(a_?)\,u_i(a_?, b)$ le gain
espéré du joueur $i$ lorsque l'ensemble des autres joueurs joue un
profil tiré selon $s_?$ et que $i$ joue $b$.  Lorsque $s$ est une
distribution de probabilités sur $A$, appelons $s[b]$ (pour $b\in A_i$
tel que $\sum_{a'\,:\, a'_i=b} s(a') > 0$) la distribution de
probabilités $s$ conditionnée à $a_i = b$ et projetée à $A_{?i}$,
c'est-à-dire concrètement la distribution qui à $a_? \in A_{?i}$
associe $s(a_?, b)/\sum_{a'_? \in A_{?i}} s(a'_?, b)$.  La condition
que $s$ soit un équilibre corrélé se réécrit alors en : $u_i(s[b], b)
\geq u_i(s[b], b')$ pour tout $b\in A_i$ tel que $\sum_{a'\,:\,
  a'_i=b} s(a') > 0$ et tout $b'\in A_i$.

Dans le cas particulier où $s$ est une distribution aux marginales
indépendantes, c'est-à-dire de la forme $s(a) = s_1(a_1) \cdots
s_N(a_N)$
(cf. \ref{remark-mixed-stragy-profile-versus-correlated-profile}), ce
qu'on a noté $s[b]$ ci-dessus est précisément la fonction qui à $a_?
\in A_{?i}$ associe le produit $s_{?i}$ des $s_j(a_j)$ pour $j\neq i$,
et la condition qu'on vient de dire est donc $u_i(s_{?i}, b) \geq
u_i(s_{?i}, b')$ pour tout $b$ dans le support de $s_i$ et tout $b'$.
D'après \ref{stupid-remark-best-mixed-strategies}, c'est justement
dire que $(s_1,\ldots,s_N)$ est un équilibre de Nash.  Autrement dit :
\emph{un équilibre de Nash est la même chose qu'un équilibre corrélé
  dans lequel les marginales se trouvent être indépendantes}.



\subsection{Jeux à somme nulle : le théorème du minimax}\label{zero-sum-games}

Plaçons nous maintenant dans le cadre d'un jeu à somme nulle,
c'est-à-dire qu'il y a $N=2$ joueurs, et qu'ils ont des gains
opposés : disons qu'on appelle $u = u_1$ le gain du joueur $1$
(« Alice »), le gain du joueur $2$ (« Bob ») étant alors $u_2 = -u$ et
n'a pas besoin d'être écrit dans le tableau.  Ainsi, Alice cherche à
\emph{maximiser} $u$ et Bob cherche à le \emph{minimiser} (puisque
maximiser $-u$ revient à minimiser $u$).

En considérant le jeu du point de vue de sa matrice de gains (où, de
nouveau, on n'a écrit que le gain d'Alice), Alice cherche à choisir
une ligne qui maximiser la valeur écrite dans le tableau et Bob
cherche à choisir une colonne qui mimimise la valeur écrite.  Mais en
général, si $u\colon A_1 \times A_2 \to \mathbb{R}$ est un tableau
fini de nombres réels, $\max_{a \in A_1} \min_{b\in A_2} u(a,b)$ ne
coïncide pas avec $\min_{b\in A_2} \max_{a \in A_1} u(a,b)$ (on a
seulement l'inégalité $\max_{a \in A_1} \min_{b\in A_2} u(a,b) \leq
\min_{b\in A_2} \max_{a \in A_1} u(a,b)$, qui sera démontrée ci-dessus
au début de la démonstration de \ref{theorem-minimax}).  L'observation
cruciale est que, quand on passe des ensembles finis $A_i$ de
stratégies pures aux simplexes $S_i$ de stratégies mixtes, on a alors
une égalité :

\begin{thm}[« du minimax », J. von Neumann, 1928]\label{theorem-minimax}
Soient $A_1,A_2$ deux ensembles finis et $S_1,S_2$ les simplexes de
stratégies mixtes (cf. \ref{definition-mixed-strategy-abst})
sur $A_1,A_2$ ; soit $u\colon A_1\times A_2 \to \mathbb{R}$ une
fonction quelconque, et on notera encore $u$ son unique extension
(explicitée en \ref{definition-mixed-strategy-gain}) en une fonction
$S_1\times S_2 \to \mathbb{R}$ affine en chaque variable.  On a
alors :
\[
\max_{x\in S_1} \min_{y\in S_2} u(x,y) =
\min_{y\in S_2} \max_{x\in S_1} u(x,y)
\]
(Ceci sous-entend notamment ces $\min$ et $\max$ existent bien.)

Plus généralement, si $S_1$ et $S_2$ sont deux convexes compacts dans
des espaces affines réels de dimension finie, et $u\colon S_1\times
S_2 \to \mathbb{R}$ une application affine en chaque variable
séparément, alors on a l'égalité ci-dessus.
\end{thm}
\begin{proof}
Plaçons-nous d'abord dans le cas (qui intéresse la théorie des jeux)
où $A_1,A_2$ sont deux ensembles finis et $S_1,S_2$ les simplexes de
stratégies mixtes.

Tout d'abord, montrons l'existence des $\min$ et $\max$ annoncés ainsi
que l'inégalité dans le sens $\max\min \leq \min\max$.  On a vu
en \ref{affine-functions-take-extrema-at-boundary} que $\min_{y\in
  S_2} u(x,y)$ existe et vaut $\min_{b\in A_2} u(x,b)$ pour tout $x\in
S_1$ : ceci montre que la fonction $x \mapsto \min_{y\in S_2} u(x,y)$
s'écrit encore $x \mapsto \min_{b\in A_2} u(x,b)$, donc il s'agit du
minimum d'un nombre \emph{fini} de fonctions continues (affines)
sur $S_1$, donc elle est encore continue et atteint ses bornes sur
l'ensemble compact $S_1$ : il existe donc $x_* \in S_1$ où cette
fonction atteint son maximum.  Soit de même $y_* \in S_2$ un point où
$y \mapsto \max_{x\in S_1} u(x,y)$ atteint son minimum.  On a alors :
\[
\max_{x\in S_1} \min_{y\in S_2} u(x,y)
= \min_{y\in S_2} u(x_*,y)
\leq u(x_*,y_*) \leq \max_{x\in S_1} u(x,y_*) =
\min_{y\in S_2} \max_{x\in S_1} u(x,y)
\]
La première relation est la définition de $x_*$, la seconde est la
définition du $\min$, la troisième la définition du $\max$, et la
quatrième est la définition de $y_*$.  (Notons que tout ceci vaut dans
n'importe quel contexte où les $\min$ et $\max$ existent, notamment
ceci prouve aussi $\max_{a \in A_1} \min_{b\in A_2} u(a,b) \leq
\min_{b\in A_2} \max_{a \in A_1} u(a,b)$ affirmé ci-dessus.)

Il s'agit maintenant de prouver l'inégalité de sens contraire.

Considérons $(x_0,y_0)$ un équilibre de Nash du jeu à somme nulle dont
la matrice de gains est donnée par $u$ pour le joueur $1$ (et $-u$
pour le joueur $2$) : on sait qu'un tel équilibre existe
d'après \ref{theorem-nash-equilibria}.  On a alors
\[
\max_{x\in S_1} \min_{y\in S_2} u(x,y)
\geq \min_{y\in S_2} u(x_0,y) = u(x_0,y_0)
= \max_{x\in S_1} u(x,y_0) \geq
\min_{y\in S_2} \max_{x\in S_1} u(x,y)
\]
Ici, la première relation est la définition du $\max$, la seconde
traduit le fait que (dans la définition de l'équilibre de Nash) Bob
fait une meilleure réponse possible à $x_0$, la troisième traduit le
fait qu'Alice fait une meilleure réponse possible à $y_0$, et la
quatrième est la définition du $\min$.

Ceci achève de montrer $\max_{x\in S_1} \min_{y\in S_2} u(x,y) =
\min_{y\in S_2} \max_{x\in S_1} u(x,y)$ lorsque $S_1$ et $S_2$ sont
deux simplexes dans $\mathbb{R}^n$, et $u\colon S_1\times S_2 \to
\mathbb{R}$ une application affine en chaque variable séparément.  Si
maintenant $S_1$ et $S_2$ sont deux polytopes, c'est-à-dire chacun
l'enveloppe d'un nombre fini de points $A_1$ et $A_2$
dans $\mathbb{R}^n$, mais plus forcément des simplexes (c'est-à-dire
qu'on ne suppose plus les points de $A_i$ affinement indépendants), la
même conclusion vaut encore : en effet, si on appelle $S'_i$ le
simplexe construit abstraitement sur $A_i$ (i.e., l'ensemble des
fonctions $s\colon A_i \to \mathbb{R}$ positives de somme $1$), et
$\varpi_i \colon S'_i \to S_i$ la fonction (affine et surjective)
envoyant $s\colon A_i \to \mathbb{R}$ positive de somme $1$ sur la
somme des $s(a)\cdot a$ (dans le $\mathbb{R}^n$ où vivent $A_i$
et $S_i$), alors la fonction $u\colon S_1\times S_2 \to \mathbb{R}$
donnée définit une fonction $u'\colon S'_1\times S'_2 \to \mathbb{R}$,
elle aussi affine en chaque variable, par $u'(x,y) = u(\varpi_1(x),
\varpi_2(y))$, et il est évident que les extrema de $u'$ sur $S'_1
\times S'_2$ coïncident avec ceux de $u$ sur $S_1 \times S_2$, donc la
conclusion qui précède, appliquée à $u'$, donne la conclusion désirée
sur $u$.

Enfin, si $S_1$ et $S_2$ sont seulement supposés être des convexes
compacts dans $\mathbb{R}^n$, on observe que pour tout $\varepsilon>0$
il existe $\Sigma_1$ et $\Sigma_2$ des polytopes contenus dans $S_1$
et $S_2$ respectivement et tels que pour tout $x\in S_1$ on ait
$\min_{y\in S_2} u(x,y) > \min_{y\in\Sigma_2} u(x,y)-\varepsilon$ et
$\max_{x\in S_1} u(x,y) < \max_{x\in\Sigma_1} u(x,y)+\varepsilon$
(explication : il est trivial que pour chaque $x$ il existe un
$\Sigma_2$ vérifiant la condition demandée, le point intéressant est
qu'un unique $\Sigma_2$ peut convenir pour tous les $x$ ; mais pour
chaque $\Sigma_2$ donné, l'ensemble des $x$ pour lesquels il convient
est un ouvert de $S_1$, qui est compact, donc un nombre fini de ces
ouverts recouvrent $S_1$, et on prend l'enveloppe convexe de la
réunion des $\Sigma_2$ en question ; on procède de même pour
$\Sigma_1$).  On a alors $\max_{x\in S_1} \min_{y\in S_2} u(x,y) >
\max_{x\in \Sigma_1} \min_{y\in \Sigma_2} u(x,y) - \varepsilon$ et une
inégalité analogue pour l'autre membre : on en déduit l'inégalité
recherchée à $2\varepsilon$ près, mais comme on peut prendre
$\varepsilon$ arbitrairement petit, on a ce qu'on voulait.
\end{proof}

Le corollaire suivant explicite la situation particulière où, en
outre, le jeu est symétrique entre les deux joueurs (c'est-à-dire que
sa matrice de gains écrite pour un seul joueur est, elle,
\emph{anti}symétrique) :

\begin{cor}\label{symmetric-zero-sum-game}
Soient $A$ un ensemble fini et $S$ le simplexe de stratégies mixtes
sur $A$ ; soit $u\colon A^2 \to \mathbb{R}$ une application
antisymétrique (i.e., $u(b,a) = -u(a,b)$), et on notera encore $u$ son
unique extension affine en chaque variable en une fonction $S^2 \to
\mathbb{R}$ affine en chaque variable.  Alors la valeur commune des
deux membres de l'égalité du théorème \ref{theorem-minimax} est $0$ :
il existe $x\in S$ tel que pour tout $y\in S$ on ait $u(x,y)\geq 0$.

Plus généralement, si $S$ est un convexe compact dans un espace affine
réel de dimension finie, et $u\colon S^2 \to \mathbb{R}$ une
application antisymétrique (i.e., $u(y,x) = -u(x,y)$) et affine en
chaque variable séparément, la même conclusion vaut.
\end{cor}
\begin{proof}
On applique le théorème : il donne $\max_{x\in S}\penalty0 \min_{y\in
  S} u(x,y) = \min_{y\in S}\penalty0 \max_{x\in S} u(x,y)$.  Mais
puisque $u$ est antisymétrique ceci s'écrit encore $\min_{y\in S}
\max_{x\in S} (-u(y,x))$, soit, en renommant les variables liées,
$\min_{x\in S}\penalty0 \max_{y\in S} (-u(x,y)) = -\max_{x\in
  S}\penalty0 \min_{y\in S} u(x,y)$.  Par conséquent, $\max_{x\in
  S}\penalty0 \min_{y\in S} u(x,y) = 0$ (il est son propre opposé), et
en prenant un $x$ qui réalise ce maximum, on a $\min_{y\in S} u(x,y) =
0$, ce qu'on voulait prouver.
\end{proof}

\smallskip

Considérons maintenant les applications du
théorème \ref{theorem-minimax} dans le cadre de la théorie des jeux.
Commençons par définir les concepts suivants :

\begin{defn}\label{definition-zero-sum-game-value-and-optimum}
Pour un jeu à somme nulle défini par une matrice de gains $u \colon
A_1 \times A_2 \to \mathbb{R}$ :
\begin{itemize}
\item La \defin[valeur (d'un jeu à somme nulle)]{valeur} est la valeur
  commune $\max_{x\in S_1} \min_{y\in S_2} u(x,y) = \min_{y\in S_2}
  \max_{x\in S_1} u(x,y)$ où $S_1,S_2$ sont les simplexes de
  stratégies mixtes des joueurs $1$ et $2$ respectivement (et $u(x,y)$
  l'espérance de gain explicitée
  en \ref{definition-mixed-strategy-gain}).
\item Une stratégie mixte du joueur $1$ qui réalise $\max_{x\in S_1}
  \min_{y\in S_2} u(x,y)$ (resp. une stratégie mixte du joueur $2$ qui
  réalise $\min_{y\in S_2} \max_{x\in S_1} u(x,y)$) est dite
  \index{optimale (stratégie)}\defin{stratégie optimale} pour le
  joueur en question.
\end{itemize}
\end{defn}

Le théorème \ref{theorem-minimax} (et, pour le dernier point, le
corollaire \ref{symmetric-zero-sum-game}) a alors les conséquences
suivantes :

\begin{prop}\label{minimax-for-games} Pour un jeu à somme nulle défini
par une matrice de gains $u \colon A_1 \times A_2 \to \mathbb{R}$ :
\begin{itemize}
\item[(i)] Un profil $(x_0,y_0) \in S_1\times S_2$ de stratégies
  mixtes est un équilibre de Nash si et seulement si $x_0$ et $y_0$
  sont chacun une stratégie optimale pour le joueur correspondant.
\item[(ii)] Tous les équilibres de Nash ont la même valeur de gain
  espéré, à savoir la valeur du jeu $v$ (pour le joueur 1).
\item[(iii)] Une stratégie mixte $x_0 \in S_1$ du joueur $1$ est une
  stratégie optimale pour celui-ci si et seulement si $u(x_0,b) \geq
  v$ pour tout $b\in A_2$ (où $v$ désigne la valeur du jeu).  Une
  stratégie mixte $y_0 \in S_2$ du joueur $2$ est une stratégie
  optimale pour celui-ci si et seulement si $u(a,y_0) \leq v$ pour
  tout $a\in A_1$.
\item[(iv)] L'ensemble $T_1$ des stratégies optimales du joueur $1$
  est un convexe, de même que l'ensemble $T_2$ des stratégies
  optimales du joueur $2$.  (L'ensemble des équilibres de Nash est
  donc aussi un convexe, à savoir $T_1 \times T_2$.)
\item[(v)] Si le jeu est symétrique, c'est-à-dire $A_2 = A_1 =: A$ et
  $u(b,a) = -u(a,b)$ (matrice de gains antisymétrique), alors la
  valeur du jeu est nulle.
\end{itemize}
\end{prop}

\begin{proof}
(i) : On a vu au cours de la preuve de \ref{theorem-minimax} que si
  $(x_0,y_0)$ est un équilibre de Nash, alors $\max_{x\in S_1}
  \min_{y\in S_2} u(x,y) \geq \min_{y\in S_2} u(x_0,y) = u(x_0,y_0) =
  \max_{x\in S_1} u(x,y_0) \geq \min_{y\in S_2} \max_{x\in S_1}
  u(x,y)$, mais comme les deux termes extrêmes sont égaux, en fait
  toutes ces inégalités sont des égalités, donc $x_0$ réalise bien le
  maximum $\max_{x\in S_1} \min_{y\in S_2} u(x,y)$ et $y_0$ le minimum
  $\min_{y\in S_2} \max_{x\in S_1} u(x,y)$, i.e., ils sont des
  stratégies optimales.  Réciproquement, si $x_*$ est une stratégie
  optimale du joueur $1$ et $y_*$ du joueur $2$, on a vu au cours de
  la même preuve qu'on a $\max_{x\in S_1} \min_{y\in S_2} u(x,y) =
  \min_{y\in S_2} u(x_*,y) \leq u(x_*,y_*) \leq \max_{x\in S_1}
  u(x,y_*) = \min_{y\in S_2} \max_{x\in S_1} u(x,y)$, et, de nouveau,
  totues ces inégalités sont des égalités, donc $x_*$ est une
  meilleure réponse possible pour le joueur $1$ contre $y_*$ et $y_*$
  est une meilleure réponse possible pour le joueur $2$ contre $x_*$,
  ce qui signifie que $(x_*,y_*)$ est un équilibre de Nash.

Le point (ii) résulte de ce qui vient d'être dit.

(iii) : Par définition, $x_0$ est optimale ssi $\min_{y\in S_2}
u(x_0,y) = v$, ou, comme on ne peut pas faire plus que le maximum,
$\min_{y\in S_2} u(x_0,y) \geq v$.  Mais on a vu
en \ref{affine-functions-take-extrema-at-boundary} que $\min_{y\in
  S_2} u(x_0,y) = \min_{b\in A_2} u(x_0,b)$, c'est-à-dire que $x_0$
est optimale si et seulement si $\min_{b\in A_2} u(x_0,b) \geq v$, ce
qui revient bien à dire $u(x_0,b) \geq v$ pour tout $b$.  Le cas de
l'autre joueur est analogue.

(iv) : Il résulte du point précédent que $T_1$ est l'intersections des
ensembles $\{x \in S_1 : u(x,b)\geq v\}$ où $b$ parcourt $A_2$ : mais
chacun de ces ensembles est convexe, donc leur intersection l'est
aussi.  Le cas de $T_2$ est analogue.  La parenthèse est une redite du
point (i).

Le point (v) n'est qu'une redite du
corollaire \ref{symmetric-zero-sum-game}.
\end{proof}

Répétons : \emph{dans un jeu à somme nulle, un profil de stratégies
  est un équilibre de Nash si et seulement si chaque joueur joue une
  stratégie optimale}, l'ensemble des stratégies optimales étant un
convexe (défini pour chaque joueur indépendamment).

L'ensemble des stratégies optimales d'un joueur donné est « très
souvent » un singleton, ce qui fait qu'on parle parfois abusivement de
« la » stratégie optimale.

Reste maintenant à expliquer comment calculer les stratégies
optimales :

\begin{algo}\label{zero-sum-games-by-linear-programming-algorithm}
Donnée une fonction $u\colon A_1 \times A_2 \to \mathbb{R}$ (avec
$A_1,A_2$ deux ensembles finis) définissant la matrice de gains pour
le joueur $1$ d'un jeu à somme nulle.  On peut calculer une stratégie
mixte optimale (cf. \ref{minimax-for-games}) pour le joueur $1$ en
résolvant, par exemple au moyen de l'algorithme du simplexe, le
problème de programmation linéaire dont les variables sont les $x_a$
pour $a \in A_1$ (les poids de la stratégie mixte) et $v$ (le gain
obtenu) cherchant à maximiser $v$ sujet aux contraintes :
\[(\forall a\in A_1)\;x_a \geq 0\]
\[\sum_{a\in A_1} x_a = 1\]
\[(\forall b\in A_2)\;v \leq u(x,b) := \sum_{a \in A_1} u(a,b)\, x_a\]
Autrement dit, il s'agit d'un problème de programmation linéaire à
$\#A_1 + 1$ variables avec des contraintes de positivité sur $\#A_1$
d'entre elles, une contrainte d'égalité et $\#A_2$ inégalités affines.
\end{algo}

(La correction de cet algorithme découle à peu près immédiatement du
point (iii) ci-dessus.)

\thingy Pour ramener ce problème à un problème de programmation
linéaire en \emph{forme normale} (maximiser $\textbf{p} x$ sous les
contraintes $\textbf{M} x \leq \textbf{q}$ et $x\geq 0$), on sépare la
variable $v$ en $v_+ - v_-$ avec $v_+,v_- \geq 0$, et le problème
devient de maximiser $v_+ - v_-$ sous les contraintes
\[v_+\geq 0,\; v_- \geq 0,\;\; (\forall a\in A_1)\;x_a \geq 0\]
\[\sum_{a\in A_1} x_a \leq 1\]
\[-\sum_{a\in A_1} x_a \leq -1\]
\[(\forall b\in A_2)\;v_+ - v_- - \sum_{a \in A_1} u(a,b)\, x_a \leq 0\]
Le problème dual (minimiser ${^{\mathrm{t}}\textbf{q}} y$ sous les
contraintes ${^{\mathrm{t}}\textbf{M}} y \geq {^\mathrm{t}\textbf{p}}$
et $y\geq 0$) est alors de minimiser $w_+ - w_-$ sous les contraintes
\[w_+\geq 0,\; w_- \geq 0,\;\; (\forall b\in A_2)\;y_b \geq 0\]
\[\sum_{b\in A_2} y_b \geq 1\]
\[-\sum_{b\in A_2} y_b \geq -1\]
\[(\forall a\in A_1)\;w_+ - w_- - \sum_{b \in A_2} u(a,b)\, y_b \geq 0\]
Il s'agit donc exactement du même problème, mais pour l'autre joueur.

Le théorème \ref{theorem-minimax} est essentiellement équivalent au
théorème de dualité pour la programmation linéaire (qui assure que si
le problème primal a un optimum $x_0$ alors le dual en a un $y_0$, et
on a égalité des optima).

Comme l'algorithme du simplexe résout simultanément le problème primal
et le problème dual, l'algorithme ci-dessus (exécuté avec l'algorithme
du simplexe) trouve simultanément une stratégie optimale pour les deux
joueurs.


%
%
%

\section{Jeux de Gale-Stewart et détermination}\label{section-gale-stewart-games}

\subsection{Définitions}

\begin{defn}\label{definition-gale-stewart-game}
Soit $X$ un ensemble non vide quelconque (à titre indicatif, les cas
$X = \{0,1\}$ et $X = \mathbb{N}$ seront particulièrement
intéressants).  Soit $A$ un sous-ensemble de $X^{\mathbb{N}}$.  Le
\defin[Gale-Stewart (jeu de)]{jeu de Gale-Stewart} $G_X(A)$ (ou
$G_X^{\mathrm{a}}(A)$, cf. \ref{remark-player-names}) est défini de la
manière suivante : Alice et Bob choisissent tour à tour un élément de
$X$ (autrement dit, Alice choisit $x_0 \in X$ puis Bob choisit $x_1
\in X$ puis Alice choisit $x_2 \in X$ et ainsi de suite).  Ils jouent
un nombre infini de tours, « à la fin » desquels la suite
$(x_0,x_1,x_2,\ldots)$ de leurs coups définit un élément de
$X^{\mathbb{N}}$ : si cet élément appartient à $A$, Alice
\defin[gain]{gagne}, sinon c'est Bob (la partie n'est jamais nulle).

Dans ce contexte, les suites finies $\underline{x} =
(x_0,\ldots,x_{\ell-1})$ d'éléments de $X$ s'appellent les
\defin[position]{positions} (y compris la suite vide $()$, qui peut
s'appeler position initiale) de $G_X(A)$, ou de $G_X$ vu que $A$
n'intervient pas ici ; leur ensemble $\bigcup_{\ell=0}^{+\infty}
X^\ell$ s'appelle parfois l'\defin[arbre d'un jeu]{arbre} du
jeu $G_X$ ; l'entier $\ell$ (tel que $\underline{x} \in X^\ell$)
s'appelle la \textbf{longueur} de $\underline{x} =
(x_0,\ldots,x_{\ell-1})$.  Une
\defin[partie|see{confrontation}]{partie} ou
\defin{confrontation}\footnote{Le mot « partie » peut malheureusement
  désigner soit un sous-ensemble soit une partie d'un jeu au sens
  défini ici : le mot « confrontation » permet d'éviter l'ambiguïté.}
de $G_X$ est une suite infinie $\dblunderline{x} =
(x_0,x_1,x_2,\ldots) \in X^{\mathbb{N}}$.
\end{defn}

\thingy\label{remark-player-names} Il peut arriver qu'on ait envie de
faire commencer la partie à Bob.  Il va de soi que ceci ne pose aucune
difficulté, il faudra juste le signaler le cas échéant.

De façon générale, sauf précision expresse du contraire, « Alice » est
le joueur qui cherche à jouer dans l'ensemble $A$ tandis que « Bob »
est celui qui cherche à jouer dans son complémentaire $B :=
X^{\mathbb{N}} \setminus A$.  Le « premier joueur » est celui qui
choisit les termes pairs de la suite, le « second joueur » est celui
qui choisit les termes impairs.

On pourra noter $G_X^{\mathrm{a}}(A)$ lorsqu'il est souhaitable
d'insister sur le fait qu'Alice joue en premier, et
$G_X^{\mathrm{b}}(A)$ lorsqu'on veut indiquer que Bob joue en
premier : formellement, le jeu $G_X^{\mathrm{b}}(A)$ est le même que
$G_X^{\mathrm{a}}(X^{\mathbb{N}}\setminus A)$ si ce n'est que les noms
des joueurs sont échangés.  On utilisera parfois $G_X(A)$ pour
désigner indifféremment $G_X^{\mathrm{a}}(A)$ ou $G_X^{\mathrm{b}}(A)$
(avec une tournure comme « quel que soit le joueur qui commence »).

Donnée une position $\underline{x}$ de longueur $\ell$, le joueur
« qui doit jouer » dans $\underline{x}$ désigne le premier joueur si
$\ell$ est paire, et le second joueur si $\ell$ est impaire.

\begin{defn}\label{definition-strategies-for-gale-stewart-games}
Pour un jeu $G_X$ comme en \ref{definition-gale-stewart-game}, une
\defin{stratégie} pour le premier joueur (resp. le second joueur) est
une fonction $\varsigma$ qui à une suite finie (=position) de longueur
paire (resp. impaire) d'éléments de $X$ associe un élément de $X$,
autrement dit une fonction $\big(\bigcup_{\ell=0}^{+\infty}
X^{2\ell}\big) \to X$ (resp. $\big(\bigcup_{\ell=0}^{+\infty}
X^{2\ell+1}\big) \to X$).

Lorsque dans une confrontation $x_0,x_1,x_2,\ldots$ de $G_X$ on a
$\varsigma((x_0,\ldots,x_{i-1})) = x_i$ pour chaque $i$ pair (y
compris $\varsigma(()) = x_0$ en notant $()$ la suite vide), on dit
que le premier joueur a joué la confrontation \textbf{selon} la
stratégie $\varsigma$ ; de même, lorsque $\tau((x_0,\ldots,x_{i-1})) =
x_i$ pour chaque $i$ impair, on dit que le second joueur a joué la
confrontation selon la stratégie $\tau$.

Si $\varsigma$ et $\tau$ sont deux stratégies pour le premier et le
second joueurs respectivement, on définit $\varsigma \ast \tau$ comme
l'unique confrontation dans laquelle premier joueur joue selon
$\varsigma$ et le second selon $\tau$ : autrement dit, $x_i$ est
défini par $\varsigma((x_0,\ldots,x_{i-1}))$ si $i$ est pair ou
$\tau((x_0,\ldots,x_{i-1}))$ si $i$ est impair.

Si on se donne une partie $A$ de $X^{\mathbb{N}}$ et qu'on convient
qu'Alice joue en premier : la stratégie $\varsigma$ pour Alice est
dite \index{stratégie gagnante}\defin[gagnante (stratégie)]{gagnante}
(dans $G_X^{\mathrm{a}}(A)$) lorsque Alice gagne toute confrontation
où elle joue selon $\varsigma$ comme premier joueur, et la stratégie
$\tau$ pour Bob est dite gagnante lorsque Bob gagne toute
confrontation où il joue selon $\tau$.  Lorsque l'un ou l'autre joueur
a une stratégie gagnante, le jeu $G_X^{\mathrm{a}}(A)$ est dit
\defin[déterminé (jeu)]{déterminé}.
\end{defn}

\thingy Il est clair que les deux joueurs ne peuvent pas avoir
simultanément une stratégie gagnante (il suffit de considérer la suite
$\varsigma \ast \tau$ où $\varsigma$ et $\tau$ seraient des stratégies
gagnantes pour les deux joueurs : elle devrait simultanément
appartenir et ne pas appartenir à $A$).

En revanche, il faut se garder de croire que les jeux $G_X(A)$ sont
toujours déterminés (on verra cependant des résultats positifs
dans \ref{section-determinacy-of-gale-stewart-games}).

\thingy\label{unshifting-notation} Introduisons la notation suivante :
si $\underline{x} := (x_0,\ldots,x_{\ell-1})$ est une suite finie
d'éléments de $X$ et si $A$ est un sous-ensemble de $X^{\mathbb{N}}$,
on notera $\underline{x}^\$ A$ l'ensemble des suites $(x_\ell,
x_{\ell+1}, \ldots) \in X^{\mathbb{N}}$ telles que
$(x_0,\ldots,x_{\ell-1}, x_\ell, \ldots)$ appartienne à $A$.
Autrement dit, il s'agit de l'image réciproque de $A$ par
l'application $X^{\mathbb{N}} \to X^{\mathbb{N}}$ qui insère
$x_0,\ldots,x_{\ell-1}$ au début de la suite.

On utilisera notamment cette notation pour une suite à un seul terme :
si $x\in X$ alors $x^\$ A$ est l'ensemble des $(x_1,x_2,x_3,\ldots)
\in X^{\mathbb{N}}$ telles que $(x,x_1,x_2,\ldots) \in A$.  (Ainsi, si
$\underline{x} := (x_0,\ldots,x_{\ell-1})$, on a $\underline{x}^\$ A =
x_{\ell-1}^{\$} \cdots x_1^{\$} x_0^{\$} A$.)

\thingy\label{gale-stewart-positions-as-games} Si $\underline{x} :=
(x_0,\ldots,x_{\ell-1})$ est une position dans un jeu de Gale-Stewart
$G_X(A)$, on peut considérer qu'elle définit un nouveau jeu
$G_X(\underline{x}, A)$ consistant à jouer \defin{à partir de là},
c'est-à-dire que les joueurs jouent pour compléter cette suite, ou, ce
qui revient au même, que leurs $\ell$ premiers coups sont imposés.

Plus précisément, on introduit le jeu $G_X^{\mathrm{a}}(\underline{x},
A)$ (resp. $G_X^{\mathrm{b}}(\underline{x}, A)$), qui est la
modification du jeu $G_X^{\mathrm{a}}(A)$
(resp. $G_X^{\mathrm{b}}(A)$) où les $\ell$ premiers coups sont
imposés par les valeurs de $\underline{x}$ (autrement dit, lorsque
$i<\ell$, au $i$-ième coup, seule la valeur $x_i$ peut être jouée).
On identifiera parfois abusivement la position $\underline{x}$ du jeu
$G_X(A)$ avec le jeu $G_X(\underline{x}, A)$ joué à partir de cette
position.

Bien sûr, plutôt qu'\emph{imposer} les $\ell$ premiers cours, on peut
aussi considérer que les joueurs jouent directement à partir du coup
numéroté $\ell$ (donc pour compléter $\underline{x}$) : les joueurs
choisissent $x_{\ell},x_{\ell+1},x_{\ell+2},\ldots$, on insère les
coups imposés $x_0,\ldots,x_{\ell-1}$ au début de la confrontation, et
on regarde si la suite tout entière appartient à $A$ pour déterminer
le gagnant.  Quitte à renuméroter les coups effectivement choisis pour
commencer à $0$, on retrouve un jeu de Gale-Stewart, défini comme suit
en utilisant la notation \ref{unshifting-notation} : à savoir, le jeu
$G_X^{\mathrm{a}}(\underline{x}, A)$ peut s'identifier au jeu
$G_X^{\mathrm{a}}(\underline{x}^\$ A)$ lorsque $\ell$ est pair (=c'est
à Alice de jouer dans la position $\underline{x}$), et
$G_X^{\mathrm{b}}(\underline{x}^\$ A)$ lorsque $\ell$ est impair
(=c'est à Bob de jouer) ; symétriquement, bien sûr,
$G_X^{\mathrm{b}}(\underline{x}, A)$ peut s'identifier au jeu
$G_X^{\mathrm{b}}(\underline{x}^\$ A)$ lorsque $\ell$ est pair, et
$G_X^{\mathrm{a}}(\underline{x}^\$ A)$ lorsque $\ell$ est impair.

\thingy\label{gale-stewart-winning-positions} On dira qu'une position
$\underline{x} = (x_0,\ldots,x_{\ell-1})$ d'un jeu de Gale-Stewart
$G_X(A)$ est \defin[gagnante (position)]{gagnante} pour Alice lorsque
Alice a une stratégie gagnante dans le jeu $G_X(\underline{x},A)$ qui
consiste à jouer à partir de cette position
(cf. \ref{gale-stewart-positions-as-games}).  On définit de même une
position gagnante pour Bob.

\medbreak

La proposition suivante est presque triviale et signifie qu'Alice (qui
doit jouer) possède une stratégie gagnante si et seulement si elle
peut jouer un coup $x$ qui l'amène à une position d'où elle (Alice) a
une stratégie gagnante, et Bob en possède une si et seulement si
n'importe quel coup $x$ joué par Alice amène à une position d'où il
(Bob) a une stratégie gagnante :
\begin{prop}\label{strategies-forall-exists-lemma}
Soit $X$ un ensemble non vide et $A \subseteq X^{\mathbb{N}}$.  Dans
le jeu de Gale-Stewart $G_X^{\mathrm{a}}(A)$, et en utilisant les
définitions faites en
\ref{unshifting-notation} et \ref{gale-stewart-positions-as-games} :
\begin{itemize}
\item Alice (premier joueur) possède une stratégie gagnante dans
  $G_X^{\mathrm{a}}(A)$ si et seulement si \emph{il existe} $x \in X$
  tel qu'elle (=Alice) possède une stratégie gagnante dans le jeu
  $G_X^{\mathrm{a}}(x, A)$ (dont on rappelle qu'il peut s'identifier
  au jeu $G_X^{\mathrm{b}}(x^{\$} A)$ défini par le sous-ensemble
  $x^{\$} A$ et où Bob joue en premier) ;
\item Bob (second joueur) possède une stratégie gagnante dans
  $G_X^{\mathrm{a}}(A)$ si et seulement si \emph{pour tout} $x \in X$
  il (=Bob) possède une stratégie gagnante dans le jeu
  $G_X^{\mathrm{a}}(x, A)$ (dont on rappelle qu'il peut s'identifier
  au jeu $G_X^{\mathrm{b}}(x^{\$} A)$ défini par le sous-ensemble
  $x^{\$} A$ et où Bob joue en premier).
\end{itemize}
(Les mêmes affirmations valent, bien sûr, en échangeant « Alice » et
    « Bob » tout du long ainsi que $G_X^{\mathrm{a}}$
    et $G_X^{\mathrm{b}}$.)
\end{prop}
\begin{proof}
La démonstration suivante ne fait que (laborieusement) formaliser
l'argument « une stratégie gagnante pour Alice détermine un premier
coup, après quoi elle a une stratégie gagnante, et une stratégie
gagnante pour Bob est prête à répondre à n'importe quel coup d'Alice
après quoi il a une stratégie gagnante » :

Si Alice (premier joueur) possède une stratégie $\varsigma$ gagnante
dans $G_X^{\mathrm{a}}(A)$, on pose $x := \varsigma(())$ le premier
coup préconisé par cette stratégie, et on définit
$\varsigma'((x_1,x_2,\ldots,x_{i-1})) =
\varsigma((x,x_1,x_2,\ldots,x_{i-1}))$ pour $i$ pair : cette
définition fait que si $(x_1,x_2,x_3,\ldots)$ est une confrontation où
Alice joue en second selon $\varsigma'$ alors $(x,x_1,x_2,x_3,\ldots)$
en est une où elle joue en premier selon $\varsigma$, donc cette suite
appartient à $A$ puisque $\varsigma$ est gagnante pour Alice
dans $G_X^{\mathrm{a}}(A)$, donc $(x_1,x_2,x_3,\ldots)$ appartient
à $x^{\$} A$, et Alice a bien une stratégie gagnante, $\varsigma'$
dans $G_X^{\mathrm{b}}(x^{\$} A)$ (où elle joue en second).

Réciproquement, si Alice possède une stratégie gagnante $\varsigma'$
dans $G_X^{\mathrm{b}}(x^{\$} A)$ (où elle joue en second), on définit
$\varsigma$ par $\varsigma(()) = x$ et
$\varsigma((x,x_1,x_2,\ldots,x_{i-1})) =
\varsigma'((x_1,x_2,\ldots,x_{i-1}))$ pour $i > 0$ pair : cette
définition fait que si $(x_0,x_1,x_2,x_3,\ldots)$ est une
confrontation où Alice joue en premier selon $\varsigma$ alors $x_0 =
x$ et $(x_1,x_2,x_3,\ldots)$ est confrontation où elle (Alice) joue en
second selon $\varsigma'$, donc cette suite appartient à $x^{\$} A$
puisque $\varsigma'$ est gagnante pour Alice second joueur
dans $G_X^{\mathrm{b}}(x^{\$} A)$, donc $(x_0,x_1,x_2,x_3,\ldots)$
appartient à $A$, et Alice a bien une stratégie gagnante, $\varsigma$,
dans $G_X^{\mathrm{a}}(A)$ (où elle joue en premier).

Si Bob (second joueur) possède une stratégie $\tau$ gagnante
dans $G_X^{\mathrm{a}}(A)$ et si $x \in X$ est quelconque, on définit
$\tau'((x_1,x_2,\ldots,x_{i-1})) = \tau((x,x_1,x_2,\ldots,x_{i-1}))$
pour $i$ impair : cette définition fait que si $(x_1,x_2,x_3,\ldots)$
est une confrontation où Bob joue en premier selon $\tau'$ alors
$(x,x_1,x_2,x_3,\ldots)$ en est une où il joue en second selon $\tau$,
donc cette suite n'appartient pas à $A$ puisque $\tau$ est gagnante
pour Bob dans $G_X^{\mathrm{a}}(A)$, donc $(x_1,x_2,x_3,\ldots)$
n'appartient pas à $x^{\$} A$, et Bob a bien une stratégie gagnante,
$\tau'$, dans $G_X^{\mathrm{b}}(x^{\$} A)$ (où il joue en premier).

Réciproquement, si pour chaque $x\in X$ Bob possède une stratégie
gagnante dans $G_X^{\mathrm{b}}(x^{\$} A)$ (où il joue en premier), on
en choisit une $\tau_x$ pour chaque $x$, et on définit $\tau$ par
$\tau((x,x_1,x_2,\ldots,x_{i-1})) = \tau_x((x_1,x_2,\ldots,x_{i-1}))$
pour $i$ impair : cette définition fait que si
$(x_0,x_1,x_2,x_3,\ldots)$ est une confrontation où Bob joue en second
selon $\tau$ alors $(x_1,x_2,x_3,\ldots)$ est confrontation où il
(Bob) joue en premier selon $\tau_{x_0}$, donc cette suite
n'appartient pas à ${x_0}^{\$} A$ puisque $\tau_{x_0}$ est gagnante
pour Bob premier joueur dans $G_X^{\mathrm{b}}({x_0}^{\$} A)$, donc
$(x_0,x_1,x_2,x_3,\ldots)$ n'appartient pas à $A$, et Bob a bien une
stratégie gagnante, $\tau$, dans $G_X^{\mathrm{a}}(A)$ (où il joue en
second).
\end{proof}

\thingy\label{strategies-forall-exists-reformulation} En utilisant la
terminologie \ref{gale-stewart-winning-positions}, la
proposition \ref{strategies-forall-exists-lemma} peut se reformuler de
la façon suivante, peut-être plus simple :
\begin{itemize}
\item une position $\underline{z}$ est gagnante pour le joueur qui
  doit jouer si et seulement si \emph{il existe} un coup $x$ menant à
  une position $\underline{z}x$ gagnante pour ce même joueur (qui est
  maintenant le joueur qui vient de jouer),
\item une position $\underline{z}$ est gagnante pour le joueur qui
  vient de jouer si et seulement si \emph{tous} les coups $x$ mènent à
  des positions $\underline{z}x$ gagnantes pour ce même joueur (qui
  est maintenant le joueur qui doit jouer).
\end{itemize}
(Dans ces affirmations, « un coup $x$ » depuis une position
$\underline{z} := (z_0,\ldots,z_{\ell-1})$ doit bien sûr se comprendre
comme menant à la position $\underline{z}x :=
(z_0,\ldots,z_{\ell-1},x)$ obtenue en ajoutant $x$ à la fin.)

\thingy On peut donc encore voir les choses comme suit : dire qu'Alice
a une stratégie gagnante dans le jeu $G_X^{\mathrm{a}}(A)$ signifie
que la position intiale $()$ est gagnante pour Alice, ce qui équivaut
(d'après ce qu'on vient de voir) à : $\exists x_0$ tel que la position
$(x_0)$ soit gagnante pour Alice ; ce qui équivaut à : $\exists x_0\,
\forall x_1$ la position $(x_0,x_1)$ est gagnante pour Alice ; ou
encore : $\exists x_0\, \forall x_1\, \exists x_2$ la position
$(x_0,x_1,x_2)$ est gagnante pour Alice ; ou encore : $\exists x_0\,
\forall x_1\, \exists x_2\, \forall x_3$ la position
$(x_0,x_1,x_2,x_3)$ est gagnante pour Alice ; et ainsi de suite.

Il peut donc être tentant de noter l'affirmation « Alice a une
  stratégie gagnante dans le jeu $G_X^{\mathrm{a}}(A)$ » par
l'écriture abusive
\[
\exists x_0\, \forall x_1\, \exists x_2\, \forall x_3\cdots
((x_0,x_1,x_2,x_3,\ldots) \in A)
\tag{$\mathscr{A}$}
\]
Cette écriture abusive n'a pas de sens mathématique, mais on peut la
\emph{définir} comme signifiant « Alice a une stratégie gagnante dans
  le jeu $G_X^{\mathrm{a}}(A)$ », tandis que « Bob a une stratégie
  gagnante dans le jeu $G_X^{\mathrm{a}}(A)$ » se noterait
\[
\forall x_0\, \exists x_1\, \forall x_2\, \exists x_3\cdots
((x_0,x_1,x_2,x_3,\ldots) \not\in A)
\tag{$\mathscr{B}$}
\]
Ces notations sont intuitivement intéressantes, mais elles présentent
l'inconvénient de laisser croire que l'affirmation $\mathscr{A}$ est
la négation de $\mathscr{B}$, ce qui n'est pas le cas en général
(comme on l'a dit, il se peut qu'aucun des joueurs n'ait de stratégie
gagnante).  Ce qui est vrai en revanche est que la négation
$\neg\mathscr{A}$ de $\mathscr{A}$ peut se réécrire (avec la notation
abusive définie ci-dessus) comme
\[
\begin{aligned}
\neg\exists x_0\, \forall x_1\, \exists x_2\, \forall x_3\cdots
((x_0,x_1,x_2,x_3,\ldots) &\in A)\\
\forall x_0\, \neg\forall x_1\, \exists x_2\, \forall x_3\cdots
((x_0,x_1,x_2,x_3,\ldots) &\in A)\\
\forall x_0\, \exists x_1\, \neg\exists x_2\, \forall x_3\cdots
((x_0,x_1,x_2,x_3,\ldots) &\in A)\\
\forall x_0\, \exists x_1\, \forall x_2\, \neg\forall x_3\cdots
((x_0,x_1,x_2,x_3,\ldots) &\in A)\\
\end{aligned}
\]
mais on ne peut pas traverser l'« infinité de quantificateurs » pour
passer à $\mathscr{B}$ sauf par exemple dans les conditions qu'on
verra en \ref{section-determinacy-of-gale-stewart-games}.


\subsection{Topologie produit}

\begin{defn}\label{definition-product-topology}
Soit $X$ un ensemble non vide.  Si $\dblunderline{x} :=
(x_0,x_1,x_2,\ldots)$ est une suite d'éléments de $X$ et
$\ell\in\mathbb{N}$, on appelle $\ell$-ième \defin{voisinage
  fondamental} de $\dblunderline{x}$, et on note
$V_\ell(\dblunderline{x})$ l'ensemble de tous les éléments
$(z_0,z_1,z_2,\ldots)$ de $X^{\mathbb{N}}$ dont les $\ell$ premiers
termes coïncident avec celles de $\dblunderline{x}$, autrement dit
$z_i = x_i$ si $i<\ell$.

On dit aussi qu'il s'agit du voisinage fondamental défini par la suite
finie $\underline{x} = (x_0,\ldots,x_{\ell-1})$ (il ne dépend
manifestement que de ces termes), et on peut le noter
$V_\ell(x_0,\ldots,x_{\ell-1})$ ou $V(x_0,\ldots,x_{\ell-1})$.

Un sous-ensemble $A \subseteq X^{\mathbb{N}}$ est dit \defin{ouvert}
[pour la topologie produit] lorsque pour tout $\dblunderline{x} \in A$
il existe un $\ell$ tel que le $\ell$-ième voisinage fondamental
$V_\ell(\dblunderline{x})$ de $\dblunderline{x}$ soit inclus dans $A$.
Autrement dit : dire que $A$ est ouvert signifie que lorsque $A$
contient une suite $\dblunderline{x} := (x_0,x_1,x_2,\ldots)$, il
existe un rang $\ell$ tel que $A$ contienne n'importe quelle suite
obtenue en modifiant la suite $\dblunderline{x}$ à partir du
rang $\ell$.

Un sous-ensemble $A \subseteq X^{\mathbb{N}}$ est dit \defin{fermé}
lorsque son complémentaire $B := X^{\mathbb{N}} \setminus A$ est
ouvert.
\end{defn}

\thingy On notera qu'il existe des parties de $X^{\mathbb{N}}$ à la
fois ouvertes et fermées : c'est le cas non seulement de $\varnothing$
et de $X^{\mathbb{N}}$, mais plus généralement de n'importe quel
voisinage fondamental $V_\ell(\dblunderline{x})$ (en effet,
$V_\ell(\dblunderline{x})$ est ouvert car si $\dblunderline{y} \in
V_\ell(\dblunderline{x})$, c'est-à-dire si $\dblunderline{y}$ coïncide avec
$\dblunderline{x}$ sur les $\ell$ premiers termes, alors toute suite
$\dblunderline{z}$ qui coïncide avec $\dblunderline{y}$ sur les $\ell$
premiers termes coïncide aussi avec $\dblunderline{x}$ dessus, et
appartient donc à $V_\ell(\dblunderline{x})$, autrement dit,
$V_\ell(\dblunderline{y})$ est inclus dans $V_\ell(\dblunderline{x})$ ; mais
$V_\ell(\dblunderline{x})$ est également fermé car si $\dblunderline{y}
\not\in V_\ell(\dblunderline{x})$, alors toute suite $\dblunderline{z}$ qui
coïncide avec $\dblunderline{y}$ sur les $\ell$ premiers termes ne
coïncide \emph{pas} avec $\dblunderline{x}$ dessus, donc n'appartient pas
à $V_\ell(\dblunderline{x})$, autrement dit $V_\ell(\dblunderline{y})$ est
inclus dans le complémentaire de $V_\ell(\dblunderline{x})$).

Il sera utile de remarquer que l'intersection de deux voisinages
fondamentaux $V,V'$ d'une même suite $\dblunderline{x}$ est encore un
voisinage fondamental de $\dblunderline{x}$ (en fait, cette intersection
est tout simplement égale à $V$ ou à $V'$).

\medbreak

L'énoncé suivant est une généralité topologique :
\begin{prop}
Soit $X$ un ensemble non vide.  Alors, dans $X^{\mathbb{N}}$ (pour la
topologie produit) :
\begin{itemize}
\item[(i)]$\varnothing$ et $X^{\mathbb{N}}$ sont ouverts,
\item[(ii)]une réunion quelconque d'ouverts est un ouvert (i.e., si
  $A_i$ est ouvert pour chaque $i\in I$ alors $\bigcup_{i\in I} A_i$
  est ouvert),
\item[(iii)]une intersection finie d'ouverts est un ouvert (i.e., si
  $A_1,\ldots,A_n$ sont ouverts alors $A_1\cap \cdots \cap A_n$ est
  ouvert).
\end{itemize}
\end{prop}
\begin{proof}
L'affirmation (i) est triviale.

Montrons (ii) : si les $A_i$ sont ouverts et si $\dblunderline{x} \in
\bigcup_{i\in I} A_i$, alors la définition d'une réunion fait qu'il
existe $i$ tel que $\dblunderline{x} \in A_i$, et comme $A_i$ est ouvert
il existe un voisinage fondamental de $\dblunderline{x}$ inclus
dans $A_i$, donc inclus dans $\bigcup_{i\in I} A_i$ : ceci montre que
$\bigcup_{i\in I} A_i$ est ouvert.

Montrons (iii) : il suffit de montrer que si $A, A'$ sont ouverts
alors $A \cap A'$ est ouvert.  Soit $\dblunderline{x} \in A \cap A'$.  Il
existe des voisinages fondamentaux $V$ et $V'$ de $\dblunderline{x}$
inclus dans $A$ et $A'$ respectivement (puisque ces derniers sont
ouverts) : alors $V \cap V'$ est un voisinage fondamental de
$\dblunderline{x}$ inclus dans $A \cap A'$ : ceci montre que $A \cap A'$
est ouvert.
\end{proof}


\subsection{Détermination des jeux ouverts}\label{section-determinacy-of-gale-stewart-games}

\thingy\label{fundamental-neighorhood-terminates-game} La remarque
suivante, bien que complètement évidente, sera cruciale : si
$(x_0,\ldots,x_{\ell-1})$ est une suite finie d'éléments de $X$ (i.e.,
une position de $G_X$) et $A$ une partie contenant le voisinage
fondamental (cf. \ref{definition-product-topology}) défini par
$(x_0,\ldots,x_{\ell-1})$, alors Alice possède une stratégie gagnante
à partir de $(x_0,\ldots,x_{\ell-1})$ dans le jeu $G_X(A)$
(cf. \ref{gale-stewart-positions-as-games}).  Mieux : quoi que fassent
l'un et l'autre joueur à partir de ce point, la partie sera gagnée par
Alice.  C'est tout simplement qu'on a fait l'hypothèse que
\emph{toute} suite commençant par $x_0,\ldots,x_{\ell-1}$ appartient
à $A$.

\begin{thm}[D. Gale \& F. M. Stewart, 1953]\label{gale-stewart-theorem}
Si $A \subseteq X^{\mathbb{N}}$ est ouvert, ou bien fermé, alors le
jeu $G_X(A)$ (qu'il s'agisse de $G_X^{\mathrm{a}}(A)$ ou
$G_X^{\mathrm{b}}(A)$) est déterminé.
\end{thm}
\begin{proof}[Première démonstration]
Il suffit de traiter le cas ouvert : le cas fermé s'en déduit
d'après \ref{remark-player-names} en passant au complémentaire,
c'est-à-dire en échangeant les deux joueurs (à condition d'avoir
traité le cas ouvert aussi le cas $G_X^{\mathrm{a}}(A)$ où Alice joue
en premier et le cas $G_X^{\mathrm{b}}(A)$ où Bob joue en premier).

Soit $A \subseteq X^{\mathbb{N}}$ ouvert.  Quel que soit le joueur qui
commence, on va montrer que si Alice (le joueur qui cherche à jouer
dans $A$) n'a pas de stratégie gagnante, alors Bob (le joueur qui
cherche à jouer dans le complémentaire) en a une.  On va définir une
stratégie $\tau$ pour Bob en évitant les positions où Alice a une
stratégie gagnante (une stratégie « défensive »).

Si $(x_0,\ldots,x_{i-1})$ est une position où c'est à Bob de jouer et
qui n'est pas gagnante pour Alice (c'est-à-dire qu'Alice n'a pas de
stratégie gagnante à partir de là,
cf. \ref{gale-stewart-winning-positions}), alors
d'après \ref{strategies-forall-exists-lemma}, (a) il existe un $x$ tel
que $(x_0,\ldots,x_{i-1},x)$ ne soit pas gagnante pour Alice :
choisissons-un tel $x$ et posons $\tau((x_0,\ldots,x_{i-1})) := x$.
Aux points où $\tau$ n'a pas été défini par ce qui vient d'être dit,
on le définit de façon arbitraire.  Par ailleurs, si
$(x_0,\ldots,x_{i-1})$ est une position où c'est à Alice de jouer et
qui n'est pas gagnante pour Alice, toujours
d'après \ref{strategies-forall-exists-lemma}, on remarque que (b) quel
que soit $y \in X$, la position $(x_0,\ldots,x_{i-1},y)$ n'est pas
gagnante pour Alice.

Si $x_0,x_1,x_2,\ldots$ est une confrontation où Bob joue
selon $\tau$, on voit par récurrence sur $i$ qu'aucune des positions
$(x_0,\ldots,x_{i-1})$ n'est gagnante pour Alice : pour $i=0$ c'est
l'hypothèse faite sur le jeu (à savoir, qu'Alice n'a pas de stratégie
gagnante depuis la position initiale), pour les positions où c'est à
Bob de jouer, c'est la construction de $\tau$ qui assure la récurrence
(cf. (a) ci-dessus), et pour les positions où c'est à Alice de jouer,
c'est le point (b) ci-dessus qui assure la récurrence.

On utilise maintenant le fait que $A$ est supposé ouvert : si
$x_0,x_1,x_2,\ldots$ appartient à $A$, alors il existe $\ell$ tel que
toute suite commençant par $x_0,\ldots,x_{\ell-1}$ appartienne à $A$.
Mais alors Alice a une stratégie gagnante à partir de la position
$(x_0,\ldots,x_{\ell-1})$
(cf. \ref{fundamental-neighorhood-terminates-game} : elle ne peut que
gagner à partir de là).  Or si Bob a joué selon $\tau$, ceci contredit
la conclusion du paragraphe précédent.  On en déduit que si Bob joue
selon $\tau$, la confrontation n'appartient pas à $A$, c'est-à-dire
que $\tau$ est gagnante pour Bob.
\end{proof}

\begin{proof}[Seconde démonstration]
Comme dans la première démonstration (premier paragraphe), on remarque
qu'il suffit de traiter le cas ouvert.  Soit $A \subseteq
X^{\mathbb{N}}$ ouvert.

On utilise la notion d'ordinaux qui sera introduite ultérieurement.
Soit $X^* := \bigcup_{\ell=0}^{+\infty} X^\ell$ l'arbre des positions
de $G_X$.

On définit les positions « gagnantes en $0$ coups pour Alice » comme
les positions $(x_0,\ldots,x_{\ell-1})$ qui définissent un voisinage
fondamental inclus dans $A$
(cf. \ref{fundamental-neighorhood-terminates-game} : quoi que les
joueurs fassent à partir de là, Alice aura gagné, et on peut
considérer qu'Alice a déjà gagné).

En supposant définies les positions gagnantes en $\alpha$ coups pour
Alice, on définit les positions « gagnantes en $\alpha+1$ coups pour
  Alice » de la façon suivante : ce sont les positions
$(x_0,\ldots,x_{\ell-1})$ où c'est à Alice de jouer et pour lesquelles
il existe un $x$ tel que $(x_0,\ldots,x_{\ell-1},x)$ soit gagnante en
$\alpha$ coups pour Alice, ainsi que les positions
$(x_0,\ldots,x_{\ell-1})$ où c'est à Bob de jouer et pour lesquels,
quel que soit $x\in X$, la position $(x_0,\ldots,x_{\ell-1},x)$ est
gagnante en $\alpha+1$ coups pour Alice (au sens où on vient de le
dire).

Enfin, si $\delta$ est un ordinal limite, en supposant définies les
positions « gagnantes en $\alpha$ coups pour Alice » pour tout $\alpha
< \delta$, on définit une position comme gagnante en $\delta$ coups
par Alice lorsqu'elle est gagnante en $\alpha$ coups pour un
certain $\alpha < \delta$.

La définition effectuée a les propriétés suivantes : (o) si une
position est gagnante en $\alpha$ coups pour Alice alors elle est
gagnante en $\alpha'$ coups pour tout $\alpha'>\alpha$, (i) si une
position où c'est à Bob de jouer est gagnante en $\alpha$ coups pour
Alice, alors tout coup (de Bob) conduit à une position gagnante en
$\alpha$ coups pour Alice, et (ii) si une position où c'est à Alice de
jouer est gagnante en $\alpha > 0$ coups pour Alice, alors il existe
un coup (d'Alice) conduisant à une position gagnante en strictement
moins que $\alpha$ coups (en fait, si $\alpha = \beta+1$ est
successeur, il existe un coup conduisant à une position gagnante en
$\beta$ coups par Alice, et si $\alpha$ est limite, la position
elle-même est déjà gagnable en strictement moins que $\alpha$ coups).

Si la position initiale $()$ est gagnante en $\alpha$ coups par Alice
pour un certain ordinal $\alpha$, alors Alice possède une stratégie
gagnante consistant à jouer, depuis une position gagnante en $\alpha$
coups, vers une position gagnante en $\beta$ coups pour un certain
$\beta < \alpha$ (ou bien $\beta = 0$), qui existe d'après (ii)
ci-dessus : comme Bob ne peut passer d'une position gagnante en
$\alpha$ coups par Alice que vers d'autres telles positions (cf. (i)),
et comme toute suite strictement décroissante d'ordinaux termine, ceci
assure à Alice d'arriver en temps fini à une position gagnante en $0$
coups.

Réciproquement, si la position initiale $()$ n'est pas gagnante en
$\alpha$ coups par Alice quel que soit $\alpha$ (appelons-la « non
  comptée »), alors Bob possède une stratégie consistant à jouer
toujours sur des telles positions non décomptées : d'après la
définition des positions gagnantes en $\alpha$ coup, quand c'est à
Alice de jouer, une position non comptée ne conduit qu'à des positions
non comptées, et quand c'est à Bob de jouer, une position non comptée
conduit à au moins une condition non comptée.  Ainsi, si Bob joue
selon cette stratégie, la confrontation $(x_0,x_1,x_2,\ldots)$ ne
passe que par des positions non comptées, et en particulier, ne passe
jamais par une position gagnante en $0$ coups par Alice, c'est-à-dire
qu'elle ne peut pas avoir un voisinage fondamental inclus dans $A$, et
comme $A$ est ouvert, elle n'appartient pas à $A$, i.e., la
confrontation est gagnée par Bob.
\end{proof}

\thingy Il ne faut pas croire que l'hypothèse « $A$ est ouvert ou bien
  fermé » est anodine : il existe des jeux $G_X^{\mathrm{a}}(A)$ qui
ne sont pas déterminés — autrement dit, même si dans toute
confrontation donnée l'un des deux joueurs gagne, aucun des deux n'a
de moyen systématique de s'en assurer.

Il ne faut pas croire pour autant que les seuls jeux déterminés soient
ceux définis par une partie ouverte.  Par exemple, il est facile de
voir que si $A$ est dénombrable, alors Bob possède une stratégie
gagnante (en effet, si $a = \{a_0,a_1,a_2,a_3,\ldots\}$, alors Bob
peut jouer au premier coup pour exclure $a_0$, c'est-à-dire jouer un
$x_1$ tel que $x_1 \neq a_{0,1}$, puis au second coup pour exclure
$a_1$, c'est-à-dire jouer un $x_3$ tel que $x_3 \neq a_{1,3}$, et
ainsi de suite $x_{2i+1} \neq a_{i,2i+1}$ : il s'agit d'un « argument
  diagonal constructif » ; l'argument fonctionne encore, quitte à
décaler les indices, si c'est Bob qui commence).

Le résultat ci-dessous généralise à la fois le
théorème \ref{gale-stewart-theorem} et ce qu'on vient de dire, et il
assez technique à démontrer :

\begin{thm}[D. A. Martin, 1975]
Si $A \subseteq X^{\mathbb{N}}$ est \defin{borélien}, c'est-à-dire
appartient à la plus petite partie de $\mathscr{P}(X^{\mathbb{N}})$
stable par complémentaire et réunions dénombrables (également appelée
\defin{tribu}) contenant les ouverts, alors le jeu $G_X(A)$ est déterminé.
\end{thm}

(Autrement dit, non seulement un ouvert et un fermé sont déterminés,
mais aussi une intersection dénombrable d'ouverts et une réunion
dénombrable de fermés, ou encore une réunion dénombrable
d'intersections dénombrables d'ouverts et une intersection dénombrable
de réunions dénombrables de fermés, « et ainsi de suite » ; les mots
« et ainsi de suite » glosent ici sur la construction des boréliens,
qui est plus complexe qu'une simple récurrence.)

\thingy Des résultats de détermination encore plus forts ont été
étudiés, et ne sont généralement pas prouvables dans la théorie des
ensembles usuelle (par exemple, l'« axiome de détermination
  projective », indémontrable dans $\mathsf{ZFC}$) ou sont même
incompatibles avec elle (l'« axiome de détermination », qui affirme
que pour toute partie $A \subseteq \{0,1\}^{\mathbb{N}}$ le jeu
$G_{\{0,1\}}(A)$ est déterminé, contredit l'axiome du choix, et a des
conséquences mathématiques remarquables comme le fait que toute partie
de $\mathbb{R}$ est mesurable au sens de Lebesgue).


\subsection{Détermination des jeux combinatoires}\label{subsection-determinacy-of-perfect-information-games}

On va définir ici rapidement les notions relatives aux jeux impartiaux
à information parfaite pour expliquer comment ces jeux peuvent se
ramener à des jeux de Gale-Stewart et comment la détermination des
jeux ouverts peut s'appliquer dans ce contexte :

\begin{defn}\label{definition-impartial-combinatorial-game}
Soit $G$ un graphe orienté (c'est-à-dire un ensemble $G$ muni d'une
relation $E$ irréflexive dont les éléments sont appelés arêtes du
graphe, cf. \ref{definitions-graphs} ci-dessous pour les définitions
générales) dont les sommets seront appelés \defin[position]{positions} de $G$,
et soit $x_0$ un sommet de $G$ qu'on appellera \index{initiale (position)}\defin{position
  initiale}.  Le \index{impartial (jeu)}\index{information parfaite (jeu à)}\defin[combinatoire (jeu)]{jeu combinatoire impartial à information
  parfaite} associé à ces données est défini de la manière suivante :
partant de $x = x_0$, Alice et Bob choisissent tour à tour un voisin
sortant de $x$, autrement dit, Alice choisit une arête $(x_0,x_1)$
de $G$, puis Bob choisit une arête $(x_1,x_2)$ de $G$, puis Alice
choisit une arête $(x_2,x_3)$, et ainsi de suite.  Si un joueur ne
peut plus jouer, il a perdu ; si la confrontation dure un temps
infini, elle est considérée comme nulle (ni gagnée ni perdue par les
joueurs).

Une \textbf{partie} ou \defin{confrontation} de ce jeu est une suite
finie ou infinie $(x_i)$ de sommets de $G$ telle que $x_0$ soit la
position initiale et que pour chaque $i\geq 1$ pour lequel $x_i$ soit
défini, ce dernier soit un voisin sortant de $x_{i-1}$.  Lorsque le
dernier $x_i$ défini l'est pour un $i$ pair, on dit que le premier
joueur \textbf{perd} et que le second \defin[gain]{gagne}, tandis que
lorsque le dernier $x_i$ défini l'est pour un $i$ impair, on dit que
le premier joueur gagne et que le second perd ; enfin, lorsque $x_i$
est défini pour tout entier naturel $i$, on dit que la confrontation
est nulle ou que les deux joueurs \defin[survie]{survivent} sans gagner.
\end{defn}

\thingy\label{combinatorial-to-gale-stewart} Pour un jeu comme
en \ref{definition-impartial-combinatorial-game}, va définir un,
ou plutôt deux, jeux de Gale-Stewart : l'intuition est que si un
joueur enfreint la « règle » du jeu (i.e., choisit un sommet qui n'est
pas un voisin sortant du sommet actuel), il a immédiatement perdu — il
n'y a manifestement pas grande différence entre avoir un jeu où un
joueur \emph{ne peut pas} faire un certain coup et un jeu où si ce
joueur fait ce coup il a immédiatement perdu (quoi qu'il se passe par
la suite).  On va définir deux jeux de Gale-Stewart plutôt qu'un parce
qu'un jeu de Gale-Stewart a forcément un gagnant (il n'y a pas de
partie nulle), donc on va définir un jeu où les parties nulles sont
comptées au bénéfice de Bob et un autre où elles sont comptées au
bénéfice d'Alice.

Autrement dit, soit $G$ un graphe orienté et $x_0 \in G$.  On pose $X
= G$ et on partitionne l'ensemble des suites à valeurs dans $X$ en
trois :
\begin{itemize}
\item l'ensemble $D$ des (confrontations nulles dans le jeu
  combinatoire, c'est-à-dire des) suites $(x_1,x_2,x_3,\ldots)$ telles
  que pour chaque $i \geq 1$ le sommet $x_i$ soit un voisin sortant
  de $x_{i-1}$ (c'est-à-dire : $(x_{i-1},x_i)$ est une arête de $G$)
  (bref, personne n'a enfreint la règle),
\item l'ensemble $A$ des (confrontations gagnées par Alice dans le jeu
  combinatoire, c'est-à-dire des) suites $(x_1,x_2,x_3,\ldots)$ telles
  qu'il existe $i \geq 1$ pour lequel $x_i$ n'est pas un voisin
  sortant de $x_{i-1}$ et que le plus petit tel $i$ soit \emph{pair}
  (i.e., Bob a enfreint la règle en premier),
\item l'ensemble $B$ des (confrontations gagnées par Bob dans le jeu
  combinatoire, c'est-à-dire des) suites $(x_1,x_2,x_3,\ldots)$ telles
  qu'il existe $i \geq 1$ pour lequel $x_i$ n'est pas un voisin
  sortant de $x_{i-1}$ et que le plus petit tel $i$ soit \emph{impair}
  (i.e., Alice a enfreint la règle en premier).
\end{itemize}
(On a choisi ici d'indicer les suites par les entiers naturels non
nuls : il va de soi que ça ne change rien à la théorie des jeux de
Gale-Stewart !  Si on préfère, on peut les faire commencer à $0$, et
mettre dans $A$ toutes les suites qui ne commencent pas par $x_0$.)

Le jeu de Gale-Stewart $G_X(A)$ est essentiellement identique au jeu
considéré en \ref{definition-impartial-combinatorial-game}, à
ceci près que les confrontations nulles sont comptées comme des gains
de Bob ; le jeu $G_X(A \cup D)$, pour sa part, est lui aussi identique
à ceci près que les nuls sont comptés comme des gains d'Alice.

\begin{lem}\label{openness-of-combinatorial-to-gale-stewart}
Avec les notations de \ref{combinatorial-to-gale-stewart}, les parties
$A$ et $B$ sont ouvertes (pour la topologie produit,
cf. \ref{definition-product-topology}).
\end{lem}
\begin{proof}
Soit $(x_1,x_2,x_3,\ldots)$ est une suite d'éléments de $X = G$.  Si
la suite appartient à $A$ alors, par définition de $A$, il existe un
$i\geq 1$ pair tel que $x_i$ ne soit pas un voisin sortant
de $x_{i-1}$ et tel que $x_j$ soit un voisin sortant de $x_{j-1}$ pour
tout $1\leq j<i$.  On en déduit que le voisinage fondamental formé de
toutes les suites qui coïncident avec $(x_1,x_2,x_3,\ldots)$ jusqu'à
$x_i$ inclus est contenu dans $A$.  La même démonstration fonctionne
pour $B$ avec $i$ impair.
\end{proof}

\begin{thm}\label{determinacy-of-perfect-information-games}
Soit $(G,x_0)$ un jeu combinatoire impartial à information parfaite
comme en \ref{definition-impartial-combinatorial-game}.  Alors
exactement l'une des trois affirmations suivantes est vraie :
\begin{itemize}
\item le premier joueur (Alice) possède une stratégie gagnante,
\item le second joueur (Bob) possède une stratégie gagnante,
\item chacun des deux joueurs possède une stratégie survivante.
\end{itemize}
(La notion de « stratégie » ici doit se comprendre comme pouvant
dépendre de l'histoire des coups joués précédemment : voir
\ref{remark-historical-versus-positional-strategies} ci-dessous.)
\end{thm}
\begin{proof}
Il est évident que les affirmations sont exclusives (si un joueur
possède une stratégie gagnante, l'autre ne peut pas posséder de
stratégie survivante, sinon on aurait une contradiction en les faisant
jouer l'une contre l'autre).

Avec les notations de \ref{combinatorial-to-gale-stewart},
d'après \ref{openness-of-combinatorial-to-gale-stewart}, les parties
$A$ et $B$ sont ouvertes, donc \ref{gale-stewart-theorem} montre que
les jeux définis par l'ouvert $A$ et le fermé $A \cup D = X\setminus
B$ sont déterminés.

Mais une stratégie gagnante d'Alice dans le jeu de Gale-Stewart défini
par $A$ est une stratégie gagnante dans le jeu combinatoire d'origine,
tandis qu'une stratégie gagnante de Bob dans ce jeu de Gale-Stewart
est une stratégie survivante dans le jeu d'origine ; ainsi, dans le
jeu d'origine, soit Alice a une stratégie gagnante soit Bob a une
stratégie survivante.

De même, une stratégie gagnante d'Alice dans le jeu défini par $A \cup
D$ est une stratégie survivante dans le jeu d'origine, tandis qu'une
stratégie gagnante de Bob dans ce jeu est une stratégie gagnante dans
le jeu d'origine ; ainsi, dans le jeu d'origine, soit Alice a une
stratégie survivante soit Bob a une stratégie gagnante.

En mettant ensemble ces deux disjonctions, on voit que l'un des trois
faits énoncés est vrai.
\end{proof}

\thingy\label{remark-historical-versus-positional-strategies} La
notion de « stratégie » implicite dans le
théorème \ref{determinacy-of-perfect-information-games} est une notion
\emph{historique} : puisque c'est ainsi qu'elles ont été définies
en \ref{definition-strategies-for-gale-stewart-games}, les stratégies
ont le droit de choisir le coup à jouer en fonction de \emph{tous les
  coups joués antérieurement}.  Il se trouve en fait qu'on a les mêmes
résultats avec des stratégies \emph{positionnelles}, c'est-à-dire, qui
ne choisissent un coup qu'en fonction du sommet $x \in G$.  C'est ce
qui va être démontré dans la section suivante.


\subsection{Détermination pour les stratégies positionnelles}\label{subsection-positional-strategies-determinacy}

\thingy Le but de la définition suivante est de formaliser, pour un
jeu combinatoire comme
en \ref{definition-impartial-combinatorial-game}, les notions de
stratégie positionnelle (dans laquelle un joueur ne choisit le coup à
jouer qu'en fonction de la position actuelle), et de stratégie
historique (dans laquelle il fait son choix en fonction de tous les
coups joués antérieurement), sachant qu'on veut montrer au final que
cette distinction a peu d'importance.  Mais la définition sur laquelle
on va vraiment travailler est formulée
en \ref{set-of-everywhere-winning-positional-strategies}, donc on peut
se contenter de lire celle-ci.

\begin{defn}
Soit $G$ un graphe orienté
(cf. \ref{definition-impartial-combinatorial-game} et \ref{definitions-graphs}).
Une \index{positionnelle (stratégie)}\defin{stratégie positionnelle} sur $G$ est une fonction
partielle $\varsigma\colon G \dasharrow G$ (i.e., une fonction définie
sur un sous-ensemble de $G$) telle que $\varsigma(x)$ soit, s'il est
défini, un voisin sortant de $x$ (s'il n'est pas défini, il faut
comprendre que le joueur abandonne la partie).  Une \index{historique (stratégie)}\defin{stratégie
  historique} sur $G$ une fonction partielle $\varsigma
\big(\bigcup_{\ell=1}^{+\infty} G^\ell\big) \dasharrow G$, où
$\big(\bigcup_{\ell=1}^{+\infty} G^\ell\big)$ désigne l'ensemble des
suites finies de $G$ de longueur non nulle (i.e., des suites
$z_0,\ldots,z_{\ell-1}$ d'éléments de $G$ avec $\ell>0$ entier) telle
que $\varsigma(z_0,\ldots,z_{\ell-1})$ soit un voisin sortant
de $z_{\ell-1}$.

Lorsque dans une confrontation $x_0,x_1,x_2,\ldots$ de $G$ (à partir
d'une position initiale $x_0$) on a $x_i = \varsigma(x_{i-1})$ pour
chaque $i\geq 1$ impair pour lequel $x_i$ est défini, on dit que le
premier joueur a joué la partie selon la stratégie positionnelle
$\varsigma$ ; tandis que si $x_i = \tau(x_{i-1})$ pour chaque $i\geq
1$ pair pour lequel $x_i$ est défini, on dit que le second joueur a
joué la partie selon la stratégie $\tau$.  Pour une stratégie
historique, il faut remplacer $x_i = \varsigma(x_{i-1})$ et $x_i =
\tau(x_{i-1})$ par $x_i = \varsigma(x_0,\ldots,x_{i-1})$ et $x_i =
\tau(x_0,\ldots,x_{i-1})$ respectivement.

La position initiale $x_0 \in G$ ayant été fixée, si $\varsigma$ et
$\tau$ sont deux stratégies (positionnelles ou historiques), on
définit $\varsigma \ast \tau$ comme la confrontation jouée lorsque le
premier joueur joue selon $\varsigma$ et le second joue selon $\tau$ :
autrement dit, $x_0$ est la position initiale, et, si $x_{i-1}$ est
défini, $x_i$ est défini par $\varsigma(x_{i-1})$ ou
$\varsigma(x_1,\ldots,x_{i-1})$ si $i$ est impair et $\tau(x_{i-1})$
ou $\tau(x_1,\ldots,x_{i-1})$ si $i$ est pair (si $x_i$ n'est pas
défini, la suite s'arrête là).

La stratégie (positionnelle ou historique) $\varsigma$ est dite
\index{stratégie gagnante}\defin[gagnante (stratégie)]{gagnante pour le premier joueur} à partir de la position
initiale $x_0$ lorsque le premier joueur gagne toute confrontation où
il joue selon $\varsigma$, c'est-à-dire que la confrontation est finie
et que le dernier $x_i$ défini l'est pour un $i$ impair.  On définit
de même une stratégie $\varsigma$ \index{stratégie survivante}\defin[survivante (stratégie)]{survivante} (c'est-à-dire,
permettant d'assurer au moins le nul) pour le premier joueur à partir
d'une position initiale $x_0$, c'est-à-dire que dans toute
confrontation où il joue selon $\varsigma$, soit la confrontation est
infinie (donc nulle) soit le dernier $x_i$ défini l'est pour un $i$
impair.
\end{defn}

\thingy La notion de stratégie pour le second joueur peut être définie
de façon analogue, bien sûr, mais la notion n'a pas beaucoup
d'intérêt : une stratégie gagnante pour le second joueur à partir de
$x_0$ est la même chose qu'une stratégie gagnante pour le premier
joueur à partir de tout voisin sortant de $x_0$.  Pour travailler avec
les stratégies positionnelles, il vaut mieux supposer qu'elles sont,
en fait, gagnantes partout où elles sont définies, ce qui amène à
faire la définition suivante :

\thingy\label{set-of-everywhere-winning-positional-strategies} Dans ce
qui suit, on va fixer un graphe orienté $G$ et on aura besoin
d'introduire l'ensemble $\mathcal{S}$ de ce qu'on pourrait appeler les
« stratégies positionnelles gagnantes partout où définies »,
c'est-à-dire des stratégies positionnelles $\varsigma$ gagnantes pour
le premier joueur à partir de n'importe quel point $x_0$ où
$\varsigma$ est défini ; autrement dit, il s'agit de l'ensemble des
fonctions partielles $\varsigma\colon G \dasharrow G$ telles que
\begin{itemize}
\item si $\varsigma(x)$ est défini alors il est un voisin sortant
  de $x$, et que
\item si $\varsigma(x_0)$ est défini et si $(x_i)$ est une suite
  (\textit{a priori} finie ou infinie) partant de $x_0$, dans laquelle
  $x_i = \varsigma(x_{i-1})$ pour $i\geq 1$ impair, et $x_i$ est
  voisin sortant de $x_{i-1}$ pour tout $i\geq 1$ [pair], alors la
  suite est de longueur finie et le dernier $x_i$ défini l'est pour un
  $i$ impair
\end{itemize}
(i.e., le premier joueur gagne n'importe quelle confrontation à partir
d'une position initiale $x_0$ du domaine de définition de $\varsigma$
et où il joue selon $\varsigma$).

L'ensemble $\mathcal{S}$ est partiellement ordonné par l'inclusion (si
$\varsigma,\tau \in \mathcal{S}$, on dit que $\varsigma$
\defin{prolonge} $\tau$ et on note $\varsigma \supseteq \tau$ ou
$\tau \subseteq \varsigma$, lorsque l'ensemble de définition de
$\varsigma$ contient celui de $\tau$ et que $\varsigma$ et $\tau$
coïncident là où $\tau$ est définie : ceci signifie bien que
$\varsigma \supseteq \tau$ en tant qu'ensembles).

\begin{lem}\label{positional-strategies-merging-lemma}
Si $\varsigma,\varsigma' \in \mathcal{S}$ avec la notation introduite
en \ref{set-of-everywhere-winning-positional-strategies}, alors il
existe $\varsigma'' \in \mathcal{S}$ qui prolonge $\varsigma$ et qui
est également définie en tout point où $\varsigma'$ l'est.
\end{lem}
\begin{proof}
Définissons $\varsigma''$ par $\varsigma''(x) = \varsigma(x)$ si
$\varsigma(x)$ est définie et $\varsigma''(x) = \varsigma'(x)$ si
$\varsigma(x)$ n'est pas définie mais que $\varsigma'(x)$ l'est.  Il
est évident que $\varsigma''$ prolonge $\varsigma$ et est également
définie en tout point où $\varsigma'$ l'est : il reste à voir que
$\varsigma'' \in \mathcal{S}$.  Mais si $x_0,x_1,x_2,\ldots$
(\textit{a priori} finie ou infinie) est une confrontation jouée par
le premier joueur selon $\varsigma''$, montrons qu'elle est
nécessairement gagnée par le premier joueur : or le premier joueur a
soit joué selon $\varsigma$ tout du long, soit selon $\varsigma'$ tout
du long, soit selon $\varsigma'$ puis $\varsigma$, mais dans tous les
cas il gagne.  De façon détaillée : soit $x_0$ est domaine de
définition de $\varsigma$ auquel cas tous les $x_i$ pairs le sont et
la confrontation est gagnée par le premier joueur ; soit $x_0$ est
dans le domaine de définition de $\varsigma'$ mais pas de $\varsigma$,
auquel cas il n'existe qu'un nombre fini de $i$ pairs tels que $x_i$
ne soit pas dans domaine de définition de $\varsigma$ (puisque
$\varsigma'$ ne peut pas donner une partie nulle) et si $j$ est le
plus grand d'entre eux, $x_{j+1}$ est défini, et soit $x_{j+2}$ n'est
pas défini (auquel cas le premier joueur a gagné) soit $x_{j+2}$ est
dans le domaine de définition de $\varsigma$ et de nouveau le premier
joueur gagne à partir de là.
\end{proof}

\begin{lem}\label{positional-strategies-union-lemma}
Si $\varsigma_i \in \mathcal{S}$ pour chaque $i\in I$ avec la notation
introduite en \ref{set-of-everywhere-winning-positional-strategies},
et si pour tous $i,j$ les fonctions $\varsigma_i$ et $\varsigma_j$
coïncident là où elles sont toutes deux définies, alors la fonction
$\bigcup_{i\in I} \varsigma_i$ (c'est-à-dire la fonction définie sur
la réunion des ensembles de définition des $\varsigma_i$ et qui
coïncide avec n'importe quel $\varsigma_i$ sur l'ensemble de
définition de celui-ci) appartient encore à $\mathcal{S}$.
\end{lem}
\begin{proof}
Si $x_0,x_1,x_2,\ldots$ (\textit{a priori} finie ou infinie) est une
confrontation jouée par le premier joueur selon $\varsigma :=
\bigcup_{i\in I} \varsigma_i$, alors en fait elle est jouée tout du
long selon $\varsigma_i$ où $i$ est n'importe quel indice tel que
$\varsigma_i(x_0)$ soit définie.  Comme $\varsigma_i \in \mathcal{S}$,
cette confrontation est gagnée par le premier joueur.
\end{proof}

Le résultat ensembliste suivant sera admis (même si on pourrait s'en
sortir en appliquant \ref{fixed-point-lemma-for-partial-functions} à
la place) :
\begin{lem}[principe maximal de Hausdorff]\label{hausdorff-maximal-principle}
Soit $\mathscr{F}$ un ensemble de parties d'un ensemble $A$.  On
suppose que $\mathscr{F}$ est non vide et que pour toute partie non
vide $\mathscr{T}$ de $\mathscr{F}$ totalement ordonnée par
l'inclusion (c'est-à-dire telle que pour $P,P' \in \mathscr{T}$ on a
soit $P \subseteq P'$ soit $P \supseteq P'$) la réunion $\bigcup_{P
  \in \mathscr{T}} P$ soit contenue dans un élément de $\mathscr{F}$.
Alors il existe dans $\mathscr{F}$ un élément $M$ maximal pour
l'inclusion (c'est-à-dire que si $P \supseteq M$ avec $P \in
\mathscr{F}$ alors $P=M$).
\end{lem}

\begin{prop}\label{existence-of-maximal-positional-strategy}
Avec la notation introduite
en \ref{set-of-everywhere-winning-positional-strategies}, il existe
$\varsigma \in \mathcal{S}$ maximal pour l'inclusion (au sens où si
$\varsigma \subseteq \varsigma'$ avec $\varsigma' \in \mathcal{S}$
alors $\varsigma = \varsigma'$) ; de plus, si $\varsigma' \in
\mathcal{S}$, alors $\varsigma$ est défini en tout point où
$\varsigma'$ l'est.
\end{prop}
\begin{proof}
L'existence de $\varsigma \in \mathcal{S}$ maximal pour l'inclusion
découle immédiatement de \ref{hausdorff-maximal-principle} en
utilisant \ref{positional-strategies-union-lemma} pour constater que
si $\mathscr{T} \subseteq \mathcal{S}$ est totalement ordonné pour
l'inclusion alors la réunion $\bigcup_{\varsigma\in\mathscr{T}}
\varsigma$ appartient à $\mathcal{S}$.

Une fois trouvé $\varsigma \in \mathcal{S}$ maximal pour l'inclusion,
si $\varsigma' \in \mathcal{S}$,
d'après \ref{positional-strategies-merging-lemma}, on peut trouver
$\varsigma''$ prolongeant $\varsigma$ et défini partout où
$\varsigma'$ l'est, et comme $\varsigma$ est maximal, on a
$\varsigma'' = \varsigma$, donc $\varsigma$ est bien défini partout où
$\varsigma'$ l'est.
\end{proof}

\thingy\label{notation-n-and-p-sets-for-combinatorial-games} En
continuant les notations introduites
en \ref{set-of-everywhere-winning-positional-strategies}, on fixe
maintenant $\varsigma \in \mathcal{S}$ maximal pour l'inclusion (dont
l'existence est garantie par la
proposition \ref{existence-of-maximal-positional-strategy}).  Soit $N$
l'ensemble des sommets $x$ de $G$ où $\varsigma(x)$ est défini (i.e.,
le domaine de définition de $\varsigma$) ; on vient de voir que $N$
est aussi l'ensemble des points où un élément quelconque de
$\mathcal{S}$ est défini, i.e., l'ensemble des positions à partir
desquelles le premier joueur a une stratégie positionnelle gagnante.
Soit $P$ l'ensemble des sommets $x\in G$ dont tous les voisins
sortants appartiennent à $N$ (y compris s'il n'y a pas de voisin
sortant, i.e., si $x$ est un puits) : clairement, $P$ est l'ensemble
des positions à partir desquelles le second joueur a une stratégie
positionnelle gagnante (quel que soit le coup du joueur adversaire, il
amènera à une position où on a une stratégie gagnante).

Enfin, on note $D := G\setminus(N\cup P)$ l'ensemble des sommets
restants.

\begin{prop}\label{positional-strategies-forall-exists-lemma}
Avec les notations introduites en
\ref{set-of-everywhere-winning-positional-strategies} et \ref{notation-n-and-p-sets-for-combinatorial-games},
\begin{itemize}
\item[(i)]un sommet $x\in G$ appartient à $N$ si et seulement si il a
  au moins un voisin sortant qui appartient à $P$,
\item[(ii)]un sommet $x\in G$ appartient à $P$ si et seulement si tous
  ses voisins sortants appartiennent à $N$.
\end{itemize}
\end{prop}
\begin{proof}
L'affirmation (ii) est la définition même de $P$ et n'a donc pas à
être prouvée.  Il s'agit donc de montrer (i).

Si $x\in N$ alors $y := \varsigma(x)$ est un voisin sortant de $x$ qui
appartient à $P$ puisque $\varsigma \in \mathcal{S}$.  Réciproquement,
si $x$ a un voisin sortant $y$ qui appartient à $P$, et si
$\varsigma(x)$ n'était pas défini, on pourrait étendre $\varsigma$ en
posant $\varsigma(x) = y$, ce qui donnerait un élément de
$\mathcal{S}$ (toute partie jouée à partir de $x$ conduit à $y$ et de
là à des points où $\varsigma$ est définie, donc est gagnée par le
premier joueur), contredisant la maximalité de $\varsigma$ ; c'est
donc que $\varsigma(x)$ était bien défini, i.e., $x\in N$.
\end{proof}

\thingy Par contraposée, sur l'ensemble $D := G\setminus(N\cup P)$ des
sommets restants de $G$, on a les propriétés suivantes : $x\in G$
appartient à $D$ si et seulement si
\begin{itemize}
\item[(i*)]tous les voisins sortants de $x$ sont dans $N\cup D$
  \emph{et}
\item[(ii*)]au moins l'un d'entre eux appartient à $D$.
\end{itemize}

On définit alors une stratégie positionnelle $\tau$ étendant
$\varsigma$ de la façon suivante : si $\varsigma(x)$ est défini (i.e.,
$x \in N$), on pose $\tau(x) = \varsigma(x)$, et si $x \in D$ on
choisit pour $\tau(x)$ un voisin sortant de $x$ qui appartienne à $D$
(lequel voisin existe d'après (ii*)).  À partir d'un sommet $x_0$
dans $D$, si l'un ou l'autre joueur joue selon $\tau$, ce joueur
survit, puisque soit son adversaire le laisse toujours dans $D$ auquel
cas le joueur considéré peut toujours jouer (selon $\tau$) en restant
dans $D$, soit son adversaire quitte $D$ et d'après (i*) joue
vers $N$, et alors le joueur considéré gagne puisqu'il joue
selon $\varsigma$.

Bref, à partir d'un sommet de $N$ le premier joueur a une stratégie
\emph{positionnelle} gagnante, à partir d'un sommet de $P$ c'est le
second joueur qui en a une, et à partir d'un sommet de $D$ les deux
joueurs ont une stratégie positionnelle survivante.  (Dans tous les
cas, on peut utiliser le $\tau$ qu'on vient de construire comme
stratégie positionnelle.)

\begin{thm}\label{positional-determinacy-of-perfect-information-games}
Soit $G$ un graphe orienté.  Quel que soit le sommet $x_0$ de $G$
choisi comme position initiale, dans le jeu combinatoire impartial à
information parfaite considéré
en \ref{definition-impartial-combinatorial-game}, exactement
l'une des trois affirmations suivantes est vraie :
\begin{itemize}
\item le premier joueur (Alice) possède une stratégie positionnelle gagnante,
\item le second joueur (Bob) possède une stratégie positionnelle gagnante,
\item chacun des deux joueurs possède une stratégie positionnelle survivante.
\end{itemize}

En particulier, l'existence d'une stratégie positionnelle gagnante ou
survivante est équivalente à celle d'une stratégie historique de même
nature.

De plus, si $N,P,D$ sont les ensembles de sommets $x_0$ à partir
desquels chacune des trois affirmations ci-dessus est vraie,
\begin{itemize}
\item un sommet $x\in G$ appartient à $N$ si et seulement si il a
  au moins un voisin sortant qui appartient à $P$,
\item un sommet $x\in G$ appartient à $P$ si et seulement si tous
  ses voisins sortants appartiennent à $N$,
\item un sommet $x\in G$ appartient à $D$ si et seulement si tous ses
  voisins sortants appartiennent à $N\cup D$ et au moins l'un d'eux
  appartient à $D$.
\end{itemize}
\end{thm}
\begin{proof}
L'existence des stratégies positionnelles a déjà été montré ci-dessus,
ainsi que les propriétés sur $N,P,D$.  L'équivalence entre stratégies
positionnelles et historiques vient du fait que toute stratégie
positionnelle peut être vue comme une stratégie historique (en
ignorant l'historique) et du fait que les trois cas sont exclusifs
aussi bien pour les stratégies historiques que positionnelles.
\end{proof}

\thingy On a en particulier redémontré le
théorème \ref{determinacy-of-perfect-information-games}, même si la
démonstration suivie ici (consistant à prendre une stratégie
positionnelle maximale pour l'inclusion) n'est peut-être pas très
éclairante.

En utilisant la notion d'ordinaux, on pourrait donner une autre
démonstration, plus explicite, du
théorème \ref{positional-determinacy-of-perfect-information-games} :
elle consiste à reprendre la seconde démonstration qui a été donnée du
théorème \ref{determinacy-of-perfect-information-games} et à constater
qu'en la modifiant à peine elle conduit maintenant à définir une
stratégie positionnelle ; un peu plus précisément, on définit par
induction sur l'ordinal $\alpha$ les positions gagnantes en $\alpha$
coups par le premier joueur et les positions gagnantes en $\alpha$
coups par le second joueur, et une stratégie gagnante consiste à jouer
d'une position gagnante en $\alpha$ coups pour le premier joueur vers
une position gagnante en $\beta<\alpha$ coups pour le second joueur
(pour les stratégies survivantes, on complète en jouant d'une position
non étiquetée vers une position non étiquetée, mais le problème de la
survie est de toute façon plus simple).

Dans le cas des jeux définis par un graphe bien-fondé
(cf. \ref{definitions-graphs}), c'est-à-dire que le nul est
impossible, la détermination est beaucoup plus simple à démontrer en
on en verra encore une nouvelle démonstration dans le cadre de la
théorie de Grundy.

Un bonus au
théorème \ref{positional-determinacy-of-perfect-information-games} est
l'affirmation suivante :

\begin{prop}\label{p-d-n-vertices-of-perfect-information-games}
Dans le contexte du
théorème \ref{positional-determinacy-of-perfect-information-games}, si
$N^*,P^*,D^*$ est une partition de $G$ en trois parties vérifiant les
trois propriétés qui ont été énoncées pour $N,P,D$ (en remplaçant
$N,P,D$ par $N^*,P^*,D^*$ respectivement), alors on a $N\subseteq N^*
\subseteq N\cup D$ et $P\subseteq P^* \subseteq P\cup D$ et bien sûr
$D\supseteq D^*$.

En particulier, si on utilise le
théorème \ref{non-well-founded-definition} ci-dessous pour définir la
\emph{plus petite} (pour l'inclusion) fonction partielle $f\colon G
\dasharrow Z := \{\mathtt{P}, \mathtt{N}\}$ telle que $f(x)$ vaille
$\mathtt{N}$ ssi $x$ a au moins un voisin sortant $y$ pour lequel
$f(y) = \mathtt{P}$, et que $f(x)$ vaille $\mathtt{P}$ ssi pour tout
voisin sortant $y$ de $x$ on a $f(y) = \mathtt{N}$, alors $f(x)$ vaut
$\mathtt{P}$ ou bien $\mathtt{N}$ ou bien est indéfinie lorsque
respectivement le premier joueur a une stratégie gagnante, le second
joueur en a une, ou les deux ont une stratégie survivante.
\end{prop}
\begin{proof}
Si $N^*,P^*,D^*$ ont les mêmes propriétés que $N,P,D$, on peut définir
une stratégie (positionnelle) consistant à jouer à partir de $N^*$
dans $P^*$ (c'est-à-dire à choisir pour chaque sommet de $N^*$ un
voisin sortant dans $P^*$) et à partir de $D^*$ dans $D^*$ : si le
premier joueur suit cette stratégie à partir d'un sommet dans $N^*$ ou
$D^*$ ne peut pas perdre puisque les propriétés des parties font que
son coup sera toujours défini : donc $N^* \cup D^* \subseteq N \cup
D$, ce qui signifie en passant au complémentaire $P^* \supseteq P$, et
si le second joueur suit cette stratégie à partir d'un sommet dans
$P^*$ ou $D^*$ il ne peut pas perdre non plus : donc $P^* \cup D^*
\subseteq P \cup D$, ce qui signifie exactement $N^* \supseteq N$.

Le deuxième paragraphe est une reformulation de la même affirmation :
la fonction $f \colon G \dasharrow Z := \{\mathtt{P}, \mathtt{N}\}$
définie par $f(x) = \mathtt{P}$ lorsque $x \in P$ et $f(x) =
\mathtt{N}$ lorsque $x \in N$ est la plus petite fonction partielle
vérifiant les propriétés qu'on a dites, i.e., la plus petite telle que
$f(x) = \Phi(x, f|_{\outnb(x)})$ avec les notations du
théorème \ref{non-well-founded-definition} et $\Phi(x, g)$ valant
$\mathtt{N}$ si $g(y)$ est défini et vaut $\mathtt{P}$ pour au moins
un $y \in \outnb(x)$ et $\mathtt{P}$ si $g(y)$ est défini et vaut
$\mathtt{N}$ pour tout $y \in \outnb(x)$ (comme $\Phi$ est cohérente
en la seconde variable, on est bien dans le cas d'application
du théorème \ref{non-well-founded-definition}, même si ici on a déjà
démontré l'existence d'une plus petite $f$).
\end{proof}



%
%
%

\section{Théorie de l'induction bien-fondée}\label{section-well-founded-induction}

Le but de cette partie est de présenter les outils fondamentaux sur
les graphes orientés bien-fondés (cf. \ref{definitions-graphs})
utiles à la théorie combinatoire des jeux impartiaux.  Il
s'agit notamment de la théorie de l'induction bien-fondée
(cf. \ref{scholion-well-founded-induction}
et \ref{scholion-well-founded-definition}).

\subsection{Graphes orientés bien-fondés}\label{subsection-well-founded-graphs}

\begin{defn}\label{definitions-graphs}
Un \defin{graphe orienté [simple]} est la donnée d'un ensemble $G$ et
d'une partie $E$ de $G^2$ ne rencontrant pas la diagonale (i.e., un
ensemble de couples $(x,y)$ avec $x\neq y$) : si on préfère, il s'agit
d'un ensemble $G$ muni d'une relation $E$ irréflexive.  Les éléments
de $G$ s'appellent \defin[sommet]{sommets} et les éléments de $E$
\defin[arête]{arêtes} de $G$, et si $(x,y) \in E$, on dit qu'il y a une
arête allant du sommet $x$ au sommet $y$, ou arête de source $x$ et de
cible $y$, ou encore que $y$ est \defin[atteindre]{atteint} par une arête de
source $x$, ou encore que $y$ est un \index{sortant (voisin)}\defin{voisin sortant} de $x$,
et on notera $\outnb(x) = \{y : (x,y) \in E\}$ l'ensemble des voisins
sortants de $x$.  Un sommet qui n'a pas de voisin sortant est
appelé \defin{puits} dans $G$.

Un tel graphe est dit \defin[fini (graphe)]{fini} lorsque $G$ est fini (il est clair
que $E$ l'est alors aussi).  Il est dit \defin[acyclique (graphe)]{acyclique} lorsqu'il
n'existe pas de suite finie (« cycle ») $x_0,\ldots,x_{n-1}$ de
sommets telle que $(x_i,x_{i+1})$ soit une arête pour chaque $0\leq
i\leq n-1$, où on convient que $x_n = x_0$.

Un graphe orienté (possiblement infini) est dit \defin[bien-fondé (graphe)]{bien-fondé} ou
\defin[progressivement fini|see{bien-fondé}]{progressivement fini} lorsqu'il n'existe pas de suite
$x_0,x_1,x_2,\ldots$ de sommets telle que $(x_i,x_{i+1})$ soit une
arête pour tout $i\in\mathbb{N}$ (i.e., aucun cycle ni chemin infini,
cf. ci-dessous).
\end{defn}

\thingy\label{finite-graph-is-well-founded-iff-acyclic} Il est
évident que tout graphe bien-fondé est acyclique (s'il
existe un cycle $x_0,\ldots,x_{n-1}$, on en déduit une suite infinie
en posant $x_i = x_{i\mod n}$) ; pour un graphe \emph{fini}, la
réciproque est vraie : en effet, s'il existe une suite infinie
$x_0,x_1,x_2,\ldots$ avec une arête de $x_i$ à $x_{i+1}$ pour
chaque $i$, il doit exister $p<q$ tels que $x_q = x_p$, et on obtient
alors un cycle $x_p,\ldots,x_{q-1}$.  En général, cependant, les
notions sont distinctes, l'exemple le plus évident (de graphe
acyclique mais mal fondé) étant sans doute celui de $\mathbb{N}$ dans
lequel on fait pointer une arête de $i$ à $i+1$ pour chaque $i$.

\begin{defn}\label{definition-accessibility-downstream}
Si $G$ est un graphe orienté on appelle \defin[accessibilité]{relation
  d'accessibilité} la clôture réflexive-transitive de la relation
donnée par les arêtes de $G$ : autrement dit, on dit que $y$ est
accessible à partir de $x$ lorsqu'il existe $x {=}
x_0,x_1,\ldots,x_{n-1},x_n {=} y$ tels que pour chaque $i$ le sommet
$x_{i+1}$ soit voisin sortant de $x_i$ (on autorise $n=0$,
c'est-à-dire que chaque sommet est toujours accessible à partir de
lui-même).

On pourra aussi introduire la relation d'\defin{accessibilité
  stricte} qui est la clôture transitive : autrement dit, la
différence avec l'accessibilité définie ci-dessus est qu'on
impose $n>0$, i.e., on ne relie pas un sommet à lui-même.

L'ensemble des sommets accessibles à partir d'un sommet $x$
s'appellera aussi l'\defin{aval} de $x$ et pourra se noter
$\downstr(x)$ (c'est donc la plus petite partie $P$ de $G$ telle que
$x\in P$ et $y\in P \limp \outnb(y)\subseteq P$,
cf. \ref{definition-downstream-closed-inductive}).  On peut considérer
l'aval de $x$ comme un sous-graphe induit de $G$ (c'est-à-dire,
considérer le graphe dont l'ensemble des sommets est l'aval de $x$ et
dont les arêtes sont celles qui partent d'un tel sommet).  On
remarquera la convention faite que $x$ appartient à son propre aval.
\end{defn}

\thingy On peut remarquer que la relation d'accessibilité sur $G$ est
antisymétrique (i.e., est une relation d'ordre partiel large) si et
seulement si $G$ est acyclique : l'accessibilité stricte est alors
l'ordre strict correspondant.  Lorsque $G$ est bien-fondé, la relation
d'accessibilité stricte est elle-même bien-fondée (au sens où le
graphe qu'elle définit est bien-fondé) : si on la voit comme une
relation d'ordre partiel ($x>y$ signifiant que $y$ est accessible à
partir de $x$), cela signifie qu'il n'y a pas de suite infinie
strictement décroissante.

Une relation d'ordre \emph{total} (strict) $>$ qui soit bien-fondée,
i.e., telle qu'il n'existe pas de suite infinie strictement
décroissante, est appelée un \defin{bon ordre}, ou définir un
ensemble \defin[bien-ordonné (ensemble)]{bien-ordonné}.

\begin{defn}\label{definition-downstream-closed-inductive}
Si $G$ est un graphe orienté, on dira qu'un ensemble $P$ de sommets de
$G$ est \defin{aval-clos} lorsqu'il vérifie la propriété suivante :
« si $x$ est dans $P$ alors tout voisin sortant de $x$ est dans $P$ »
(soit $x\in P \limp \outnb(x)\subseteq P$ ; ou de façon équivalente,
« tout sommet accessible à partir d'un sommet de $P$ est encore
  dans $P$ »,).

Réciproquement, on dira qu'un ensemble $P$ de sommets de $G$ est
\defin{aval-inductif} lorsqu'il vérifie la propriété suivante : « si
  $x \in G$ est tel que tout voisin sortant de $x$ appartient à $P$,
  alors $x$ lui-même appartient à $P$ » (i.e. « $P$ contient tout
  sommet dont tous les voisins sortants sont dans $P$ »,
soit $\outnb(x)\subseteq P \limp x\in P$).
\end{defn}

\thingy\label{trivial-remark-downstream} Il est clair qu'une
intersection ou réunion quelconque d'ensembles aval-clos est encore
aval-close.  L'aval de $x$ (ensemble des sommets accessibles
depuis $x$) est toujours aval-clos, c'est même la plus petite partie
aval-close contenant $x$, et il est facile de se convaincre qu'un
ensemble de sommets est aval-clos si et seulement si il est une
réunion d'avals.

Pour ce qui est des ensembles aval-inductifs, il est clair qu'une
intersection quelconque d'ensembles aval-inductifs est aval-inductive.
Leur nature, au moins dans un graphe bien-fondé, va être précisée dans
ce qui suit, et ceci justifiera le terme d'« aval-inductif ».

\begin{prop}[induction bien-fondée]\label{well-founded-induction}
Pour un graphe orienté $G$, les affirmations suivantes sont
équivalentes :

\begin{itemize}
\item[(*)]$G$ est bien-fondé.
\item[(\dag)]Tout ensemble \emph{non vide} $N$ de sommets de $G$
  contient un sommet $x \in N$ qui est un puits pour $N$,
  c'est-à-dire qu'il n'existe aucune arête $(x,y)$ de $G$ avec $y \in
  N$ (i.e., aucun voisin sortant de $x$ n'appartient à $N$).
\item[(\ddag)]Si une partie $P\subseteq G$ vérifie la propriété
  suivante « si $x \in G$ est tel que tout voisin sortant de $x$
    appartient à $P$, alors $x$ lui-même appartient à $P$ » (i.e.,
  « $P$ est aval-inductif »,
  cf. \ref{definition-downstream-closed-inductive}), alors $P = G$.
\end{itemize}
\end{prop}
\begin{proof}
L'équivalence entre (\dag) et (\ddag) est complètement formelle : elle
s'obtient en posant $P = G\setminus N$ ou réciproquement $N =
G\setminus P$, en passant à la contraposée, et en passant aussi à la
contraposée à l'intérieur de la propriété d'être aval-inductif (entre
guillemets dans (\ddag)), et encore une fois dans la prémisse de cette
propriété (« tout voisin sortant de $x$ appartient à $P$ » équivaut à
« aucun voisin sortant de $x$ n'appartient à $N$ », i.e., « $x$ est un
  puits pour $N$ »).

Pour montrer que (\dag) implique (*), il suffit d'appliquer (\dag) à
l'ensemble $N := \{x_0,x_1,x_2,\ldots\}$ des sommets d'une suite telle
qu'il y ait une arête de $x_i$ à $x_{i+1}$.

Pour montrer que (*) implique (\dag), on suppose que $N$ est un
ensemble non-vide de sommets sans puits [i.e., puits pour $N$] : comme
$N$ est non-vide, on choisit $x_0 \in N$, et comme $x_0$ n'est pas un
puits on peut choisir $x_1 \in N$ atteignable à partir de $x_0$ par
une arête, puis $x_2 \in N$ atteignable à partir de $x_1$ et ainsi de
suite — par récurrence (et par l'axiome du choix [dépendants]), on
construit ainsi une suite $(x_i)$ de sommets telle qu'il y ait une
arête de $x_i$ à $x_{i+1}$.
\end{proof}

La définition (*) choisie pour un graphe bien-fondé est la plus
compréhensible, mais en un certain sens la définition (\ddag) est « la
  bonne » (en l'absence de l'axiome du choix, il faut utiliser (\dag)
ou (\ddag), et en mathématiques constructives il faut
utiliser (\ddag)).  En voici une traduction informelle :

\begin{scho}\label{scholion-well-founded-induction}
Pour montrer une propriété $P$ sur les sommets d'un graphe bien-fondé,
on peut supposer (comme « hypothèse d'induction »), lorsqu'il s'agit
de montrer que $x$ a la propriété $P$, que cette propriété est déjà
acquise pour tous les voisins sortants de $x$.
\end{scho}

Exactement comme le principe de récurrence sur les entiers naturels,
le principe d'induction bien-fondée peut servir non seulement à
démontrer des propriétés sur les graphes bien-fondés, mais aussi à
définir des fonctions dessus :

\begin{thm}[définition par induction bien-fondée]\label{well-founded-definition}
Soit $G$ un graphe orienté bien-fondé et $Z$ un ensemble quelconque.
Notons $\outnb(x) = \{y : (x,y) \in E\}$ l'ensemble des
voisins sortants d'un sommet $x$ de $G$ (i.e., des $y$ atteints par une
arête de source $x$).

Appelons $\mathscr{F}$ l'ensemble des couples $(x,f)$ où $x\in G$ et
$f$ une fonction de l'ensemble des voisins sortants de $x$ vers $Z$
(autrement dit, $\mathscr{F}$ est $\bigcup_{x \in G} \big(\{x\}\times
Z^{\outnb(x)}\big)$).  Soit enfin $\Phi\colon \mathscr{F} \to Z$
une fonction quelconque.  Alors il existe une unique fonction $f\colon
G \to Z$ telle que pour tout $x \in G$ on ait
\[
f(x) = \Phi(x,\, f|_{\outnb(x)})
\]
\end{thm}
\begin{proof}
Montrons d'abord l'unicité : si $f$ et $f'$ vérifient toutes les deux
la propriété annoncée, soit $P$ l'ensemble des sommets $x$ de $G$ tels
que $f(x) = f'(x)$.  Si $x \in G$ est tel que $\outnb(x)
\subseteq P$, c'est-à-dire que $f|_{\outnb(x)} =
f'|_{\outnb(x)}$, alors $f(x) = \Phi(x,\,
f|_{\outnb(x)}) = \Phi(x,\, f'|_{\outnb(x)}) = f'(x)$,
autrement dit, $x\in P$.  La phrase précédente affirme précisément que
$P$ vérifie la propriété entre guillemets dans (\ddag)
de \ref{well-founded-induction}, et d'après la proposition en
question, on a donc $P = G$, c'est-à-dire $f = f'$.  Ceci montre
l'unicité.

Pour montrer l'existence, on considère l'ensemble $\mathfrak{E}$ des
fonctions $e\colon H\to Z$ définies sur une partie aval-close $H
\subseteq G$ et telles que pour tout $e(x) = \Phi(x, e|_{\outnb(x)})$
pour tout $x\in H$.  Si $e,e' \in \mathfrak{E}$ alors $e$ et $e'$
coïncident là où toutes deux sont définies, comme le montre l'unicité
qu'on a montrée (appliquée à $e$ et $e'$ sur l'ensemble aval-clos $H
\cap H'$ de définition commun de $e$ et $e'$).  En particulier, la
réunion [des graphes] de tous les $e\in\mathfrak{E}$ définit encore un
élément $f$ de $\mathfrak{E}$, maximal pour le prolongement.  Soit $P$
l'ensemble des $x \in G$ où $f$ est définie.  Si $P$ contient (i.e.,
$f$ est définie sur) tous les voisins sortants d'un certain $x\in G$,
alors $f$ est nécessairement définie aussi en $x$, sans quoi on
pourrait l'y prolonger par $f(x) = \Phi(x,\, f|_{\outnb(x)})$, ce qui
contredirait la maximalité de $f$.  Par induction bien-fondée, on en
conclut $P = G$, c'est-à-dire que $f$ est définie sur $G$ tout entier.
C'est ce qu'on voulait.
\end{proof}

Ce théorème est difficile à lire.  En voici une traduction
informelle :

\begin{scho}\label{scholion-well-founded-definition}
Pour définir une fonction $f$ sur un graphe bien-fondé, on peut
supposer, lorsqu'on définit $f(x)$, que $f$ est déjà défini (i.e.,
connu) sur tous les voisins sortants de $x$ : autrement dit, on
peut librement utiliser la valeur de $f(y)$ sur chaque sommet $y$
voisin sortant de $x$, dans la définition de $f(x)$.
\end{scho}

Voici un exemple d'application de la définition par induction
bien-fondée :

\begin{defn}\label{definition-well-founded-rank}
Soit $G$ un graphe orienté bien-fondé dans lequel chaque sommet n'a
qu'un nombre fini de voisins sortants (par exemple, un graphe fini
acyclique).  En utilisant le théorème \ref{well-founded-definition},
on définit alors une fonction $\rk\colon G \to \mathbb{N}$ par
$\rk(x) = \max\{\rk(y) : y\in\outnb(x)\} + 1$ où il est convenu que
$\max\varnothing = -1$ ; formellement, c'est-à-dire qu'on pose
$\Phi(x, r) = \max\{r(y) : y\in\outnb(x)\} + 1$ avec $\Phi(x, r) = 0$
si $x$ est un puits, et qu'on appelle $\rk$ la fonction telle que
$\rk(x) = \Phi(x, \rk|_{\outnb(x)})$ dont l'existence et l'unicité
sont garanties par le théorème.  Cette fonction $\rk$ s'appelle la
\defin[rang]{fonction rang} sur $G$, on dit que $\rk(x)$ est le rang (ou
rang bien-fondé) d'un sommet $x$.  (Voir
aussi \ref{ordinal-valued-rank-and-grundy-function} pour une
généralisation.)
\end{defn}

\thingy Autrement dit, un sommet de rang $0$ est un puits,
un sommet de rang $1$ est un sommet non-puits dont tous les
voisins sortants sont terminaux, un sommet de rang $2$ est un sommet dont
tous les voisins sortants sont de rang $\leq 1$ mais et au moins un est de
rang exactement $1$, et ainsi de suite.

Il revient au même de définir le rang de la manière suivante : le rang
$\rk(x)$ d'un sommet $x$ d'un graphe orienté bien-fondé est la plus
grande longueur possible d'un chemin orienté partant de $x$,
c'est-à-dire, le plus grand $n$ tel qu'il existe une suite
$x_0,x_1,\ldots,x_n$ telle que $x_0 = x$ et que pour chaque $i$ le
sommet $x_{i+1}$ soit atteint par une arête de source $x_i$.

Voici un autre exemple de définition par induction bien-fondée :

\begin{defn}\label{definition-grundy-function}
Soit $G$ un graphe orienté bien-fondé dans lequel chaque sommet n'a
qu'un nombre fini de voisins sortants.  En utilisant le
théorème \ref{well-founded-definition}, on définit alors une fonction
$\gr\colon G \to \mathbb{N}$ par $\gr(x) = \mex\{\gr(y) :
y\in\outnb(x)\}$ où, si $S\subseteq\mathbb{N}$, on note \index{mex}$\mex S
:= \min(\mathbb{N}\setminus S)$ pour le plus petit entier naturel
\emph{n'appartenant pas} à $S$ ; formellement, c'est-à-dire qu'on pose
$\Phi(x, g) = \mex\{g(y) : y\in\outnb(x)\}$ et qu'on appelle
$\gr$ la fonction telle que $\gr(x) = \Phi(x,
\gr|_{\outnb(x)})$ dont l'existence et l'unicité sont
garanties par le théorème.  Cette fonction $\gr$ s'appelle la
\defin[Grundy (fonction de)]{fonction de Grundy} sur $G$, on dit que $\gr(x)$ est la
valeur de Grundy d'un sommet $x$.  (Voir
aussi \ref{ordinal-valued-rank-and-grundy-function} pour une
généralisation.)
\end{defn}

\thingy En particulier, un sommet de valeur de Grundy $0$ est un
sommet qui n'a que des sommets de valeur de Grundy $>0$ comme voisins
sortants (ceci inclut le cas particulier d'un puits), tandis qu'un
sommet de valeur de Grundy $>0$ est un sommet ayant au moins un sommet
de valeur de Grundy $0$ comme voisin sortant.

On verra que la notion de fonction de Grundy, et particulièrement le
fait que la valeur soit nulle ou pas, a énormément d'importance dans
l'étude de la théorie des jeux impartiaux.  On verra aussi comment la
définir sans l'hypothèse que chaque sommet n'a qu'un nombre fini de
voisins sortants (mais ce ne sera pas forcément un entier naturel :
cf. \ref{ordinal-valued-rank-and-grundy-function}).  En attendant,
peut se passer de cette hypothèse pour définir isolément l'ensemble
des sommets de valeur de Grundy $0$ :

\begin{defn}\label{definition-grundy-kernel}
Soit $G$ un graphe orienté bien-fondé.  En utilisant le
théorème \ref{well-founded-definition}, on définit alors une partie $P
\subseteq G$ telle que $x \in P$ ssi $\outnb(x) \cap P =
\varnothing$ ; formellement, c'est-à-dire que pour $f\colon \outnb(x)
\to \{\mathtt{P},\mathtt{N}\}$ on définit $\Phi(x, g)$ comme valant
$\mathtt{N}$ si $g$ prend la valeur $\mathtt{P}$, et $\mathtt{P}$ si
$g$ vaut constamment $\mathtt{N}$ (y compris si $g$ est la fonction
vide), et qu'on appelle $f$ la fonction telle que $f(x) = \Phi(x,
f|_{\outnb(x)})$ dont l'existence et l'unicité sont garanties par le
théorème, et enfin on pose $P = \{x \in G : f(x) = \mathtt{P}\}$.

Les éléments de $P$ (i.e., les $x$ tels que $f(x) = \mathtt{P}$)
seront appelés les \defin[P-position]{P-sommets} (ou P-positions) de $G$, tandis
que les éléments du complémentaire $G\setminus P$ (i.e., les $x$ tels
que $f(x) = \mathtt{N}$) seront appelés \defin[N-position]{N-sommets} (ou
N-positions) de $G$ : ainsi, \emph{un P-sommet est un sommet dont tous
  les voisins sortants sont des N-sommets, et un N-sommet est un
  sommet qui a au moins un P-sommet pour voisin sortant}.
\end{defn}

\thingy\label{discussion-n-and-p-vertices} Dans un jeu combinatoire comme exposé en
\ref{introduction-graph-game} et/ou
\ref{definition-impartial-combinatorial-game}, si le graphe est
bien-fondé, les P-sommets sont les positions du jeu à partir
desquelles le joueur Précédent (=second joueur) a une stratégie
gagnante, tandis que les N-sommets sont celles à partir desquelles le
joueur suivant (`N' comme « Next », =premier joueur) a une stratégie
gagnante (consistant, justement, à jouer vers un P-sommet) : ceci
résulte de \ref{strategies-forall-exists-lemma}
ou \ref{strategies-forall-exists-reformulation} (ou encore
de \ref{positional-strategies-forall-exists-lemma} dans le cadre plus
général de la
section \ref{subsection-positional-strategies-determinacy}).

On peut donc résumer ainsi la stratégie gagnante « universelle » pour
n'importe quel jeu combinatoire impartial à connaissance parfaite dont
le graphe est bien-fondé : \emph{jouer vers un P-sommet}, après quoi
l'adversaire sera obligé de jouer vers un N-sommet (si tant est qu'il
peut jouer), et on pourra de nouveau jouer vers un P-sommet, et ainsi
de suite (et comme le graphe est bien-fondé, le jeu termine forcément
en temps fini, et celui qui a joué pour aller vers les P-sommets aura
gagné puisqu'il peut toujours jouer).

\thingy Pour illustrer la technique de démonstration par induction
bien-fondée, montrons que si $\gr$ est la fonction de Grundy
introduite en \ref{definition-grundy-function} et si $P$ est la partie
introduite en \ref{definition-grundy-kernel}, alors on a $x \in P$ si
et seulement si $\gr(x) = 0$, i.e., les P-sommets sont ceux pour
lesquels la fonction de Grundy est nulle.

Par induction bien-fondé (cf. \ref{scholion-well-founded-induction}),
on peut supposer la propriété (« $x \in P$ si et seulement si
  $\gr(x) = 0$ ») déjà connue pour tous les voisins sortants $y$ du
sommet $x$ où on cherche à la vérifier.  Si $x \in P$, cela signifie
par définition de $P$ que tous ses voisins sortants $y$ de $x$ sont
dans $G \setminus P$, et d'après l'hypothèse d'induction cela signifie
$\gr(y) \neq 0$ pour tout $y \in \outnb(x)$, c'est-à-dire
$\mex\{\gr(y) : y\in\outnb(x)\} = 0$ (puisque $\mex S = 0$ si et
seulement si $0 \not\in S$), ce qui signifie $\gr(x) = 0$.  Toutes
ces reformulations étaient des équivalences : on a bien montré que $x
\in P$ équivaut à $\gr(x) = 0$, ce qui conclut l'induction.

\thingy\label{ordinal-valued-rank-and-grundy-function} En anticipant
sur la notion d'ordinaux introduite dans la
partie \ref{section-ordinals}, la fonction rang peut se généraliser à
n'importe quel graphe bien-fondé (sans hypothèse de nombre fini de
voisins sortants), à condition d'autoriser des valeurs ordinales
quelconques : précisément, on définit $\rk(x) = \sup\{\rk(y)+1 :
y\in\outnb(x)\}$ (avec cette fois $\sup\varnothing = 0$ pour garder
$\rk(x) = 0$ si $x$ est un puits) c'est-à-dire que $\rk(x)$ est le
plus petit ordinal strictement supérieur à $\rk(y)$ pour tout
$y\in\outnb(x)$.

De même, la fonction de Grundy peut se généraliser à n'importe quel
graphe bien-fondé en définissant $\gr(x) = \mex\{\gr(y) :
y\in\outnb(x)\}$ où \index{mex}$\mex S$ désigne le plus petit ordinal
\emph{n'appartenant pas} à $S$
(voir \ref{definition-grundy-function-again} pour une écriture
formelle de cette définition).  Il reste vrai (avec la même
démonstration) que $\gr(x)$ est nul si et seulement si $x$ est un
P-sommet.

On peut donc résumer ainsi la stratégie gagnante « universelle » pour
n'importe quel jeu combinatoire impartial à connaissance parfaite dont
le graphe est bien-fondé : \emph{jouer de façon à annuler la fonction
  de Grundy} (c'est-à-dire, jouer vers un P-sommet), après quoi
l'adversaire sera obligé de jouer vers un sommet dont la valeur de
Grundy est non-nulle, et on pourra de nouveau jouer vers zéro, et
ainsi de suite (et comme le graphe est bien-fondé, le jeu termine
forcément en temps fini, et celui qui a joué pour annuler la fonction
de Grundy aura gagné puisqu'il peut toujours jouer).

\thingy\label{well-founded-induction-algorithm} Lorsque $G$ est
\emph{fini} et bien-fondé, i.e., est un graphe fini acyclique
(cf. \ref{finite-graph-is-well-founded-iff-acyclic}), la fonction $f$
définie par le théorème \ref{well-founded-definition} peut se calculer
algorithmiquement de la façon suivante (si tant est que $\Phi$
elle-même est calculable) :
\begin{itemize}
\item utiliser un algorithme de \defin{tri topologique} (en ayant
  inversé le sens des arêtes pour se ramener à la convention usuelle)
  pour trouver une numérotation des sommets de $G$ telle que pour
  toute arête $(x,y)$ de $G$ le sommet $y$ précède $x$ dans la
  numérotation ;
\item parcourir tous les éléments $x\in G$ dans l'ordre de cette
  numérotation, et définir $f(x) = \Phi(x, f|_{\outnb(x)})$, qui aura
  toujours un sens car tous les éléments de $\outnb(x)$ auront été
  parcourus avant $x$.
\end{itemize}
Si le tri topologique a été effectué par parcours en largeur, il est
compatible avec la fonction rang (et inversement, tout tri raffinant
la fonction rang est un tri topologique et peut s'obtenir par parcours
en largeur).  La notion de rang ordinal
(cf. \ref{ordinal-valued-rank-and-grundy-function}) est donc une forme
de généralisation transfinie du tri topologique.

Un algorithme évitant le tri topologique, au prix d'une plus grande
complexité, mais ayant le bénéfice de fonctionner dans une situation
plus générale (graphes non nécessairement bien-fondés/acycliques),
sera exposé dans la section suivante
(cf. \ref{non-well-founded-induction-algorithm}).



\subsection{Généralisations aux graphes non nécessairement bien-fondés}

\begin{defn}\label{definition-wfpart}
L'ensemble des sommets d'un graphe orienté dont l'aval est bien-fondé,
autrement dit, l'ensemble des sommets $x$ tels qu'il n'existe pas de
suite $x_0,x_1,x_2,\ldots$ de sommets où $x_0 = x$ et où pour
chaque $i$ le sommet $x_{i+1}$ est atteint par une arête de
source $x_i$, est appelé la \defin{partie bien-fondée} du graphe.
\end{defn}

\thingy\label{trivial-remark-wfpart} Il sera utile pour la suite de
remarquer que la partie bien-fondée de $G$ est à la fois aval-close et
aval-inductive (car on peut construire une suite infinie $x_0, x_1,
x_2\ldots$, avec $x_{i+1}$ voisin sortant de $x_i$, commençant par un
$x_0$ donné si et seulement si on peut le faire en commençant pour un
certain voisin sortant $x_1$ de ce $x_0$).

\begin{prop}\label{wfpart-is-smallest-inductive}
Si $G$ est un graphe orienté non supposé bien-fondé, la partie
bien-fondée de $G$ est la plus petite (pour l'inclusion) partie $P$
aval-inductive de $P$ (i.e., vérifiant la propriété « si $x \in
  G$ est tel que tout voisin sortant de $x$ appartient à $P$, alors
  $x$ lui-même appartient à $P$ »,
cf. \ref{definition-downstream-closed-inductive}).
\end{prop}
\begin{proof}
La plus petite partie aval-inductive de $G$ a bien un sens, car
l'intersection de toutes les parties aval-inductives est encore
aval-inductive (cf. \ref{trivial-remark-downstream}).

Si $x$ est un sommet de $G$ et $\downstr(x)$ désigne son aval
(considéré comme sous-graphe induit de $G$), il est clair que pour
toute partie $P$ aval-inductive de $G$, la partie $P \cap \downstr(x)$
de $\downstr(x)$ est aval-inductive dans ce dernier (le point
important étant que les voisins sortants d'un sommet de $\downstr(x)$
dans ce dernier sont les mêmes que ceux dans $G$).  En particulier, si
$\downstr(x)$ est bien-fondé (c'est-à-dire, si $x$ appartient à la
partie bien-fondée de $G$), alors $x$ appartient à toute partie
aval-inductive de $G$.

Mais réciproquement, la partie bien-fondée de $G$ est elle-même
aval-inductive (car si $\downstr(y)$ est bien-fondé pour tout voisin
sortant de $x$, il est clair que $\downstr(x)$ est aussi bien-fondé,
cf. \ref{trivial-remark-wfpart}), donc un sommet qui appartient à
toute partie aval-inductive de $G$ est, en particulier, dans la partie
bien-fondée de $G$.
\end{proof}

\thingy Pour pouvoir énoncer le théorème suivant, on aura besoin de
faire les rappels, définitions et remarques suivants.  Une
\defin[partielle (fonction)]{fonction partielle} $A \dasharrow Z$, où $A$ est un ensemble
quelconque, n'est rien d'autre qu'une fonction définie $D \to Z$ sur
une partie $D\subseteq A$ de $Z$ (appelée \defin[définition (ensemble de)]{ensemble de
  définition} de la partie).  Si $f,g\colon A \dasharrow Z$ sont deux
fonctions partielles, on dit que $f$ \defin{prolonge} $g$ et on note
$f \supseteq g$ ou $g\subseteq f$, lorsque l'ensemble de définition
$D_f$ de $f$ contient celui $D_g$ de $g$ et que $f$ et $g$ coïncident
sur $D_f$ (ceci signifie bien que $f \supseteq g$ en tant qu'ensembles
si on identifie une fonction avec son graphe).  Il s'agit bien sûr là
d'un ordre partiel (sur l'ensemble des fonctions partielles $A
\dasharrow Z$).

Enfin, si $\Phi$ est une fonction partielle elle-même définie sur
l'ensemble des fonctions partielles $A \dasharrow Z$ (cet ensemble est
$\bigcup_{D\subseteq A} Z^D$, si on veut), on dit que $\Phi$ est
\defin[cohérente (fonction)]{cohérente} lorsque $\Phi(f) = \Phi(g)$ à chaque fois que $f$
prolonge $g$ et que $\Phi(g)$ est définie (autrement dit, une fois que
$\Phi$ est définie sur une fonction partielle $g$, elle est définie
sur tout prolongement de $g$ et y prend la même valeur que sur $g$ ;
intuitivement, il faut s'imaginer que si $g$ apporte assez
d'information pour décider la valeur de $\Phi(g)$, toute information
supplémentaire reste cohérente avec cette valeur).

\begin{thm}\label{non-well-founded-definition}
Soit $G$ un graphe orienté et $Z$ un ensemble quelconque.
Notons $\outnb(x) = \{y : (x,y) \in E\}$ l'ensemble des
voisins sortants d'un sommet $x$ de $G$ (i.e., des $y$ atteints par une
arête de source $x$).

Appelons $\mathscr{F}$ l'ensemble des couples $(x,f)$ où $x\in G$ et
$f$ une fonction \emph{partielle} de l'ensemble des voisins sortants
de $x$ vers $Z$ (autrement dit, $\mathscr{F}$ est $\bigcup_{x \in G}
\bigcup_{D \subseteq \outnb(x)} \big(\{x\}\times Z^D\big)$).  Soit
enfin $\Phi\colon \mathscr{F} \dasharrow Z$ une fonction partielle
cohérente en la deuxième variable, c'est-à-dire telle que $\Phi(x,f) =
\Phi(x,g)$ dès que $f \supseteq g$ et que $\Phi(x,g)$ est définie.
Alors il existe une plus petite (au sens du prolongement) fonction
partielle $f\colon G \dasharrow Z$ telle que pour tout $x \in G$ on
ait
\[
f(x) = \Phi(x,\, f|_{\outnb(x)})
\]
(au sens où chacun des deux membres est défini ssi l'autre l'est, et
dans ce cas ils ont la même valeur).

Si $\Phi(x,g)$ est défini à chaque fois que $g$ est totale, alors la
fonction $f$ qu'on vient de décrire est définie \emph{au moins} sur la
partie bien-fondée de $G$.
\end{thm}
\begin{proof}
Soit $\mathscr{D}$ l'ensemble des fonctions partielles $f \colon G
\dasharrow Z$.  Pour $f \in \mathscr{D}$, on définit $\Psi(f)$ comme
la fonction partielle $x \mapsto \Phi(x, f|_{\outnb(x)})$.  Remarquons
que si $f$ prolonge $g$ dans $\mathscr{D}$ alors $\Psi(f)$
prolonge $\Psi(g)$ (c'est une traduction de la cohérence supposée
sur $\Phi$).  Le lemme suivant (appliqué à $X=G$, les autres notations
étant inchangées) permet de conclure à l'existence de $f$.

Pour ce qui est de la dernière affirmation, on procède par induction
bien-fondée sur la partie bien-fondée de $G$ : si $f|_{\outnb(x)}$ est
totale, i.e., si $f$ est définie sur chaque voisin sortant de $x$,
alors l'hypothèse faite sur $\Phi$ assure que $f(x)$ est définie, et
l'induction bien-fondée ((\ddag) de \ref{well-founded-induction}
appliqué à l'intersection $P$ de la partie bien-fondée de $G$ et de
l'ensemble de définition de $f$) montre alors que $f$ est définie
partout sur la partie bien-fondée de $G$.
\end{proof}

La démonstration du théorème repose crucialement sur le lemme suivant,
dont on va donner deux démonstrations :
\begin{lem}\label{fixed-point-lemma-for-partial-functions}
Soient $X$ et $Z$ deux ensembles quelconques : notons $\mathscr{D}$
l'ensemble des fonctions partielles $X\dasharrow Z$, qu'on verra comme
des parties de $X\times Z$ ne contenant jamais deux couples $(x,z_1)$
et $(x,z_2)$ avec la même première coordonnée.  (Lorsque $f\supseteq
g$, on dit que « $f$ prolonge $g$ ».)

Soit $\Psi \colon \mathscr{D} \to \mathscr{D}$ une fonction (totale !)
vérifiant : $\Psi$ est \emph{croissante} pour l'inclusion,
c'est-à-dire que si $f$ prolonge $g$, alors $\Psi(f)$
prolonge $\Psi(g)$.  Alors il existe une plus petite (au sens du
prolongement) fonction partielle $f \in \mathscr{D}$ telle que
$\Psi(f) = f$.
\end{lem}
\begin{proof}[Première démonstration]
Montrons d'abord que \emph{si} il existe une fonction partielle $f \in
\mathscr{D}$ telle que $\Psi(f) = f$, ou même simplement $\Psi(f)
\subseteq f$, alors il en existe une plus petite.  Pour cela, il
suffit de considérer l'intersection $h$ de toutes les $f$ telles que
$\Psi(f) \subseteq f$ (considérées comme des parties de $X\times Z$) :
dès lors qu'il existe au moins un $f \in \mathscr{D}$ tel que $\Psi(f)
\subseteq f$, cette intersection $h$ est bien définie et est bien un
élément de $\mathscr{D}$.  Si $\Psi(f) \subseteq f$ alors $h \subseteq
f$ (puisque $h$ est l'intersection des $f$), donc $\Psi(h) \subseteq
\Psi(f)$ (par croissance de $\Psi$), donc $\Psi(h) \subseteq f$ (par
transitivité), et comme ceci est vrai pour tous les $f$ dont
l'intersection est $h$, on a finalement $\Psi(h) \subseteq h$ ; mais
la croissance de $\Psi$ donne alors aussi $\Psi(\Psi(h)) \subseteq
\Psi(h)$, et du coup $\Psi(h)$ fait partie des $f$ qu'on a
intersectées pour former $h$, et on a ainsi $h \subseteq \Psi(h)$ ;
bref, $\Psi(h) = h$, et $h$ est à la fois le plus petit élément $f \in
\mathscr{D}$ vérifiant $\Psi(f) \subseteq f$ (de par sa construction)
et le plus petit vérifiant $\Psi(f) = f$ (puisqu'il vérifie cette
propriété).

Reste à montrer qu'il existe bien une fonction partielle $f$ telle que
$\Psi(f) = f$, ou même simplement $\Psi(f) \subseteq f$.  Pour cela,
on introduit l'ensemble $\mathscr{E}$ des $f \in \mathscr{D}$ qui
vérifient $\Psi(f) \supseteq f$ (inclusion dans le sens opposé du
précédent !).  Notons que $\mathscr{E}$ n'est pas vide puisque
$\varnothing \in \mathscr{E}$ (où $\varnothing$ est la fonction vide,
définie nulle part).

Soit maintenant $\mathfrak{M}$ l'ensemble des applications (totales !)
$T\colon\mathscr{E}\to\mathscr{E}$ qui vérifient (i) $T(f) \supseteq
f$ pour tout $f\in \mathscr{E}$ et (ii) si $f \supseteq g$ alors $T(f)
\supseteq T(g)$.  Ainsi $\id_{\mathscr{E}} \in \mathfrak{M}$
(trivialement) et $\Psi \in \mathfrak{M}$ (par définition de
$\mathscr{E}$ et par croissance de $\Psi$) ; et si $T, T' \in
\mathfrak{M}$ on a $T'\circ T \in \mathfrak{M}$ (en notant $T'\circ T$
la composée).  L'observation suivante sera cruciale : si $g \in
\mathscr{E}$ et $T, T' \in \mathfrak{M}$, alors on a à la fois
$(T'\circ T)(g) \supseteq T(g)$ (d'après (i) pour $T'$) et $(T'\circ
T)(g) \supseteq T'(g)$ (d'après (i) pour $T$ et (ii) pour $T'$).

Affirmation : si $g \in \mathscr{E}$ alors la réunion des $T(g)$ pour
tous les $T\in\mathfrak{M}$ est, en fait, une fonction partielle.  En
effet, l'observation faite ci-dessus montre que si $T, T' \in
\mathfrak{M}$ alors les fonctions partielles $T(g)$ et $T'(g)$ sont
toutes deux restrictions d'une même fonction partielle $(T'\circ
T)(g)$, donc il ne peut pas y avoir de conflit entre leurs valeurs (au
sens où si toutes les deux sont définies en un $x\in X$, elles y
coïncident) — c'est exactement ce qui permet de dire que la réunion
est encore une fonction partielle.  Notons $U(g) :=
\bigcup_{T\in\mathfrak{M}} T(g)$ cette réunion.  On a au moins $U(g)
\in \mathscr{D}$.  Mais en fait, comme $U(g)$ prolonge tous les
$T(g)$, la croissance de $\Psi$ assure que $\Psi(U(g))$ prolonge tous
les $\Psi(T(g))$, qui prolongent eux-mêmes les $T(g)$ (puisque $T(g)
\in \mathscr{E}$), bref $\Psi(U(g)) \supseteq U(g)$ et ainsi $U(g) \in
\mathscr{E}$.

Mais alors $U \in \mathfrak{M}$ (on vient de voir que $U$ est une
fonction $\mathscr{E}\to\mathscr{E}$, et les propriétés (i) et (ii)
sont claires).  En particulier, $\Psi\circ U \in \mathfrak{M}$, donc
$(\Psi\circ U)(g)$ fait partie des $T(g)$ dont $U(g)$ est la réunion,
et on a donc $(\Psi\circ U)(g) \subseteq U(g)$, l'inclusion réciproque
ayant déjà été vue (et de toute façon on n'en a pas besoin).  On a
donc bien trouvé une fonction partielle $f := U(\varnothing)$ telle
que $\Psi(f) \subseteq f$ (même $\Psi(f) = f$).
\end{proof}

\begin{proof}[Seconde démonstration]
On utilise la notion d'ordinaux.  On pose $f_0 = \varnothing$, et par
induction sur l'ordinal $\alpha$ on définit $f_{\alpha+1} =
\Psi(f_\alpha)$ et si $\delta$ est un ordinal limite alors $f_\delta =
\bigcup_{\gamma<\delta} f_\gamma$.  On montre simultanément par
induction sur $\alpha$ que $f_\alpha$ est bien définie, est une
fonction partielle, et, grâce à la croissance de $\Psi$, prolonge
$f_\beta$ pour chaque $\beta<\alpha$ (c'est ce dernier point qui
permet de conclure que $\bigcup_{\gamma<\delta} f_\gamma$ est une
fonction partielle lorsque $\delta$ est un ordinal limite : la réunion
d'une famille totalement ordonnée pour l'inclusion de fonctions
partielles est encore une fonction partielle).  Les inclusions
$f_\beta \subseteq f_\alpha$ pour $\beta<\alpha$ ne peuvent pas être
toutes strictes sans quoi on aurait une surjection d'un ensemble sur
la classe des ordinaux.  Il existe donc $\tau$ tel que $f_{\tau+1} =
f_\tau$, c'est-à-dire $\Psi(f_\tau) = f_\tau$.  D'autre part, si
$\Psi(h) = h$, alors par induction sur $\alpha$ on montre $f_\alpha
\subseteq h$ pour chaque $\alpha$ (l'étape successeur étant que si
$f_\alpha \subseteq h$ alors $f_{\alpha+1} = \Psi(f_\alpha) \subseteq
\Psi(h) = h$), donc en particulier $f_\tau \subseteq h$, et $f_\tau$
est bien le plus petit $f$ tel que $\Psi(f) = f$.
\end{proof}

\thingy\label{non-well-founded-induction-algorithm} Lorsque $G$ est
\emph{fini}, la fonction $f$ définie par le
théorème \ref{non-well-founded-definition} (et \textit{a fortiori} par
le théorème \ref{well-founded-definition}) peut se calculer
algorithmiquement de la façon suivante :
\begin{itemize}
\item initialement, poser $f(x) = \mathtt{undefined}$ pour tout $x\in
  G$ ;
\item parcourir tous les éléments $x\in G$ et, si $f(y)$ est défini
  pour suffisamment de voisins sortants $y$ de $x$ pour imposer
  $f(x)$, autrement dit, si $\Phi(x, f|_{\outnb(x)})$ est défini,
  alors définir $f(x)$ à cette valeur (noter qu'elle ne changera
  jamais ensuite du fait de la cohérence de $\Phi$) ; et
\item répéter l'opération précédente jusqu'à ce que $f$ ne change
  plus.
\end{itemize}
La démonstration donnée avec les ordinaux correspond exactement à cet
algorithme, à ceci près que « répéter l'opération » devra peut-être se
faire de façon transfinie.

\thingy Les différentes définitions par induction bien-fondée
proposées en exemple dans la
section \ref{subsection-well-founded-graphs}, c'est-à-dire la notion
de rang bien-fondé (cf. \ref{definition-well-founded-rank}), de
fonction de Grundy (cf. \ref{definition-grundy-function}) et de
P-sommets et N-sommets (cf. \ref{definition-grundy-kernel}), et même
la généralisation ordinale des deux premiers
(cf. \ref{ordinal-valued-rank-and-grundy-function}), peuvent se
généraliser à des graphes non nécessairement bien-fondés en utilisant
le théorème \ref{non-well-founded-definition} : il faudra juste
admettre que ce soient des fonctions \emph{partielles}, non définies
sur certains sommets (voire, tous les sommets) qui ne sont pas dans la
partie bien-fondée du graphe (cf. \ref{definition-wfpart}).
Néanmoins, ces définitions ne présentent qu'un intérêt limité : le
rang bien-fondé, par exemple, s'avère être défini exactement sur la
partie bien-fondée de $G$ (vérification facile) et coïncider avec le
rang bien-fondé sur ce sous-graphe, et pour ce qui est des P-sommets
et N-sommets, on obtient exactement la fonction évoquée
en \ref{p-d-n-vertices-of-perfect-information-games}.



\subsection{Écrasement transitif}

\begin{defn}\label{definition-transitive-collapse}
Soit $G$ un graphe orienté bien-fondé.  En utilisant le
théorème \ref{well-founded-definition} (modulo la remarque qui suit),
on définit alors une fonction $f$ sur $G$ par $f(x) = \{f(y) :
y\in\outnb(x)\}$.  L'image $f(G)$ de $G$ par la fonction $f$ (c'est-à-dire
l'ensemble des $f(x)$ pour $x\in G$) s'appelle l'\defin{écrasement
  transitif} ou \defin{écrasement de Mostowski} de $G$, tandis que
$f$ s'appelle la fonction d'écrasement, et la valeur $f(x)$ (qui n'est
autre que l'écrasement transitif de l'aval de $x$ vu comme un graphe
orienté) s'appelle l'écrasement transitif du sommet $x$.

On considérera l'écrasement $f(G)$ de $G$ comme un graphe orienté, en
plaçant une arête de $u$ vers $v$ lorsque $v \in u$ ; autrement dit,
lorsque $v = f(y)$ et $u = f(x)$ pour certains $x,y$ de $G$ tels qu'il
existe une arête de $x$ vers $y$.
\end{defn}

\thingy En particulier, un puits a pour écrasement $\varnothing$, un
sommet qui n'a pour voisins sortants que des sommets terminaux a pour
écrasement $\{\varnothing\}$, un sommet qui n'a pour voisins sortants
que de tels sommets a pour écrasement $\{\{\varnothing\}\}$ tandis que
s'il a aussi des sommets terminaux pour voisins sortants ce sera
$\{\varnothing,\penalty0 \{\varnothing\}\}$, et ainsi de suite.

La terminologie « transitif » fait référence au fait qu'un ensemble
$E$ est dit transitif lorsque $v \in u\in E$ implique $v \in E$ (de
façon équivalente, $E \subseteq \mathscr{P}(E)$ où $\mathscr{P}(E)$
est l'ensemble des parties de $E$) : c'est le cas de $f(G)$ ici.

Il y a une subtilité ensembliste dans la définition ci-dessus, c'est
qu'on ne peut pas donner \textit{a priori} un ensemble $Z$ dans lequel
$f$ prend sa valeur : il faut en fait appliquer une généralisation
de \ref{well-founded-definition} où $Z$ est remplacé par l'univers de
tous les ensembles : nous ne rentrerons pas dans ces subtilités, et
admettrons qu'il existe bien une unique fonction $f$ sur $G$ qui
vérifie $f(x) = \{f(y) : y\in\outnb(x)\}$ pour chaque $x\in G$.

\begin{defn}
Un graphe orienté $G$ est dit \defin[extensionnel (graphe)]{extensionnel} lorsque deux
sommets $x$ et $x'$ ayant le même ensemble de voisins sortants ($\outnb(x)
= \outnb(x')$) sont égaux.
\end{defn}

Pour bien comprendre et utiliser la définition ci-dessus, il est
pertinent de rappeler que \emph{deux ensembles sont égaux si et
  seulement si ils ont les mêmes éléments} (\defin[extensionalité (axiome)]{axiome
  d'extensionalité}).

\begin{prop}\label{extensional-iff-collapse-injective}
Un graphe orienté bien-fondé est extensionnel si et seulement si sa
fonction d'écrasement $f$ définie
en \ref{definition-transitive-collapse} est injective.
\end{prop}
\begin{proof}
Si $f$ est injective et si $\outnb(x) = \outnb(x')$ alors en
particulier $f(x) = f(x')$ (puisque $f(x)$ est définie comme l'image
par $f$ de $\outnb(x)$), donc $x = x'$.  Ceci montre que $G$ est
extensionnel.

Réciproquement, $G$ extensionnel et montrons que $f$ est injective,
c'est-à-dire que $f(x) = f(x')$ implique $x = x'$.  On va procéder par
induction bien-fondée sur le graphe $G^2$ dont les sommets sont les
couples $(x,x')$ de sommets de $G$, avec une arête de $(x,x')$ vers
$(y,y')$ lorsqu'il y en a de $x$ vers $y$ \emph{et} de $x'$
vers $y'$ : il est clair que $G^2$ est bien-fondé ; soit $P$
l'ensemble des sommets $(x,x')$ de $G^2$ tels que $f(x) = f(x')$
implique $x=x'$ (autrement dit, l'ensemble des sommets tels que
$(x,x')$ tels que $x=x'$ \emph{ou bien} $f(x) \neq f(x')$).  Soit
$(x,x')$ un sommet de $G^2$ dont tous les voisins sortants vérifient
l'hypothèse d'induction (i.e., appartiennent à $P$) : on suppose $f(x)
= f(x')$ et on veut montrer $x = x'$ pour pouvoir conclure $P = G^2$.
Or $f(x) \subseteq f(x')$ signifie que tout $f(y) \in f(x)$ appartient
à $f(x')$, c'est-à-dire que pour tout voisin sortant $y$ de $x$ il
existe un voisin sortant $y'$ de $x'$ pour lequel $f(y) = f(y')$, et
l'hypothèse d'induction montre alors $y = y'$ : ainsi, $\outnb(x)
\subseteq \outnb(x')$, et par symétrie, $f(x) = f(x')$ montre
$\outnb(x) = \outnb(x')$ donc, par extensionalité de $G$, on a $x =
x'$ comme on le voulait.
\end{proof}

\thingy Si $G$ est un graphe orienté (non nécessairement bien-fondé),
et si $\sim$ est une relation d'équivalence sur l'ensemble des sommets
de $G$, on peut définir un \emph{graphe quotient} $G/\sim$ dont les
sommets sont les classes d'équivalences pour $\sim$ de sommets de $G$
et dont les arêtes sont les couples $(\bar x,\bar y)$ de classes
telles qu'il existe une arête $(x,y)$ (i.e., de $x$ vers $y$) dans $G$
avec $\bar x$ classe de $x$ et $\bar y$ celle de $y$ ; si ce graphe
est extensionnel, on peut dire abusivement que $\sim$ l'est :
concrètement, cela signifie que pour tous $x,x' \in G$, si pour chaque
voisin sortant $y$ de $x$ il existe un voisin sortant $y'$ de $x'$
avec $y \sim y'$ et que la même chose vaut en échangeant $x$ et $x'$,
alors $x \sim x'$.  Une intersection quelconque de relations
d'équivalence extensionnelles sur $G$ est encore une relation
d'équivalence extensionnelle, donc il existe une plus petite relation
d'équivalence extensionnelle $\equiv$ sur $G$, c'est-à-dire un plus
grand quotient $G/\equiv$ de $G$ qui soit extensionnel (« plus grand »
au sens où tout quotient de $G$ par une relation d'équivalence se
factorise à travers ce quotient $G/\equiv$).

Le contenu essentiel de la
proposition \ref{extensional-iff-collapse-injective} est que
l'écrasement transitif $f(G)$ d'un graphe $G$ bien-fondé réalise ce
plus grand quotient extensionnel $G/\equiv$ : la relation $f(x) =
f(x')$ sur $G$ est précisément la plus petite relation d'équivalence
extensionnelle $\equiv$ sur $G$ (en effet, la relation $f(x) = f(x')$
est évidemment extensionnelle, donc contient $\equiv$ par définition
de celle-ci, mais l'écrasement de $G/\equiv$ est le même que celui
de $G$, et comme la fonction d'écrasement est injective sur
$G/\equiv$, on a bien $f(x) = f(x')$ ssi $x\equiv x'$).


%
%
%

\section{Introduction aux ordinaux}\label{section-ordinals}

\subsection{Présentation informelle}\label{subsection-informal-description-of-ordinals}

\thingy Les \index{ordinal}ordinaux sont une sorte de nombres, totalement ordonnés et
même « bien-ordonnés », qui généralisent les entiers naturels en
allant « au-delà de l'infini » : les entiers naturels
$0,1,2,3,4,\ldots$ sont en particulier des ordinaux (ce sont les plus
petits), mais il existe un ordinal qui vient après eux, à
savoir $\omega$, qui est lui-même suivi de
$\omega+1,\omega+2,\omega+3,\ldots$, après quoi vient $\omega\cdot 2$
(ou simplement $\omega 2$), et beaucoup d'autres choses.

\begin{center}
\begin{tikzpicture}
\begin{scope}[line width=1.5pt,cap=round,join=round]
% x = 10*(1-0.8^k*(1-0.2*(1-0.8^n)))
% y = 2*0.8^k*0.8^n
\draw (0.00000,2.00000) -- (0.00000,-2.00000);
\draw (0.40000,1.60000) -- (0.40000,-1.60000);
\draw (0.72000,1.28000) -- (0.72000,-1.28000);
\draw (0.97600,1.02400) -- (0.97600,-1.02400);
\draw (1.18080,0.81920) -- (1.18080,-0.81920);
\draw (1.34464,0.65536) -- (1.34464,-0.65536);
\draw (1.47571,0.52429) -- (1.47571,-0.52429);
\draw (1.58057,0.41943) -- (1.58057,-0.41943);
\draw (1.66446,0.33554) -- (1.66446,-0.33554);
\draw (1.73156,0.26844) -- (1.73156,-0.26844);
\draw[fill] (1.78525,0.21475) -- (1.78525,-0.21475) -- (2.00000,0.00000) -- (1.78525,0.21475);
\draw (2.00000,1.60000) -- (2.00000,-1.60000);
\draw (2.32000,1.28000) -- (2.32000,-1.28000);
\draw (2.57600,1.02400) -- (2.57600,-1.02400);
\draw (2.78080,0.81920) -- (2.78080,-0.81920);
\draw (2.94464,0.65536) -- (2.94464,-0.65536);
\draw (3.07571,0.52429) -- (3.07571,-0.52429);
\draw (3.18057,0.41943) -- (3.18057,-0.41943);
\draw (3.26446,0.33554) -- (3.26446,-0.33554);
\draw (3.33156,0.26844) -- (3.33156,-0.26844);
\draw (3.38525,0.21475) -- (3.38525,-0.21475);
\draw[fill] (3.42820,0.17180) -- (3.42820,-0.17180) -- (3.60000,0.00000) -- (3.42820,0.17180);
\draw (3.60000,1.28000) -- (3.60000,-1.28000);
\draw (3.85600,1.02400) -- (3.85600,-1.02400);
\draw (4.06080,0.81920) -- (4.06080,-0.81920);
\draw (4.22464,0.65536) -- (4.22464,-0.65536);
\draw (4.35571,0.52429) -- (4.35571,-0.52429);
\draw (4.46057,0.41943) -- (4.46057,-0.41943);
\draw (4.54446,0.33554) -- (4.54446,-0.33554);
\draw (4.61156,0.26844) -- (4.61156,-0.26844);
\draw (4.66525,0.21475) -- (4.66525,-0.21475);
\draw (4.70820,0.17180) -- (4.70820,-0.17180);
\draw[fill] (4.74256,0.13744) -- (4.74256,-0.13744) -- (4.88000,0.00000) -- (4.74256,0.13744);
\draw (4.88000,1.02400) -- (4.88000,-1.02400);
\draw (5.08480,0.81920) -- (5.08480,-0.81920);
\draw (5.24864,0.65536) -- (5.24864,-0.65536);
\draw (5.37971,0.52429) -- (5.37971,-0.52429);
\draw (5.48457,0.41943) -- (5.48457,-0.41943);
\draw (5.56846,0.33554) -- (5.56846,-0.33554);
\draw (5.63556,0.26844) -- (5.63556,-0.26844);
\draw (5.68925,0.21475) -- (5.68925,-0.21475);
\draw (5.73220,0.17180) -- (5.73220,-0.17180);
\draw (5.76656,0.13744) -- (5.76656,-0.13744);
\draw[fill] (5.79405,0.10995) -- (5.79405,-0.10995) -- (5.90400,0.00000) -- (5.79405,0.10995);
\draw (5.90400,0.81920) -- (5.90400,-0.81920);
\draw (6.06784,0.65536) -- (6.06784,-0.65536);
\draw (6.19891,0.52429) -- (6.19891,-0.52429);
\draw (6.30377,0.41943) -- (6.30377,-0.41943);
\draw (6.38766,0.33554) -- (6.38766,-0.33554);
\draw (6.45476,0.26844) -- (6.45476,-0.26844);
\draw (6.50845,0.21475) -- (6.50845,-0.21475);
\draw (6.55140,0.17180) -- (6.55140,-0.17180);
\draw (6.58576,0.13744) -- (6.58576,-0.13744);
\draw (6.61325,0.10995) -- (6.61325,-0.10995);
\draw[fill] (6.63524,0.08796) -- (6.63524,-0.08796) -- (6.72320,0.00000) -- (6.63524,0.08796);
\draw (6.72320,0.65536) -- (6.72320,-0.65536);
\draw (6.85427,0.52429) -- (6.85427,-0.52429);
\draw (6.95913,0.41943) -- (6.95913,-0.41943);
\draw (7.04302,0.33554) -- (7.04302,-0.33554);
\draw (7.11012,0.26844) -- (7.11012,-0.26844);
\draw[fill] (7.16381,0.21475) -- (7.16381,-0.21475) -- (7.37856,0.00000) -- (7.16381,0.21475);
\draw (7.37856,0.52429) -- (7.37856,-0.52429);
\draw (7.48342,0.41943) -- (7.48342,-0.41943);
\draw (7.56730,0.33554) -- (7.56730,-0.33554);
\draw (7.63441,0.26844) -- (7.63441,-0.26844);
\draw (7.68810,0.21475) -- (7.68810,-0.21475);
\draw[fill] (7.73105,0.17180) -- (7.73105,-0.17180) -- (7.90285,0.00000) -- (7.73105,0.17180);
\draw (7.90285,0.41943) -- (7.90285,-0.41943);
\draw (7.98673,0.33554) -- (7.98673,-0.33554);
\draw (8.05384,0.26844) -- (8.05384,-0.26844);
\draw (8.10753,0.21475) -- (8.10753,-0.21475);
\draw (8.15048,0.17180) -- (8.15048,-0.17180);
\draw[fill] (8.18484,0.13744) -- (8.18484,-0.13744) -- (8.32228,0.00000) -- (8.18484,0.13744);
\draw (8.32228,0.33554) -- (8.32228,-0.33554);
\draw (8.38939,0.26844) -- (8.38939,-0.26844);
\draw (8.44307,0.21475) -- (8.44307,-0.21475);
\draw[fill] (8.48602,0.17180) -- (8.48602,-0.17180) -- (8.65782,0.00000) -- (8.48602,0.17180);
\draw (8.65782,0.26844) -- (8.65782,-0.26844);
\draw (8.71151,0.21475) -- (8.71151,-0.21475);
\draw (8.75446,0.17180) -- (8.75446,-0.17180);
\draw[fill] (8.78882,0.13744) -- (8.78882,-0.13744) -- (8.92626,0.00000) -- (8.78882,0.13744);
\draw (8.92626,0.21475) -- (8.92626,-0.21475);
\draw (8.96921,0.17180) -- (8.96921,-0.17180);
\draw (9.00357,0.13744) -- (9.00357,-0.13744);
\draw[fill] (9.03106,0.10995) -- (9.03106,-0.10995) -- (9.14101,0.00000) -- (9.03106,0.10995);
\draw (9.14101,0.17180) -- (9.14101,-0.17180);
\draw (9.17537,0.13744) -- (9.17537,-0.13744);
\draw[fill] (9.20285,0.10995) -- (9.22484,-0.08796) -- (9.31281,0.00000) -- (9.20285,0.10995);
\draw[fill] (9.31281,0.13744) -- (9.31281,-0.13744) -- (9.45024,0.00000) -- (9.31281,0.13744);
\draw[fill] (9.45024,0.10995) -- (9.45024,-0.10995) -- (9.56020,0.00000) -- (9.45024,0.10995);
\draw[fill] (9.56020,0.08796) -- (9.56020,-0.08796) -- (9.64816,0.00000) -- (9.56020,0.08796);
\draw[fill] (9.64816,0.07037) -- (9.64816,-0.07037) -- (9.71853,0.00000) -- (9.64816,0.07037);
\draw[fill] (9.71853,0.05629) -- (9.71853,-0.05629) -- (9.77482,0.00000) -- (9.71853,0.05629);
\draw[fill] (9.77482,0.04504) -- (9.77482,-0.04504) -- (9.81986,0.00000) -- (9.77482,0.04504);
\draw[fill] (9.81986,0.03603) -- (9.81986,-0.03603) -- (9.85588,0.00000) -- (9.81986,0.03603);
\draw[fill] (9.85588,0.02882) -- (9.85588,-0.02882) -- (9.88471,0.00000) -- (9.85588,0.02882);
\draw[fill] (9.88471,0.02306) -- (9.88471,-0.02306) -- (10.00000,0.00000) -- (9.88471,0.02306);
\end{scope}
\node[anchor=north] at (0.00000,-2.00000) {$0$};
\node[anchor=north] at (0.40000,-1.60000) {$1$};
\node[anchor=north] at (0.72000,-1.28000) {$2$};
\node[anchor=north] at (0.97600,-1.02400) {$3$};
\node[anchor=north] at (2.00000,-1.60000) {$\omega$};
\node[anchor=north] at (2.32000,-1.28000) {$\scriptscriptstyle \omega+1$};
\node[anchor=north] at (3.60000,-1.28000) {$\omega2$};
\node[anchor=north] at (4.88000,-1.02400) {$\omega3$};
\end{tikzpicture}
\\{\footnotesize (Une rangée de $\omega^2$ allumettes.)}
\end{center}

\thingy Les ordinaux servent à mesurer la taille des ensembles
bien-ordonnés (c'est-à-dire, les ensembles totalement ordonnés dans
lesquels il n'existe pas de suite infinie strictement décroissante)
exactement comme les entiers naturels servent à mesurer la taille des
ensembles finis.

\thingy On pourra ajouter les ordinaux, et les multiplier, et même
élever un ordinal à la puissance d'un autre, mais il n'y aura pas de
soustraction ($\omega-1$ n'a pas de sens, en tout cas pas en tant
qu'ordinal, parce que $\omega$ est le plus petit ordinal infini).  Les
ordinaux ont de nombreux points en commun avec les entiers naturels
(l'addition est associative, la multiplication aussi, on peut les
écrire en binaire, etc.), mais aussi des différences importantes
(l'addition n'est pas commutative : on a $1+\omega = \omega$ mais
$\omega+1 > \omega$).

\thingy Comme la récurrence pour les entiers naturels, il y a sur les
ordinaux (ou de façon équivalente, sur les ensembles bien-ordonnés) un
principe d'\emph{induction transfinie}
(cf. \ref{scholion-transfinite-induction}), qui est en fait
l'application directe à eux du principe d'induction bien-fondée : son
énoncé est essentiellement le même que le principe parfois appelé de
« récurrence forte » pour les entiers naturels, c'est-à-dire que :
\begin{center}
Pour montrer une propriété sur tous les ordinaux $\alpha$ on peut
faire l'hypothèse d'induction qu'elle est déjà connue pour les
ordinaux $<\alpha$ au moment de la montrer pour $\alpha$.
\end{center}

Comme la récurrence sur les entiers naturels, et comme l'induction
bien-fondée dont c'est un cas particulier, l'induction transfinie
permet soit de \emph{démontrer} des propriétés sur les ordinaux, soit
de \emph{définir} des fonctions sur ceux-ci :
\begin{center}
Pour définir une fonction sur tous les ordinaux $\alpha$ on peut faire
l'hypothèse d'induction qu'elle est déjà définie pour les
ordinaux $<\alpha$ au moment de la définir pour $\alpha$.
\end{center}

\thingy\label{decreasing-sequences-of-ordinals-terminate}
Ce qui importe surtout pour la théorie des jeux est le fait suivant :
\begin{center}
\emph{toute suite strictement décroissante d'ordinaux est finie}
\end{center}
(généralisation du fait que toute suite strictement d'entiers naturels
est finie).  À cause de ça, les ordinaux peuvent servir à « mesurer »
toutes sortes de processus qui terminent à coup sûr en temps fini, ou
à généraliser les entiers naturels pour toutes sortes de processus qui
terminent à coup sûr en temps fini mais pas en un nombre d'étapes
borné \textit{a priori}.

Par exemple, on peut imaginer que le dessin de la figure ci-dessus
(figurant les ordinaux $<\omega^2$) représente une rangée d'allumettes
qu'on pourrait utiliser dans un jeu de nim
(cf. \ref{introduction-nim-game}) : si on convient que les allumettes
doivent être effacées \emph{par la droite}, ce qui revient à diminuer
strictement l'ordinal qui les compte (initialement $\omega^2$), la
ligne sera toujours vidée en temps fini même si les joueurs essaient
de la faire durer le plus longtemps possible (le premier coup va faire
tomber l'ordinal $\omega^2$ à $\omega\cdot k + n$ avec
$k,n\in\mathbb{N}$, après quoi les coups suivants l'amèneront au plus
à $\omega\cdot k + n'$ avec $n'<n$ qui va finir par tomber à $0$, puis
on tombe à $\omega\cdot k' + m$ avec $k'<k$, et en continuant ainsi on
finit forcément par retirer toutes les allumettes).

Plus formellement, quel que soit l'ordinal $\alpha$, l'ensemble
$\{\beta : \beta<\alpha\}$ des ordinaux plus petits, vu comme un
graphe pour la relation $>$ (i.e., on fait pointer une arête orientée
de chaque ordinal $\beta$ vers chaque ordinal strictement plus petit),
est bien-fondé, ou de façon équivalente, bien-ordonné.

\thingy\label{ordinal-counting-genies-story}
Voici une façon imagée d'y penser qui peut servir à faire le
lien avec la théorie des jeux : imaginons un génie qui exauce des vœux
en nombre limité (les vœux eux-mêmes sont aussi limités et ne
permettent certainement pas de faire le vœu d'avoir plus de vœux —
peut-être qu'on ne peut que souhaiter un paquet de carambars, ou de
transformer son ennemi en crapaud, ou d'annuler une transformation en
crapaud qu'on aurait soi-même subie, ou des choses de ce genre).  Si
le génie est prêt à exaucer $3$ vœux, on peut imaginer qu'à la fin de
chaque vœu qu'on prononce on doive dire « maintenant, il me reste
$n$ vœux » avec $n$ strictement inférieur à la valeur antérieure
(initialement $3$).

Cette définition se généralise aux ordinaux : un génie qui exauce
$\alpha$ vœux est un génie qui demande qu'on formule un vœu et qu'on
choisisse un ordinal $\beta < \alpha$, après quoi le vœu est exaucé et
le génie se transforme en un génie qui exauce $\beta$ vœux.

Ainsi, pour un génie qui exauce $\omega$ vœux on devra, lors du tout
premier vœu qu'on formule, décider quel nombre de vœux il reste, ce
nombre étant un \emph{entier naturel}, aussi grand qu'on le souhaite —
mais fini.  Cela peut sembler sans importance (si on a de toute façon
autant de vœux que l'on souhaite, même $N = 10^{1000}$, peu importe
qu'on doive choisir un nombre dès le début).  Mais comparons avec un
génie qui exauce $\omega+1$ vœux : pour celui-ci, lors du premier vœu
que l'on formule, on pourra décider qu'il reste $\omega$ vœux et c'est
au vœu suivant qu'on devra redescendre à un entier naturel (et le
choisir).  La différence entre avoir $\omega$ et $\omega+1$ vœux
apparaîtra si on imagine un combat entre Aladdin et Jafar où Jafar
utilise des vœux pour transformer Aladdin en crapaud et Aladdin pour
redevenir humain : si Jafar a initialement $\omega$ vœux et Aladdin
aussi, Jafar transforme Aladdin en crapaud et choisit qu'il lui reste
$N$ vœux avec $N$ fini, alors Aladdin redevient humain choisit qu'il
lui reste aussi au moins $N$ vœux, et au final il est sauf ; alors que
si Jafar a initialement $\omega+1$ vœux et Aladdin seulement $\omega$,
Jafar transforme Aladdin en crapaud et tombe à $\omega$, puis Aladdin
est obligé de choisir un $N$ fini en formulant le vœu de redevenir
humain, et Jafar peut choisir au moins $N$ vœux et gagne le combat
(ainsi que quelques paquets de carambars).

\thingy\label{introduction-von-neumann-ordinals}
La construction moderne des ordinaux, introduite par
J. von Neumann en 1923, est mathématiquement très élégante mais peut-être
d'autant plus difficile à comprendre qu'elle est subtile :
\begin{center}
\index{von Neumann (ordinal de)}\emph{un ordinal est l'ensemble des ordinaux strictement plus petits que lui}
\end{center}
— ainsi, l'entier $0$ est défini comme l'ensemble vide $\varnothing$
(puisqu'il n'y a pas d'ordinaux plus petits que lui), l'entier $1$ est
défini comme l'ensemble $\{0\} = \{\varnothing\}$ ayant pour seul
élément $0$ (puisque $0$ est le seul ordinal plus petit que $1$),
l'entier $2$ est défini comme $\{0,1\} =
\{\varnothing,\{\varnothing\}\}$, l'entier $3$ comme $\{0,1,2\} =
\{\varnothing,\{\varnothing\}, \penalty0
\{\varnothing,\{\varnothing\}\}\}$, et ainsi de suite, et l'ordinal
$\omega$ est défini comme l'ensemble $\mathbb{N} = \{0,1,2,3,\ldots\}$
de tous les entiers naturels, puis $\omega+1$ comme l'ensemble $\omega
\cup \{\omega\}$ des entiers naturels auquel on a ajouté le seul
élément $\omega$, et plus généralement $\alpha+1 = \alpha \cup
\{\alpha\}$.  Formellement, un ordinal est l'écrasement transitif
(cf. \ref{definition-transitive-collapse}) d'un ensemble bien-ordonné
(i.e., totalement ordonné bien-fondé).

Cette définition a certains avantages, par exemple la borne supérieure
d'un ensemble $S$ d'ordinaux est simplement la réunion
$\bigcup_{\alpha\in S} \alpha$ de(s éléments de) $S$.  Néanmoins, elle
n'est pas vraiment nécessaire à la théorie des ordinaux, et nous
tâcherons d'éviter d'en dépendre.  Mais il faut au moins retenir une
idée :

\thingy Pour tout ensemble $S$ d'ordinaux, il existe un ordinal qui est
plus grand que tous les éléments de $S$ ; il existe même un \emph{plus
  petit} ordinal plus grand que tous les éléments de $S$,
c'est-à-dire, une \emph{borne supérieure} de $S$.  Ce fait est la clé
de l'inexhaustibilité des ordinaux : quelle que soit la manière dont
on essaie de rassembler des ordinaux en un ensemble, on peut trouver
un ordinal strictement plus grand qu'eux (en particulier, les ordinaux
ne forment pas un ensemble, pour un peu la même raison que l'ensemble
de tous les ensembles n'existe pas : il est « trop gros » pour tenir
dans un ensemble).

Pour dire les choses différemment avec un slogan peut-être un peu
approximatif :
\begin{center}
À chaque fois qu'on a construit les ordinaux jusqu'à un certain point,
on crée un nouvel ordinal qui vient juste après tous ceux-là.
\end{center}

\thingy Pour aider à comprendre comment les choses commencent, et en
partant de l'idée générale que
\begin{center}
\emph{comprendre un ordinal, c'est comprendre tous les ordinaux
  strictement plus petits que lui, et comment ils s'ordonnent}
\end{center}
voici comment s'arrangent les plus petits ordinaux.

Après les entiers naturels $0,1,2,3,\ldots$ vient l'ordinal $\omega$
puis $\omega+1,\omega+2,\omega+3$ et ainsi de suite, après quoi
viennent $\omega 2, \omega 2+1, \omega 2+2,\ldots$ qui sont suivis de
$\omega 3$ et le même mécanisme recommence.  Les ordinaux $\omega k +
n$ pour $k,n\in\mathbb{N}$ sont ordonnés par l'ordre lexicographique
donnant plus de poids à $k$ : l'ordinal qui vient immédiatement après
(c'est-à-dire, leur ensemble, si on utilise la construction de
von Neumann) est $\omega\cdot\omega = \omega^2$, qui est suivi de
$\omega^2+1,\omega^2+2,\ldots$ et plus généralement des $\omega^2 +
\omega k + n$, qui sont eux-mêmes suivis de $\omega^2 \cdot 2$.  Les
$\omega^2 \cdot n_2 + \omega \cdot n_1 + n_0$ sont ordonnés par
l'ordre lexicographique donnant plus de poids à $n_2$, puis à $n_1$
puis à $n_0$.  L'ordinal qui vient immédiatement après tous ceux-ci
est $\omega^3$.

En itérant ce procédé, on fabrique de même $\omega^4$, puis $\omega^5$
et ainsi de suite : les $\omega^r \cdot n_r + \cdots + \omega \cdot
n_1 + n_0$ (c'est-à-dire en quelque sorte des polynômes en $\omega$ à
coefficients dans $\mathbb{N}$) sont triés par ordre lexicographique
en donnant plus de poids aux coefficients $n_i$ pour $i$ grand (et en
identifiant bien sûr un cas où $n_r = 0$ par celui où il est omis : il
s'agit de l'ordre lexicographique sur les suites d'entiers nulles à
partir d'un certain rang).  L'ordinal qui vient immédiatement après
est $\omega^\omega$, puis on a tous les $\omega^\omega + \omega^r
\cdot n_r + \cdots + \omega \cdot n_1 + n_0$ jusqu'à $\omega^\omega
\cdot 2$, et de même $\omega^\omega \cdot 3$, etc., jusqu'à
$\omega^\omega \cdot \omega = \omega^{\omega+1}$.

En répétant $\omega$ fois toute cette séquence, on obtient
$\omega^{\omega+2}$, puis de nouveau $\omega^{\omega+3}$ et ainsi de
suite : après quoi vient $\omega^{\omega 2}$, et on voit comment on
peut continuer de la sorte.

\thingy Plus généralement, tout ordinal va s'écrire de façon unique
sous la forme $\omega^{\gamma_s} n_s + \cdots + \omega^{\gamma_1} n_1$
où $\gamma_s > \cdots > \gamma_1$ sont des ordinaux et
$n_s,\ldots,n_1$ sont des entiers naturels non nuls (si un $n_i$ est
nul il convient de l'omettre) : il s'agit d'une sorte d'écriture en
« base $\omega$ » de l'ordinal, appelée \defin[Cantor (forme normale de)]{forme normale de
  Cantor}.  On compare deux formes normales de Cantor en comparant le
terme dominant (le plus à gauche, i.e., $\omega^{\gamma_s} n_s$ dans
les notations qui viennent d'être données, ce qui se fait lui-même en
comparant les $\gamma_s$ et sinon, les $n_s$), et s'ils sont égaux, en
comparant le suivant et ainsi de suite.

La forme normale de Cantor ne permet cependant pas de « comprendre »
tous les ordinaux, car il existe des ordinaux tels que $\varepsilon =
\omega^\varepsilon$.  Le plus petit d'entre eux est noté
$\varepsilon_0$ et est la limite de
$\omega,\omega^\omega,\omega^{\omega^\omega},\omega^{\omega^{\omega^\omega}},\ldots$

\thingy On retiendra qu'il existe trois sortes d'ordinaux, cette
distinction étant souvent utile dans les inductions :
\begin{itemize}
\item l'ordinal nul $0$, qui est souvent un cas spécial,
\item les ordinaux « successeurs », c'est-à-dire ceux qui sont de la
  forme $\beta+1$ pour $\beta$ un ordinal plus petit (de façon
  équivalente, il y a un plus grand ordinal strictement plus petit),
\item les ordinaux qui sont la borne supérieure (ou « limite ») des
  ordinaux strictement plus petits et qu'on appelle, pour cette
  raison, « ordinaux limites » (autrement dit, $\delta$ est limite
  lorsque pour tout $\beta<\delta$ il existe $\beta'$ avec
  $\beta<\beta'<\delta$).
\end{itemize}
À titre d'exemple, l'ordinal $0$, l'ordinal $42$ et l'ordinal $\omega$
sont des exemples de ces trois cas.  D'autres ordinaux limites sont
$\omega\cdot 2$, ou $\omega^2$, ou encore $\omega^\omega +
\omega^3\cdot 7$, ou bien $\omega^{\omega+1}$ ; en revanche,
$\omega^\omega + 1$ ou $\omega^{\omega 2} + 1729$ sont
successeurs.  Dans la forme normale de Cantor, un ordinal est
successeur si et seulement si le dernier terme (le plus à droite) est
un entier naturel non nul.

\thingy\label{introduction-nimbers-and-numbers}
Les ordinaux vont servir à définir différents jeux qui, pris
isolément, sont extrêmement peu intéressants, mais qui ont la vertu de
permettre de « mesurer » d'autres jeux : ces jeux ont en commun que,
partant d'un ordinal $\alpha$, l'un ou l'autre joueur, ou les deux,
ont la possibilité de le faire décroître (strictement), c'est-à-dire
de le remplacer par un ordinal $\beta < \alpha$ strictement plus petit
— comme expliqué en \ref{decreasing-sequences-of-ordinals-terminate},
ce processus termine forcément.  Dans le cadre esquissé
en \ref{introduction-graph-game}, on a trois jeux associés à un
ordinal $\alpha$ :
\begin{itemize}
\item Un jeu \emph{impartial}, c'est-à-dire que les deux joueurs ont
  les mêmes options à partir de n'importe quelle position $\beta \leq
  \alpha$, à savoir, les ordinaux $\beta' < \beta$ — autrement dit,
  les deux joueurs peuvent décroître l'ordinal.  Dans le cadre
  de \ref{introduction-graph-game}, le graphe a pour sommets les
  ordinaux $\beta \leq \alpha$ avec une arête (« verte », i.e.,
  utilisable par tout le monde) reliant $\beta$ à $\beta'$ lorsque
  $\beta'<\beta$.  Il s'agit du jeu de nim
  (cf. \ref{introduction-nim-game}) avec une seule ligne d'allumettes
  ayant initialement $\alpha$ allumettes (les allumettes sont bien
  ordonnées et doivent être retirées \emph{par la droite} dans un
  dessin comme au début de cette section).  Ce jeu s'appelle parfois
  le \index{nimbre}« nimbre » associé à l'ordinal $\alpha$.
\item Deux jeux \emph{partisans} (=partiaux), où un joueur n'a aucun
  coup possible (il a donc immédiatement perdu si c'est à son tour de
  jouer, ce qui rend le jeu, pris isolément, encore plus inintéressant
  que le précédent) : un jeu « bleu » ou « positif », dans lequel seul
  le joueur « bleu » (également appelé « gauche », « Blaise »...) peut
  jouer, exactement comme dans le jeu impartial ci-dessus, tandis que
  l'autre joueur ne peut rien faire, et un jeu « rouge » ou
  « négatif », dans lequel seul le joueur « rouge » (également appelé
  « droite », « Roxane »...) peut jouer tandis que l'autre ne peut
  rien faire.  Dans le cadre de \ref{introduction-graph-game}, le
  graphe a pour sommets les ordinaux $\beta \leq \alpha$ avec une
  arête reliant $\beta$ à $\beta'$ lorsque $\beta'<\beta$, ces arêtes
  étant toutes bleues ou toutes rouges selon le jeu considéré.  Il
  s'agit d'un jeu qui correspond à un certain avantage du joueur bleu,
  respectivement rouge, à rapprocher de l'histoire
  \ref{ordinal-counting-genies-story} ci-dessus.  Le jeu bleu est
  parfois appelé le « nombre surréel » associé à l'ordinal $\alpha$,
  tandis que le rouge est l'opposé du bleu.
\end{itemize}



\subsection{Ensembles bien-ordonnés et induction transfinie}

\thingy\label{definition-well-ordered-set}
Un ensemble \defin[ordonné (ensemble)]{[partiellement] ordonné} est un ensemble
muni d'une relation $>$ (d'ordre \emph{strict}) qui soit à la fois
\begin{itemize}
\item irréflexive ($x>x$ n'est jamais vrai quel que soit $x$), et
\item transitive ($x>y$ et $y>z$ entraînent $x>z$).
\end{itemize}
Une telle relation est automatiquement antisymétrique ($x>y$ et $y>x$
ne sont jamais simultanément vrais pour $x\neq y$).  On peut tout
aussi bien définir un ensemble partiellement ordonné en utilisant
l'ordre \emph{large} $\geq$ ($x\geq y$ étant défini par $x>y$ ou
$x=y$, ou symétriquement $x>y$ étant défini par $x\geq y$ et $x\neq
y$), réflexive, antisymétrique et transitive.  On note bien sûr $x\leq
y$ pour $y\geq x$ et $x<y$ pour $y>x$.

Un ensemble partiellement ordonné est dit \defin[totalement ordonné (ensemble)]{totalement ordonné}
lorsque pour tous $x\neq y$ on a soit $x>y$ soit $y>x$.

%% TODO: éclaircir le fait que dans ce qui suit, « bien-fondé » se
%% comprend pour le graphe reliant $x$ à $y$ ssi $x>y$ (voir aussi la
%% précédente occurrence du terme « bien-ordonné »).

Un ensemble totalement ordonné bien-fondé $W$ est dit
\defin[bien-ordonné (ensemble)]{bien-ordonné}.  D'après \ref{well-founded-induction}, ceci
peut se reformuler de différentes façons :
\begin{itemize}
\item[(*)]$W$ est un ensemble totalement ordonné dans lequel il
  n'existe pas de suite infinie strictement décroissante.
\item[(\dag)]$W$ est un ensemble (partiellement) ordonné dans lequel
  toute partie \emph{non vide} $N$ a un plus petit élément.
\item[(\ddag)]$W$ est totalement ordonné, et si une partie $P\subseteq
  W$ vérifie la propriété suivante « si $x \in W$ est tel que tout
    élément strictement plus petit que $x$ appartient à $P$, alors $x$
    lui-même appartient à $P$ » (on pourra dire « $P$ est inductif »),
  alors $P = W$.
\end{itemize}
(Dans (\dag), « partiellement ordonné » suffit car si $\{x,y\}$ a un
plus petit élément c'est bien qu'on a $x<y$ ou $y>x$.)

\begin{scho}[principe d'induction transfinie]\label{scholion-transfinite-induction}
Pour montrer une propriété $P$ sur les éléments d'un ensemble
bien-ordonné $W$, on peut supposer (comme « hypothèse d'induction »),
lorsqu'il s'agit de montrer que $x$ a la propriété $P$, que cette
propriété est déjà connue de tous les éléments strictement plus petits
que $x$.
\end{scho}

\begin{prop}\label{downward-closed-subsets-of-well-ordered-sets}
Soit $W$ un ensemble bien-ordonné et $S \subseteq W$ tel que $u < v$
avec $v \in S$ implique $u \in S$ (on peut dire que $S$ est
« aval-clos »).  Alors soit $S = W$ soit il existe $x\in W$ tel que $S
= \precs(x) := \{y\in W : y<x\}$.
\end{prop}
\begin{proof}
Si $S \neq W$, soit $x$ le plus petit élément de $W$ qui n'est pas
dans $S$.  Si $y < x$ alors $y \in S$ par minimalité de $x$.  Si $y
\geq x$ alors on a $y \not\in S$ car le contraire ($y \in S$)
entraînerait $x \in S$ d'après l'hypothèse faite sur $S$.  On a donc
montré que $y \in S$ si et seulement si $y < x$, c'est-à-dire
précisément $S = \precs(x)$.
\end{proof}

Le théorème suivant est un cas particulier du
théorème \ref{well-founded-definition} :

\begin{thm}[définition par induction transfinie]\label{transfinite-definition}
Soit $W$ un ensemble bien-ordonné et $Z$ un ensemble quelconque.
Notons $\precs(x) = \{y : y<x\}$ l'ensemble des éléments strictement
plus petits que $x$ dans $W$.

Appelons $\mathscr{F}$ l'ensemble des couples $(x,f)$ où $x\in W$ et
$f$ une fonction de $\precs(x)$ vers $Z$ (autrement dit, $\mathscr{F}$
est $\bigcup_{x \in W} \big(\{x\}\times Z^{\precs(x)}\big)$).  Soit
enfin $\Phi\colon \mathscr{F} \to Z$ une fonction quelconque.  Alors
il existe une unique fonction $f\colon W \to Z$ telle que pour tout $x
\in W$ on ait
\[
f(x) = \Phi(x,\, f|_{\precs(x)})
\]
\end{thm}

\begin{scho}\label{scholion-transfinite-definition}
Pour définir une fonction $f$ sur un ensemble bien-ordonné, on peut
supposer, lorsqu'on définit $f(x)$, que $f$ est déjà défini (i.e.,
connu) sur tous les éléments strictement plus petits que $x$ :
autrement dit, on peut librement utiliser la valeur de $f(y)$
pour $y<x$, dans la définition de $f(x)$.
\end{scho}



\subsection{Comparaison d'ensembles bien-ordonnés, et ordinaux}

\thingy Avant d'énoncer les résultats suivants, faisons une remarque
évidente et une définition.  La remarque est que si $W$ est
bien-ordonné et $E \subseteq W$ est un sous-ensemble de $W$, alors $E$
est lui-même bien ordonné (pour l'ordre induit) ; ceci s'applique en
particulier à $\precs(x) = \{y : y<x\}$.  La définition est qu'une
fonction $f$ entre ensembles ordonnés est dite \defin[croissante (fonction)]{croissante}
lorsque $x \leq y$ implique $f(x) \leq f(y)$, et \defin[strictement croissante (fonction)]{strictement
  croissante} lorsque $x < y$ implique $f(x) < f(y)$, ce qui entre des
ensembles totalement ordonnés signifie exactement la même chose que
« injective et croissante ».

\begin{prop}
Si $W$ est un ensemble bien-ordonné, et si $f\colon W\to W$ est
strictement croissante, alors $x \leq f(x)$ pour tout $x\in W$.
\end{prop}
\begin{proof}
Montrons par induction transfinie que $x \leq f(x)$.  Par hypothèse
d'induction, on peut supposer $y \leq f(y)$ pour tout $y<x$.
Supposons par l'absurde $f(x) < x$.  Alors l'hypothèse d'induction
appliquée à $y := f(x)$ donne $f(x) \leq f(f(x))$, tandis que la
stricte croissance de $f$ appliquée à $f(x) < x$ donne $f(f(x)) <
f(x)$.  On a donc une contradiction.
\end{proof}

\begin{cor}
Si $W$ est bien-ordonné, la seule bijection croissante $W \to W$ est
l'identité.
\end{cor}
\begin{proof}
Si $f\colon W\to W$ est une bijection croissante, la proposition
précédente appliquée à $f$ montre $x \leq f(x)$ pour tout $x \in W$,
mais appliquée à $f^{-1}$ elle montre $x \leq f^{-1}(x)$ donc $f(x)
\leq x$ et finalement $f(x) = x$.
\end{proof}

\begin{cor}\label{uniqueness-of-isomorphisms-between-well-ordered-sets}
Si $W,W'$ sont deux ensembles bien-ordonnés, il existe \emph{au plus
  une} bijection croissante $W \to W'$ (i.e., s'il en existe une, elle
est unique).
\end{cor}
Une telle bijection peut s'appeler un \textbf{isomorphisme} d'ensemble
bien-ordonnés, et on peut dire que $W$ et $W'$ ont \defin[isomorphes
  (ensembles bien-ordonnés)]{isomorphes} lorsqu'il existe un
isomorphisme entre eux (on conviendra en
cf. \ref{definition-of-ordinals} ci-dessous qu'on le note $\#W =
\#W'$).
\begin{proof}
Si $f,g\colon W\to W'$ sont deux bijections croissantes, appliquer le
corollaire précédent à la composée de l'une et de la réciproque de
l'autre.
\end{proof}

\begin{cor}\label{uniqueness-of-initial-segment-isomorphic-to-a-well-ordered-set}
Si $W$ est un ensemble bien-ordonné, $x\in W$ et $\precs(x) = \{y :
y<x\}$, alors il n'existe pas de fonction strictement croissante $W
\to \precs(x)$.
\end{cor}
\begin{proof}
Une telle fonction serait en particulier une fonction strictement
croissante $W \to W$, donc vérifierait $x \leq f(x)$ d'après la
proposition, ce qui contredit $f(x) \in \precs(x)$.
\end{proof}

Le théorème suivant assure que donnés deux ensembles bien-ordonnés, il
y a moyen de les comparer :
\begin{thm}\label{comparison-of-well-ordered-sets}
Si $W,W'$ sont deux ensembles bien-ordonnés, alors exactement l'une
des affirmations suivantes est vraie :
\begin{itemize}
\item il existe une bijection croissante $f\colon W \to \precs(y)$
  avec $y\in W'$,
\item il existe une bijection croissante $f\colon \precs(x) \to W'$
  avec $x\in W$,
\item il existe une bijection croissante $f\colon W \to W'$.
\end{itemize}
(Dans chaque cas, la bijection est automatiquement unique
d'après \ref{uniqueness-of-isomorphisms-between-well-ordered-sets}.
De plus, $y$ est unique dans le premier cas et $x$ l'est dans le
second,
d'après \ref{uniqueness-of-initial-segment-isomorphic-to-a-well-ordered-set}.)
\end{thm}
\begin{proof}
Les affirmations sont exclusives d'après le corollaire précédent.
Plus précisément, s'il existe une fonction strictement croissante
$f\colon W \to W'$, en composant à droite par $f$ on voit qu'il ne
peut pas exister de fonction strictement croissante $W' \to
\precs(x)$, donc le premier ou le dernier cas excluent celui du
milieu, mais par symétrie les deux derniers excluent le premier et les
trois cas sont bien exclusifs.

Considérons maintenant l'ensemble des couples $(x,y) \in W\times W'$
tels qu'il existe une bijection croissante (forcément unique !) entre
$\precs_W(x)$ et $\precs_{W'}(y)$ : d'après ce qu'on vient
d'expliquer, pour chaque $x$ il existe au plus un $y$ tel qu'il y ait
une telle bijection croissante, et pour chaque $y$ il existe au plus
un $x$.  On peut donc voir cet ensemble de couples comme (le graphe
d')une fonction partielle injective $\varphi\colon W \dasharrow W'$
(autrement dit, $\varphi(x)$ est l'unique $y$, s'il existe, tel qu'il
existe une bijection croissante entre $\precs_W(x)$
et $\precs_{W'}(y)$).

Si $x_0 \leq x$ et si $\varphi(x) =: y$ est défini, une bijection
croissante $f\colon \precs_W(x) \to \precs_{W'}(y)$ peut se
restreindre à $\precs_W(x_0)$ et il est clair que son image est
exactement $\precs_{W'}(f(x_0))$, ce qui montre que $\varphi(x_0)$ est
définie et vaut $f(x_0) < y = \varphi(x)$, notamment $\varphi$ est
strictement croissante.
D'après \ref{downward-closed-subsets-of-well-ordered-sets}, l'ensemble
de définition de $\varphi$ est soit $W$ tout entier soit de la forme
$\precs_W(x)$ pour un $x\in W$.  Symétriquement, son image est soit
$W'$ tout entier soit de la forme $\precs_{W'}(y)$ pour un $y\in W'$.
Et comme on vient de le voir, $\varphi$ est une bijection croissante
entre l'un et l'autre.  Ce ne peut pas être une bijection croissante
entre $\precs_W(x)$ et $\precs_{W'}(y)$ sinon $(x,y)$ lui-même serait
dans $\varphi$, une contradiction.  C'est donc que $\varphi$ est
exactement une bijection comme un des trois cas annoncés.
\end{proof}

L'affirmation suivante est une trivialité, mais peut-être utile à
écrire explicitement :
\begin{prop}\label{triviality-on-comparison-of-initial-segments-in-well-ordered-sets}
Soit $X$ un ensemble bien-ordonné : si $w,w' \in X$ et qu'on pose $W =
\precs_X(w)$ et $W' = \precs_X(w')$, les trois cas du
théorème \ref{comparison-of-well-ordered-sets} se produisent
exactement lorsque $w < w'$, resp. $w > w'$, resp. $w = w'$.
\end{prop}
\begin{proof}
C'est évident : si $w < w'$ alors l'identité fournit une bijection
croissante $W \to \precs_{W'}(w)$, et de même dans les autres cas.
\end{proof}

\begin{defn}\label{definition-of-ordinals}
Soient $W,W'$ deux ensembles bien-ordonnés.  On notera $\#W < \#W'$,
resp. $\#W > \#W'$, resp. $\#W = \#W'$, dans les trois cas du
théorème \ref{comparison-of-well-ordered-sets}.  Autrement dit, $\#W =
\#W'$ signifie qu'il existe une bijection croissante $W \to W'$
(unique
d'après \ref{uniqueness-of-isomorphisms-between-well-ordered-sets}),
ce qui définit une relation d'équivalence entre ensembles
bien-ordonnés, et on note $\#W < \#W'$ lorsque $\#W = \#\precs(y)$
pour un $y \in W'$, le théorème \ref{comparison-of-well-ordered-sets}
assurant qu'il s'agit d'une relation d'ordre total entre les classes
d'équivalence qu'on vient de définir.  La
proposition \ref{triviality-on-comparison-of-initial-segments-in-well-ordered-sets}
assure que si $w,w'$ sont deux éléments d'un même ensemble
bien-ordonné, alors $\#\precs(w) < \#\precs(w')$ se produit si et
seulement si $w < w'$, et de même en remplaçant le signe $<$ par
$=$ ou $>$.

La classe d'équivalence\footnote{Pour être parfaitement rigoureux, à
  cause de subtilités ensemblistes, on ne peut pas vraiment définir
  des classes d'équivalence de façon usuelle dans ce contexte, d'où
  l'intérêt de la définition suivante (ordinaux de von Neumann).}
$\#W$ pour la relation $\#W = \#W'$ s'appelle l'\defin{ordinal}
de $W$.  Par abus de notation, si $w$ est un élément d'un ensemble
bien-ordonné, on peut noter $\#w$ pour $\#\precs(w)$ (autrement dit,
on associe un ordinal non seulement à un ensemble bien-ordonné, mais
aussi à un élément d'un ensemble bien-ordonné).

Si on préfère éviter la définition par classe d'équivalence, on peut
aussi définir $\#W$ comme l'écrasement transitif
(cf. \ref{definition-transitive-collapse}) de $W$ (\defin[von Neumann
  (ordinal de)]{ordinal de von Neumann}), à savoir $\#W = \{\#x : x\in
W\}$ où $\#x = \{\#y : y<x\}$, cette définition ayant bien un sens par
induction transfinie (\ref{transfinite-definition}
et \ref{scholion-transfinite-definition}).

On appelle $\omega$ l'ordinal $\#\mathbb{N}$ de l'ensemble des entiers
naturels, et on identifie tout entier naturel $n$ à l'ordinal de
$\precs(n) = \{0,\ldots,n-1\}$ dans $\mathbb{N}$.
\end{defn}

\thingy Les deux façons de définir les ordinaux reviennent
essentiellement au même : en effet, s'il y a une bijection croissante
$W \to W'$ (forcément unique), alors les écrasements transitifs de $W$
et $W'$ coïncident, et réciproquement, si les écrasements transitifs
de $W$ et $W'$ coïncident, on définit une bijection croissante $W \to
W'$ en envoyant $x \in W$ sur l'unique $y \in W'$ tel que
$\precs_W(x)$ et $\precs_{W'}(y)$ aient le même écrasement transitif
(on peut au préalable montrer le lemme suivant : si $W$ est
bien-ordonné et $y < x$ dans $W$ alors l'écrasement transitif de
$\precs(y)$, qui est un élément de celui de $\precs(x)$, ne lui est
pas égal — ceci résulte d'une induction transfinie sur $x$).

Les ordinaux de von Neumann ont l'avantage d'être des ensembles
bien-définis et de vérifier $\beta < \alpha$ si et seulement si $\beta
\in \alpha$ ; ils ont comme inconvénient d'être peut-être plus
difficiles à visualiser.  Mais même si on n'identifie pas $\alpha =
\#W$ à l'ensemble des ordinaux strictement plus petits, il est
important de garder à l'esprit que l'ensemble des ordinaux strictement
plus petits est $\{\#\precs(x) : x \in W\}$ (par définition de
l'ordre !), et que $\alpha = \#\{\beta < \alpha\}$ (idem).  Même si
nous éviterons de supposer explicitement que les ordinaux sont
construits à la façon de von Neumann, il arrivera souvent qu'on dise
« un élément de $\alpha$ » pour parler d'un ordinal strictement plus
petit que $\alpha$ (cela peut être considéré comme un abus de
langage).

\bigbreak

Le théorème \ref{comparison-of-well-ordered-sets} a la conséquence
importante suivante sur les ordinaux :

\begin{thm}\label{sets-of-ordinals-are-well-ordered}
Tout ensemble d'ordinaux est bien-ordonné : deux ordinaux sont
toujours comparables (on a toujours $\beta<\alpha$ ou $\beta>\alpha$
ou $\beta=\alpha$), et il n'existe pas de suite infinie strictement
décroissante d'ordinaux.

Autrement dit : dans tout ensemble non vide d'ordinaux il y en a un
plus petit.
\end{thm}
\begin{proof}
Le théorème \ref{comparison-of-well-ordered-sets} signifie exactement
que les ordinaux sont \emph{totalement} ordonnés.  Reste à expliquer
qu'ils sont bien-ordonnés, c'est-à-dire, qu'il n'existe pas de suite
infinie strictement décroissante $\alpha_0 > \alpha_1 > \alpha_2 >
\cdots$.  Mais si on avait une telle suite, en appelant $W$ un
ensemble bien-ordonné tel que $\#W = \alpha_0$, chaque $\alpha_i$
suivant s'écrit $\#\precs(w_i)$ pour un $w_i \in W$, et
d'après \ref{triviality-on-comparison-of-initial-segments-in-well-ordered-sets}
on devrait avoir une suite strictement décroissante $w_1 > w_2 >
\cdots$ dans $W$, ce qui contredit le fait que $W$ est bien-ordonné.

La dernière affirmation vient de l'équivalence entre (*) et (\dag)
dans \ref{definition-well-ordered-set}.
\end{proof}

\begin{prop}\label{sup-and-strict-sup-of-sets-of-ordinals}
Tout ensemble $S$ d'ordinaux a une borne supérieure : autrement dit,
il existe un ordinal $\sup S$ qui est le plus petit majorant (large)
de $S$ (i.e., le plus petit ordinal $\alpha$ tel que $\beta\leq\alpha$
pour tout $\beta \in S$), et un ordinal $\sup^+ S$ qui est le plus
petit majorant strict de $S$ (i.e., le plus petit ordinal $\alpha$ tel
que $\beta<\alpha$ pour tout $\beta \in S$).

(On verra plus loin que $\sup^+ S = \sup\{\beta+1 \colon \beta \in
S\}$, donc cette notion n'est pas vraiment nouvelle.)
\end{prop}
\begin{proof}
D'après ce qu'on vient de voir (dernière affirmation
de \ref{sets-of-ordinals-are-well-ordered}), il suffit de montrer
qu'il existe un majorant strict de $S$.  Quitte à remplacer $S$ par sa
réunion avec l'ensemble des ordinaux inférieurs à un ordinal
quelconque de $S$ (pour les ordinaux de von Neumann, ceci revient à
remplacer $S$ par $S \cup \bigcup_{\alpha\in S} \alpha$), on peut
supposer que (*) si $\alpha \in S$ et $\beta < \alpha$ alors $\beta
\in S$.  On vient de voir que $S$ est bien-ordonné : si $\alpha =
\#S$, montrons qu'il s'agit d'un majorant strict de $S$ ; or si $\beta
\in S$, on a $\beta = \#\precs_S(\beta)$ d'après l'hypothèse (*) qu'on
vient d'assurer, et la définition de l'ordre sur les ordinaux donne
$\beta<\alpha$ : ainsi, $\alpha$ est bien un majorant strict comme
voulu.
\end{proof}

\thingy Une conséquence de cette proposition est qu'il n'y a pas
d'ensemble de tous les ordinaux (car si $S$ était un tel ensemble, il
aurait un majorant strict, qui par définition ne peut pas appartenir
à $S$) : c'est le \defin[Burali-Forti (paradoxe de)]{« paradoxe » de Burali-Forti} ; le mot
« paradoxe » fait référence à une conception ancienne de la théorie
des ensembles, mais selon les fondements modernes des mathématiques,
ce phénomène n'a rien de paradoxal (intuitivement, il y a trop
d'ordinaux pour pouvoir tenir dans un ensemble, de même qu'il n'y a
pas d'ensemble de tous les ensembles).  Ces subtilités ensemblistes ne
poseront pas de problème dans la suite de ces notes, il faut juste
reconnaître leur existence.



\subsection{Ordinaux successeurs et limites}

\thingy On appelle \defin[successeur (ordinal)]{successeur} d'un ordinal $\alpha$ le plus
petit ordinal strictement supérieur à $\alpha$ (qui existe d'après la
proposition \ref{sup-and-strict-sup-of-sets-of-ordinals} : si on veut,
c'est $\sup^+\{\alpha\}$) : il est facile de voir que cet ordinal est
fabriqué en ajoutant un unique élément à la fin d'un ensemble
bien-ordonné d'ordinal $\alpha$ ; on le note $\alpha+1$, et on a
$\beta \leq \alpha$ si et seulement si $\beta < \alpha+1$.
Réciproquement, tout ordinal ayant un plus grand élément (i.e.,
l'ordinal d'un ensemble bien-ordonné ayant un plus grand élément) est
un successeur : en effet, si $W$ a un plus grand élément $x$, alors
$\#W$ est le successeur de $\#\precs(x)$.

\thingy On distingue maintenant trois sortes d'ordinaux :
\begin{itemize}
\item l'ordinal \defin[nul (ordinal)]{nul} $0 = \#\varnothing$, mis à part de tous
  les autres,
\item les ordinaux \defin[successeur (ordinal)]{successeurs}, c'est-à-dire ceux qui ont un plus
  grand élément (au sens expliqué ci-dessus),
\item les autres, qu'on appelle ordinaux \defin[limite (ordinal)]{limites}.
\end{itemize}

La terminologie d'ordinaux « limites » s'explique ainsi : si $\delta$
est un ordinal non nul qui n'est pas successeur, cela signifie que
pour chaque $\beta<\delta$ il existe $\beta'$ avec
$\beta<\beta'<\delta$ (puisque $\beta$ n'est pas le plus grand élément
de $\delta$).  Ceci permet de dire que $\sup\{\beta < \alpha\} =
\sup^+\{\beta < \alpha\}$ (de façon générale, on a $\sup^+\{\beta <
\alpha\} = \alpha$ par définition), et on va définir la notion de
limite ainsi :

\thingy Si $\delta$ est un ordinal limite et $f$ est une fonction
\emph{croissante} définie sur les ordinaux strictement plus petits
que $\delta$ et à valeurs ordinales, on appelle \defin{limite} de $f$
en $\delta$ la valeur $\sup\{f(\xi) : \xi<\delta\}$.  On pourra la
noter $\lim_{\xi\to\delta} f(\xi)$ ou simplement $\lim_\delta f$.  (Il
s'agit bien d'une limite pour une certaine topologie : la topologie de
l'ordre ; plus exactement, c'est une limite car pour tous $\beta_1 <
\lim_\delta f < \beta_2$, il existe $\xi_0$ tel que $\beta_1 < f(\xi)
< \beta_2$ si $\xi_0 \leq \xi < \delta$.)  Si $f$ est aussi définie en
$\delta$ et que $f(\delta) = \lim_\delta f$, on dit que $f$ est
\defin[continue (fonction ordinale)]{continue} en $\delta$.

Ainsi, si $\delta$ est un ordinal limite, on peut écrire $\delta =
\lim_{\xi\to\delta} \xi$ (et réciproquement, si $f$ est
\emph{strictement} croissante, alors $\lim_{\xi\to\delta} f(\xi)$ est
forcément un ordinal limite).

À titre d'exemple, si $(u_n)$ est une suite croissante d'entiers
naturels, sa limite en tant que fonction ordinale $\omega \to \omega$
est soit un entier naturel (lorsque la suite est bornée, donc
constante à partir d'un certain rang) soit $\omega$ (lorsque la suite
n'est pas bornée).  Notamment, $\lim_{n\to\omega} 2^n = \omega$ (ce
qui permettra de dire que $2^\omega = \omega$ quand on aura défini cet
objet).


\subsection{Somme, produit et exponentielle d'ordinaux}

Les résultats de cette section seront admis (ils ne sont pas très
difficiles à montrer — presque toujours par induction transfinie —
mais seraient trop longs à traiter en détails).

\thingy\label{definition-sum-of-ordinals} Il existe deux façons
équivalentes de définir la somme $\alpha+\beta$ de deux ordinaux.

La première façon consiste à prendre un ensemble bien-ordonné $W$ tel
que $\alpha = \#W$ et un ensemble bien-ordonné $W'$ tel que $\beta =
\#W'$, et définir $\alpha + \beta := \#W''$ où $W''$ est l'ensemble
bien-ordonné qui est réunion disjointe de $W$ et $W'$ avec l'ordre qui
place $W'$ \emph{après} $W$, c'est-à-dire formellement $W'' := \{(0,w)
: w\in W\} \cup \{(1,w') : w'\in W'\}$ ordonné en posant $(i,w_1) <
(i,w_2)$ ssi $w_1 < w_2$ et $(0,w) < (1,w')$ quels que soient $w\in W$
et $w' \in W'$ (il est facile de voir qu'il s'agit bien d'un bon
ordre).

Autrement dit, intuitivement, une rangée de $\alpha+\beta$ allumettes
s'obtient en ajoutant $\beta$ allumettes \emph{après} (i.e., \emph{à
  droite d'})une rangée de $\alpha$ allumettes.

La seconde façon consiste à définir $\alpha+\beta$ par induction
transfinie sur $\beta$ (le \emph{second} terme) :
\begin{itemize}
\item $\alpha + 0 = \alpha$,
\item $\alpha + (\beta+1) = (\alpha+\beta) + 1$ (cas successeur),
\item $\alpha + \delta = \lim_{\xi\to\delta} (\alpha+\xi)$ si $\delta$
  est limite.
\end{itemize}

Nous ne ferons pas la vérification du fait que ces définitions sont
bien équivalentes, qui n'est cependant pas difficile (il s'agit de
vérifier que la première définition vérifie bien les clauses
inductives de la seconde).

\thingy Quelques propriétés de l'addition des ordinaux sont les
suivantes :
\begin{itemize}
\item l'addition est associative, c'est-à-dire que
  $(\alpha+\beta)+\gamma = \alpha+(\beta+\gamma)$ (on notera donc
  simplement $\alpha+\beta+\gamma$ et de même quand il y a plus de
  termes) ;
\item l'ordinal nul est neutre à gauche comme à droite, c'est-à-dire
  que $0+\alpha = \alpha = \alpha+0$ ;
\item le successeur de $\alpha$ est $\alpha + 1$ ;
\item l'addition n'est pas commutative en général : par exemple,
  $1+\omega = \omega$ (en décalant d'un cran toutes les allumettes)
  alors que $\omega + 1 > \omega$ ;
\item l'addition est croissante en chaque variable, et même
  strictement croissante en la seconde (si $\alpha\leq\alpha'$ alors
  $\alpha+\beta \leq \alpha'+\beta$, et si $\beta<\beta'$ alors
  $\alpha+\beta < \alpha+\beta'$ ; le fait que $0<1$ mais $0+\omega =
  1+\omega$ explique qu'il n'y a pas croissante stricte en la première
  variable) ;
\item l'addition est continue en la seconde variable (c'est exactement
  ce que dit le cas limite dans la définition par induction
  transfinie) ;
\item lorsque $\alpha \leq \alpha'$, il existe un unique $\beta$ tel
  que $\alpha' = \alpha + \beta$ (certains auteurs le notent $-\alpha
  + \alpha'$ : on prendra garde au fait qu'il s'agit d'une
  soustraction \emph{à gauche}) ;
\item comme conséquence de l'une des deux propriétés précédentes : si
  $\alpha + \beta = \alpha + \beta'$ alors $\beta = \beta'$
  (simplification \emph{à gauche} des sommes ordinales).
\end{itemize}

\thingy On pourrait aussi définir des sommes de séries d'ordinaux, ces
séries étant elles-mêmes indicées par d'autres ordinaux (le cas des
séries ordinaires étant le cas où l'ensemble d'indices est $\omega$).
Précisément, si $\alpha_\iota$ est un ordinal pour tout $\iota <
\gamma$ (avec $\gamma$ un autre ordinal), on peut définir
$\sum_{\iota<\gamma} \alpha_\iota$ par induction transfinie
sur $\gamma$ :
\begin{itemize}
\item $\sum_{\iota<0} \alpha_\iota = 0$ (somme vide !),
\item $\sum_{\iota<\gamma+1} \alpha_\iota = \big(\sum_{\iota<\gamma}
  \alpha_\iota\big) + \alpha_\gamma$ (cas successeur),
\item $\sum_{\iota<\delta} \alpha_\iota = \lim_{\xi\to\delta}
  \sum_{\iota<\xi} \alpha_\iota$ (cas limite).
\end{itemize}
Ainsi, dans le cas d'une série indicée par les entiers naturels,
$\sum_{n<\omega} \alpha_n$ est la limite $n\to\omega$ de la suite
croissante d'ordinaux $\alpha_0 + \cdots + \alpha_{n-1}$ (limite qui
existe toujours en tant qu'ordinal).

Cette notion de somme peut servir à définir le produit, $\alpha\beta =
\sum_{\iota<\beta} \alpha$, mais on va le redéfinir de façon plus
simple :

\thingy\label{definition-product-of-ordinals}
Il existe deux façons équivalentes de définir le produit
$\alpha\cdot\beta$ (ou $\alpha\beta$) de deux ordinaux.

La première façon consiste à prendre un ensemble bien-ordonné $W$ tel
que $\alpha = \#W$ et un ensemble bien-ordonné $W'$ tel que $\beta =
\#W'$, et définir $\alpha \cdot \beta := \#(W\times W')$ où $W \times
W'$ est l'ensemble bien-ordonné qui est le produit cartésien de $W$ et
$W'$ avec l'ordre lexicographique donnant plus de poids à $W'$,
c'est-à-dire $(w_1,w_1') < (w_2,w_2')$ ssi $w_1' < w_2'$ ou bien $w_1'
= w_2'$ et $w_1 < w_2$ (il est facile de voir qu'il s'agit bien d'un
bon ordre).

Autrement dit, intuitivement, une rangée de $\alpha\beta$ allumettes
s'obtient en prenant une rangée de $\beta$ allumettes et en y
\emph{remplaçant} chaque allumette par une rangée de $\alpha$
allumettes.

La seconde façon consiste à définir $\alpha\beta$ par induction
transfinie sur $\beta$ (le \emph{second} facteur) :
\begin{itemize}
\item $\alpha \cdot 0 = 0$,
\item $\alpha \cdot (\beta+1) = (\alpha\cdot\beta) + \alpha$ (cas successeur),
\item $\alpha \cdot \delta = \lim_{\xi\to\delta} (\alpha\cdot\xi)$
  si $\delta$ est limite.
\end{itemize}

Nous ne ferons pas la vérification du fait que ces définitions sont
bien équivalentes, qui n'est cependant pas difficile (il s'agit de
vérifier que la première définition vérifie bien les clauses
inductives de la seconde).

\thingy Quelques propriétés de la multiplication des ordinaux sont les
suivantes :
\begin{itemize}
\item la multiplication est associative, c'est-à-dire que
  $(\alpha\beta)\gamma = \alpha(\beta\gamma)$ (on notera donc
  simplement $\alpha\beta\gamma$ et de même quand il y a plus de
  facteurs) ;
\item l'ordinal nul est absorbant à gauche comme à droite, c'est-à-dire
  que $0\cdot\alpha = 0 = \alpha\cdot0$ ;
\item l'ordinal $1$ est neutre à gauche comme à droite, c'est-à-dire
  que $1\cdot\alpha = \alpha = \alpha\cdot1$ ;
\item la multiplication est distributive \emph{à droite} sur
  l'addition, c'est-à-dire que $\alpha(\beta+\gamma) = \alpha\beta +
  \alpha\gamma$ (en particulier, $\alpha\cdot 2 = \alpha+\alpha$) ;
\item la multiplication n'est pas commutative en général : par exemple,
  $2\cdot\omega = \omega$ (en doublant chaque allumette)
  alors que $\omega \cdot 2 > \omega$ ;
\item la distributivité à gauche ne vaut pas en général : par exemple,
  $(1+1)\cdot\omega = 2\cdot\omega = \omega$ n'est pas égal à
  $\omega+\omega = \omega\cdot 2$ ;
\item la multiplication est croissante en chaque variable, et même
  strictement croissante en la seconde lorsque la première est non
  nulle (si $\alpha\leq\alpha'$ alors $\alpha\cdot\beta \leq
  \alpha'\cdot\beta$, et si $\beta<\beta'$ et $\alpha>0$ alors
  $\alpha\cdot\beta < \alpha\cdot\beta'$) ;
\item la multiplication est continue en la seconde variable (c'est
  exactement ce que dit le cas limite dans la définition par induction
  transfinie) ;
\item \defin{division euclidienne} : pour tout $\alpha$ (ici appelé
  dividende) et tout $\beta>0$ (ici appelé diviseur) il existe
  $\gamma$ (ici appelé quotient) et $\rho<\beta$ (ici appelé reste)
  uniques tels que $\alpha = \beta\gamma + \rho$ (on prendra garde au
  fait qu'il s'agit d'une division \emph{à gauche}) ;
\item comme conséquence de l'une des deux propriétés précédentes : si
  $\beta\gamma = \beta\gamma'$ avec $\beta>0$, alors $\gamma =
  \gamma'$ (simplification \emph{à gauche} des produits ordinaux).
\end{itemize}

À titre d'exemple concernant la division euclidienne, tout ordinal
$\alpha$ peut s'écrire de façon unique comme $\alpha = \omega\gamma +
r$ avec $r$ un entier naturel : on a alors $r>0$ si et seulement si
$r$ est successeur (les ordinaux limites sont donc exactement les
$\omega\gamma$ avec $\gamma>0$) ; ce $r$ sera le « chiffre des
unités » de l'écriture de $\alpha$ en forme normale de Cantor
($\xi_{(0)}$ dans la notation de \ref{base-tau-writing-of-ordinals}).
On peut aussi écrire tout ordinal $\alpha$ de façon unique comme
$\alpha = 2\gamma + r$ avec $r$ valant $0$ ou $1$ : on peut dire que
$\alpha$ est « pair » ou « impair » selon le cas (à titre d'exemple,
$\omega$ est pair car $\omega = 2\cdot\omega$) ; ce $r$ sera de même
le « chiffre des unités » de l'écriture binaire.

\thingy On pourrait aussi définir des produits d'ordinaux, ces
produits étant eux-mêmes indicés par d'autres ordinaux (le cas des
produits infinis ordinaires étant le cas où l'ensemble d'indices
est $\omega$).  Précisément, si $\alpha_\iota$ est un ordinal non nul
pour tout $\iota < \gamma$ (avec $\gamma$ un autre ordinal), on peut
définir $\prod_{\iota<\gamma} \alpha_\iota$ par induction transfinie
sur $\gamma$ :
\begin{itemize}
\item $\prod_{\iota<0} \alpha_\iota = 1$ (produit vide !),
\item $\prod_{\iota<\gamma+1} \alpha_\iota = \big(\prod_{\iota<\gamma}
  \alpha_\iota\big) \cdot \alpha_\gamma$ (cas successeur),
\item $\prod_{\iota<\delta} \alpha_\iota = \lim_{\xi\to\delta}
  \prod_{\iota<\xi} \alpha_\iota$ (cas limite),
\end{itemize}
et on peut étendre au cas où certains ordinaux sont nuls en décrétant
que, dans ce cas, le produit est nul (évidemment).

Ainsi, dans le cas d'un produit d'ordinaux non nuls indicé par les
entiers naturels, $\prod_{n<\omega} \alpha_n$ est la limite
$n\to\omega$ de la suite croissante d'ordinaux $\alpha_0 \cdots
\alpha_{n-1}$ (limite qui existe toujours en tant qu'ordinal).

Cette notion de produit peut servir à définir le produit,
$\alpha^\beta = \prod_{\iota<\beta} \alpha$, mais on va le redéfinir
de façon plus simple :

\thingy Il existe deux façons équivalentes de définir l'exponentielle
$\alpha^\beta$ de deux ordinaux.

La première façon consiste à prendre un ensemble bien-ordonné $W$ tel
que $\alpha = \#W$ et un ensemble bien-ordonné $W'$ tel que $\beta =
\#W'$, et définir $\alpha ^ \beta := \#(W^{(W')})$ où $W^{(W')}$ est
l'ensemble des fonctions $W' \to W$ \emph{prenant presque partout la
  valeur $0$}, c'est-à-dire partout sauf en un nombre fini de points
la plus petite valeur de $W$, cet ensemble étant muni de l'ordre
lexicographique donnant plus de poids aux grandes composantes de la
fonction, c'est-à-dire que $f < g$ lorsque le plus grand $w' \in W'$
tel que $f(w') \neq g(w')$ vérifie en fait $f(w') < g(w')$ (on peut
vérifier qu'il s'agit bien d'un bon ordre).

La seconde façon consiste à définir $\alpha^\beta$ par induction
transfinie sur $\beta$ (l'exposant) :
\begin{itemize}
\item $\alpha ^ 0 = 1$,
\item $\alpha ^ {\beta+1} = (\alpha^\beta) \cdot \alpha$ (cas successeur),
\item $\alpha ^ \delta = \lim_{\xi\to\delta} \alpha^\xi$
  si $\delta$ est limite.
\end{itemize}

Nous ne ferons pas la vérification du fait que ces définitions sont
bien équivalentes, qui n'est cependant pas difficile (il s'agit de
vérifier que la première définition vérifie bien les clauses
inductives de la seconde).

\thingy Quelques propriétés de l'exponentiation des ordinaux sont les
suivantes :
\begin{itemize}
\item pour tout $\beta$, on a $1^\beta = 1$ ;
\item pour tout $\beta>0$, on a $0^\beta = 0$ (en revanche, $0^0=1$) ;
\item on a $\alpha^{\beta+\gamma} = \alpha^\beta \cdot \alpha^\gamma$ ;
\item on a $\alpha^{\beta\gamma} = (\alpha^\beta)^\gamma$ ;
\item l'exponentiation est croissante en chaque variable (si on écarte
  $0^0 = 1$), et même strictement croissante en l'exposant ($\beta$)
  lorsque la base $\alpha$ est $>1$ ;
\item l'exponentiation est continue en l'exposant variable (c'est
  exactement ce que dit le cas limite dans la définition par induction
  transfinie).
\end{itemize}

Il peut être éclairant de vérifier $2^\omega = \omega$ avec les deux
définitions de l'exponentiation.  Selon la définition avec des
ensembles bien-ordonnés, $2^\omega = \#(\{0,1\}^{(\mathbb{N})})$ où
$\{0,1\}^{(\mathbb{N})}$ est l'ensemble des suites de $0$ et de $1$
dont presque tous les termes sont des $0$, ordonnées par l'ordre
lexicographique donnant le plus de poids aux valeurs lointaines de la
suite : or de telles suites peuvent se voir comme des écritures
binaire (écrites à l'envers, i.e., en commençant par le poids faible),
et se comparent comme des écritures binaires, si bien que
$\{0,1\}^{(\mathbb{N})}$ est isomorphe, en tant qu'ensemble
bien-ordonné, à l'ensemble $\mathbb{N}$ des naturels, et son ordinal
est bien $\omega$.  Selon la définition inductive, $2^\omega$ est la
limite de $2^n$ pour $n\to\omega$, c'est-à-dire la borne supérieure de
$\{2^0,2^1,2^2,2^3,\ldots\}$, or cette borne supérieure est $\omega$
(ce ne peut pas être plus, parce que tous les $2^n$ sont des entiers
naturels donc majorés par $\omega$, et ce ne peut pas être moins car
aucun ordinal $<\omega$, i.e., aucun entier naturel, ne majore tous
les $2^n$).

\thingy\label{base-tau-writing-of-ordinals} Soient $\alpha,\tau$ des
ordinaux avec $\tau>1$ (dans la pratique, on ne s'intéressera guère
qu'à $\tau = 2$ et $\tau = \omega$) : alors il existe une unique
écriture
\[
\alpha = \tau^{\gamma_s} \xi_s + \cdots + \tau^{\gamma_1} \xi_1
\]
où $\gamma_s > \cdots > \gamma_1$ et $\xi_s,\ldots,\xi_1$ tous non
nuls et strictement inférieurs à $\tau$, ou, ce qui revient au même
(mais en changeant $\xi_i$ en $\xi_{(\gamma_i)}$),
\[
\alpha = \cdots + \tau^\iota \xi_{(\iota)} + \cdots + \tau \xi_{(1)} + \xi_{(0)}
\]
où les $\xi_{(\iota)}$ sont tous nuls sauf un nombre fini (ce qui rend
finie la somme ci-dessus) et tous $<\tau$.

Deux telles expressions se comparent par l'ordre lexicographique
donnant le plus de poids aux puissances élevées de $\tau$.  On parle
d'écriture de $\alpha$ \defin[écriture en base $\tau$]{en base $\tau$} : on dit que les
$\xi_{(\iota)}$ sont les \emph{chiffres} de cette écriture.  On
souligne que les chiffres sont \emph{tous nuls sauf un nombre fini}
(ce qui permet de les comparer lexicographiquement).

Les deux cas les plus importants sont $\tau=2$ et $\tau=\omega$ : le
cas $\tau=2$ correspond à l'\index{binaire (écriture)}\defin{écriture binaire} d'un ordinal,
c'est-à-dire son écriture comme somme décroissante finie de puissances
de $2$ distinctes, et le cas $\tau=\omega$ s'appelle écriture en
\defin[Cantor (forme normale de)]{forme normale de Cantor}, c'est-à-dire comme somme
décroissante finie de puissances de $\omega$.

La forme normale de Cantor est la manière usuelle d'écrire les
ordinaux (par exemple, $\omega$, $\omega+7$, $\omega\cdot 5$,
$\omega^2$, $\omega^\omega$ ou encore $\omega^{\omega\cdot 2}$ sont
des formes normales de Cantor, fussent-elles un peu dégénérées ; un
exemple moins dégénéré serait $\omega^{\omega 3}\cdot 7 +
\omega^{\omega+5}\cdot 42 + \omega^3 + 666$) ; on va expliquer au
paragraphe suivant que la forme normale de Cantor itérée (c'est-à-dire
appliquée aux exposants $\gamma_i$ eux-mêmes, et à leurs exposants, et
ainsi de suite) permet de « comprendre » et de manipuler
informatiquement un ordinal $\varepsilon_0$ passablement grand.
L'écriture binaire est moins souvent utilisée pour les ordinaux, et
son rapport avec la forme normale de Cantor sera expliqué
en \ref{binary-versus-cantor-normal-form}.

\thingy La forme normale de Cantor (ou plus exactement, la forme
normale de Cantor \emph{itérée}, c'est-à-dire appliquée récursivement
aux exposants de la forme normale de Cantor) permet de comprendre, et
de manipuler informatiquement, les ordinaux strictement inférieurs à
un certain ordinal \index{epsilon 0 (ordinal)}appelé $\varepsilon_0$.

Plus exactement, on peut donner la définition suivante :
un ordinal $<\varepsilon_0$ est \emph{soit} un entier
naturel $n$ (qui pourra aussi s'écrire $\omega^0\cdot n$ si on le
souhaite), qu'on compare comme on compare usuellement les entiers
naturels, \emph{soit} une écriture de la forme $\omega^{\gamma_s} n_s
+ \cdots + \omega^{\gamma_1} n_1$ où $\gamma_s > \cdots > \gamma_1$
sont eux-mêmes des ordinaux $<\varepsilon_0$ (précédemment définis),
et les $n_i$ sont des entiers naturels non nuls, et de telles
expressions se comparent par ordre lexicographique (i.e., comparer
d'abord $\gamma_s$, ou, s'ils sont égaux, comparer les $n_s$, ou,
s'ils sont égaux, comparer le terme suivant, etc.).

À titre d'exemple, $\alpha := \omega^{\omega^{\omega 3}\cdot 7 +
  \omega^{\omega+5}\cdot 42}\cdot 1729 + \omega^{\omega^{\omega
    2}\cdot 1000 + \omega + 33}\cdot 299\,792\,458$ est un ordinal
$<\varepsilon_0$ (écrit en forme normale de Cantor) car le premier
exposant, $\omega^{\omega 3}\cdot 7 + \omega^{\omega+5}\cdot 42$, est
lui-même un ordinal $<\varepsilon_0$ écrit en forme normale de Cantor
et supérieur au second exposant, $\omega^{\omega 2}\cdot 1000 + \omega
+ 33$.  L'ordinal $\alpha$ est plus grand que $\beta :=
\omega^{\omega^{\omega 3}\cdot 7 + \omega^{\omega+4}\cdot 333}\cdot
2016 + \omega^{\omega^{\omega 3}\cdot 5 + \omega + 33}\cdot
299\,792\,458$ car (une fois qu'on a vérifié, de même, que $\beta$ est
correctement écrit) le premier exposant de $\alpha$, à savoir
$\omega^{\omega 3}\cdot 7 + \omega^{\omega+5}\cdot 42$, est supérieur
au premier exposant de $\beta$, à savoir $\omega^{\omega 3}\cdot 7 +
\omega^{\omega+4}\cdot 333$ (cette comparaison se fait elle-même en
comparant le premier exposant, dans les deux cas $\omega\cdot 3$, puis
son coefficient, dans les deux cas $7$, puis l'exposant suivant, et
c'est là qu'on constate que $\omega+5$ dépasse $\omega+4$).

Formellement, on peut définir $\varepsilon_0$ comme l'ordinal (le $\#$
au sens de \ref{definition-of-ordinals}) de l'ensemble bien-ordonné
des écritures qu'on vient de présenter pour l'ordre qu'on vient de
dire.  Cet ordinal est la limite de
$\omega,\omega^\omega,\omega^{\omega^\omega},\omega^{\omega^{\omega^\omega}},\ldots$
c'est-à-dire le plus petit ordinal tel que $\varepsilon =
\omega^\varepsilon$.

\thingy (\textbf{Digression.}  Plus généralement, on appelle
$\varepsilon_\alpha$ le $\alpha$-ième ordinal tel que $\varepsilon =
\omega^\varepsilon$ : si on préfère, $\varepsilon_{\beta+1}$ est la
limite de $\varepsilon_\beta, \varepsilon_\beta^{\varepsilon_\beta},
\varepsilon_\beta^{\varepsilon_\beta^{\varepsilon_\beta}}, \ldots$, et
lorsque $\delta$ est limite, $\varepsilon_\delta$ est la limite des
$\varepsilon_\xi$ pour $\xi\to\delta$ ; on pourrait utiliser ce genre
de notations jusqu'à
$\varepsilon_{\varepsilon_{\varepsilon_{\varepsilon_{...}}}}$ qui est le
plus petit ordinal tel que $\zeta = \varepsilon_\zeta$.  Mais il ne
faut pas s'imaginer qu'on puisse espérer épuiser les ordinaux par ce
genre d'opérations : tous les ordinaux que nous venons de mentionner
sont \emph{dénombrables}, c'est-à-dire qu'ils sont les ordinaux de
certains bons ordres sur $\mathbb{N}$ (ou un ensemble fini), et toutes
les opérations $(\alpha,\beta) \mapsto \alpha+\beta$, $(\alpha,\beta)
\mapsto \alpha\beta$, $(\alpha,\beta) \mapsto \alpha^\beta$ et même
$\alpha \mapsto \varepsilon_\alpha$ envoient les ordinaux dénombrables
sur les ordinaux dénombrables.  Or il existe des ordinaux
indénombrables, le plus petit étant appelé $\omega_1$, qui est donc la
borne supérieure des ordinaux dénombrables, ou, dans la construction
de von Neumann, l'ensemble des ordinaux dénombrables : cet ordinal a
la propriété que \emph{toute suite strictement croissante (indicée par
  les entiers naturels) à valeurs dans [les ordinaux plus petits
    que] $\omega_1$ est bornée}, c'est-à-dire qu'on ne peut jamais le
fabriquer comme limite d'une suite.)

\thingy Expliquons rapidement pourquoi la forme normale de Cantor
permet de calculer la somme ou le produit de deux ordinaux.

Pour l'addition, à part le fait qu'elle est associative, on a
simplement besoin de savoir que :
\begin{itemize}
\item si $\gamma < \gamma'$ alors $\omega^\gamma\cdot n +
  \omega^{\gamma'}\cdot n' = \omega^{\gamma'} \cdot n'$ (autrement
  dit, les plus grandes puissances de $\omega$ absorbent les plus
  petites situées à leur gauche),
\item et bien sûr, $\omega^\gamma\cdot n + \omega^\gamma \cdot n' =
  \omega^\gamma \cdot (n+n')$ par associativité à droite (on peut donc
  regrouper deux puissances égales).
\end{itemize}

À titre d'exemple, la somme de $\omega^{\omega 2}\cdot 2 +
\omega^7\cdot 3$ et $\omega^{\omega 2}\cdot 5 + 12$ dans cet ordre
vaut $\omega^{\omega 2}\cdot 7 + 12$, et dans l'autre sens elle vaut
$\omega^{\omega 2}\cdot 7 + \omega^7\cdot 3$.

Pour la multiplication, comme elle est distributive à droite et
associative, il suffit de savoir calculer le produit de
$\omega^{\gamma_s} n_s + \cdots + \omega^{\gamma_1} n_1$ (en forme
normale de Cantor, avec les $\gamma_i$ strictement décroissants et les
$n_i$ strictement positifs) par $\omega^{\gamma'}$ \emph{à droite}, et
par $n' > 0$ (entier) lui aussi \emph{à droite}.  Or on a :
\begin{itemize}
\item $(\omega^{\gamma_s} n_s + \cdots + \omega^{\gamma_1} n_1) \times
  \omega^{\gamma'} = \omega^{(\gamma_s + \gamma')}$ dès que
  $\gamma'>0$, et
\item $(\omega^{\gamma_s} n_s + \cdots + \omega^{\gamma_1} n_1) \times
  n' = \omega^{\gamma_s} n_s n' + \cdots + \omega^{\gamma_1} n_1$ dès
  que $n'>0$ (i.e., seul le coefficient le plus à gauche est multiplié
  par $n'$, les autres sont inchangés).
\end{itemize}

À titre d'exemple, le produit de $\omega^{\omega 2}\cdot 2 +
\omega^7\cdot 3$ et $\omega^{\omega 2}\cdot 5 + 12$ dans cet ordre
vaut $\omega^{\omega 4}\cdot 5 + \omega^{\omega 2}\cdot 24 +
\omega^7\cdot 3$, et dans l'autre sens il vaut $\omega^{\omega 4}\cdot
2 + \omega^{\omega 2 + 7}\cdot 3$.

\thingy\label{binary-versus-cantor-normal-form}
Puisque $\omega = 2^\omega$ (et par conséquent,
$\omega^\gamma\, 2^c = 2^{\omega\gamma + c}$), l'écriture binaire d'un
ordinal s'obtient en remplaçant chaque chiffre $n$ (un entier naturel)
dans sa forme normale de Cantor par l'écriture binaire de $n$ (somme
de $2^c$ pour des entiers naturels $c$ distincts).  À titre d'exemple,
\[
\begin{array}{rllll}
&\omega^{\omega^2\cdot 3}\cdot 5 &\strut + \omega^{\omega+1}\cdot 8
&\strut + \omega^{17}\cdot 6 &\strut + 12\\
=& 2^{\omega^3\cdot 3 + 2} + 2^{\omega^3\cdot 3} &\strut + 2^{\omega^2+\omega+3}
&\strut + 2^{\omega\cdot 17+2} + 2^{\omega\cdot 17+1} &\strut + 2^3 + 2^2
\end{array}
\]

\subsection{Retour sur le jeu de l'hydre}\label{subsection-hydra-game-again}

\thingy À tout arbre fini enraciné $T$ on peut associer un
ordinal $o(T) <\varepsilon_0$ par récurrence sur la profondeur de
l'arbre, de la façon suivante : si $T_1,\ldots,T_r$ sont les
sous-arbres ayant pour racine les fils de la racine, triés de façon
que $o(T_1) \geq \cdots \geq o(T_r)$, alors on pose $o(T) =
\omega^{o(T_1)} + \cdots + \omega^{o(T_r)}$ (autrement dit, quitte à
regrouper les puissances identiques, ceci donne la forme normale de
Cantor de $o(T)$).

À titre d'exemple, dans le dessin
\begin{center}
\begin{tikzpicture}[baseline=0]
\draw[very thin] (-1.5,0) -- (1.5,0);
\begin{scope}[every node/.style={circle,fill,inner sep=0.5mm}]
\node (P0) at (0,0) {};
\node (P1) at (0,1) {};
\node (P2) at (-1,2) {};
\node (P3) at (0,2) {};
\node (P4) at (1,2) {};
\node (P5) at (0.5,3) {};
\node (P6) at (1.5,3) {};
\end{scope}
\begin{scope}[line width=1.5pt]
\draw (P0) -- (P1);
\draw (P1) -- (P2);
\draw (P1) -- (P3);
\draw (P1) -- (P4);
\draw (P4) -- (P5);
\draw (P4) -- (P6);
\end{scope}
\node[anchor=west] at (P6) {$x$};
\node[anchor=west] at (P4) {$y$};
\node[anchor=west] at (P1) {$z$};
\end{tikzpicture}
\end{center}
le sous-arbre partant de $x$ a pour valeur $0$ (il n'a aucun fils),
celui partant de $y$ a pour valeur $\omega^0 + \omega^0 = 2$, celui
partant de $z$ a pour valeur $\omega^2 + 2$, et l'arbre tout entier a
pour valeur $\omega^{\omega^2 + 2}$.

Il est facile de se convaincre que si on remplace un $T_i$ par un
autre arbre $T'_i$ avec $o(T'_i) < o(T_i)$ alors l'arbre $T'$
résultant vérifie $o(T') < o(T)$.  Ceci marche donc à n'importe quel
niveau de remplacement.

\thingy Rappelons maintenant les règles du \index{hydre (jeu de
  l')}jeu de l'hydre qui ont été données
en \ref{introduction-hydra-game} : Hercule choisit une \emph{tête} de
l'hydre, c'est-à-dire une feuille $x$ de l'arbre, et la décapite en la
supprimant de l'arbre.  L'hydre se reproduit alors de la façon
suivante : soit $y$ le nœud parent de $x$ dans l'arbre, et $z$ le nœud
parent de $y$ (grand-parent de $x$, donc) : si l'un ou l'autre
n'existe pas, rien ne se passe (l'hydre passe son tour) ; sinon,
l'hydre choisit un entier naturel $n$ (aussi grand qu'elle veut) et
attache à $z$ autant de nouvelles copies de $y$ (mais sans la tête $x$
qui a été décapitée) qu'elle le souhaite.  Si on appelle $T_y$ le
sous-arbre de l'hydre partant du nœud $y$ avant décapitation et $T'_y$
le même après décapitation, on a $o(T'_y) < o(T_y)$ (puisqu'on a
décapité une tête), donc $\omega^{o(T'_y)}\cdot n < \omega^{o(T_y)}$
\emph{quel que soit $n$} entier naturel.  Ainsi, si on appelle $T_z$
et $T'_z$ les sous-arbres partant de $z$ avant et après décapitation +
reproduction, on voit que $T'_z$ est obtenu en remplaçant
$\omega^{o(T_y)}$ dans son écriture en forme normale de Cantor par un
terme $\omega^{o(T'_y)}\cdot n$ strictement plus petit.  On a donc
$o(T_z') < o(T_z)$, et cette inégalité stricte vaut encore pour
l'arbre tout entier.

On voit donc que si on considère la suite des ordinaux associés aux
différentes formes de l'hydre au cours d'une partie du jeu de l'hydre,
cette suite est \emph{strictement décroissante}, et d'après
\ref{decreasing-sequences-of-ordinals-terminate}
ou \ref{sets-of-ordinals-are-well-ordered}, le jeu termine donc
toujours en temps fini.



%
%
%

\section{Jeux combinatoires impartiaux à information parfaite}\label{section-combinatorial-impartial-games}

\subsection{Récapitulations}

\thingy On a introduit en \ref{introduction-graph-game} et plus
formellement en \ref{definition-impartial-combinatorial-game} la
notion de \emph{jeu combinatoire impartial} (à information parfaite)
associé à un graphe $G$ muni d'un sommet initial $x_0$ (c'est le jeu
où, partant de $x = x_0$, chacun des deux joueurs choisit à son tour
un voisin sortant de $x$ : si l'un des joueurs ne peut pas jouer, il a
perdu, tandis que si la confrontation se continue indéfiniment, elle
est considérée comme nulle) ; on a vu
en \ref{determinacy-of-perfect-information-games} que ces jeux sont
déterminés.

On fera dans toute cette partie l'hypothèse supplémentaire que le
graphe $G$ est bien-fondé (cf. \ref{definitions-graphs}), c'est-à-dire
qu'aucune confrontation ne peut être nulle (=durer indéfiniment) : il
arrive forcément un point où l'un des joueurs ne peut plus jouer (et a
donc perdu), et la détermination signifie qu'un des joueurs a une
stratégie \emph{gagnante}.  On peut qualifier un tel jeu de
\index{bien-fondé (jeu)}« bien-fondé » ou de \defin[terminant
  (jeu)]{terminant}.

\thingy\label{combinatorial-positions-as-games} Il va de soi qu'une
position $x$ d'un jeu combinatoire impartial $G$ peut elle-même être
considérée comme un jeu combinatoire impartial : le jeu joué \defin{à
  partir de là}, à savoir celui dont $x$ est la position initiale (on
avait fait une remarque semblable pour les jeux de Gale-Stewart
en \ref{gale-stewart-positions-as-games}) et dont l'ensemble des
positions est $G$ ou, mieux (cf. paragraphe suivant) l'aval de $x$.
On se permettra souvent d'identifier ainsi une position et un jeu.

Pour éviter des discussions inutiles, on fera aussi souvent
implicitement l'hypothèse, en parlant d'un jeu combinatoire impartial
$(G,x_0)$, que tout sommet de $G$ est accessible depuis $x_0$
(cf. \ref{definition-accessibility-downstream} ; i.e., il n'y a pas de
positions « inutiles » car inaccessibles) : on peut toujours se
ramener à cette hypothèse en remplaçant $G$ par l'aval de $x_0$,
c'est-à-dire, en supprimant les positions inaccessibles.

On notera parfois abusivement un jeu $G$ au lieu de la donnée
$(G,x_0)$ (c'est d'autant plus critiquable que $x_0$ est peut-être
plus important que $G$ comme expliqué aux deux paragraphes
précédents ; néanmoins, ce n'est pas vraiment grave car $x_0$ peut
être retrouvé, si on a fait l'hypothèse du paragraphe précédent, comme
le seul sommet à partir duquel tout sommet est accessible).

\smallbreak

On rappelle la définition de la fonction de Grundy
(généralisant \ref{definition-grundy-function} à des valeurs non
nécessairement finies, comme expliqué
en \ref{ordinal-valued-rank-and-grundy-function}) :

\begin{defn}\label{definition-grundy-function-again}
Soit $G$ un graphe orienté bien-fondé.  On appelle \defin[Grundy
  (fonction de)]{fonction de Grundy} sur $G$ la fonction $\gr$
définie sur $G$ et à valeurs ordinales définie inductivement (en
utilisant le théorème \ref{well-founded-definition}) par $\gr(x) =
\mex\{\gr(y) : y\in\outnb(x)\}$ (où $\mex S$ désigne le plus petit
ordinal n'appartenant pas à $S$, qui existe d'après
\ref{sets-of-ordinals-are-well-ordered}
et \ref{sup-and-strict-sup-of-sets-of-ordinals}).  On appelle valeur
(ou fonction) de Grundy du jeu combinatoire impartial (terminant)
associé à $(G,x_0)$ la valeur $\gr(x_0)$.
\end{defn}

On rappelle qu'on a vu (cf. \ref{discussion-n-and-p-vertices}
à \ref{ordinal-valued-rank-and-grundy-function}) que le second joueur
(=joueur Précédent) a une stratégie gagnante si et seulement si la
valeur de Grundy est $0$, tandis que le premier joueur (=joueur
suivant (Next)) en a une si et seulement si elle est $\neq 0$.  (On
parlera de \index{P-position}P-positions et de
\index{N-position}N-positions selon que la valeur de Grundy est nulle
ou non.)  La stratégie « universelle » consiste toujours à \emph{jouer
  de façon à annuler la valeur de Grundy} du sommet vers lequel on
joue (i.e., jouer vers les P-positions).

(On notera parfois $G\doteq 0$ pour dire que le second joueur a une
stratégie gagnante, i.e., la valeur de Grundy de $G$ est $0$, i.e., la
position initiale de $G$ est une P-position ; et $G\fuzzy 0$ pour dire
que le premier joueur a une stratégie gagnante, i.e., la valeur de
Grundy de $G$ est $\neq 0$, i.e., la position initiale de $G$ est une
N-position.  Ces notations seront surtout utiles dans la
section \ref{section-combinatorial-partizan-games}.)

La signification exacte de la valeur de Grundy (au-delà du fait
qu'elle soit nulle ou non nulle) est plus mystérieuse.  Pour
l'éclaircir, on va introduire les \emph{nimbres} (déjà évoqués
en \ref{introduction-nimbers-and-numbers}) :

\begin{defn}\label{definition-nimber}
Soit $\alpha$ un ordinal.  On appelle alors \defin{nimbre} associé
à $\alpha$, et on note $*\alpha$, le jeu combinatoire impartial dont
\begin{itemize}
\item l'ensemble des positions est l'ensemble des ordinaux $\beta \leq
  \alpha$ (c'est-à-dire, si on utilise la construction de von Neumann
  des ordinaux, cf. \ref{introduction-von-neumann-ordinals},
  l'ensemble $\alpha+1$),
\item la relation d'arête (définissant le graphe) est $>$,
  c'est-à-dire que les voisins sortants de $\beta\leq\alpha$ sont les
  ordinaux $\beta'<\beta$, et
\item la position initiale est $\alpha$.
\end{itemize}
Autrement dit, il s'agit du jeu où, partant de l'ordinal $\beta =
\alpha$, chaque joueur peut diminuer l'ordinal $\beta$, c'est-à-dire
le remplacer par un ordinal $\beta' < \beta$ de son choix (ce jeu
termine en temps fini d'après
\ref{decreasing-sequences-of-ordinals-terminate}
ou \ref{sets-of-ordinals-are-well-ordered}).

On dira aussi que $*\alpha$ est le jeu constituée d'« une rangée de
  $\alpha$ allumettes » au nim (si $\alpha$ est fini, cela coïncide
bien avec ce qu'on a dit en \ref{introduction-nim-game}).  En accord
avec la remarque \ref{combinatorial-positions-as-games}, on peut
identifier les positions du nimbre $*\alpha$ avec les jeux $*\beta$
pour $\beta\leq\alpha$.
\end{defn}

\begin{prop}\label{grundy-of-nimbers-triviality}
La valeur de Grundy du nimbre $*\alpha$ est $\alpha$.
\end{prop}
\begin{proof}
On procède par induction transfinie
(cf. \ref{scholion-transfinite-induction}) : si on sait que
$\gr(*\beta) = \beta$ pour tout $\beta<\alpha$, alors
$\gr(*\alpha)$ est le plus petit ordinal différent de $\beta$
pour $\beta<\alpha$, c'est donc exactement $\alpha$.
\end{proof}


\subsection{Somme de nim}\label{subsection-nim-sum}

\begin{defn}\label{definition-nim-sum-of-games}
Soient $G_1,G_2$ deux jeux combinatoires impartiaux dont on note
$x_1,x_2$ les positions initiales et $E_1,E_2$ les relations d'arêtes.
On appelle \defin{somme de nim} (ou simplement « somme ») de $G_1$ et
$G_2$, et on note $G_1 \oplus G_2$ le jeu combinatoire impartial dont
\begin{itemize}
\item l'ensemble des positions est $G_1 \times G_2$,
\item la relation d'arête (définissant le graphe) est $(E_1 \times
  \id_{G_2}) \cup (\id_{G_1} \times E_2)$, c'est-à-dire que les
  voisins sortants de $(y_1,y_2) \in G_1 \times G_2$ sont les
  $(z_1,y_2)$ avec $z_1$ voisin sortant de $y_1$ ainsi que les
  $(y_1,z_2)$ avec $z_2$ voisin sortant de $y_1$, et
\item la position initiale est $(x_1,x_2)$.
\end{itemize}
\end{defn}

\thingy Autrement dit, jouer à $G_1 \oplus G_2$ signifie que chaque
joueur a, lorsque son tour vient (depuis la position $(y_1,y_2)$), le
choix entre jouer dans $G_1$ (c'est-à-dire aller en $(z_1,y_2)$ avec
$z_1$ voisin sortant de $y_1$ dans $G_1$) \emph{ou exclusif} jouer
dans $G_2$ (c'est-à-dire aller en $(y_1,z_2)$ avec $z_2$ voisin
sortant de $y_2$ dans $G_2$).

Il est clair que la somme de nim est commutative et associative, au
sens où les jeux $G_1 \oplus G_2$ et $G_2 \oplus G_1$ sont isomorphes
(=les graphes sont isomorphes par un isomorphisme qui envoie la
position initiale de l'un sur celle de l'autre, en l'occurrence
$(y_1,y_2) \mapsto (y_2,y_1)$) et de même pour $(G_1 \oplus G_2)
\oplus G_3$ et $G_1 \oplus (G_2 \oplus G_3)$.  On peut donc parler
sans souci du jeu $G_1 \oplus \cdots \oplus G_n$ ou $\bigoplus_{i=1}^n
G_i$ pour la somme de nim d'un nombre fini de jeux.  Il s'agit
simplement du jeu où chaque joueur, lorsque son tour vient, a la
faculté de joueur un coup dans un et un seul des $G_i$.

Notamment, si $\alpha_1,\ldots,\alpha_n$ sont des ordinaux, le
\defin[nim (jeu de)]{jeu de nim} ayant $n$ rangées avec ces nombres
(« nimbres » !) d'allumettes est défini comme le jeu
$\bigoplus_{i=1}^n(*\alpha_i)$, c'est-à-dire le jeu où chaque joueur,
lorsque son tour vient, a la faculté de décroître un et un seul
des $\alpha_i$.

\begin{prop}\label{nim-sum-is-well-founded}
Si $G_1,G_2$ sont bien-fondés, alors $G_1\oplus G_2$ est bien-fondé.
\end{prop}
\begin{proof}
S'il existe une suite infinie $x_{(i)} = (x_{(i),1}, x_{(i),2})$ de
sommets de $G_1 \oplus G_2$ avec $x_{(i+1)}$ voisin sortant
de $x_{(i)}$, alors soit l'ensemble des $i$ tels que $x_{(i+1),1}$
soit voisin sortant de $x_{(i),1}$ soit celui des $i$ tels que
$x_{(i+1),2}$ soit voisin sortant de $x_{(i),2}$ doit être infini :
dans les deux cas on a une contradiction.  (Autre démonstration : la
relation d'accessibilité sur $G_1\oplus G_2$ définit l'ordre partiel
produit, qui est inclus, i.e., plus faible, que l'ordre
lexicographique, et ce dernier est bien-ordonné
d'après \ref{definition-product-of-ordinals}.)
\end{proof}

\begin{defn}\label{definition-nim-sum-of-ordinals}
Soient $\alpha_1,\alpha_2$ deux ordinaux.  On appelle \defin{somme de
  nim} de $\alpha_1$ et $\alpha_2$ et on note $\alpha_1 \oplus
\alpha_2$ l'ordinal défini inductivement
(cf. \ref{scholion-transfinite-definition}) par
\begin{align*}
\alpha_1 \oplus \alpha_2 = \mex\big(
& \{\beta_1\oplus\alpha_2 : \beta_1 < \alpha_1\}\\
\cup\, & \{\alpha_1\oplus\beta_2 : \beta_2 < \alpha_2\}\big)
\end{align*}
Autrement dit, il s'agit du plus petit ordinal qui n'est ni de la
forme $\beta_1\oplus\alpha_2$ pour $\beta_1 < \alpha_1$ ni de la forme
$\alpha_1\oplus\beta_2$ pour $\beta_2 < \alpha_2$ ; cette définition a
bien un sens d'après \ref{nim-sum-is-well-founded}.  Encore autrement
(en utilisant \ref{grundy-of-nimbers-triviality} et une induction
transfinie évidente), il s'agit de la valeur de Grundy du jeu
$(*\alpha_1) \oplus (*\alpha_2)$.
\end{defn}

\thingy Pour comprendre cette définition, le mieux est de calculer
quelques valeurs.  Voici le tableau des $n_1 \oplus n_2$ pour $n_1 <
16$ et $n_2 < 16$ :
{\small\[
\begin{array}{c|cccccccccccccccc}
\oplus&0&1&2&3&4&5&6&7&8&9&10&11&12&13&14&15\\\hline
0&0&1&2&3&4&5&6&7&8&9&10&11&12&13&14&15\\
1&1&0&3&2&5&4&7&6&9&8&11&10&13&12&15&14\\
2&2&3&0&1&6&7&4&5&10&11&8&9&14&15&12&13\\
3&3&2&1&0&7&6&5&4&11&10&9&8&15&14&13&12\\
4&4&5&6&7&0&1&2&3&12&13&14&15&8&9&10&11\\
5&5&4&7&6&1&0&3&2&13&12&15&14&9&8&11&10\\
6&6&7&4&5&2&3&0&1&14&15&12&13&10&11&8&9\\
7&7&6&5&4&3&2&1&0&15&14&13&12&11&10&9&8\\
8&8&9&10&11&12&13&14&15&0&1&2&3&4&5&6&7\\
9&9&8&11&10&13&12&15&14&1&0&3&2&5&4&7&6\\
10&10&11&8&9&14&15&12&13&2&3&0&1&6&7&4&5\\
11&11&10&9&8&15&14&13&12&3&2&1&0&7&6&5&4\\
12&12&13&14&15&8&9&10&11&4&5&6&7&0&1&2&3\\
13&13&12&15&14&9&8&11&10&5&4&7&6&1&0&3&2\\
14&14&15&12&13&10&11&8&9&6&7&4&5&2&3&0&1\\
15&15&14&13&12&11&10&9&8&7&6&5&4&3&2&1&0
\end{array}
\]}
Chaque case est calculée en prenant la \emph{plus petite valeur qui ne
  figure pas déjà plus à gauche dans la ligne ou plus haut dans la
  colonne}.

\thingy\label{remark-on-dealing-with-mex} Dans ce qui suit, on
utilisera souvent implicitement le raisonnement suivant : pour montrer
que $\mex S = \alpha$, il suffit de montrer que (1) tout élément de
$S$ est différent de $\alpha$, et que (2) tout ordinal $<\alpha$
appartient à $S$.

\begin{prop}\label{nim-sum-is-commutative}
L'opération $\oplus$ est commutative sur les ordinaux.
\end{prop}
\begin{proof}[Première démonstration]
Par induction transfinie sur $\alpha_1$ et $\alpha_2$, on prouve
$\alpha_2\oplus\alpha_1 = \alpha_1\oplus\alpha_2$ : en effet,
$\alpha_2\oplus\alpha_1 = \mex (\{\alpha_2\oplus\beta_1:
\beta_1<\alpha_1\} \cup \{\beta_2\oplus\alpha_1: \beta_2<\alpha_2\})$,
et par hypothèse d'induction ceci vaut $\mex (\{\beta_1\oplus\alpha_2:
\beta_1<\alpha_1\} \cup \{\alpha_1\oplus\beta_2: \beta_2<\alpha_2\}) =
\alpha_1\oplus\alpha_2$.
\end{proof}
\begin{proof}[Seconde démonstration]
Cela résulte de la commutativité de $\oplus$ sur les jeux et de
l'observation que $\alpha_1\oplus\alpha_2 = \gr(*\alpha_1 \oplus
*\alpha_2)$.
\end{proof}

\begin{prop}\label{zero-is-neutral-for-nim-sum}
L'ordinal $0$ est neutre pour $\oplus$.
\end{prop}
\begin{proof}[Première démonstration]
Par induction sur $\alpha$, on prouve $\alpha \oplus 0 = \alpha$ : en
effet, $\alpha \oplus 0 = \mex \{\beta\oplus 0: \beta<\alpha\}$, et
par hypothèse d'induction ceci vaut $\mex \{\beta: \beta<\alpha\} =
\alpha$.
\end{proof}
\begin{proof}[Seconde démonstration]
Cela résulte de l'observation que $\alpha\oplus 0 = \gr(*\alpha_1
\oplus *0)$ et du fait que $*0$ est le jeu trivial ayant une seule
position, si bien que $G\oplus *0 \cong G$ pour n'importe quel $G$.
\end{proof}

\begin{prop}\label{nim-sum-cancellative}
Pour tous $\alpha_1,\alpha_2,\alpha_2'$, si $\alpha_1\oplus\alpha_2 =
\alpha_1\oplus\alpha_2'$, alors $\alpha_2=\alpha_2'$.
\end{prop}
\begin{proof}
Si ce n'est pas le cas, supposons sans perte de généralité que
$\alpha_2'<\alpha_2$.  Alors $\alpha_1\oplus\alpha_2$ est le $\mex$ d'un
ensemble contenant $\alpha_1\oplus\alpha_2'$, donc il ne peut pas lui
être égal, d'où une contradiction.
\end{proof}

\begin{thm}\label{nim-sum-for-games-versus-ordinals}
Si $G_1,G_2$ sont deux jeux combinatoires impartiaux bien-fondés ayant
valeurs de Grundy respectivement $\alpha_1,\alpha_2$, alors la valeur
de Grundy de $G_1\oplus G_2$ est $\alpha_1\oplus\alpha_2$.
\end{thm}
\begin{proof}
On procède par induction bien-fondée sur les positions de $G_1\oplus
G_2$ (ce qui est justifié d'après \ref{nim-sum-is-well-founded}) : il
s'agit de montrer que $\gr(y_1 \oplus y_2) = \gr(y_1) \oplus \gr(y_2)$
(en notant $y_1\oplus y_2$ la position $(y_1,y_2)$ de $G_1\oplus
G_2$), et on peut supposer ce fait déjà connu lorsque l'un de $y_1$ ou
exclusif $y_2$ est remplacé par un voisin sortant.  Or $\gr(y_1 \oplus
y_2)$ est le plus petit ordinal qui n'est ni de la forme $\gr(z_1
\oplus y_2)$ (avec $z_1$ voisin sortant de $y_1$) ni de la forme
$\gr(y_1 \oplus z_2)$ (avec $z_2$ voisin sortant de $y_2$),
c'est-à-dire, par hypothèse d'induction, ni de la forme $\gr(z_1)
\oplus \gr(y_2)$ ni de la forme $\gr(y_1) \oplus \gr(z_2)$.  Mais
quand $z_1$ parcourt les voisins sortants de $y_1$, les ordinaux
$\gr(z_1)$ sont tous distincts de $\gr(y_1)$ et parcourent au moins
tous les ordinaux strictement plus petits que lui (c'est la définition
de $\gr$) : par conséquent, les $\gr(z_1) \oplus \gr(y_2)$ sont tous
distincts de $\gr(y_1) \oplus \gr(y_2)$ (on a
utilisé \ref{nim-sum-cancellative}) et parcourent au moins tous les
$\beta_1 \oplus \gr(y_2)$ pour $\beta_1 < \gr(y_1)$, et de même, les
$\gr(y_1) \oplus \gr(z_2)$ sont tous distincts de $\gr(y_1) \oplus
\gr(y_2)$ et parcourent au moins tous les $\gr(y_2) \oplus \beta_2$
pour $\beta_2 < \gr(y_2)$ : ceci montre bien
(cf. \ref{remark-on-dealing-with-mex}) que le mex de toutes ces
valeurs est $\gr(y_1) \oplus \gr(y_2)$, c'est-à-dire $\gr(y_1 \oplus
y_2) = \gr(y_1) \oplus \gr(y_2)$.
\end{proof}

\begin{prop}\label{nim-sum-associative}
L'opération $\oplus$ est associative.
\end{prop}
\begin{proof}[Première démonstration]
Par induction sur $\alpha_1$, $\alpha_2$ et $\alpha_3$, on prouve
$(\alpha_1\oplus\alpha_2) \oplus \alpha_3 = \alpha_1 \oplus
(\alpha_2\oplus\alpha_3)$.  Pour cela, il suffit de prouver que tout
ordinal $\xi$ strictement inférieur à l'un des deux membres est
différent de l'autre.  Par symétrie, supposons $\xi <
(\alpha_1\oplus\alpha_2) \oplus \alpha_3$ et on veut prouver $\xi \neq
\alpha_1 \oplus (\alpha_2\oplus\alpha_3)$ : alors on a soit $\xi =
\gamma \oplus \alpha_3$ avec $\gamma < \alpha_1\oplus\alpha_2$ qui
peut lui-même s'écrire (A) $\gamma = \beta_1\oplus\alpha_2$ où
$\beta_1<\alpha_1$ ou bien (B) $\gamma = \alpha_1\oplus\beta_2$ où
$\beta_2<\alpha_2$, soit (C) $\xi = (\alpha_1\oplus\alpha_2) \oplus
\alpha_3'$ où $\alpha_3'<\alpha_3$.  Dans le cas (A), en utilisant
l'hypothèse d'induction, on a maintenant $\xi =
(\beta_1\oplus\alpha_2) \oplus \alpha_3 = \beta_1 \oplus
(\alpha_2\oplus\alpha_3)$ et d'après \ref{nim-sum-cancellative} ceci
est différent de $\alpha_1 \oplus (\alpha_2\oplus\alpha_3)$.  Les cas
(B) et (C) sont tout à fait analogues.
\end{proof}
\begin{proof}[Seconde démonstration]
Cela résulte de l'associativité de $\oplus$ sur les jeux et de
l'observation que $(\alpha_1\oplus\alpha_2)\oplus\alpha_3 =
\gr(((*\alpha_1) \oplus (*\alpha_2)) \oplus (*\alpha_3))$ qui
utilise \ref{nim-sum-for-games-versus-ordinals}.
\end{proof}

\begin{prop}\label{nim-sum-has-characteristic-two}
Pour tout $\alpha$ on a $\alpha\oplus\alpha = 0$.
\end{prop}
\begin{proof}[Première démonstration]
Par induction sur $\alpha$, on prouve $\alpha\oplus\alpha = 0$.  Or on
a $\alpha\oplus\alpha = \mex\{\beta\oplus\alpha:\beta<\alpha\}$, donc
il suffit de montrer $\beta\oplus\alpha \neq 0$ pour tout
$\beta<\alpha$.  Mais si $\beta\oplus\alpha = 0$, puisque l'hypothèse
d'induction assure $\beta\oplus\beta = 0$,
avec \ref{nim-sum-cancellative} on conclut $\alpha=\beta$, d'où une
contradiction.
\end{proof}
\begin{proof}[Seconde démonstration]
Le second joueur a une stratégie gagnante au jeu
$(*\alpha)\oplus(*\alpha)$ consistant à reproduire chaque coup de son
adversaire sur l'autre ligne (assurant que la position reste toujours
de la forme $(*\beta)\oplus(*\beta)$).  On a donc
$\gr((*\alpha)\oplus(*\alpha))$, c'est-à-dire $\alpha \oplus \alpha =
0$.
\end{proof}

\begin{prop}\label{nim-sum-multiple-sum}
La somme de nim $\alpha_1\oplus\cdots\oplus\alpha_n$ est le plus petit
ordinal qui n'est pas de la forme
$\alpha_1\oplus\cdots\oplus\beta_i\oplus \cdots\oplus\alpha_n$, où
exactement un des $\alpha_i$ a été remplacé par un ordinal $\beta_i$
strictement plus petit.
\end{prop}
\begin{proof}[Première démonstration]
On procède par récurrence sur $n$.  Tout d'abord, en
utilisant \ref{nim-sum-cancellative}, on voit que chaque $\alpha_1
\oplus \cdots \oplus \beta_i \oplus \cdots \oplus \alpha_n$ est
différent de $\alpha_1 \oplus \cdots \oplus \alpha_n$.  D'autre part,
si $\xi < \alpha_1 \oplus \cdots \oplus \alpha_n$, alors quitte à
écrire $\alpha_1 \oplus \cdots \oplus \alpha_n = \alpha' \oplus
\alpha_n$ avec $\alpha' := \alpha_1 \oplus \cdots \oplus
\alpha_{n-1}$, on a soit $\xi = \gamma \oplus \alpha_n$ où $\gamma <
\alpha'$, soit $\xi = \alpha' \oplus \beta_n$ où $\beta_n <
\alpha_n$ ; le second cas a bien la forme $\xi = \alpha_1 \oplus
\cdots \oplus \alpha_{n-1} \oplus \beta_n$ voulue, et dans le premier
cas l'hypothèse de récurrence permet d'écrire $\gamma = \alpha_1
\oplus \cdots \oplus \beta_i \oplus \cdots \oplus \alpha_{n-1}$ avec
$\beta_i < \alpha_i$ (où $i\leq n-1$), d'où $\xi = \alpha_1 \oplus
\cdots \oplus \beta_i \oplus \cdots \oplus \alpha_n$.  Ceci prouve
bien que $\alpha_1 \oplus \cdots \oplus \alpha_n$ est le plus petit
ordinal qui ne soit pas de la forme $\alpha_1 \oplus \cdots \oplus
\beta \oplus \cdots \oplus \alpha_n$, comme souhaité.
\end{proof}
\begin{proof}[Seconde démonstration]
En appliquant \ref{nim-sum-for-games-versus-ordinals} (et une
récurrence immédiate sur $n$), on voit que
$\alpha_1\oplus\cdots\oplus\alpha_n =
\gr((*\alpha_1)\oplus\cdots\oplus(*\alpha_n))$.  Or les positions de
$(*\alpha_1)\oplus\cdots\oplus(*\alpha_n)$ sont justement les
$(*\alpha_1)\oplus\cdots\oplus(*\beta_i)\oplus
\cdots\oplus(*\alpha_n)$ pour $\beta_i < \alpha_i$, et la conclusion
découle de la définition de $\gr$.
\end{proof}

\begin{prop}\label{nim-sum-of-powers-of-two-is-ordinary-sum}
Soient $\gamma_1 > \cdots > \gamma_r$ des ordinaux : alors la somme
usuelle (cf. \ref{definition-sum-of-ordinals}) $2^{\gamma_1} + \cdots
+ 2^{\gamma_r}$ des puissances de $2$ correspondantes coïncide avec la
somme de nim $2^{\gamma_1} \oplus \cdots \oplus 2^{\gamma_r}$.
\end{prop}
\begin{proof}
On procède par induction transfinie sur $2^{\gamma_1} + \cdots +
2^{\gamma_r}$.  On veut montrer que $\alpha := 2^{\gamma_1} + \cdots +
2^{\gamma_r}$ est égal à $2^{\gamma_1} \oplus \cdots \oplus
2^{\gamma_r}$, c'est-à-dire, d'après \ref{nim-sum-multiple-sum}, qu'il
s'agit du plus petit ordinal qui n'est pas de la forme $2^{\gamma_1}
\oplus \cdots \oplus \beta_i \oplus \cdots \oplus 2^{\gamma_r}$ pour
un certain $\beta_i < 2^{\gamma_i}$.  Pour ceci
(cf. \ref{remark-on-dealing-with-mex}), il suffit de montrer que
(1) $\alpha$ n'est pas de la forme qu'on vient de dire, et que
(2) tout ordinal $<\alpha$ l'est.

Or dans la somme de nim $2^{\gamma_1} \oplus \cdots \oplus \beta_i
\oplus \cdots \oplus 2^{\gamma_r}$ avec $\beta_i < 2^{\gamma_i}$, en
écrivant $\beta_i$ en binaire et en utilisant l'hypothèse d'induction,
on a $\beta_i = 2^{\gamma'_1} \oplus \cdots \oplus 2^{\gamma'_s}$ où
$\gamma_i > \gamma'_1 > \cdots > \gamma'_s$.  Avec les propriétés déjà
démontrées sur la somme de nim (commutativité, associativité, et
surtout \ref{nim-sum-has-characteristic-two}), on peut supprimer les
puissances de $2$ qui apparaissent deux fois et on obtient ainsi une
somme de nim de puissances de $2$ distinctes qui fait intervenir les
puissances $2^{\gamma_1}$ à $2^{\gamma_{i-1}}$ et certaines puissances
strictement plus petites que $2^{\gamma_i}$.  En utilisant de nouveau
l'hypothèse d'induction, cette somme se revoit comme une écriture
binaire, et puisqu'il n'y a pas $2^{\gamma_i}$ dedans, elle est
strictement plus petite que $\alpha$ (on rappelle que les écritures
binaires des ordinaux se comparent lexicographiquement).  Ceci
montre (1).

Pour ce qui est de (2), considérons un ordinal $\alpha'<\alpha$, qui
s'écrit donc en binaire comme une somme de la forme $2^{\gamma_1} +
\cdots + 2^{\gamma_{i-1}} + 2^{\gamma''_1} + \cdots + 2^{\gamma''_s}$
avec $\gamma_i > \gamma''_1 > \cdots > \gamma''_s$.  L'hypothèse
d'induction permet de réécrire cette somme comme une somme de nim.
Quitte à ajouter $2^{\gamma_{i+1}},\ldots,2^{\gamma_r}$ dans la somme
et annuler les puissances de $2$ qui apparaissent deux fois, on voit
qu'on peut écrire $\alpha'$ sous la forme $2^{\gamma_1} \oplus \cdots
\oplus \beta_i \oplus \cdots \oplus 2^{\gamma_r}$ où $\beta_i$ est une
somme de nim de puissances de $2$ toutes strictement plus petites
que $2^{\gamma_i}$.  En utilisant de nouveau l'hypothèse d'induction,
cette somme de nim est une somme usuelle, c'est-à-dire une écriture
binaire, et $\beta_i < 2^{\gamma_i}$ puisqu'il n'y a pas de
$2^{\gamma_i}$ dans l'écriture binaire en question.  Ceci prouve (2).
\end{proof}

Récapitulons ce qu'on a montré :

\begin{thm}\label{summary-of-grundy-theory}
La somme de nim d'ordinaux (définie
en \ref{definition-nim-sum-of-ordinals}) peut se calculer en écrivant
les ordinaux en question en binaire et en effectuant le \emph{ou
  exclusif} de ces écritures binaires (c'est-à-dire que le coefficient
devant chaque puissance de $2$ donnée vaut $1$ lorsque exactement l'un
des coefficients des nombres ajoutés vaut $1$, et $0$ sinon).

La valeur de Grundy de la somme de nim de jeux combinatoires
impartiaux bien-fondés se calcule donc comme le \emph{ou exclusif} des
valeurs de Grundy des composantes.  Notamment, la valeur de Grundy
d'un jeu de nim est le \emph{ou exclusif} des nombres d'allumettes des
différentes lignes.

Enfin, la valeur de Grundy d'un jeu combinatoire impartial
bien-fondé $G$ peut se voir comme l'unique ordinal $\alpha$ tel que le
second joueur ait une stratégie gagnante dans $G \oplus *\alpha$.
\end{thm}
\begin{proof}
L'affirmation du premier paragraphe résulte
de \ref{nim-sum-of-powers-of-two-is-ordinary-sum} et
de \ref{nim-sum-has-characteristic-two} (ainsi que de la commutativité
et l'associativité, \textit{ad lib.}) pour annuler les puissances
de $2$ identiques.

L'affirmation du second paragraphe est une reformulation
de \ref{nim-sum-for-games-versus-ordinals} (et pour le jeu de nim,
cf. aussi \ref{grundy-of-nimbers-triviality}).

Enfin, l'affirmation du troisième paragraphe en résulte d'après
\ref{nim-sum-cancellative} (pour l'unicité de $\alpha$)
et \ref{nim-sum-has-characteristic-two}.
\end{proof}

\thingy On savait déjà que dire que la valeur de Grundy d'un jeu
combinatoire impartial bien-fondé $G$ vaut $0$ signifie que le second
joueur a une stratégie gagnante.  On voit maintenant que dire que
cette valeur vaut $1$ signifie que le second joueur a une stratégie
gagnante dans le jeu modifié où les joueurs peuvent passer un tour
exactement une fois dans le jeu (le premier joueur qui utilise cette
faculté la consomme pour tout le monde, et le jeu n'est fini qu'une
fois ce tour consommé) : en effet, c'est simplement une reformulation
du jeu $G \oplus *1$.

\thingy On peut définir une relation d'équivalence $\doteq$ entre jeux
combinatoires impartiaux bien-fondés par : $G \doteq H$ lorsque le
second joueur possède une stratégie gagnante au jeu $G\oplus H$
(c'est-à-dire que [la position initiale de] $G\oplus H$ est une
P-position, ou encore que sa valeur de Grundy est nulle).

Selon le modèle de \ref{conway-equivalence-on-partizan-games} plus
bas, on peut montrer que $\doteq$ est une relation d'équivalence et
qu'elle est compatible avec la somme de nim (c'est-à-dire que si
$G\doteq G'$ et $H\doteq H'$ alorsq $G\oplus H \doteq G'\oplus H'$) :
le point crucial est que si le second joueur possède une stratégie
gagnante à un jeu $G_0$ alors le joueur qui a une stratégie gagnante à
$G\oplus G_0$ est le même qu'à $G$ (comparer
avec \ref{basic-facts-on-sum-of-partizan-games} ci-dessous).

Cette relation $\doteq$ sera particulièrement pertinente pour les jeux
partisans (dont la théorie est esquissée dans la
section \ref{section-combinatorial-partizan-games} ci-dessous), mais
dans le cas des jeux impartiaux considérés ici, elle est complètement
contenue dans la valeur de Grundy : on a $G\doteq H$ si et seulement
si $\gr(G) = \gr(H)$ (en effet, $G\doteq H$ équivaut par définition au
fait que le second joueur ait une stratégie gagnante à $G\oplus H$,
c'est-à-dire $\gr(G\oplus H) = 0$, ce qui
d'après \ref{nim-sum-for-games-versus-ordinals} signifie $\gr(G)
\oplus \gr(H) = 0$, autrement dit,
cf. \ref{nim-sum-has-characteristic-two}, que $\gr(G) = \gr(H)$).
Bref, les « classes d'équivalence » pour la relation $\doteq$ sont
précisément données par les ordinaux à travers la valeur de Grundy.

\thingy À côté de la somme de nim définie
en \ref{definition-nim-sum-of-ordinals}, il existe aussi une opération
de \index{nim (produit de)|see{produit de nim}}\defin{produit de nim}
définie inductivement par
\[
\alpha\otimes\beta := \mex\Big\{(\alpha'\otimes\beta)
\oplus (\alpha\otimes\beta') \oplus (\alpha'\otimes\beta')
: \alpha'<\alpha,\text{~et~}\beta'<\beta\Big\}
\]

On renvoie à l'exercice \ref{game-for-nim-product} pour une
interprétation de cette opération en termes de jeux, et
à \ref{inductions-on-nim-product} pour plus d'informations sur cette
opération : on y montre notamment qu'elle est commutative et
distributive sur la somme de nim, et en fait, on peut voir que les
deux opérations de nim munissent la classe de tous les ordinaux d'une
structure de corps commutatif (de caractéristique $2$ et
algébriquement clos).



%
%
%

\section{Notions sur les combinatoires partisans à information parfaite}\label{section-combinatorial-partizan-games}

\subsection{Jeux partisans, ordre, et somme}

\begin{defn}\label{definition-partizan-combinatorial-game}
Soit $G$ un graphe orienté dont l'ensemble $E$ des arêtes est réunion
de deux sous-ensembles $L$ (les arêtes \emph{bleues}) et $R$ (les
arêtes \emph{rouges}) et soit $x_0$ un sommet de $G$.  Le
\index{partisan (jeu)}\index{information parfaite (jeu
  à)}\defin[combinatoire (jeu)]{jeu combinatoire partisan à
  information parfaite} associé à ces données est défini de la manière
suivante : partant de $x = x_0$, Blaise (=joueur Bleu, =joueur gauche,
=Left) et Roxane (=joueur Rouge, =joueur droite, =Right) choisissent
tour à tour un voisin sortant de $x$ avec la contrainte supplémentaire
que chacun doit suivre une arête de sa couleur, autrement dit, Blaise
choisit une arête bleue $(x_0,x_1) \in L$ de $G$, puis Roxane choisit
une arête rouge $(x_1,x_2) \in R$ de $G$, puis Blaise choisit une
arête $(x_2,x_3) \in L$, et ainsi de suite.  Si un joueur ne peut plus
jouer, il a perdu ; si la confrontation dure un temps infini, elle est
considérée comme nulle (ni gagnée ni perdue par les joueurs).  On peut
considérer le jeu où Blaise commence (qu'on vient de décrire) ou celui
où Roxane commence (exactement analogue : Roxane choisit $(x_0,x_1)
\in R$ puis Blaise choisit $(x_1,x_2) \in L$, etc.).

On dira que $y$ est un « voisin sortant bleu » ou une « option
  gauche » de $x$ lorsque $(x,y) \in L$, et de même un « voisin
  sortant rouge » ou une « option droite » de $x$ lorsque $(x,y) \in
R$.

Une arête à la fois rouge et bleue (i.e., empruntable par les deux
joueurs) sera dite \textbf{verte}\footnote{Il serait sans doute plus
  logique de la qualifier de violette...}.

Les jeux impartiaux sont identifiés aux jeux partisans pour lesquels
$L=R=E$ (i.e., toutes les arêtes sont vertes : les joueurs ont
toujours les mêmes ensembles d'options).
\end{defn}

\thingy On fera toujours l'hypothèse que le graphe est bien-fondé,
c'est-à-dire que la partie nulle est impossible.  Les mêmes remarques
qu'en \ref{combinatorial-positions-as-games} s'appliquent pour les
jeux partisans.

\thingy On ne précise pas si Blaise ou Roxane joue en premier.  Dans
chacun des cas, il résulte des techniques de la
section \ref{subsection-determinacy-of-perfect-information-games} (ou
de \ref{partizan-games-as-impartial-games} plus bas) que l'un des deux
joueurs a une stratégie gagnante.  Il résulte que quatre cas
(exclusifs et exhaustifs) sont possibles pour un jeu combinatoire
partisan terminant (=bien-fondé) $G$ :
\begin{itemize}
\item Blaise a une stratégie gagnante, qui que soit le joueur qui
  commence : dans ce cas, on dira que le jeu est \defin[positif
    (jeu)]{strictement positif} et on notera $G > 0$ ;
\item Roxane a une stratégie gagnante, qui que soit le joueur qui
  commence : dans ce cas, on dira que le jeu est \defin[négatif
    (jeu)]{strictement négatif} et on notera $G < 0$ ;
\item le second joueur a une stratégie gagnante, quel qu'il soit :
  dans ce cas, on dira que le jeu est \defin[nul (jeu)]{nul} (au sens
  de Conway) et on notera $G \doteq 0$ ;
\item le premier joueur a une stratégie gagnante, quel qu'il soit :
  dans ce cas, on dira que le jeu est \defin[flou (jeu)]{flou} et on
  notera $G \fuzzy 0$.
\end{itemize}
Pour motiver la convention de considérer les jeux où le \emph{second}
joueur a une stratégie gagnante comme nuls, on rappelle que dans le
cas où $G$ est impartial, ce sont les jeux dont la valeur de Grundy
est nulle (et notamment que ces jeux ont la propriété qu'on peut les
ignorer dans une somme de nim).

On définit aussi $G \geq 0$ (le jeu est alors qualifié de positif)
lorsque $G > 0$ ou bien $G \doteq 0$ : concrètement, cela signifie
donc que Blaise a une stratégie gagnante à condition qu'il joue en
\emph{second}.  De même, $G \leq 0$ signifie que Roxane a une
stratégie gagnante à condition qu'elle joue en second.  On note
parfois $G \vartriangleleft 0$ pour dire que $G<0$ ou bien $G\fuzzy 0$
(c'est-à-dire la négation de $G > 0$ : Blaise a une stratégie gagnante
à condition qu'il joue en \emph{premier}), et de même $G
\vartriangleright 0$ pour $G>0$ ou $G\fuzzy 0$.

\thingy On définit également l'\defin{opposé} d'un jeu combinatoire
partisan $G$ comme le jeu $-G$ (ou $\ominus G$) dans lequel les arêtes
bleues et rouges sont échangées (i.e., $L$ et $R$ sont échangés).  Il
va de soi que $-G$ est strictement positif, resp. strictement négatif,
resp. nul, resp. flou, selon que $G$ est strictement négatif,
resp. strictement positif, resp. nul, resp. flou.

\begin{defn}\label{definition-partizan-sum-of-games}
Soient $G_1,G_2$ deux jeux combinatoires partisans dont on note
$x_1,x_2$ les positions initiales et $E_1,E_2$ les relations d'arêtes,
avec $L_1,L_2$ les ensembles d'arêtes bleues et $R_1,R_2$ les rouges.
On appelle \defin{somme disjonctive}, ou simplement « somme », de
$G_1$ et $G_2$, et on note $G_1 + G_2$ (ou $G_1 \oplus G_2$) le jeu
combinatoire partisan dont
\begin{itemize}
\item l'ensemble des positions est $G_1 \times G_2$,
\item la relation d'arête (définissant le graphe) est $(E_1 \times
  \id_{G_2}) \cup (\id_{G_1} \times E_2)$, c'est-à-dire que les
  voisins sortants de $(y_1,y_2) \in G_1 \times G_2$ sont les
  $(z_1,y_2)$ avec $z_1$ voisin sortant de $y_1$ ainsi que les
  $(y_1,z_2)$ avec $z_2$ voisin sortant de $y_1$, 
\item la couleur des arêtes vient de $G_1$ et de $G_2$, c'est-à-dire
  que l'ensemble des arêtes bleues est $(L_1 \times \id_{G_2}) \cup
  (\id_{G_1} \times L_2)$ (formé des arêtes $((y_1,y_2),(z_1,y_2))$
  avec $(y_1,z_1)$ dans $L_1$ et des $((y_1,y_2),(y_1,z_2))$ avec
  $(y_2,z_2)$ dans $L_2$) et l'ensemble des arêtes rouges est de même
  $(R_1 \times \id_{G_2}) \cup (\id_{G_1} \times R_2)$, et
\item la position initiale est $(x_1,x_2)$.
\end{itemize}
\end{defn}

\thingy Autrement dit, comme pour les jeux impartiaux, jouer à $G_1 +
G_2$ signifie que chaque joueur a, lorsque son tour vient (depuis la
position $(y_1,y_2)$), le choix entre jouer dans $G_1$ (c'est-à-dire
aller en $(z_1,y_2)$ avec $z_1$ voisin sortant de $y_1$ dans $G_1$ de
la couleur du joueur) \emph{ou exclusif} jouer dans $G_2$
(c'est-à-dire aller en $(y_1,z_2)$ avec $z_2$ voisin sortant de $y_2$
dans $G_2$ de la couleur du joueur).  La somme de nim est le cas
particulier de la somme disjonctive appliquée aux jeux impartiaux.

Comme pour les jeux impartiaux, la somme est commutative et
associative (à isomorphisme près).

Toute la cohérence de la théorie est fondée sur la proposition
élémentaire suivante :

\begin{prop}\label{basic-facts-on-sum-of-partizan-games}
La somme de deux jeux (combinatoires partisans bien-fondés) nuls est
nulle.  La somme de deux jeux strictement positifs, ou d'un
strictement positif et d'un nul, est strictement positive.  La somme
de deux jeux strictement négatifs, ou d'un strictement négatif et d'un
nul, est strictement négative.  La somme d'un jeu flou et d'un jeu nul
est floue.  (En revanche, la somme de deux jeux flous n'est pas
nécessairement floue.)

Enfin, la somme d'un jeu et de son opposé est nulle.
\end{prop}
\begin{proof}
Supposons que le second joueur ait une stratégie gagnante dans chacun
de $G_1$ et $G_2$ : il en a aussi une à $G_1 + G_2$, consistant à
répliquer à chaque coup de son adversaire dans la même composante que
celui-ci a joué, selon la stratégie gagnante dans cette composante.
Ceci montre que si $G_1 \doteq 0$ et $G_2 \doteq 0$ alors $G_1 + G_2
\doteq 0$.

Supposons que Blaise ait une stratégie gagnante dans $G_1$ (i.e., $G_1
> 0$) et qu'il en ait encore une dans $G_2$ à condition de jouer en
second (i.e., $G_2 \geq 0$).  Alors il en a une dans $G_1 \oplus G_2$,
consistant à répliquer à chaque coup de Roxane dans la même composante
qu'elle a joué, selon la stratégie gagnante qu'il a dans la composante
en question ; et s'il doit jouer le premier coup, il jouera selon la
stratégie gagnante dans $G_1$.  Ceci montre que si $G_1 > 0$ et $G_2
\geq 0$ alors $G_1 + G_2 > 0$.

L'affirmation concernant les jeux négatifs est évidemment symétrique.

Supposons que le premier joueur ait une stratégie gagnante dans $G_1$
(i.e., $G_1 \fuzzy 0$) et que le second en ait une dans $G_2$ (i.e.,
$G_2 \doteq 0$).  Alors le premier joueur a une dans $G_1 \oplus G_2$,
consistant à jouer au premier coup dans $G_1$ selon la stratégie
gagnante de celui-ci, puis à répliquer à chaque coup de son adversaire
dans la même composante que celui-ci a joué, selon la stratégie
gagnante dans cette composante.  Ceci montre que si $G_1 \fuzzy 0$ et
$G_2 \doteq 0$ alors $G_1 + G_2 \fuzzy 0$.

(La somme de deux jeux flous n'est pas forcément floue puisque
$*\alpha$ est flou pour tout ordinal $\alpha>0$, or
$(*\alpha)\oplus(*\alpha) \doteq 0$ comme on l'a vu
en \ref{nim-sum-has-characteristic-two} ou comme il résulte de la
dernière affirmation.)

Enfin, le second joueur a une stratégie gagnante dans $G + (-G)$
consistant à répliquer chaque coup de son adversaire dans la
composante opposée à celle où celui-ci vient de jouer.  Ceci montre
que $G + (-G) \doteq 0$.
\end{proof}

\thingy On définit fort logiquement $G \doteq H$, resp. $G > H$,
resp. $G \geq H$, resp. $G\fuzzy H$, resp. etc., lorsque $G + (-H)$
est $\doteq 0$, resp. $>0$, resp. $\geq 0$, resp. $\fuzzy 0$,
resp. etc.  La relation $G \doteq H$ se lit « $G$ et $H$ sont égaux au
  sens de Conway », la relation $G > H$ (resp. $G < H$) se lit
« $G$ est strictement supérieur (resp. inférieur) à $H$ », la relation
$G \geq H$ (resp. $G \leq H$) se lit « $G$ est supérieur
  (resp. inférieur) ou égal à $H$ », la relation $G \fuzzy H$ se lit
« $G$ est confus par rapport à $H$ ».

Intuitivement, il faut comprendre $G>H$ comme signifiant que « Blaise
  a strictement plus d'avantage dans $G$ que dans $H$ » (i.e., $G$ est
strictement plus favorable à Blaise que $H$).

Sous ces conditions :

\begin{prop}\label{conway-equivalence-on-partizan-games}
La relation $\doteq$ est une relation d'équivalence.  Elle est
compatible à la somme (c'est-à-dire que si $G \doteq G'$ et $H \doteq
H'$ alors $G + H \doteq G' + H'$) ainsi qu'aux relations $>$, $\geq$,
$\fuzzy$, etc. (c'est-à-dire que si $G \doteq G'$ et $H \doteq H'$ et
$G > H$ alors $G' > H'$ et de même en remplaçant « $>$ » par une des
autres relations).

Les jeux combinatoires partisans bien-fondés modulo la relation
$\doteq$ forment un groupe abélien partiellement ordonné par la
relation $>$.
\end{prop}
\begin{proof}
Tout est complètement formel à partir de ce qu'on a dit
en \ref{basic-facts-on-sum-of-partizan-games}.  Le fait que $G \doteq
G$ signifie exactement que $G + (-G) \doteq 0$, ce qu'on a vu ; le
fait que $G \doteq G'$ implique $G' \doteq G$ vient de ce que l'opposé
d'un jeu nul est nul ; le fait que $G \doteq G'$ et $G' \doteq G''$
impliquent $G \doteq G''$ vient de ce que la somme de deux jeux nuls
est nulle.  La compatibilité à la somme vient aussi de ce que la somme
de jeux nuls est nulle (remarquer que $G' \doteq G + (G' + (-G))$).
La compatibilité à l'ordre vient de ce que la somme d'un jeu nul et
d'un jeu strictement positif est strictement positif.  Le fait que $>$
est une relation d'ordre vient de ce que la somme de deux jeux
strictement positifs est strictement positive (pour la transitivité).

Les propriétés de groupe sont claires : on a vu l'associativité et la
commutativité (à isomorphisme près, donc \textit{a fortiori} à
$\doteq$ près), et comme $G + (-G) \doteq 0$, on a bien l'existence
d'un symétrique.
\end{proof}

\thingy La théorie des jeux combinatoires partisans bien-fondés
développée par John H. Conway s'intéresse essentiellement aux jeux
\emph{modulo} la relation $\doteq$.  C'est-à-dire, aux propriétés des
jeux ou opérations dessus qui sont préservées par cette relation
d'équivalence : comme on vient de le voir, le « signe » du jeu, ou la
somme disjonctive en sont ; un exemple d'une opération qui \emph{n'est
  pas} compatible à $\doteq$ est l'« impartialisation » du jeu définie
en \ref{partizan-games-as-impartial-games} ci-dessous, ou une
opération plus simple consistant à rendre vertes toutes les arêtes du
jeu — il n'est pas difficile de trouver des exemples de jeux tels que
$G \doteq G'$ mais que cette relation ne vaille plus après
l'opération.

On peut appeler \defin[valeur (d'un jeu combinatoire partisan)]{valeur}
d'un jeu $G$ la classe de $G$ pour la relation $\doteq$.  La
proposition ci-dessus affirme donc que les \emph{valeurs} des jeux
combinatoires partisans bien-fondés forment un groupe abélien
partiellement ordonné.  Cette notion de valeur peut être considérée
comme une généralisation de la fonction de Grundy (puisque
d'après \ref{summary-of-grundy-theory}, pour deux jeux impartiaux
$H,H'$, la relation $H\doteq H'$ équivaut à $H \oplus H' \doteq 0$
c'est-à-dire $\gr(H\oplus H') = 0$ c'est-à-dire $\gr(H) \oplus \gr(H')
= 0$, c'est-à-dire $\gr(H) = \gr(H')$), et les \emph{nimbres} peuvent
être considérés comme les valeurs particulières des jeux impartiaux
(sur lesquelles la loi de groupe est le « ou exclusif » des
représentations binaires).  La structure du groupe des valeurs des
jeux combinatoires partisans (bien-fondés) généraux est cependant
beaucoup plus difficile à comprendre.  Les \emph{nombres surréels}
évoqués en \ref{subsection-surreal-numbers} permettent d'éclaircir un
petit peu cette structure.


\subsection{Lien entre jeux partisans et jeux impartiaux}

\thingy\label{partizan-games-as-impartial-games} Le point suivant
mérite d'être éclairci : \emph{on peut toujours décrire un jeu
  combinatoire partisan à partir d'un jeu impartial ou deux}.

En effet, à partir d'un jeu combinatoire partisan $(G,x_0)$, on peut
fabriquer un graphe $\tilde G$ dont les sommets sont l'ensemble $G
\times \{\mathrm{L},\mathrm{R}\}$ des couples formés d'un sommet
de $G$ et d'une étiquette $\mathrm{L}$ ou $\mathrm{R}$ (qui sert à
retenir « le joueur qui doit maintenant jouer »), avec une arête entre
$(x,\mathrm{L})$ et $(y,\mathrm{R})$ dans $\tilde G$ lorsque $(x,y)$
appartient à l'ensemble $L$ des arêtes bleues de $G$ et une arête
entre $(x,\mathrm{R})$ et $(y,\mathrm{L})$ dans $\tilde G$ lorsque
$(x,y)$ appartient à l'ensemble $R$ des arêtes rouges de $G$.
Autrement dit, le jeu $\tilde G$ retient dans sa position à quel
joueur il incombe de jouer, et permet à partir d'une position de ce
genre de suivre une arête de la couleur correspondante dans $G$ (après
quoi ce sera à l'autre joueur de jouer).  Selon qu'on prend pour
position initiale $(x_0,\mathrm{L})$ ou $(x_0,\mathrm{R})$, on obtient
deux jeux combinatoires impartiaux (l'un correspondant à « commencer à
  jouer en tant que Blaise » et l'autre à « commencer à jouer en tant
  que Roxane ») : on les notera, disons, $\tilde G_{\mathrm{L}}$ et
$\tilde G_{\mathrm{R}}$.  De plus, si $G$ est bien-fondé, $\tilde G$
l'est aussi.

Cette construction fait que les affirmations « Blaise a une stratégie
  gagnante dans [le jeu partisan] $G$ s'il joue en premier » et « le
  premier joueur a une stratégie gagnante dans [le jeu
    impartial] $\tilde G_{\mathrm{L}}$ » sont équivalentes presque par
définition, et de même « Blaise a une stratégie gagnante dans [le jeu
    partisan] $G$ s'il joue en second » est équivalent à « le second
  joueur a une stratégie gagnante dans [le jeu impartial] $\tilde
  G_{\mathrm{R}}$ » et de même en échangeant simultanément Blaise et
Roxane et $L$ et $R$.  Bref, on a les équivalences suivantes :
\begin{center}
\begin{tabular}{ccccc}
$G\doteq 0$&ssi&$\tilde G_{\mathrm{L}}\doteq 0$&et&$\tilde G_{\mathrm{R}}\doteq 0$\\
$G>0$&ssi&$\tilde G_{\mathrm{L}}\fuzzy 0$&et&$\tilde G_{\mathrm{R}}\doteq 0$\\
$G<0$&ssi&$\tilde G_{\mathrm{L}}\doteq 0$&et&$\tilde G_{\mathrm{R}}\fuzzy 0$\\
$G\fuzzy 0$&ssi&$\tilde G_{\mathrm{L}}\fuzzy 0$&et&$\tilde G_{\mathrm{R}}\fuzzy 0$\\
$G\geq 0$&ssi&&&$\tilde G_{\mathrm{R}}\doteq 0$\\
$G\leq 0$&ssi&$\tilde G_{\mathrm{L}}\doteq 0$&&\\
\end{tabular}
\end{center}
où lorsque $H$ est un jeu combinatoire impartial on a écrit $H\doteq
0$ pour dire que sa position initiale est une P-position (si on
préfère, $\gr(H) = 0$) et $H\fuzzy 0$ pour dire qu'elle est une
N-position (si on préfère, $\gr(H) \neq 0$).

\thingy À cause de la remarque du point précédent, on peut se demander
quel est l'intérêt de l'étude des jeux combinatoires partisans :
plutôt qu'étudier $G$, autant étudier les deux jeux impartiaux $\tilde
G_{\mathrm{L}}$ (correspondant à faire commencer Blaise) et $\tilde
G_{\mathrm{R}}$ (correspondant à faire commencer Roxane) au moyen de
la théorie des jeux combinatoires impartiaux.  La raison pour laquelle
cette approche ne marche pas est que \emph{la construction $\tilde G$
  ne se comporte pas bien vis-à-vis de la somme}, c'est-à-dire que
$\widetilde{G+H}$ ne coïncide pas du tout avec $\tilde G \oplus \tilde
H$.

Pour s'en rendre compte et comprendre la différence, le mieux est de
considérer la somme de deux copies du jeu $G$ des échecs\footnote{Les
  échecs ne sont peut-être pas un jeu bien-fondé au sens où nous
  l'entendons, et d'ailleurs on peut discuter la mathématisation des
  règles exactes, mais on veut juste donner une idée.} :
\begin{itemize}
\item Si on fait la somme $G + G$ de deux copies des échecs comme jeux
  \emph{partisans}, alors Blaise et Roxane ont chacun une couleur et
  gardent la même couleur tout du long du jeu : chacun, à son tour,
  peut jouer sur l'un ou l'autre échiquier, mais jouera toujours de la
  même couleur (du coup, sur un échiquier donné, il se peut qu'une
  couleur joue plusieurs fois de suite, ce qui n'est normalement pas
  permis par les règles des échecs).
\item Si on fait la somme $\tilde G \oplus \tilde G$ de deux copies
  des échecs comme jeux \emph{impartiaux} construits avec la
  construction $\tilde{{\cdot}}$ (c'est-à-dire en déplaçant
  l'information du joueur à qui il incombe de jouer dans la
  « position » des échecs), alors Blaise et Roxane jouent sur deux
  échiquiers et choisissent celui sur lequel ils vont faire un coup,
  mais ce coup sera fait de la couleur qui doit jouer sur cet
  échiquier-là (i.e., la couleur opposée à celle du joueur qui a joué
  en dernier sur cet échiquier-là) : du coup, Blaise et Roxane n'ont
  pas une couleur bien définie chacun, mais en contrepartie, la partie
  jouée sur un échiquier donné sera une partie d'échecs normale
  (alternant les deux couleurs).  En fait, comme $\tilde G \oplus
  \tilde G$ est forcément un jeu nul (puisque c'est la somme de deux
  jeux impartiaux égaux, cf. \ref{nim-sum-has-characteristic-two}), le
  second joueur a une stratégie gagnante ici (consistant, dans les
  faits, à faire jouer son adversaire contre lui-même !).
\end{itemize}

On voit bien qu'il s'agit de jeux très différents, et la première
construction (la somme disjonctive de jeux partisans) est plus
naturelle si on doit étudier quel joueur a une avance sur lequel.

\thingy Même si la construction $\tilde G$ n'est pas compatible avec
la somme comme on vient de l'expliquer, on peut se rappeler que, pour
toute position $x$ d'un jeu impartial $H$, on a $x \doteq 0$ si et
seulement si $y \fuzzy 0$ pour tout voisin sortant $y$ de $x$, et
inversement $x \fuzzy 0$ si et seulement si $y \doteq 0$ pour un
voisin sortant $y$ de $x$ (c'est une reformulation
de \ref{definition-grundy-kernel}).  On en déduit :

\begin{prop}
Soit $G$ un jeu combinatoire partisan bien-fondé.  Alors pour toute
option $x$ de $G$ (identifiée au jeu ayant $x$ pour position
initiale) :
\begin{itemize}
%% \item On a $x \doteq 0$ si et seulement si : $\ell \vartriangleleft 0$
%%   pour tout voisin sortant \emph{bleu} $\ell$ de $x$, et $r
%%   \vartriangleright 0$ pour tout voisin sortant \emph{rouge}
%%   $r$ de $x$.
%% \item On a $x > 0$ si et seulement si : $\ell \geq 0$ pour au moins un
%%   voisin sortant \emph{bleu} $\ell$ de $x$, et $r \vartriangleright 0$
%%   pour tout voisin sortant \emph{rouge} $r$ de $x$.
%% \item On a $x < 0$ si et seulement si : $\ell \vartriangleleft 0$ pour
%%   tout voisin sortant \emph{bleu} $\ell$ de $x$, et $r \leq 0$ pour au
%%   moins un voisin sortant \emph{rouge} $r$ de $x$.
%% \item On a $x \fuzzy 0$ si et seulement si : $\ell \geq 0$ pour au moins un
%%   voisin sortant \emph{bleu} $\ell$ de $x$, et $r \leq 0$ pour au
%%   moins un voisin sortant \emph{rouge} $r$ de $x$.
\item On a $x \geq 0$ si et seulement si : $r \vartriangleright 0$
  pour tout voisin sortant \emph{rouge} $r$ de $x$.
\item On a $x \leq 0$ si et seulement si : $\ell \vartriangleleft 0$
  pour tout voisin sortant \emph{bleu} $\ell$ de $x$.
\item On a $x \vartriangleright 0$ si et seulement si : $\ell \geq 0$
  pour au moins un voisin sortant \emph{bleu} $\ell$ de $x$.
\item On a $x \vartriangleleft 0$ si et seulement si : $r \leq 0$ pour
  au moins un voisin sortant \emph{rouge} $r$ de $x$.
\end{itemize}
Autrement dit : une position positive est une position depuis laquelle
Roxane ne peut jouer que vers des positions non négatives
(=strictement positives ou floues) ; et une position négative est une
position depuis laquelle Blaise ne peut jouer que vers des positions
non positives (=strictement négatives ou floues).
\end{prop}
\begin{proof}
On a par exemple $x \geq 0$ ssi et seulement si $\tilde x_{\mathrm{R}}
\doteq 0$, soit si et seulement si tout $\tilde y_{\mathrm{L}} =
(y,\mathrm{L})$ voisin sortant de $\tilde x_{\mathrm{R}} =
(x,\mathrm{R})$ vérifie $\tilde y_{\mathrm{L}} \fuzzy 0$, c'est-à-dire
$y \vartriangleright 0$ pour tout $y$ voisin sortant rouge de $x$.
Les autres cas sont analogues.
\end{proof}


\subsection{Les nombres surréels (une esquisse)}\label{subsection-surreal-numbers}

\thingy Parmi les jeux combinatoires partisans (ou leurs valeurs,
c'est-à-dire ces mêmes jeux vus modulo $\doteq$), il en est de
particulièrement importants qui sont appelés « nombres surréels » (ou
simplement « nombres ») : ils sont \emph{totalement} ordonnés (entre
deux nombres surréels $x,x'$, on peut avoir $x<x'$ ou $x\doteq x'$ ou
$x>x'$ mais jamais $x\fuzzy x'$), ils forment eux aussi un groupe
abélien (et même un corps !), et on peut s'en servir pour comparer les
jeux combinatoires partisans (bien-fondés) généraux.

Les nombres surréels sont par ailleurs remarquables en ce qu'ils
généralisent \emph{à la fois} les ordinaux et les nombres réels (et
contiennent des éléments surprenants comme $\omega-42$ ou
$\omega+\sqrt{2}$ ou $2\pi\omega$ ou $1/\omega$ ou $\sqrt{\omega}$
ou $\omega^{\sqrt{5}}$).

\begin{defn}
Soit $\alpha$ un ordinal et $\sigma\colon\{\beta : \beta<\alpha\} \to
\{+,-\}$ une fonction quelconque définie sur les ordinaux $<\alpha$ et
à valeurs dans $\{+,-\}$ (on dira que $\sigma$ est une \defin{suite de
  signes} et que $\alpha$ est sa \textbf{longueur}).  Le
\index{surréel (nombre)}\defin{nombre surréel} associé à ces données
est le jeu combinatoire partisan bien-fondé dont
\begin{itemize}
\item l'ensemble des positions est l'ensemble des ordinaux $\beta \leq
  \alpha$,
\item la relation d'arête (définissant le graphe) est $>$,
  c'est-à-dire que les voisins sortants de $\beta\leq\alpha$ sont les
  ordinaux $\beta'<\beta$,
\item l'arête $(\beta,\beta')$ est colorée en bleu si $\sigma(\beta')
  = +$ et en rouge si $\sigma(\beta') = -$, et
\item la position initiale est $\alpha$.
\end{itemize}
Lorsque la suite de signes $\sigma$ est constamment égale à $+$, le
nombre surréel défini est appelé nombre surréel associé à
l'ordinal $\alpha$.
\end{defn}

\thingy Autrement dit, on peut considérer qu'on a affaire à une rangée
de $\alpha$ allumettes, mais cette fois-ci elles sont coloriées en
bleu ou rouge : les allumettes doivent être retirées par la droite, et
le joueur qui joue doit avoir la même couleur que l'allumette retirée
la plus à gauche (par exemple, si $\sigma(0)=+$, l'allumette la plus à
gauche est bleue et seul Blaise a le droit de retirer toutes les
allumettes d'un coup ; si $\sigma(1)=-$, l'allumette suivante est
rouge et Roxane a le droit de retirer toutes les allumettes à droite
de celle-là incluse, etc.).  Au moins pour un ordinal $\alpha$ fini,
ce jeu peut être vu comme un cas particulier du jeu de Hackenbush
introduit en \ref{introduction-hackenbush}, pour un graphe formé d'une
seule tour verticale, en disposant les allumettes les unes sur les
autres plutôt qu'en ligne (dans ce cas, $\sigma(0)$ donne la couleur
de celle qui est le plus en bas et qui supporte toutes les autres,
$\sigma(1)$ donne la couleur de la suivante, etc.).

Le cas particulier introduit en \ref{introduction-nimbers-and-numbers}
sous le nom de nombre surréel associé à l'ordinal $\alpha$ est bien le
cas où $\sigma$ est constamment $+$ (seul le joueur bleu peut jouer à
décroître l'ordinal, l'autre joueur ne peut jamais rien faire).

L'opposé d'un nombre surréel défini par sa suite de signes s'obtient
en changeant tous les signes de la suite.

\thingy Il est évident qu'un nombre surréel défini par une suite de
signes $\sigma$ est strictement positif lorsque $\sigma(0)=+$ (Blaise
peut gagner en un seul coup en retirant immédiatement toutes les
allumettes, c'est-à-dire en jouant vers la position $0$), et
strictement négatif lorsque $\sigma(0)=-$.  Le seul nombre surréel qui
n'est ni strictement positif ni strictement négatif est $0$, défini
par l'ordinal $\alpha=0$ et la suite de signes vide.  Autrement dit,
le signe d'un nombre surréel est donné par le premier signe de la
suite de signes.  Un nombre surréel n'est jamais flou.

En fait, on peut se convaincre que les nombres surréels sont
totalement ordonnés par l'ordre lexicographique sur leurs suites de
signes : si $x$ est défini par $\sigma$ de longueur $\alpha$ et $x'$
par $\sigma'$ de longueur $\alpha'$, et si on appelle $\gamma$ la
longueur commune entre $\sigma$ et $\sigma'$ (c'est-à-dire le plus
grand ordinal $\leq\max(\alpha,\alpha')$ tel que $\sigma(\beta) =
\sigma'(\beta)$ si $\beta<\gamma$), alors on a $x<x'$ si et seulement
si $\sigma(\gamma) < \sigma(\gamma')$ où on convient que $- <
\mathrm{nd} < +$ avec $\mathrm{nd}$ signifiant « non défini »
(c'est-à-dire le cas où $\sigma$ ou $\sigma'$ a justement comme
longueur $\gamma$, de sorte que la valeur en $\gamma$ n'est pas
définie).  Par exemple, l'ordre sur tous les nombres surréels de
longueur $\leq 2$ représentés par leurs suites de signes est donné par
$(--) < (-) < (-+) < () < (+-) < (+) < (++)$ (en fait, il s'agit des
nombres $-2 < -1 < -\frac{1}{2} < 0 < \frac{1}{2} < 1 < 2$).
L'égalité au sens de Conway entre deux nombres surréels ne peut se
produire que s'ils ont la même longueur et la même suite de signes.

\thingy On peut montrer que la somme (i.e., la somme disjonctive, en
tant que jeux) de deux nombres surréels est égale au sens de Conway
(i.e., a la même valeur) qu'un (unique) nombre surréel : autrement
dit, les nombres surréels forment pour l'addition un sous-groupe des
jeux partisans.

À titre d'exemple, le jeu $x$ défini par la suite de signes $(+-)$
(c'est-à-dire de longueur $2$ avec $\sigma(0)=+$ et $\sigma(1)=-$ :
autrement dit, Blaise peut jouer vers le jeu nul $0 = ()$ et Roxane
peut jouer vers le jeu $1 = (+)$) vérifie $x+x \doteq 1$ : en effet,
on peut se convaincre que le second joueur quel qu'il soit a une
stratégie gagnante dans le jeu formé de la somme disjonctive de deux
copies de $(+-)$ et d'une copie de $-1 = (-)$ (l'opposé de
l'ordinal $1$) : par exemple, si Blaise commence à jouer dans $(+-) +
(+-) + (-)$, il joue efface forcément un des signes $+$ et le $-$ qui
suit, ce qui laisse $(+-) + (-)$ et Roxane gagne en jouant vers $(+) +
(-)$ ; tandis que si Roxane commence dans $(+-) + (+-) + (-)$, elle va
jouer soit vers $(+) + (+-) + (-)$ soit vers $(+-) + (+-)$, dans le
premier cas Blaise gagne en jouant $(+) + (-)$, dans le second c'est
encore plus facile.  On dira donc que $x$ correspond au nombre réel
$\frac{1}{2}$.

\thingy Il n'est cependant pas évident de calculer la suite de signes
de $x+x'$ à partir de celles de $x$ et de celle de $x'$.  (Pour donner
un exemple, ajouter $1$ à un nombre surréel directement à partir de sa
suite de signes se fait de la façon — bien compliquée — suivante : on
commence par sauter tous les blocs de signes identiques dont la
longueur est multiple de $\omega$ ; si on arrive ainsi au bout de la
suite ou bien que le signe suivant est un $+$, on insère un $+$ à cet
endroit-là ; s'il y a au moins deux signes $-$, on en retire un ; s'il
y a un unique signe $-$, soit c'est le dernier de la suite auquel cas
on le retire, soit il est lui-même suivi d'un signe $+$, auquel cas on
remplace cette combinaison $-+$ par $+-$.)

Les ordinaux se voient, comme on l'a déjà dit, comme les nombres
surréels dont la suite de signe n'a que des signes $+$ (l'addition sur
les nombres surréels ne redonne pas exactement l'addition usuelle sur
les ordinaux, puisque cette dernière n'est pas commutative, mais elle
n'est pas très éloignée : en fait, l'addition qu'on obtient est
l'addition terme à terme des formes normales de Cantor, c'est-à-dire
l'addition des mêmes puissances de $\omega$, opération également
appelée « somme naturelle » des ordinaux).

Les nombres dyadiques (ceux de la forme $\frac{p}{2^k}$) se voient
comme les nombres surréels dont la suite de signes est de longueur
\emph{finie} (par exemple $\frac{1}{2} = (+-)$).  Les nombres réels se
voient comme les nombres surréels dont la suite de signes et de
longueur $\leq\omega$ et qui ne se terminent pas par une infinité de
signes tous égaux (pour obtenir la suite de signes d'un nombre réel
strictement positif et qui n'est pas dyadique, on commence par mettre
un nombre de $+$ égal à sa partie entière, puis la séquence $+-$, puis
l'écriture binaire de la partie fractionnaire du nombre en remplaçant
$1$ par $+$ et $0$ par $-$).

\thingy On peut même définir une multiplication sur les nombres
surréels qui font d'eux un corps totalement ordonné et « réel-clos »
(c'est-à-dire que, comme sur les réels, tout polynôme de degré impair
a une racine et que tout élément positif a une racine carrée).  Cette
multiplication peut même être définie entre un nombre réel et [la
  valeur d']un jeu partisan.





%
%
%

\section{Exercices}\label{section-exercises}

(Les sections de cette partie sont numérotées de la même manière que
les parties de l'ensemble, et font appel aux notions correspondantes.)

\medbreak

\setbox0=\vbox\bgroup
\subsection{Introduction et typologie}
\egroup

\subsection{Jeux en forme normale}

\exercice\label{normal-form-game-exercise-two-by-two}

Alice joue contre Bob un jeu dans lequel elle choisit une option parmi
deux possibles appelées U et V, et Bob choisit une option parmi deux
appelées X et Y (les modalités du choix varient selon les questions
ci-dessous) : les gains d'\emph{Alice} (c'est-à-dire, la fonction
qu'elle cherche à maximiser) sont donnés par le tableau ci-dessous, en
fonction de son choix (colonne de gauche) et de celui de Bob (ligne du
haut) :

\begin{center}
\begin{tabular}{r|ccc}
$\downarrow$A, B$\rightarrow$&X&Y\\\hline
U&$3$&$1$\\
V&$0$&$4$\\
\end{tabular}
\end{center}

\smallbreak

(1) On suppose que Bob fait son choix \emph{après} Alice, et en
connaissant le choix d'Alice, et qu'il cherche à minimiser le gain
d'Alice (i.e., le gain de Bob est l'opposé de celui d'Alice).  Comment
Alice a-t-elle intérêt à jouer et comment Bob répondra-t-il ?  Quelle
est le gain d'Alice dans ce cas ?

\begin{corrige}
Si Alice choisit U, Bob répondra par Y et le gain d'Alice sera $1$ ;
si Alice choisit V, Bob répondra par X et le gain d'Alice sera $0$.
Comme Alice veut maximiser son gain, elle a intérêt à choisir U, et
Bob répondra par Y, et le gain d'Alice sera $1$ dans ce cas.
\end{corrige}

\smallbreak

(2) On suppose maintenant que Bob fait son choix \emph{avant} Alice,
et qu'Alice connaîtra le choix de Bob ; on suppose toujours que Bob
cherche à minimiser le gain d'Alice.  Que fera Bob et comment Alice
répondra-t-elle ?  Quelle est le gain d'Alice dans ce cas ?

\begin{corrige}
Si Bob choisit X, Alice répondra par U et le gain d'Alice sera $3$ ;
si Bob choisit Y, Alice répondra par V et le gain d'Alice sera $4$.
Comme Bob veut minimiser le gain d'Alice, il a intérêt à choisir X, et
Alice répondra par U, et le gain d'Alice sera $3$ dans ce cas.
\end{corrige}

\smallbreak

(3) On suppose qu'Alice et Bob font leur choix séparément, sans
connaître le choix de l'autre, et toujours que Bob cherche à minimiser
le gain d'Alice.  Comment ont-ils intérêt à faire leurs choix ?  Quel
est le gain (espéré) d'Alice dans ce cas ?

\begin{corrige}
On a affaire à un jeu à somme nulle écrit en forme normale :
l'algorithme \ref{zero-sum-games-by-linear-programming-algorithm} nous
indique qu'on obtient la stratégie optimale d'Alice en résolvant le
problème de programmation linéaire suivant :
\[
\left\{
\begin{aligned}
\text{maximiser\ }v&\text{\ avec}\\
p_U \geq 0\;, \;\; p_V \geq 0&\\
p_U + p_V &= 1\\
v - 3p_U - 0p_V &\leq 0\;\;\text{(X)}\\
v - 1p_U - 4p_V &\leq 0\;\;\text{(Y)}\\
\end{aligned}
\right.
\]
On peut l'écrire sous forme normale en réécrivant $v = v_+ - v_-$ avec
$v_+,v_- \geq 0$, mais on gagne un petit peu en remarquant que $v$
sera forcément positif puisque tous les gains du tableau le sont, donc
on peut ajouter la contrainte $v \geq 0$.  Pour la même raison, on
peut transformer la contrainte d'égalité $p_U + p_V = 1$ en
l'inégalité $p_U + p_V \leq 1$.  Une application de l'algorithme du
simplexe donne finalement l'optimum $v = 2$ atteint pour $p_U =
\frac{2}{3}$ et $p_V = \frac{1}{3}$, avec pour le dual $q_X =
\frac{1}{2}$ et $q_Y = \frac{1}{2}$ (les inégalités (X) et (Y) sont
toutes les deux saturées).

Autrement dit, Alice joue les options U et V avec probabilités
$\frac{2}{3}$ et $\frac{1}{3}$, Bob réplique avec les options X et Y
de façon équiprobable, et le gain espéré d'Alice est $2$, qui est la
valeur du jeu à somme nulle en forme normale considéré ici.
\end{corrige}

\smallbreak

(4) On suppose maintenant que Bob cherche à \emph{maximiser} le gain
d'Alice (i.e., il n'est plus son adversaire comme dans les questions
(1), (2) et (3), mais son allié).  On cherche à déterminer quels sont
les équilibres de Nash possibles.  On notera $(p_U,p_V, q_X,q_Y)$ un
profil de stratégies mixtes général, où $p_U,p_V$ (positifs de
somme $1$) sont les poids des deux options d'Alice (=probabilités
qu'elle les joue), et $q_X,q_Y$ (positifs, également de somme $1$) les
poids des deux options de Bob.  On va discuter selon le support des
stratégies (i.e., selon les ensembles d'options qui ont un poids
strictement positif).\spaceout (a) Quels sont les équilibres de Nash
évidents en stratégies pures ?  Expliquer pourquoi ce sont bien les
seuls équilibres de Nash où l'un des deux joueurs a une stratégie
pure.\spaceout (b) Calculer ce que doivent valoir $p_U$ et $p_V$ dans
un équilibre de Nash où $q_X > 0$ et $q_Y > 0$ (i.e., les options X et
Y de Bob sont dans le support), et ce que doivent valoir $q_X$ et
$q_Y$ dans un équilibre de Nash où $p_U > 0$ et $p_V > 0$ (i.e., les
options U et V d'Alice sont dans le support).  Ces contraintes
donnent-elles effectivement un équilibre de Nash ?\spaceout
(c) Conclure quant à l'ensemble des équilibres de Nash du jeu
considéré.

\begin{corrige}
(a) Deux équilibres de Nash sont évidents : si Alice joue (purement) U
  et Bob joue (purement) X, aucun n'a intérêt à changer puisque $3$
  est maximum sur la ligne et sur la colonne ; si Alice joue
  (purement) V et Bob joue (purement) Y, aucun n'a intérêt à changer
  puisque $4$ est maximum sur la ligne et sur la colonne.  Les gains
  d'Alice (et de Bob) dans chacun de ces cas sont donc $3$ et $4$
  respectivement.

Il s'agit là de l'ensemble des équilibres de Nash où l'un des deux
joueurs a une stratégie pure : par exemple, si Alice joue purement U,
Bob ne peut que répondre par purement X puisque le gain $3$ est
strictement plus grand que tout autre gain sur la ligne (i.e., donner
un poids non nul à une autre option de Bob ne pourrait que diminuer le
gain).

(b) Si $q_X > 0$ et $q_Y > 0$, les stratégies pures X et Y de Bob
donnent forcément le même gain, car si l'une d'elle donnait un gain
strictement supérieure à l'autre, Bob aurait intérêt à augmenter le
poids $q$ correspondant et améliorerait ainsi strictement sa réponse.
Autrement dit, l'espérance de gain contre la stratégie pure X,
c'est-à-dire $3 p_U$, est égale à l'espérance de gain contre la
stratégie pure Y, soit $p_U + 4 p_V$.  On a donc $3 p_U = p_U + 4
p_V$, et comme aussi $p_U + p_V = 1$ on trouve $(p_U, p_V) =
(\frac{2}{3}, \frac{1}{3})$ en résolvant le système (soit la même
stratégie mixte que trouvée en (3), ce qui n'est pas un hasard vu que
le signe des gains de Bob n'est pas du tout intervenu dans le
raisonnement).  De même, si $p_U > 0$ et $p_V > 0$, on a $3 q_X + q_Y
= 4 q_Y$, et comme aussi $q_X + q_Y = 1$ on trouve $(q_X, q_Y) =
(\frac{1}{2}, \frac{1}{2})$ en résolvant le système (même remarque).

Bref, on a un équilibre de Nash \emph{potentiel} donné par $(p_U,p_V,
q_X,q_Y) = (\frac{2}{3}, \frac{1}{3}, \frac{1}{2}, \frac{1}{2})$,
c'est-à-dire exactement le même profil que dans la question (3) quand
les joueurs étaient adversaires.  Ceci peut surprendre, mais le fait
est qu'aucun des deux joueurs n'a (strictement) intérêt à changer
unilatéralement sa stratégie puisque les deux options qui se
présentent à lui sont indifférentes (elles ont le même gain espéré)
compte tenu du comportement de l'autre.  On a donc bien trouvé un
troisième équilibre de Nash.

(c) Pour résumer, on a trois équilibres de Nash récapitulés par le
tableau (triés par ordre de gain espéré décroissant) :
\begin{center}
\vskip-\baselineskip
\begin{tabular}{cc|cc|c}
$p_U$&$p_V$&$q_X$&$q_Y$&gain\\\hline
$0$&$1$&$0$&$1$&$4$\\
$1$&$0$&$1$&$0$&$3$\\
$\frac{2}{3}$&$\frac{1}{3}$&$\frac{1}{2}$&$\frac{1}{2}$&$2$\\
\end{tabular}
\end{center}
et on a prouvé que c'étaient bien les seuls.
\end{corrige}

%
%
%

\exercice\label{normal-form-game-exercise-three-by-three}

Alice joue contre Bob un jeu dans lequel elle choisit une option parmi
trois possibles appelées U, V et W, et Bob choisit une option parmi
trois appelées X, Y et Z (les modalités du choix varient selon les
questions ci-dessous) : les gains d'\emph{Alice} (c'est-à-dire, la
fonction qu'elle cherche à maximiser) sont donnés par le tableau
ci-dessous, en fonction de son choix (colonne de gauche) et de celui
de Bob (ligne du haut) :

\begin{center}
\begin{tabular}{r|ccc}
$\downarrow$A, B$\rightarrow$&X&Y&Z\\\hline
U&$6$&$0$&$4$\\
V&$0$&$6$&$4$\\
W&$2$&$2$&$7$\\
\end{tabular}
\end{center}

\smallbreak

(1) On suppose que Bob fait son choix \emph{après} Alice, et en
connaissant le choix d'Alice, et qu'il cherche à minimiser le gain
d'Alice (i.e., le gain de Bob est l'opposé de celui d'Alice).  Comment
Alice a-t-elle intérêt à jouer et comment Bob répondra-t-il ?  Quelle
est le gain d'Alice dans ce cas ?

\begin{corrige}
Si Alice choisit U, Bob répondra par Y et le gain d'Alice sera $0$ ;
si Alice choisit V, Bob répondra par X et le gain d'Alice sera $0$ ;
si Alice choisit W, Bob répondra par X ou Y indifféremment et le gain
d'Alice sera $2$.  Comme Alice veut maximiser son gain, elle a intérêt
à choisir W, et Bob répondra par X ou Y indifféremment, et le gain
d'Alice sera $2$ dans ce cas.
\end{corrige}

\smallbreak

(2) On suppose maintenant que Bob fait son choix \emph{avant} Alice,
et qu'Alice connaîtra le choix de Bob ; on suppose toujours que Bob
cherche à minimiser le gain d'Alice.  Que fera Bob et comment Alice
répondra-t-elle ?  Quelle est le gain d'Alice dans ce cas ?

\begin{corrige}
Si Bob choisit X, Alice répondra par U et le gain d'Alice sera $6$ ;
si Bob choisit Y, Alice répondra par V et le gain d'Alice sera $6$ ;
si Bob choisit Z, Alice répondra par W indifféremment et le gain
d'Alice sera $7$.  Comme Bob veut minimiser le gain d'Alice, il a
intérêt à choisir X ou Y, et Alice répondra par U ou V respectivement,
et le gain d'Alice sera $6$ dans ce cas.
\end{corrige}

\smallbreak

(3) On suppose que Bob joue \emph{aléatoirement}, en choisissant de
façon équiprobable (i.e., avec probabilité $\frac{1}{3}$) chacune des
trois options (X, Y et Z) qui s'offrent à lui.  Alice a connaissance
de ce fait mais ne connaît pas le résultat du tirage : comment
a-t-elle intérêt à jouer ?  Quelle est le gain (espéré) d'Alice dans
ce cas ?

\begin{corrige}
Le choix de Bob étant purement aléatoire, le gain espéré d'Alice est
donné par la combinaison convexe correspondante des colonnes du
graphe : si elle choisit U, son gain espéré est $\frac{1}{3}\times 6 +
\frac{1}{3}\times 0 + \frac{1}{3}\times 4 = \frac{10}{3} = 3 +
\frac{1}{3}$, si elle choisit V, son gain espéré est
$\frac{1}{3}\times 0 + \frac{1}{3}\times 6 + \frac{1}{3}\times 4 =
\frac{10}{3} = 3 + \frac{1}{3}$, et si elle choisit W, son gain espéré
est $\frac{1}{3}\times 2 + \frac{1}{3}\times 2 + \frac{1}{3}\times 7 =
\frac{11}{3} = 3 + \frac{2}{3}$.  Alice a donc intérêt à choisir W,
pour un gain espéré de $3 + \frac{2}{3}$.
\end{corrige}

\smallbreak

(4) On suppose qu'Alice et Bob font leur choix séparément, sans
connaître le choix de l'autre, et toujours que Bob cherche à minimiser
le gain d'Alice.  Comment ont-ils intérêt à faire leurs choix ?  Quel
est le gain (espéré) d'Alice dans ce cas ?

\begin{corrige}
On a affaire à un jeu à somme nulle écrit en forme normale :
l'algorithme \ref{zero-sum-games-by-linear-programming-algorithm} nous
indique qu'on obtient la stratégie optimale d'Alice en résolvant le
problème de programmation linéaire suivant :
\[
\left\{
\begin{aligned}
\text{maximiser\ }v&\text{\ avec}\\
p_U \geq 0\;, \;\; p_V \geq 0\;, \;\; p_W \geq 0&\\
p_U + p_V + p_W &= 1\\
v - 6p_U - 0p_V - 2p_W &\leq 0\;\;\text{(X)}\\
v - 0p_U - 6p_V - 2p_W &\leq 0\;\;\text{(Y)}\\
v - 4p_U - 4p_V - 7p_W &\leq 0\;\;\text{(Z)}\\
\end{aligned}
\right.
\]
On peut l'écrire sous forme normale en réécrivant $v = v_+ - v_-$ avec
$v_+,v_- \geq 0$, mais on gagne un petit peu en remarquant que $v$
sera forcément positif puisque tous les gains du tableau le sont, donc
on peut ajouter la contrainte $v \geq 0$.  Une application de
l'algorithme du simplexe donne finalement l'optimum $v = 3$ atteint
pour $p_U = \frac{1}{2}$ et $p_V = \frac{1}{2}$ et $p_W = 0$, avec
pour le dual $q_X = \frac{1}{2}$ et $q_Y = \frac{1}{2}$ et $q_Z = 0$
(les inégalités (X) et (Y) sont saturées et (Z) ne l'est pas).

Autrement dit, Alice joue les options U et V de façon équiprobable,
Bob réplique avec les options X et Y de façon équiprobable, et le gain
espéré d'Alice est $3$, qui est la valeur du jeu à somme nulle en
forme normale considéré ici.
\end{corrige}

\smallbreak

(5) On suppose maintenant que Bob cherche à \emph{maximiser} le gain
d'Alice (i.e., il n'est plus son adversaire comme dans les questions
(1), (2) et (4), mais son allié).  On cherche à déterminer quels sont
les équilibres de Nash possibles.  On notera $(p_U,p_V,p_W,
q_X,q_Y,q_Z)$ un profil de stratégies mixtes général, où $p_U,p_V,p_W$
(positifs de somme $1$) sont les poids des trois options d'Alice
(=probabilités qu'elle les joue), et $q_X,q_Y,q_Z$ (positifs,
également de somme $1$) les poids des trois options de Bob.  On va
discuter selon le support des stratégies (i.e., selon les ensembles
d'options qui ont un poids strictement positif).\spaceout (a) Pour
commencer, quelles sont les équilibres de Nash évidents en stratégies
pures ?  Expliquer pourquoi ce sont bien les seuls équilibres de Nash
où l'un des deux joueurs a une stratégie pure.\spaceout (b) Rappeler
pourquoi dans un équilibre de Nash où $q_X > 0$ et $q_Y > 0$ (i.e.,
les options X et Y sont dans le support), on a $6 p_U + 2 p_W = 6 p_V
+ 2 p_W$ (les stratégies pures X et Y de Bob sont forcément
indifférentes compte tenu de la stratégie d'Alice), et comment on
obtient des égalités de ce genre pour toute autre hypothèse sur le
support.\spaceout (c) Montrer que l'hypothèse $q_X,q_Y,q_Z > 0$
(toutes les options de Bob sont dans le support) conduit à une valeur
impossible pour $p_U,p_V,p_W$, et en déduire qu'aucun équilibre de
Nash n'a les trois options de Bob dans le support.  Montrer ensuite
qu'aucun équilibre de Nash n'a les trois options d'Alice dans le
support (procéder de même et se ramener à ce qu'on vient de
montrer).\spaceout (d) Trouver les équilibres de Nash avec
$q_X,q_Y>0$, en étudiant chacun des cas $p_U=0$, $p_V=0$ et
$p_W=0$.\spaceout (e) Trouver les équilibres de Nash avec $q_X,q_Z>0$,
en étudiant chacun des cas $p_U=0$, $p_V=0$ et $p_W=0$.\spaceout
(f) Conclure quant à l'ensemble des équilibres de Nash du jeu
considéré.

\begin{corrige}
(a) Trois équilibres de Nash sont évidents : si Alice joue
  (purement) U et Bob joue (purement) X, aucun n'a intérêt à changer
  puisque $6$ est maximum sur la ligne et sur la colonne ; si Alice
  joue (purement) V et Bob joue (purement) Y, aucun n'a intérêt à
  changer puisque $6$ est maximum sur la ligne et sur la colonne ; et
  si Alice joue (purement) W et Bob joue (purement) Z, aucun n'a
  intérêt à changer puisque $7$ est maximum sur la ligne et sur la
  colonne.  Les gains d'Alice (et de Bob) dans chacun de ces trois cas
  sont donc $6$, $6$ et $7$ respectivement.

Il s'agit là de l'ensemble des équilibres de Nash où l'un des deux
joueurs a une stratégie pure : par exemple, si Alice joue purement U,
Bob ne peut que répondre par purement X puisque le gain $6$ est
strictement plus grand que tout autre gain sur la ligne (i.e., donner
un poids non nul à une autre option de Bob ne pourrait que diminuer le
gain).

(b) Si $q_X > 0$ et $q_Y > 0$, les stratégies pures X et Y de Bob
donnent forcément le même gain, car si l'une d'elle donnait un gain
strictement supérieure à l'autre, Bob aurait intérêt à augmenter le
poids $q$ correspondant et améliorerait ainsi strictement sa réponse.
Autrement dit, l'espérance de gain contre la stratégie pure X,
c'est-à-dire $6 p_U + 2 p_W$, est égale à l'espérance de gain contre
la stratégie pure Y, soit $6 p_V + 2 p_W$.  De même, si $q_X > 0$ et
$q_Z > 0$, on a $6 p_U + 2 p_W = 4 p_U + 4 p_V + 7 p_W$, et ainsi de
suite.  Ceci vaut aussi dans l'autre sens : si $p_U > 0$ et $p_V > 0$
alors $6 q_X + 4 q_Z = 6 q_Y + 4 q_Z$ par exemple (les stratégies
pures U et V d'Alice sont indifférentes).

(c) Si $q_X,q_Y,q_Z > 0$, d'après ce qu'on vient de dire en (b), on a
$6 p_U + 2 p_W = 6 p_V + 2 p_W = 4 p_U + 4 p_V + 7 p_W$, et comme
aussi $p_U + p_V + p_W = 1$ on trouve $(p_U, p_V, p_W) = (\frac{5}{8},
\frac{5}{8}, -\frac{1}{4})$ en résolvant le système.  Une de ces
valeurs est strictement négative, donc un tel équilibre de Nash n'est
pas possible.  Bref, ceci montre que tout équilibre de Nash omet
forcément une des options X, Y ou Z de Bob.

De même, si $p_U,p_V,p_W > 0$, d'après ce qu'on a dit en (b), on a $6
q_X + 4 q_Z = 6 q_Y + 4 q_Z = 2 q_X + 2 q_Y + 7 q_Z$, ce qui avec $q_X
+ q_Y + q_Z = 1$ se résout en $(q_X, q_Y, q_Z) = (\frac{3}{8},
\frac{3}{8}, \frac{1}{4})$.  Ces valeurs sont \textit{a priori}
possibles, mais comme on a montré ci-dessus que $q_X,q_Y,q_Z$ ne
peuvent pas être simultanément strictement positifs, cette possibilité
est exclue.  Bref, ceci montre que tout équilibre de Nash omet
forcément une des options U, V ou W d'Alice.

(d) Si $q_X,q_Y > 0$, on a $6 p_U + 2 p_W = 6 p_V + 2 p_W$, et on a
toujours $p_U + p_V + p_W = 1$.  On résout séparément chacun des
systèmes obtenu en ajoutant une des équations $p_U = 0$, $p_V = 0$ et
$p_W = 0$ (on sait qu'on a forcément une de ces égalités en vertu de
la conclusion de la question (c)).  Pour $p_U = 0$ ou $p_V = 0$ on
trouve $(p_U, p_V, p_W) = (0,0,1)$, c'est-à-dire une stratégie pure,
qu'on a déjà traitée.  Pour $p_W = 0$, on trouve $(p_U, p_V, p_W) =
(\frac{1}{2},\frac{1}{2},0)$, ce qui implique à son tour $(q_X, q_Y,
q_Z) = (\frac{1}{2},\frac{1}{2},0)$ (en résolvant $6 q_X + 4 q_Z = 6
q_Y + 4 q_Z$ avec $q_Z = 0$ et $q_X + q_Y + q_Z = 1$), or ceci n'est
pas un équilibre de Nash car le gain espéré est $3$ (pour Alice comme
pour Bob, bien sûr) et Bob fera mieux en répondant Z (pour un gain
espéré de $4$).

(e) Si $q_X,q_Z > 0$, on a $6 p_U + 2 p_W = 4 p_U + 4 p_V + 7 p_W$, et
on a toujours $p_U + p_V + p_W = 1$.  On résout séparément chacun des
systèmes obtenu en ajoutant une des équations $p_U = 0$, $p_V = 0$ et
$p_W = 0$ (on sait qu'on a forcément une de ces égalités en vertu de
la conclusion de la question (c)).  Pour $p_U = 0$, on trouve une
solution impossible (un des poids est strictement négatif).  Pour $p_W
= 0$, on trouve $(p_U, p_V, p_W) = (\frac{2}{3},\frac{1}{3},0)$, ce
qui implique à son tour $(q_X, q_Y, q_Z) = (0,0,1)$, c'est-à-dire une
stratégie pure, qu'on a déjà traitée.  Reste $p_V = 0$ : on trouve
alors $(p_U, p_V, p_W) = (\frac{5}{7},0,\frac{2}{7})$, ce qui implique
à son tour $(q_X, q_Y, q_Z) = (\frac{3}{7},0,\frac{4}{7})$, donnant un
gain espéré de $\frac{34}{7}$ ; on vérifie que la stratégie pure V
d'Alice donne un gain strictement inférieure contre celle
$(\frac{3}{7},0,\frac{4}{7})$ de Bob (à savoir, $\frac{16}{7}$) et que
la stratégie pure Y de Bob donne un gain strictement inférieur contre
celle $(\frac{5}{7},0,\frac{2}{7})$ d'Alice (en l'occurrence,
$\frac{4}{7}$), donc aucun des joueurs n'a intérêt à changer
unilatéralement sa stratégie (les deux options du support sont
indifférentes, et la troisième diminue strictement le gain).  Bref, on
a bien trouvé un équilibre de Nash en stratégies mixtes.

Comme le jeu est symétrique par échange simultané des options U et V
d'Alice et X et Y de Bob, on a aussi traité le cas $q_Y,q_Z > 0$.
Bref, on a traité les six possibilités de support de la stratégie de
Bob (les stratégies pures en (a), celles ayant les trois options dans
le support en (c), et celles ayant deux options dans le support en
(d) et (e)).

(f) Les cinq équilibres de Nash trouvés sont récapitulés par le
tableau (triés par ordre de gain espéré décroissant) :
\begin{center}
\begin{tabular}{ccc|ccc|c}
$p_U$&$p_V$&$p_W$&$q_X$&$q_Y$&$q_Z$&gain\\\hline
$0$&$0$&$1$&$0$&$0$&$1$&$7\phantom{\strut+\frac{0}{1}}$\\
$1$&$0$&$0$&$1$&$0$&$0$&$6\phantom{\strut+\frac{0}{1}}$\\
$0$&$1$&$0$&$0$&$1$&$0$&$6\phantom{\strut+\frac{0}{1}}$\\
$\frac{5}{7}$&$0$&$\frac{2}{7}$&$\frac{3}{7}$&$0$&$\frac{4}{7}$&$4+\frac{6}{7}$\\
$0$&$\frac{5}{7}$&$\frac{2}{7}$&$0$&$\frac{3}{7}$&$\frac{4}{7}$&$4+\frac{6}{7}$\\
\end{tabular}
\end{center}
et on a prouvé que c'étaient bien les seuls.
\end{corrige}


%
%
%

\exercice\label{three-player-game-exercise}

On considère le jeu en forme normale suivant : \emph{trois} joueurs
(Alice, Bob et Charlie, par exemple) choisissent indépendamment les
uns des autres un élément de l'ensemble $\{\mathtt{0},\mathtt{1}\}$ :
\begin{itemize}
\item s'ils ont tous les trois choisi la même option, ils gagnent
  tous $0$,
\item si l'un d'entre eux a choisi une option différente des deux
  autres, celui qui a choisi cette option gagne $2$ et chacun des deux
  autres gagne $-1$.
\end{itemize}

Il sera utile de remarquer que les joueurs ont des rôles complètement
symétriques, et que les options sont également symétriques.

(Attention, même si la somme des gains des trois joueurs vaut
toujours $0$, ce n'est pas un « jeu à somme nulle » au sens classique,
car ces derniers ne sont définis que pour \emph{deux} joueurs.)

\smallbreak

(1) Écrire le tableau des gains du jeu considéré.  (On choisira une
façon raisonnable de présenter un tableau à trois entrées, par exemple
comme plusieurs tableaux à deux entrées mis côte à côte.)

\begin{corrige}
On fait deux tableaux, l'un pour le cas où Alice joue $\mathtt{0}$,
\begin{center}
\begin{tabular}{r|cc}
$\mathtt{0}$A, $\downarrow$B, C$\rightarrow$&$\mathtt{0}$&$\mathtt{1}$\\\hline
$\mathtt{0}$&$0,0,0$&$-1,-1,+2$\\
$\mathtt{1}$&$-1,+2,-1$&$+2,-1,-1$\\
\end{tabular}
\end{center}
et l'autre pour le cas où Alice joue $\mathtt{1}$,
\begin{center}
\begin{tabular}{r|cc}
$\mathtt{1}$A, $\downarrow$B, C$\rightarrow$&$\mathtt{0}$&$\mathtt{1}$\\\hline
$\mathtt{0}$&$+2,-1,-1$&$-1,+2,-1$\\
$\mathtt{1}$&$-1,-1,+2$&$0,0,0$\\
\end{tabular}
\end{center}
Chacune des entrées doit bien sûr lister trois nombres, pour les gains
d'Alice, Bob et Charlie respectivement.
\end{corrige}

\smallbreak

Si $p \in [0;1]$, on notera simplement $p$ la stratégie mixte d'un
joueur qui consiste à choisir l'option $\mathtt{1}$ avec
probabilité $p$, et l'option $\mathtt{0}$ avec probabilité $1-p$.

(2) Vérifier que l'espérance de gain d'Alice si elle joue selon la
stratégie mixte $p$ tandis que Bob joue selon la stratégie mixte $q$
et Charlie selon la stratégie mixte $r$ vaut : $-2pq -2pr +4qr + 2p -
q -r$.  (Ici, $p,q,r$ sont trois réels entre $0$ et $1$.)

\begin{corrige}
Si Alice joue $\mathtt{0}$, son espérance de gain est $-q(1-r) -
(1-q)r + 2qr$ d'après le premier tableau donné en réponse à la
question précédente, soit $4qr - q - r$.  Si Alice joue $\mathtt{1}$,
son espérance de gain vaut $2(1-q)(1-r) -q(1-r) - (1-q)r = 4qr - 3q -
3r + 2$.  Si elle joue $p$, son espérance de gain vaut $1-p$ fois $4qr
- q - r$ plus $p$ fois $4qr - 3q - 3r + 2$, ce qui vaut l'expression
$-2pq -2pr +4qr + 2p - q -r$ annoncée.
\end{corrige}

\smallbreak

(3) On se demande à quelle condition sur la stratégie mixte $q$ jouée
par Bob et la stratégie mixte $r$ jouée par Charlie les options
$\mathtt{0}$ et $\mathtt{1}$ d'Alice sont indifférentes pour elle
(c'est-à-dire, lui apportent la même espérance de gain).  Montrer que
c'est le cas si et seulement si $q + r = 1$.

\begin{corrige}
On cherche à quelle condition la valeur $4qr - q - r$ (qui se retrouve
en substituant $0$ à $p$ dans $-2pq -2pr +4qr + 2p - q -r$) est égale
à $4qr - 3q - 3r + 2$ (obtenue en mettant $p$ à $1$).  La différence
entre les deux vaut $2 - 2q - 2r$, qui est donc nulle si et seulement
si $q+r = 1$, comme annoncé.
\end{corrige}

\smallbreak

(4) Déduire de la question (3) que si un profil $(p,q,r)$ de
stratégies mixtes est un équilibre de Nash et que $0<p<1$ alors
$q+r=1$.

\begin{corrige}
Si $(p,q,r)$ est un équilibre de Nash en stratégies mixtes et si
$0<p<1$, c'est-à-dire si Alice pondère effectivement ses deux
stratégies pures, c'est que les gains espérés qu'elles lui apportent
sont égaux (si l'une était strictement meilleure que l'autre, Alice
aurait strictement intérêt à ne jouer que celle-là), c'est-à-dire
$q+r=1$ comme on vient de le voir.
\end{corrige}

\smallbreak

(5) En déduire tous les équilibres de Nash $(p,q,r)$ du jeu (on pourra
distinguer des cas selon que $p=0$, $p=1$ ou $0<p<1$ et de même pour
$q$ et $r$ ; la symétrie doit permettre de simplifier le travail).

\begin{corrige}
Considérons un équilibre de Nash $(p,q,r)$.  On a vu en (4) que si
l'un des trois nombres n'est ni $0$ ni $1$, la somme des deux autres
vaut nécessairement $1$.

(A) Si au moins deux des trois nombres sont strictement entre $0$ et
$1$, disons sans perte de généralité que $p$ et $q$ le sont.  Alors
$q+r=1$ et $p+r=1$, ce qui donne $p=q$.  Mais le fait que $r=1-q$ avec
$0<q<1$ implique que $0<r<1$.  On a donc aussi $p+q=1$, ce qui
implique $p=q=\frac{1}{2}$ et du coup $r=\frac{1}{2}$ puisque le
raisonnement est complètement symétrique.  Or il est clair que
$(\frac{1}{2},\frac{1}{2},\frac{1}{2})$ est bien un équilibre de Nash
(toutes les options deviennent indifférentes pour tout le monde).

(B) Si un seul des trois nombres est strictement entre $0$ et $1$,
disons sans perte de généralité que $p$ l'est.  Alors $q+r=1$, et
comme $q$ et $r$ doivent valoir chacun $0$ ou $1$, les seules
possibilités sont $(p,0,1)$ et symétriquement $(p,1,0)$.  Vérifions
que $(p,0,1)$ constitue bien un équilibre de Nash, y compris si $p=0$
ou $p=1$ (le cas $(p,1,0)$ étant bien sûr symétrique) : dans
$(p,0,1)$, Alice a un gain espéré de $-1$ qui ne varie pas selon $p$ ;
Bob y a un gain espéré de $3p-1$, qui est supérieur ou égal au gain
$-3p+2$ qu'il espère obtenir en changeant d'option, et le cas de
Charlie est exactement symétrique.  On a donc bien affaire à un
équilibre de Nash.

(C) Enfin, si $p,q,r$ dont tous dans $\{0,1\}$, on a déjà vu en (B)
que si deux valent $0$ et un vaut $1$ ou le contraire, on a affaire à
un équilibre de Nash.  Reste le cas de $(0,0,0)$ ou $(1,1,1)$, et ce
ne sont certainement pas des équilibres de Nash car les trois joueurs
ont intérêt à changer unilatéralement de stratégie.

Finalement, on a trouvé comme équilibre de Nash : le point isolé
$(\frac{1}{2},\frac{1}{2},\frac{1}{2})$, et l'hexagone formé des
segments paramétrés par $(p,0,1)$, $(1,0,r)$, $(1,q,0)$, $(p,1,0)$,
$(0,1,r)$ et $(0,q,1)$ où la seule variable prend une valeur
quelconque dans $[0;1]$ (intuitivement, ce sont des équilibres où deux
joueurs sont en position de gagner, et le troisième, qui est en
position de « faiseur de roi », ne peut pas gagner mais choisit au
hasard lequel des deux autres gagne).
\end{corrige}

\smallbreak

(6) Dans cette question, on modifie le jeu : plutôt que faire leurs
choix indépendamment, les joueurs les font et les déclarent
successivement (Alice, puis Bob, puis Charlie).  (a) Que va faire Bob
si Alice choisit $\mathtt{0}$ ?  (b) Informellement, expliquer qui est
avantagé ou désavantagé par cette modification de la règle.

\begin{corrige}
(a) Si Alice choisit $\mathtt{0}$, Bob a bien sûr intérêt à
  choisir $\mathtt{1}$ (car s'il choisit lui-même $\mathtt{0}$,
  Charlie choisira $\mathtt{1}$ et Bob a le pire gain possible
  de $-1$).  Charlie sera alors dans la situation de choisir qui
  d'Alice ou de Bob gagne sans pouvoir gagner lui-même (il est
  « faiseur de roi »).  La situation est complètement symétrique si
  Alice choisit $\mathtt{1}$, et le choix d'Alice est complètement
  indifférent puisque les deux options sont équivalentes de son point
  de vue.

(b) On peut donc dire que Charlie est désavantagé par le fait de jouer
  en dernier : il ne pourra pas gagner, seulement choisir lequel des
  deux autres joueurs gagne.  (Ceci est un peu paradoxal quand on se
  rappelle que dans un jeu à \emph{deux} joueurs à somme nulle, on ne
  peut qu'être avantagé par le fait d'avoir connaissance de l'option
  choisie par l'adversaire.)
\end{corrige}

%
%
%

\subsection{Jeux de Gale-Stewart et détermination}

\exercice

On considère dans cet exercice une variante des jeux de Gale-Stewart
où au lieu qu'un joueur « gagne » et l'autre « perde », il va leur
être attribué un gain réel comme dans les jeux à somme nulle.  Plus
exactement, ce type de jeu est décrit comme suit.  On fixe un ensemble
$X$ (non vide) et une fonction $u \colon X^{\mathbb{N}} \to
\mathbb{R}$ dite fonction de gain d'Alice (la fonction de gain de Bob
serait $-u$).  Alice et Bob choisissent tour à tour un élément de $X$
comme pour un jeu de Gale-Stewart, et jouent un nombre infini de
coups, « à la fin » desquels la suite $\dblunderline{x} =
(x_0,x_1,x_2,\ldots)$ définit un élément $\dblunderline{x} \in
X^{\mathbb{N}}$.  Le gain d'Alice est alors $u(\dblunderline{x})$ (le but
d'Alice est de le maximiser) et le gain de Bob est $-u(\dblunderline{x})$
(le but de Bob est de maximiser cette quantité, c'est-à-dire de
minimiser $u(\dblunderline{x})$).

\smallbreak

(1) Expliquer pourquoi ce type de jeu peut être considéré comme une
généralisation des jeux de Gale-Stewart.

\begin{corrige}
Si $G_X(A)$ où $A \subseteq X^{\mathbb{N}}$ est un jeu de Gale-Stewart
défini par un sous-ensemble $A$ de $X^{\mathbb{N}}$, on peut définir
une fonction $u \colon X^{\mathbb{N}} \to \mathbb{R}$ par
$u(\dblunderline{x}) = 1$ si $x\in A$ et $u(\dblunderline{x}) = -1$
si $x\not\in A$, c'est-à-dire utiliser la valeur $+1$ pour un gain
d'Alice et $-1$ pour un gain de Bob.  Les buts d'Alice et de Bob dans
le jeu défini par cette fonction sont alors exactement les mêmes que
dans le jeu $G_X(A)$.
\end{corrige}

\smallbreak

(2) On suppose désormais que $u \colon X^{\mathbb{N}} \to \mathbb{R}$
est \emph{continue}, ce qui signifie par définition que pour tout
$\dblunderline{x} \in X^{\mathbb{N}}$ et tout $\varepsilon > 0$ réel, il
existe $\ell$ tel que pour tout $y \in V_\ell(\dblunderline{x})$ on ait
$|u(\dblunderline{y}) - u(\dblunderline{x})| < \varepsilon$.  Montrer que
quel que soit $v \in\mathbb{R}$ l'ensemble des $\dblunderline{x}$ tels
que $u(\dblunderline{x}) > v$ est ouvert, et de même pour l'ensemble des
$\dblunderline{x}$ tels que $u(\dblunderline{x}) < v$.

\begin{corrige}
Fixons un réel $v$.  Montrons que l'ensemble $A_v$ des $\dblunderline{x}$
tels que $u(\dblunderline{x}) > v$ est ouvert.  Si $\dblunderline{x}$ est
dans $A_v$, i.e., vérifie $u(\dblunderline{x}) > v$, alors la définition
de la continuité appliquée à ce $\dblunderline{x}$ et à $\varepsilon :=
u(\dblunderline{x}) - v$ assure qu'il existe $\ell$ tel que pour tout $y
\in V_\ell(\dblunderline{x})$ on ait $|u(\dblunderline{y}) -
u(\dblunderline{x})| < \varepsilon$, et cette dernière inégalité implique
$u(\dblunderline{y}) > u(\dblunderline{x}) - \varepsilon = v$, donc on a
encore $u(\dblunderline{y}) > v$, autrement dit $\dblunderline{y} \in A_v$.
On a donc montré que pour tout $\dblunderline{x}$ dans $A_v$ il existe
$\ell$ tel que tout $\dblunderline{y}$ dans $V_\ell(\dblunderline{x})$ soit
encore dans $A_v$, c'est exactement dire que $A_v$ est ouvert.  Le cas
de l'ensemble $B_v$ des $\dblunderline{x}$ tels que $u(\dblunderline{x}) <
v$ est exactement analogue (on prendra $\varepsilon := v -
u(\dblunderline{x})$).
\end{corrige}

\smallbreak

(3) (a) En déduire que pour tout $v$ réel, soit Alice possède une
stratégie lui garantissant un gain $\geq v$, soit Bob possède une
stratégie lui garantissant qu'Alice aura un gain $\leq v$
(c'est-à-dire que lui, Bob, aura un gain $\geq -v$).\spaceout
(b) Montrer que les deux ne peuvent se produire simultanément que pour
\emph{au plus un} réel $v$.  Un tel $v$ s'appelle la \textbf{valeur}
du jeu.

\begin{corrige}
(a) Quel que soit $v$ réel, on vient de voir que l'ensemble $A_v$ des
  $\dblunderline{x}$ tels que $u(\dblunderline{x}) > v$ est ouvert.  La
  détermination des jeux de Gale-Stewart ouverts
  (théorème \ref{gale-stewart-theorem}) implique donc que soit Alice a
  une stratégie lui garantissant de tomber dans $A_v$, c'est-à-dire un
  gain $>v$, et \textit{a fortiori $\geq v$}, soit Bob a une stratégie
  qui lui garantit de tomber dans le complémentaire de $A_v$,
  c'est-à-dire garantissant qu'Alice aura un gain $\leq v$.

(b) S'il existait deux réels $v' < v$ tels qu'Alice possède une
  stratégie lui garantissant un gain $\geq v$ et que Bob possède une
  stratégie lui garantissant qu'Alice aura un gain $\leq v'$, on
  aurait une contradiction lorsque ces deux joueurs jouent leurs
  stratégies respectives.
\end{corrige}

\smallbreak

(4) Indépendamment de tout ce qui précède, montrer que si $w$ est un
réel et $I_0$ et $I_1$ deux parties de $\mathbb{R}$ vérifiant (i) si
$v \in I_0$ et $v' < v$ alors $v' \in I_0$ et de même pour $I_1$, et
(ii) tout $v < w$ appartient à $I_0 \cup I_1$, alors en fait soit tout
$v<w$ appartient à $I_0$ soit tout $v<w$ appartient à $I_1$.  (On
pourra considérer à quel $I_a$ appartient $w - \frac{1}{n}$ et faire
varier $n$.)

\begin{corrige}
Pour chaque entier $n>0$, on a $w - \frac{1}{n} \in I_0 \cup I_1$
d'après (ii), disons donc $w - \frac{1}{n} \in I_{a(n)}$ avec $a(n)
\in \{0,1\}$.  Soit il existe une infinité de valeurs de $n$ telles
que $a(n)$ vaille $0$, soit $a(n)$ vaut constamment $1$ à partir d'un
certain point — et en particulier, il existe une infinité de valeurs
de $n$ telles que $a(n)$ vaille $1$.  Bref, dans tous les cas, il
existe $b \in \{0,1\}$ tel que $a(n)$ prenne une infinité de fois la
valeur $b$, c'est-à-dire $w - \frac{1}{n} \in I_b$ une infinité
de $n$, autrement dit, pour des $n$ arbitrairement grands.  Mais si
$v<w$, on peut trouver $n$ plus grand que $1/(w-v)$ tel que $w -
\frac{1}{n} \in I_b$ d'après ce qu'on vient de dire, or $v <
w-\frac{1}{n}$, et la propriété (i) montre alors que $v \in I_b$ : on
a bien montré que tout $v<w$ appartient à $I_b$.
\end{corrige}

\smallbreak

(5) On suppose désormais que $X = \{0,1\}$ (et toujours que $u$ est
continue).  Le but de cette question est de montrer (pour $w \in
\mathbb{R}$) que si Alice possède pour tout $v < w$ une stratégie lui
garantissant un gain $\geq v$, alors en fait Alice possède une
stratégie lui garantissant un gain $\geq w$ : autrement dit, si Alice
peut se garantir n'importe quel gain $v < w$ alors elle peut se
garantir le gain $w$ lui-même.  Pour abréger, on dira qu'une position
est « $v$-gagnante [pour Alice] » si Alice possède une stratégie lui
garantissant un gain $\geq v$ à partir de cette position (on notera
qu'une position $v$-gagnante est évidemment $v'$-gagnante pour
tout $v' < v$) ; et on dira qu'une position est « presque $w$-gagnante
[pour Alice] » si elle est $v$-gagnante pour tout $v < w$.

On cherche donc à montrer, en s'inspirant de la démonstration du
théorème de Gale-Stewart, que si la position initiale (ou, en fait,
une position quelconque) est presque $w$-gagnante, alors elle est
$w$-gagnante.\spaceout (a) Montrer que si une position
$(z_0,\ldots,z_{i-1})$ où c'est à Alice de jouer est presque
$w$-gagnante, alors \emph{il existe} $x \in X$ (un coup d'Alice) tel
que $(z_0,\ldots,z_{i-1},x)$ soit presque $w$-gagnante (on pourra
utiliser (4)).\spaceout (b) Montrer que si une position
$(z_0,\ldots,z_{i-1})$ où c'est à Bob de jouer est presque
$w$-gagnante, alors \emph{pour tout} $x \in X$ (coup de Bob), la
position $(z_0,\ldots,z_{i-1},x)$ est presque $w$-gagnante.\spaceout
(c) Montrer que si $\dblunderline{x} \in X^{\mathbb{N}}$ et
$u(\dblunderline{x}) < w$ alors il existe $\ell$ tel que
$(x_0,\ldots,x_{\ell-1})$ ne soit pas presque $w$-gagnante.\spaceout
(d) En déduire (en construisant une stratégie d'Alice qui cherche à
rester dans des positions presque $w$-gagnantes) que si la position
initiale est presque $w$-gagnante, alors elle est $w$-gagnante.

\begin{corrige}
(a) Si $\underline{z} := (z_0,\ldots,z_{i-1})$ est une position où
  c'est à Alice de jouer qui est presque $w$-gagnante, elle est
  $v$-gagnante pour tout $v < w$, autrement dit, pour chaque $v < w$,
  Alice possède une stratégie qui lui garantit un gain $\geq v$.
  D'après \ref{strategies-forall-exists-reformulation}, cela signifie
  que pour chaque $v < w$, il existe un coup $x \in X$ d'Alice tel que
  la position $\underline{z}x = (z_0,\ldots,z_{i-1},x)$ soit
  $v$-gagnante.  Pour $x \in \{0,1\}$, soit $I_x$ l'ensemble des $v$
  tels que la position $\underline{z}x = (z_0,\ldots,z_{i-1},x)$ soit
  $v$-gagnante : la propriété (i) de la question (4) a été observée
  dans la question, et la propriété (ii) est exactement la phrase
  précédente.  On déduit de la question (4) qu'il existe $x \in
  \{0,1\}$ tel que la position $\underline{z}x =
  (z_0,\ldots,z_{i-1},x)$ soit $v$-gagnante pour tout $v<w$, autrement
  dit, soit presque $w$-gagnante, ce qu'on voulait montrer.

(b) Si $\underline{z} := (z_0,\ldots,z_{i-1})$ est une position où
  c'est à Bob de jouer qui est presque $w$-gagnante, elle est
  $v$-gagnante pour tout $v < w$, autrement dit, pour chaque $v < w$,
  Alice possède une stratégie qui lui garantit un gain $\geq v$.
  D'après \ref{strategies-forall-exists-reformulation}, cela signifie
  que pour chaque $v < w$, et pour chaque coup $x \in X$ de Bob, la
  position $\underline{z}x = (z_0,\ldots,z_{i-1},x)$ est $v$-gagnante.
  En permutant les quantificateurs : pour chaque coup $x \in X$ de Bob
  et pour chaque $v < w$, la position $\underline{z}x$ est
  $v$-gagnante.  C'est bien dire que pour chaque coup $x\in X$ de Bob,
  la position $\underline{z}x$ est presque $w$-gagnante, ce qu'on
  voulait montrer.

(c) Si $u(\dblunderline{x}) < w$, et si $v$ est choisi tel que
  $u(\dblunderline{x}) < v < w$, alors par continuité de $u$
  (cf. question (2)), il existe $\ell$ tel que $u$ soit $<v$ sur
  $V_\ell(\dblunderline{x})$, autrement dit, aucune confrontation qui
  prolonge $x_0,\ldots,x_{\ell-1}$ ne peut donner un gain $\geq v$ à
  Alice, et en particulier, la position $x_0,\ldots,x_{\ell-1}$ n'est
  pas $v$-gagnante, donc elle n'est pas presque $w$-gagnante.

(d) Si $(x_0,\ldots,x_{i-1})$ est une position où c'est à Alice de
  jouer et qui est presque $w$-gagnante, alors d'après (a) il existe
  un $x$ tel que $(x_0,\ldots,x_{i-1},x)$ soit presque $w$-gagnante :
  choisissons-un tel $x$ et posons $\tau((x_0,\ldots,x_{i-1})) := x$ :
  d'après (b), quel que soit $y \in X$, la position
  $(x_0,\ldots,x_{i-1},x,y)$ est presque $w$-gagnante (et c'est de
  nouveau à elle de jouer).  Aux points où $\tau$ n'a pas été défini
  par ce qui vient d'être dit, on le définit de façon arbitraire.

Si $\dblunderline{x} = (x_0,x_1,x_2,\ldots)$ est une confrontation où
Alice joue selon $\tau$, on voit par récurrence sur $i$ que chacune
des positions $(x_0,\ldots,x_{i-1})$ est presque $w$-gagnante si la
position initiale l'est : pour $i=0$ c'est l'hypothèse faite (à savoir
que la position initiale est presque $w$-gagnante), pour les positions
où c'est à Alice de jouer, c'est la construction de $\tau$ qui assure
la récurrence (cf. (a) ci-dessus), et pour les positions où c'est à
Bob de jouer, c'est la question (b) qui assure la récurrence.  Mais la
question (c) assure qu'on a alors $u(\dblunderline{x}) \geq w$.  On a
donc bien défini une stratégie qui garantit à Alice un gain $\geq w$.
\end{corrige}

\smallbreak

(6) (On suppose toujours que $X = \{0,1\}$ que $u$ est
continue.)\spaceout (a) Montrer que si $u$ est bornée, alors la valeur
du jeu existe, autrement dit, il existe $v$ tel qu'Alice possède une
stratégie lui garantissant un gain $\geq v$ et que Bob possède une
stratégie lui garantissant qu'Alice aura un gain $\leq v$.  (On pourra
considérer la borne supérieure $w$ des $v$ pour lesquels Alice possède
une stratégie lui garantissant un gain $\geq v$, et appliquer la
question (5).)\spaceout (b) En considérant une fonction comme
$\arctan\circ u$ ou $\tanh\circ u$, montrer qu'on peut s'affranchir de
l'hypothèse « $u$ est bornée ».\spaceout (bonus) Déduire de ce qui
précède que toute fonction continue $u\colon \{0,1\}^{\mathbb{N}} \to
\mathbb{R}$ est bornée et atteint ses bornes (on pourra considérer un
jeu où les coups de Bob sont purement et simplement ignorés).

\begin{corrige}
(a) Soit $I$ l'ensemble des $v$ pour lesquels Alice possède une
  stratégie lui garantissant un gain $\geq v$ (i.e., pour lesquels la
  position initiale est $v$-gagnante, dans la terminologie de (4)).
  Si $v$ est strictement supérieur à un majorant de $u$ (qui existe
  car $u$ est supposée bornée), il est évident qu'il est impossible
  d'obtenir un gain $\geq v$, donc $v$ majore strictement $I$.  Si $v$
  est strictement inférieur à un minorant de $u$, il est évident qu'il
  est impossible de \emph{ne pas} obtenir un gain $\geq v$, donc
  $v\in I$.  Soit $w := \sup I$, qui existe puisque $I$ est non vide
  et majoré.  Si $v < w$, alors il existe $v' \in I$ tel que $v<v'<w$,
  ce qui implique que $v \in I$ (puisqu'une stratégie garantissant un
  gain $\geq v'$ garantit \textit{a fortiori} un gain $\geq v$).  On
  voit donc qu'Alice a une stratégie lui garantissant un gain $\geq v$
  pour tout $v < w$.  La question (5) montre alors qu'Alice a une
  stratégie lui garantissant un gain $\geq w$.  \textit{A contrario},
  si $v > w$, alors Alice n'a pas de stratégie lui garantissant un
  gain $\geq v$, donc Bob a une stratégie lui garantissant qu'Alice
  aura un gain $\leq v$ (d'après la question (3)(a)).  En appliquant
  l'analogue pour Bob de la question (5) (qui s'en déduit en
  échangeant les joueurs, c'est-à-dire en changeant $u$ en $-u$), Bob
  a lui aussi une stratégie lui garantissant qu'Alice aura un
  gain $\leq w$.  On a bien prouvé que $w$ est la valeur du jeu.

(b) Si $u$ n'est plus supposée bornée, soit $\tilde u = h\circ u$ où
  $h\colon\mathbb{R}\to\mathbb{R}$ est une fonction continue bornée et
  strictement croissante (par exemple $\arctan$ ou $\tanh$).  Alors
  $\tilde u$ est continue et bornée, donc la question précédente
  montre que le jeu qu'elle définit a une valeur $\tilde w$.  Cette
  valeur ne peut pas être strictement supérieure à toute valeur de $h$
  (vu qu'Alice a une stratégie lui garantissant un gain $\geq\tilde
  w$) ni strictement inférieure à toute valeur de $h$ (vu que Bob a
  une stratégie lui garantissant qu'Alice aura un gain $\leq\tilde
  w$).  C'est donc (en utilisant les valeurs intermédiaires) que
  $\tilde w = h(w)$ pour un certain $w$.  Alors il est clair qu'une
  stratégie garantissant à Alice un gain $\geq\tilde w$ dans le jeu
  défini par $\tilde u$ lui garantit un gain $\geq w$ dans le jeu
  défini par $u$, et de même pour Bob : ceci montre que $w$ est la
  valeur du jeu défini par $u$.

(bonus) Si $u$ est une fonction continue $\{0,1\}^{\mathbb{N}} \to
  \mathbb{R}$, en considérant la fonction $\hat u\colon
  \{0,1\}^{\mathbb{N}} \to \mathbb{R}$ qui ignore les coups joués par
  Bob et calcule $u$ sur les coups joués par Alice (i.e.,
  formellement, $\hat u(\dblunderline{x}) = u(\check{\dblunderline{x})}$ où
  $\check x_i = x_{2i}$), il est clair qu'Alice a une stratégie lui
  garantissant un gain $\geq v$ si et seulement si $u$ prend une
  valeur $\geq v$.  Comme on vient de voir qu'il existe un plus grand
  tel $v$ (la valeur du jeu), cela signifie bien que $u$ est bornée et
  atteint ses bornes.
\end{corrige}

\smallbreak

(7) Soit $u \colon \{0,1\}^{\mathbb{N}} \to \mathbb{R}$ qui à
$(x_0,x_1,x_2,\ldots)$ associe $\sum_{i=0} x_i 2^{-i-1}$ (le nombre réel
dont la représentation binaire est donnée par $0$ virgule la suite
des $x_i$).  Vérifier que $u$ est continue et calculer la valeur du
jeu qu'elle définit (quelle est la stratégie optimale pour Alice et
pour Bob ?).

\begin{corrige}
La fonction $u$ est continue car si $\varepsilon < 2^{-\ell}$
  alors la valeur $u(\dblunderline{x})$ est définie à $\varepsilon$ près
  par la donnée des $\ell$ premiers termes de la
  suite $\dblunderline{x}$.  Il est évident qu'Alice a intérêt à ne jouer
  que des $1$ (jouer autre chose ne ferait que diminuer son gain) et
  Bob que des $0$.  La valeur du jeu est donc $u(0,1,0,1,0,1,\ldots) =
  \frac{1}{3}$.
\end{corrige}


\setbox0=\vbox\bgroup
\subsection{Théorie de l'induction bien-fondée}
\egroup
\bigbreak


\subsection{Introduction aux ordinaux}

\exercice

Ranger les ordinaux suivants par ordre croissant :
\spaceout $\omega^{\omega+1} + \omega^\omega\cdot 33$ ;
\spaceout $\omega\cdot 3 + 42$ ;
\spaceout $\omega^{\omega+1} + \omega + 33$ ;
\spaceout $\omega^{\omega+2} + \omega^\omega$ ;
\spaceout $\omega^2\cdot 42 + 1000$ ;
\spaceout $\omega^2 + \omega$ ;
\spaceout $\omega^2\cdot 42 + \omega$ ;
\spaceout $\omega^{\omega^2 + 1}$ ;
\spaceout $\omega^{\omega^{(\omega\cdot 2)}}$ ;
\spaceout $\omega^{\omega^\omega} + 1$ ;
\spaceout $\omega^{\omega+1} + \omega\cdot 33$ ;
\spaceout $\omega^{\omega^2}$ ;
\spaceout $\omega^{\omega^2 + 1} + \omega^{\omega\cdot 2}\cdot 1000$ ;
\spaceout $\omega^{\omega^2 + \omega}$ ;
\spaceout $\omega\cdot 3$ ;
\spaceout $\omega^{(\omega^\omega\cdot 2)}$ ;
\spaceout $\omega^{\omega^3}$ ;
\spaceout $\omega^{\omega+1} + 1000$ ;
\spaceout $\omega^{\omega+2}$ ;
\spaceout $\omega^{\omega+1}\cdot 2$ ;
\spaceout $\omega\cdot 2 + 1729$ ;
\spaceout $\omega^2 + 1000$ ;
\spaceout $42$ ;
\spaceout $\omega^{\omega^2 + \omega} + \omega^{\omega^2 + 1}$ ;
\spaceout $\omega^{\omega\cdot 2}\cdot 1000$ ;
\spaceout $\omega^2\cdot 42$ ;
\spaceout $\omega^{\omega+1} + \omega^2\cdot 33$ ;
\spaceout $\omega^2$ ;
\spaceout $\omega$ ;
\spaceout $\omega^{\omega+1}$ ;
\spaceout $\omega^{\omega^2 + 1} + \omega^{\omega^2}\cdot 42$ ;
\spaceout $\omega^{\omega\cdot 2} + \omega^{\omega+2}$ ;
\spaceout $\omega^{\omega^2 + \omega} + \omega^{\omega^2} + \omega^{\omega+1}$ ;
\spaceout $\omega^{\omega^{(\omega^2)}}$ ;
\spaceout $\omega^{\omega^2 + 1}\cdot 2$ ;
\spaceout $\omega^2 + \omega\cdot 42$ ;
\spaceout $\omega + 42$ ;
\spaceout $\omega^{\omega^2\cdot 2}$ ;
\spaceout $\omega^{\omega\cdot 2 + 42}$ ;
\spaceout $\omega^{\omega\cdot 2}$ ;
\spaceout $\omega\cdot 2$ ;
\spaceout $\omega^{\omega+1} + \omega^2 + 33$ ;
\spaceout $\omega^{\omega^{(\omega+1)}}$ ;
\spaceout $\omega^\omega$ ;
\spaceout $\omega^{\omega^2 + \omega + 1}$ ;
\spaceout $\omega^{\omega^\omega}\cdot 2$ ;
\spaceout $\omega^{\omega^\omega}$ ;
\spaceout $0$ ;
\spaceout $\omega^{\omega^2} + \omega^{\omega+1}$ ;
\spaceout $\omega^{(\omega^\omega + 1)}$.

\begin{corrige}
On vérifie que tous ces ordinaux sont écrits en forme normale de
Cantor (et les exposants de $\omega$ aussi, etc.).  On les compare
donc en comparant à chaque fois la plus grande puissance de $\omega$.

Dans l'ordre croissant : \spaceout $0$ ;
\spaceout $42$ ;
\spaceout $\omega$ ;
\spaceout $\omega + 42$ ;
\spaceout $\omega\cdot 2$ ;
\spaceout $\omega\cdot 2 + 1729$ ;
\spaceout $\omega\cdot 3$ ;
\spaceout $\omega\cdot 3 + 42$ ;
\spaceout $\omega^2$ ;
\spaceout $\omega^2 + 1000$ ;
\spaceout $\omega^2 + \omega$ ;
\spaceout $\omega^2 + \omega\cdot 42$ ;
\spaceout $\omega^2\cdot 42$ ;
\spaceout $\omega^2\cdot 42 + 1000$ ;
\spaceout $\omega^2\cdot 42 + \omega$ ;
\spaceout $\omega^\omega$ ;
\spaceout $\omega^{\omega+1}$ ;
\spaceout $\omega^{\omega+1} + 1000$ ;
\spaceout $\omega^{\omega+1} + \omega + 33$ ;
\spaceout $\omega^{\omega+1} + \omega\cdot 33$ ;
\spaceout $\omega^{\omega+1} + \omega^2 + 33$ ;
\spaceout $\omega^{\omega+1} + \omega^2\cdot 33$ ;
\spaceout $\omega^{\omega+1} + \omega^\omega\cdot 33$ ;
\spaceout $\omega^{\omega+1}\cdot 2$ ;
\spaceout $\omega^{\omega+2}$ ;
\spaceout $\omega^{\omega+2} + \omega^\omega$ ;
\spaceout $\omega^{\omega\cdot 2}$ ;
\spaceout $\omega^{\omega\cdot 2} + \omega^{\omega+2}$ ;
\spaceout $\omega^{\omega\cdot 2}\cdot 1000$ ;
\spaceout $\omega^{\omega\cdot 2 + 42}$ ;
\spaceout $\omega^{\omega^2}$ ;
\spaceout $\omega^{\omega^2} + \omega^{\omega+1}$ ;
\spaceout $\omega^{\omega^2 + 1}$ ;
\spaceout $\omega^{\omega^2 + 1} + \omega^{\omega\cdot 2}\cdot 1000$ ;
\spaceout $\omega^{\omega^2 + 1} + \omega^{\omega^2}\cdot 42$ ;
\spaceout $\omega^{\omega^2 + 1}\cdot 2$ ;
\spaceout $\omega^{\omega^2 + \omega}$ ;
\spaceout $\omega^{\omega^2 + \omega} + \omega^{\omega^2} + \omega^{\omega+1}$ ;
\spaceout $\omega^{\omega^2 + \omega} + \omega^{\omega^2 + 1}$ ;
\spaceout $\omega^{\omega^2 + \omega + 1}$ ;
\spaceout $\omega^{\omega^2\cdot 2}$ ;
\spaceout $\omega^{\omega^3}$ ;
\spaceout $\omega^{\omega^\omega}$ ;
\spaceout $\omega^{\omega^\omega} + 1$ ;
\spaceout $\omega^{\omega^\omega}\cdot 2$ ;
\spaceout $\omega^{(\omega^\omega + 1)}$ ;
\spaceout $\omega^{(\omega^\omega\cdot 2)}$ ;
\spaceout $\omega^{\omega^{(\omega+1)}}$ ;
\spaceout $\omega^{\omega^{(\omega\cdot 2)}}$ ;
\spaceout et enfin\spaceout $\omega^{\omega^{(\omega^2)}}$.
\end{corrige}



%
%
%

\exercice

(a) Que vaut $(\omega+1) + (\omega+1)$ ?

(b) Plus généralement, que vaut $(\omega+1) + \cdots + (\omega+1)$
avec $n$ termes $\omega+1$ (où $n$ est un entier naturel $\geq 1$) ?

(c) En déduire ce que vaut $(\omega+1)\cdot n$.

(d) En déduire ce que vaut $(\omega+1)\cdot \omega$.

(e) En déduire ce que vaut $(\omega+1)\cdot(\omega+1)$.

(f) En déduire ce que vaut $(\omega+1)^2$.

\begin{corrige}
(a) On a $(\omega+1) + (\omega+1) = \omega + 1 + \omega + 1 = \omega +
  (1 + \omega) + 1 = \omega + \omega + 1 = \omega\cdot 2 + 1$.

(b) En procédant de même, on voit que dans la somme de $n$ termes
  $\omega + 1$, chaque $1$ est absorbé par le $\omega$ qui
  \emph{suit}, sauf le dernier $1$ qui demeure : la somme vaut
  donc $\omega\cdot n + 1$.

(c) Quel que soit l'ordinal $\alpha$, la somme $\alpha + \cdots +
  \alpha$ avec $n$ termes $\alpha$ vaut $\alpha\cdot n$ (ceci se voit
  soit par une récurrence immédiate sur $n$ avec la définition par
  induction de la multiplication, soit en utilisant la distributivité
  à droite, c'est-à-dire $\alpha\cdot n = \alpha\cdot(1 + \cdots + 1)
  = \alpha + \cdots + \alpha$).  On a donc $(\omega+1)\cdot n =
  \omega\cdot n + 1$.

(d) L'ordinal $(\omega+1)\cdot \omega$ est donc la limite
  (c'est-à-dire la borne supérieure) des $(\omega+1)\cdot n =
  \omega\cdot n + 1$ pour $n\to\omega$.  Cette borne supérieure
  vaut $\omega^2$ : en effet, $\omega^2 \geq \omega\cdot n + 1$ pour
  chaque $n<\omega$, mais inversement, si $\gamma < \omega^2$, on a
  $\gamma < \omega\cdot n$ pour un certain $n$ (par exemple en
  utilisant le fait que $\omega^2 = \omega\cdot\omega$ est elle-même
  la limite des $\omega\cdot n$, c'est-à-dire le plus petit ordinal
  supérieur ou égal à eux), et en particulier $\gamma < \omega\cdot n
  + 1$ ; ou, si on préfère, quel que soit $n$ on a $\omega\cdot n \leq
  \omega\cdot n + 1 \leq \omega\cdot (n + 1)$ où $\omega\cdot n$ et
  $\omega\cdot (n+1)$ ont la même limite $\omega^2$ quand
  $n\to\omega$, d'où il résulte que $\omega\cdot n + 1$ aussi.

(e) On a $(\omega+1)\cdot (\omega+1) = (\omega+1)\cdot \omega +
  (\omega + 1) = \omega^2 + \omega + 1$.

(f) On a toujours $\alpha^2 = \alpha\cdot\alpha$, donc $(\omega+1)^2 =
  \omega^2 + \omega + 1$ comme on vient de le montrer.
\end{corrige}



%
%
%

\exercice

(a) Que vaut $(\omega 2) \cdot (\omega 2)$ ?

(b) Plus généralement, que vaut $(\omega 2) \cdots (\omega 2)$ avec
$n$ facteurs $\omega 2$ (où $n$ est un entier naturel $\geq 1$) ?

(c) En déduire ce que vaut $(\omega 2)^n$.

(d) En déduire ce que vaut $(\omega 2)^\omega$.  Comparer avec
$\omega^\omega \cdot 2^\omega$.

(e) En déduire ce que vaut $(\omega 2)^{\omega+n}$ pour $n\geq 1$
entier naturel.

(f) En déduire ce que vaut $(\omega 2)^{\omega 2}$.

\begin{corrige}
(a) On a $(\omega 2) \cdot (\omega 2) = \omega \cdot 2 \cdot \omega
  \cdot 2 = \omega \cdot (2 \cdot \omega) \cdot 2 = \omega \cdot
  \omega \cdot 2 = \omega^2 \cdot 2$.

(b) En procédant de même, on voit que dans le produit de $n$ facteurs
  $\omega 2$, chaque $2$ est absorbé par le $\omega$ qui \emph{suit},
  sauf le dernier $2$ qui demeure : le produit vaut donc $\omega^n
  \cdot 2$.

(c) Quel que soit l'ordinal $\alpha$, le produit $\alpha \cdots
  \alpha$ avec $n$ facteurs $\alpha$ vaut $\alpha^n$ (ceci se voit
  soit par une récurrence immédiate sur $n$ avec la définition par
  induction de l'exponentiation, soit en écrivant $\alpha^n =
  \alpha^{1+\cdots+1} = \alpha \cdots \alpha$).  On a donc $(\omega
  2)^n = \omega^n \cdot 2$.

(d) L'ordinal $(\omega 2)^\omega$ est la limite (c'est-à-dire la borne
  supérieure) des $\omega^n \cdot 2$ pour $n\to\omega$.  Cette borne
  supérieure vaut $\omega^\omega$ : en effet, $\omega^\omega \geq
  \omega^n \cdot 2$ pour chaque $n<\omega$, mais inversement, si
  $\gamma < \omega^\omega$, on a $\gamma < \omega^n$ pour un
  certain $n$ (par exemple en utilisant le fait que $\omega^\omega$
  est lui-même la limite des $\omega^n$, c'est-à-dire le plus petit
  ordinal supérieur ou égal à eux), et en particulier $\gamma <
  \omega^n \cdot 2$ ; ou, si on préfère, quel que soit $n$ on a
  $\omega^n \leq \omega^n \cdot 2 \leq \omega^{n+1}$ où $\omega^n$ et
  $\omega^{n+1}$ ont la même limite $\omega^\omega$ quand
  $n\to\omega$, d'où il résulte que $\omega^n \cdot 2$ aussi.

  Bref, $(\omega 2)^\omega = \omega^\omega$.  En revanche,
  $\omega^\omega \cdot 2^\omega = \omega^\omega \cdot \omega =
  \omega^{\omega+1}$ est strictement plus grand.

(e) On a $(\omega 2)^{\omega + n} = (\omega 2)^\omega \cdot (\omega
  2)^n = \omega^\omega \cdot \omega^n \cdot 2$ d'après les questions
  précédentes, donc ceci vaut $\omega^{\omega+n} \cdot 2$.

(f) L'ordinal $(\omega 2)^{\omega 2}$ est la limite des
  $\omega^{\omega+n} \cdot 2$ pour $n\to\omega$, et le même
  raisonnement qu'en (d) montre que cette limite est
  $\omega^{\omega+\omega} = \omega^{\omega 2}$.  Bref, $(\omega
  2)^{\omega 2} = \omega^{\omega 2}$.
\end{corrige}



%
%
%

\exercice

On dit qu'un ordinal $\alpha$ est \textbf{infini} lorsque
$\alpha\geq\omega$.  Montrer qu'un ordinal est infini si et seulement
si $1+\alpha = \alpha$.

\begin{corrige}
Si $\alpha$ est infini, on a $\alpha \geq \omega$, donc il existe un
unique ordinal $\beta$ tel que $\alpha = \omega + \beta$.  On a alors
$1 + \alpha = 1 + (\omega + \beta) = (1 + \omega) + \beta = \omega +
\beta = \alpha$.

Si, en revanche, $\alpha$ est fini, c'est-à-dire $\alpha < \omega$,
alors $\alpha$ est un entier naturel, et comme l'addition ordinale sur
les entiers naturels coïncide avec l'addition usuelle sur ceux-ci, on
a $1 + \alpha > \alpha$.
\end{corrige}



%
%
%

\exercice

On rappelle que si $\alpha' \geq \alpha$ sont deux ordinaux, il existe
un unique $\beta$ tel que $\alpha' = \alpha + \beta$.\spaceout (a) En
déduire que si $\gamma < \gamma'$ alors $\omega^\gamma +
\omega^{\gamma'} = \omega^{\gamma'}$ (on pourra utiliser la conclusion
de l'exercice précédent).\spaceout (b) Expliquer pourquoi
$\omega^{\gamma'} + \omega^\gamma$, lui, est strictement plus grand
que $\omega^{\gamma'}$ et $\omega^\gamma$.

\begin{corrige}
(a) Si $\gamma < \gamma'$, il existe $\beta$ tel que $\gamma' = \gamma
  + \beta$, si bien qu'on a $\omega^\gamma + \omega^{\gamma'} =
  \omega^\gamma + \omega^{\gamma + \beta} = \omega^\gamma +
  \omega^\gamma \cdot \omega^\beta = \omega^\gamma (1 +
  \omega^\beta)$.  La conclusion voulue découle donc du fait que $1 +
  \omega^\beta = \omega^\beta$ : or ceci résulte de l'exercice
  précédent (on a $\beta \neq 0$ puisque $\gamma' \neq \gamma$, donc
  $\beta \geq 1$, donc $\omega^\beta \geq \omega$).

(b) On a $\omega^\gamma > 0$ donc $\omega^{\gamma'} + \omega^\gamma >
  \omega^{\gamma'}$ (par stricte croissance de la somme en la variable
  de droite), et comme $\omega^{\gamma'} > \omega^\gamma$, la somme
  est également $> \omega^\gamma$.  (On pouvait aussi invoquer la
  comparaison des formes normales de Cantor.)
\end{corrige}



%
%
%

\exercice

(A) (1) Que vaut $2^{\omega+1}$ ?\spaceout (2) Que vaut
$2^{\omega^2}$ ?\spaceout (3) Expliquer pourquoi $\omega^\omega =
\omega\cdot \omega^\omega$.  En déduire ce que vaut
$2^{\omega^\omega}$.\spaceout (À chaque fois, on écrira les ordinaux
demandés en forme normale de Cantor.)

(B) On suppose que $\varepsilon = \omega^\varepsilon$.\spaceout
(1) Que vaut $\varepsilon^\varepsilon$ ?\spaceout (2) Que vaut
$\varepsilon^{\varepsilon^\varepsilon}$ (on pourra utiliser un des
deux exercices précédents) ?\spaceout (À chaque fois, plusieurs
écritures sont possibles.)

\begin{corrige}
(A) (1) On a $2^{\omega+1} = 2^\omega\cdot 2^1 = \omega\cdot
  2$.\spaceout (2) On a $2^{\omega^2} = 2^{\omega\cdot\omega} =
  (2^\omega)^\omega = \omega^\omega$.\spaceout (3) On a $\omega\cdot
  \omega^\omega = \omega^1 \cdot \omega^\omega = \omega^{1+\omega} =
  \omega^\omega$.  On en déduit que $2^{\omega^\omega} =
  2^{\omega\cdot \omega^\omega} = (2^\omega)^{\omega^\omega} =
  \omega^{\omega^\omega}$.

(B) (1) On a $\varepsilon^\varepsilon =
  (\omega^\varepsilon)^\varepsilon = \omega^{\varepsilon^2}$ ou, si on
  préfère, $\omega^{\omega^{\varepsilon\cdot 2}}$.\spaceout  (2) On a
  $\varepsilon^{\varepsilon^\varepsilon} =
  (\omega^\varepsilon)^{\varepsilon^\varepsilon} = \omega^{\varepsilon
    \cdot \varepsilon^\varepsilon} = \omega^{\varepsilon^{1 +
      \varepsilon}}$.  Or $1 + \varepsilon = \varepsilon$ d'après un
  des exercices précédents (parce que $\varepsilon$ est infini ou
  parce que la somme est $\omega^0 + \omega^\varepsilon$), donc
  $\varepsilon^{\varepsilon^\varepsilon} =
  \omega^{\varepsilon^\varepsilon}$.  D'après la sous-question
  précédente, c'est aussi $\omega^{\omega^{\varepsilon^2}}$ ou encore
  $\omega^{\omega^{\omega^{\varepsilon\cdot 2}}}$.
\end{corrige}

%
%
%

\subsection{Jeux combinatoires à information parfaite}

\exercice

Soit $k$ un entier naturel fixé.  On considère le jeu suivant : une
position du jeu consiste en un certain nombre (fini) de jetons placés
sur des cases étiquetées par les ordinaux.  Par exemple, il pourrait y
avoir trois jetons sur la case $42$ et un sur la case $\omega$ ; ou
bien douze jetons sur la case $0$.  Un coup du jeu consiste à retirer
un jeton d'une case, disons $\alpha$, et le remplacer par exactement
$k$ jetons situées sur des cases $<\alpha$ quelconques (y compris
plusieurs fois la même).  Par exemple, s'il y a un jeton sur la
case $3$ et si $k=7$, on peut le remplacer par quatre jetons sur la
case $2$ et trois sur la case $0$.  (Le nombre de jetons présents dans
le jeu augmente donc de $k-1$ à chaque coup joué.)  On remarquera que
les jetons sur la case $0$ ne peuvent plus être retirés ou servir de
quelque manière que ce soit (en tout cas si $k\geq 1$) : on pourra
dire que la case $0$ est la « défausse » des jetons.  Le jeu se
termine lorsque tous les jetons sont sur la case $0$ (=dans la
défausse), car il n'est alors plus possible de jouer.  Les joueurs
(Alice et Bob) jouent à tour de rôle et celui qui ne peut plus jouer a
perdu.

(1) Montrer que le jeu termine toujours en temps fini.  (On pourra par
exemple coder la position sous la forme d'un ordinal écrit en forme
normale de Cantor et montrer qu'il décroît strictement.)

\begin{corrige}
À une position du jeu ayant $n_i$ jetons sur la case $\alpha_i$ on
peut associer l'ordinal $\omega^{\alpha_1} n_1 + \cdots +
\omega^{\alpha_s} n_s$ où les $\alpha_i$ ont été triés de façon à
avoir $\alpha_1 > \cdots > \alpha_s$.  Un coup consistant à remplacer
$\omega^{\alpha_i} n_i$ par $\omega^{\alpha_i} (n_i-1)$ en même temps
qu'on ajoute un nombre fini ($k$) de termes strictement inférieurs à
$\omega^{\alpha_i}$ : ceci fait donc décroître strictement l'ordinal
en question, et en particulier, la partie doit terminer en temps fini.
\end{corrige}

\smallbreak

(2) Dans le cas particulier où $k=1$, expliquer pourquoi le jeu est
simplement une reformulation du jeu de nim.

\begin{corrige}
Lorsque $k=1$, un coup consiste simplement à déplacer un jeton vers
une case d'ordinal strictement plus petit ; on peut identifier la
position ayant $n_i$ jetons sur la case $\alpha_i$ à un partie de nim
ayant $n_i$ lignes avec $\alpha_i$ allumettes : le coup consistant à
déplacer un jeton de la case $\alpha_i$ vers la case $\alpha_i' <
\alpha_i$ peut se voir comme un coup de nim consistant à diminuer le
nombre d'allumettes de la ligne qui en a $\alpha_i$ pour qu'il en
reste $\alpha'_i$.  Les jeux sont donc complètement équivalents.
\end{corrige}

\smallbreak

(3) Dans le cas général, montrer qu'une position du jeu peut se voir
comme somme de nim de positions ayant un seul jeton.  Que peut-on dire
des positions ayant plusieurs jetons sur la même case ?  Expliquer
comment on pourrait modifier les règles, sans changer vraiment le jeu,
pour qu'il n'y ait jamais plusieurs jetons sur la même case.

\begin{corrige}
Il n'y a aucune interaction entre les jetons dans le jeu.  En
particulier, jouer à une somme de nim de deux positions du jeu
considéré dans cet exercice revient au même que de jouer à la position
ayant la réunion de ces deux ensembles de jetons.  Par exemple, la
somme de nim de deux positions, l'une ayant un unique jeton sur la
case $\alpha$ et l'autre ayant un unique jeton sur la case $\alpha'$
revient au même que la position ayant deux jetons, un sur la case
$\alpha$ et l'autre sur la case $\alpha'$ (quitte à se rappeler si un
jeton est un descendant de celui de la case $\alpha$ ou de celui de la
case $\alpha'$, on peut transformer toute position ou partie d'un jeu
en une position ou partie de l'autre), et la même chose vaut avec plus
de deux jetons.  (Si on veut être extrêmement rigoureux, « revenir au
  même » dans ce paragraphe signifie que les jeux en question ont le
même écrasement transitif, cf. \ref{definition-transitive-collapse},
ce qui implique notamment qu'ils ont même valeur de Grundy.)

En particulier, deux jetons sur la même case peuvent être considérés
comme s'annulant (puisque la somme de nim de deux jeux égaux donne un
jeu nul, c'est-à-dire dont la valeur de Grundy est nulle,
cf. \ref{nim-sum-has-characteristic-two}) : concrètement, ajouter deux
jetons sur la même case ne changera jamais rien, dès qu'un joueur joue
sur l'un, son adversaire pourra reproduire le même coup sur l'autre.

On peut donc modifier la règle du jeu sans le changer
substantiellement (concrètement, sans en changer la valeur de Grundy)
en décrétant que si on jeton tombe sur une case déjà occupée par un
autre jeton, les deux s'annulent — si bien qu'au final il n'y a jamais
plus qu'un jeton sur une case donnée.
\end{corrige}

\smallbreak

(4) Montrer que la valeur de Grundy d'un état du jeu est la somme de
nim sur tous les jetons du jeu d'un valeur $f_k(\alpha)$ où $\alpha$
est la case où se trouve le jeton.

\begin{corrige}
Si $x$ est une position quelconque du jeu, $N$ son nombre de jetons et
$\alpha_1,\ldots,\alpha_N$ les cases sur lesquelles se trouvent les
jetons, on vient de voir que $x$ est équivalent à la somme de nim des
positions $u_{\alpha_i}$ où $u_{\alpha}$ désigne la position ayant un
unique jeton sur la case $\alpha$.
D'après \ref{summary-of-grundy-theory}, on en déduit que $\gr(x) =
\bigoplus_{i=1}^N \gr(u_{\alpha_i})$, et en posant $f_k(\alpha) =
\gr(u_\alpha)$, on a montré ce qui était demandé : $\gr(x) =
\bigoplus_{i=1}^N f_k(\alpha_i)$.
\end{corrige}

\smallbreak

(5) Donner une définition inductive directe de la fonction $f_k$ (sans
faire référence à un jeu).  Que vaut $f_k(0)$ (pour $k\geq 1$) ?

\begin{corrige}
Comme $f_k(\alpha)$ est la valeur de Grundy de la position $u_\alpha$
ayant un unique jeton sur la case $\alpha$, la définition de la
fonction de Grundy fait que $f_k(\alpha)$ est le $\mex$ de toutes les
valeurs de Grundy des options (=voisins sortants) de $u_\alpha$,
c'est-à-dire des coups qu'on peut faire à partir de là.  La règle du
jeu est que ce sont les positions ayant $k$ jetons sur cases
numérotées $<\alpha$, et on a vu que la valeur de Grundy d'une telle
position est la somme de nim des $f_k$ correspondants.  On a donc la
définition inductive
\[
f_k(\alpha) = \mex\Big\{\bigoplus_{i=1}^k f_k(\beta_i)
: \beta_1,\ldots,\beta_k < \alpha \Big\}
\]
c'est-à-dire que $f_k(\alpha)$ est le plus petit ordinal qui ne peut
pas s'écrire comme somme de nim de $k$ valeurs de $f_k$ sur des
ordinaux strictement plus petits.

Pour $\alpha=0$, notamment, on a $f_k(0) = 0$ d'après la définition
ci-dessus (il s'agit du $\mex$ de l'ensemble vide) ou simplement en se
rappelant que la position $u_0$ ayant un jeton dans la défausse ne
permet pas de faire de coup.
\end{corrige}

\smallbreak

(6) Pour $k=1$, que vaut $f_1(\alpha)$ ?  Et pour $k=0$, que vaut
$f_0(\alpha)$ ?

\begin{corrige}
Pour $k=1$, la définition inductive trouvée à la question précédente
donne $f_1(\alpha) = \mex\{f_1(\beta) : \beta<\alpha\}$ : une
induction transfinie immédiate montre $f_1(\alpha) = \alpha$ (en
effet, si on suppose que $f_1(\beta) = \beta$ pour tout
$\beta<\alpha$, alors la définition montre que $f_1(\alpha)$ vaut
$\mex\{\beta : \beta<\alpha\} = \alpha$).  Si on préfère, comme on a
vu que le jeu pour $k=1$ est simplement le jeu de nim, on
invoque \ref{grundy-of-nimbers-triviality} pour affirmer $f_1(\alpha)
= \alpha$.

Pour $k=0$, la définition inductive donne $f_0(\alpha) = \mex\{0\}$
(la somme de nim vide vaut $0$), donc $f_0(\alpha) = 1$ pour
tout $\alpha$.  Si on préfère, comme le jeu pour $k=0$ consiste
simplement à ce que chaque joueur tour à tour retire un jeton, il est
évident que seule compte la parité du nombre total $N$ de jetons, et
que la valeur de Grundy de ce jeu est égale à $0$ ou $1$ selon que ce
nombre $N$ total de jetons est pair ou impair.
\end{corrige}

\smallbreak

(7) Pour tout $k\geq 1$, montrer que $f_k$ est une fonction
strictement croissante.

\begin{corrige}
Si $\alpha < \alpha'$ alors $f_k(\alpha')$ est le $\mex$ d'un ensemble
$\{\bigoplus_{i=1}^k f_k(\beta_i) : \beta_1,\ldots,\beta_k <
\alpha'\}$ qui contient celui $\{\bigoplus_{i=1}^k f_k(\beta_i) :
\beta_1,\ldots,\beta_k < \alpha\}$ dont $f_k(\alpha)$ est le $\mex$.
Mais par ailleurs, $f_k(\alpha') \neq f_k(\alpha)$ puisqu'on peut
prendre $\beta_1 = \alpha$ et $\beta_2,\ldots,\beta_k = 0$ (et qu'on a
vu $f_k(0) = 0$).  On a donc $f_k(\alpha') > f_k(\alpha)$.
\end{corrige}

\smallbreak

(8) Calculer les premières valeurs de $f_2(n)$ (pour $n$ entier
naturel).  En considérant le nombre $n-1$ et notamment la parité du
nombre de $1$ dans l'écriture binaire de $n-1$, et/ou la parité du
nombre de $1$ dans l'écriture binaire de $f_2(n)$, formuler une
conjecture quant à la valeur de $f_2(n)$.  Montrer cette conjecture.

\begin{corrige}
On a vu que $f_2(n)$ est défini inductivement (i.e., récursivement)
comme le plus petit entier naturel qui n'est pas de la forme $f_2(n_1)
\oplus f_2(n_2)$ avec $n_1,n_2<n$.  Pour faciliter les calculs, on
peut aussi remarquer que (pour $n>0$) c'est le plus petit entier
naturel \emph{non nul} qui n'est pas sous la forme $f_2(n_1) \oplus
f_2(n_2)$ avec $n_1 < n_2 < n$ (en effet, le cas $n_1 = n_2$ donne de
toute façon zéro, et dans les autres cas, par symétrie on peut prendre
$n_1 < n_2$) ; ou encore, si on préfère séparer le cas $n_1=0$ du
reste : le plus petit entier naturel non nul qui n'est pas ni sous la
forme $f_2(n')$ pour $0<n'<n$ ni sous la forme $f_2(n_1) \oplus
f_2(n_2)$ pour $0<n_1<n_2<n$.  On calcule ainsi de proche en proche :

\begin{center}
\begin{tabular}{r|ccccccccccc}
$n$&$0$&$1$&$2$&$3$&$4$&$5$&$6$&$7$&$8$&$9$&$10$\\\hline
$n-1$&&$0$&$1$&$2$&$3$&$4$&$5$&$6$&$7$&$8$&$9$\\\hline
$f_2(n)$&$0$&$1$&$2$&$4$&$7$&$8$&$11$&$13$&$14$&$16$&$19$\\
\end{tabular}
\end{center}

En considérant la parité du nombre de $1$ dans l'écriture binaire de
$n-1$ et/ou de $f_2(n)$, on constate expérimentalement que (A) les
valeurs $f_2(n)$ pour $n>0$ ont un nombre impair de $1$ dans leur
écriture binaire, et sont même exactement tous ces nombres rangés par
ordre croissant, et (B) $f_2(n)$ vaut $2(n-1)$ ou $2(n-1)+1$ selon que
$n-1$ a un nombre impair ou pair de $1$ dans son écriture binaire.  Il
est facile de voir que ces affirmations sont équivalentes : les
nombres ayant un nombre impair de $1$ dans leur écriture binaire
s'obtiennent, dans l'ordre, en prenant les écritures binaires des
entiers (ici $n-1$) et en ajoutant le bit $0$ ou $1$ à la fin
(c'est-à-dire en calculant $2(n-1)$ ou $2(n-1)+1$) pour que le nombre
total de bits $1$ soit impair, ce qui donne l'équivalence entre
(A) et (B).

(Note : certains appellent « odieux » — un jeu de mot sur « odd » en
anglais — les entiers naturels ayant un nombre impair de $1$ dans leur
écriture binaire : voir par exemple \url{http://oeis.org/A000069}
trouvé en entrant les premières valeurs de $f_2(n)$ dans l'OEIS.)

Montrons les affirmations (A) et (B) conjecturées ci-dessus, en
procédant par récurrence sur $n$, autrement dit, montrons par
récurrence sur $n>0$ que $f_2(n)$ parcourt dans l'ordre tous les
nombres, dits « odieux », ayant un nombre impair de $1$ dans leur
écriture binaire.  On a remarqué ci-dessus que $f_2(n)$ est le plus
petit entier naturel non nul qui n'est pas ni sous la forme $f_2(n')$
pour $0<n'<n$ ni sous la forme $f_2(n_1) \oplus f_2(n_2)$ pour
$0<n_1<n_2<n$.  Pour montrer que $f_2(n)$ est le plus petit nombre
odieux $m$ strictement supérieur aux $f_2(n')$ pour $n'<n$, il suffit
donc de montrer que (a) les $f_2(n_1) \oplus f_2(n_2)$ ne sont jamais
odieux, et (b) ils parcourent au moins tous les nombres non-odieux
inférieurs à $m$.  Autrement dit, il suffit de montrer que : (a) la
somme de nim de deux nombres odieux n'est pas odieuse, et que (b) tout
nombre qui n'est pas odieux est la somme de nim de deux nombres odieux
plus petits que lui.  Or ces deux affirmations sont claires sur la
représentation binaire.  Ceci démontre donc la conjecture.
\end{corrige}


%
%
%

\exercice\label{game-for-nim-product}

On considère le jeu suivant : une position du jeu consiste en un
certain nombre fini de jetons placés sur un damier possiblement
transfini dont les cases étiquetées par un couple $(\alpha,\beta)$
d'ordinaux (on dira que $\alpha$ est [le numéro de] la ligne de la
case et $\beta$ [le numéro de] la colonne).  Plusieurs jetons peuvent
se trouver sur la même case sans effet particulier (ils s'empilent).

Un coup du jeu consiste à faire l'opération suivante : le joueur qui
doit jouer choisit un jeton du jeu, disons sur la case
$(\alpha,\beta)$, et il choisit aussi arbitrairement $\alpha' <
\alpha$ (i.e., une ligne située plus haut) et $\beta' < \beta$ (i.e.,
une colonne située plus à gauche) : il retire alors le jeton choisi de
la case $(\alpha,\beta)$ et le remplace par \emph{trois} jetons, sur
les cases $(\alpha',\beta)$, $(\alpha,\beta')$ et $(\alpha',\beta')$.
(Par exemple, un coup valable consiste à remplacer un jeton sur la
case $(42,7)$ par trois sur les cases $(18,7)$, $(42,5)$ et $(18,5)$.)
Le nombre de jetons augmente donc de $2$ à chaque coup.

\begin{center}
\vskip-\baselineskip
\begin{tikzpicture}
\draw[step=.5cm,help lines] (0,-2) grid (3.5,0);
\begin{scope}[every node/.style={circle,fill,inner sep=1.2mm}]
\node at (1.25,-0.75) {};
\node at (1.25,-1.75) {};
\node at (2.75,-0.75) {};
\node[fill=gray] at (2.75,-1.75) {};
\end{scope}
\node[anchor=east] at (0,-0.75) {$\alpha'$};
\node[anchor=east] at (0,-1.75) {$\alpha$};
\node[anchor=south] at (1.25,-0.1) {$\beta'$};
\node[anchor=south] at (2.75,-0.1) {$\beta$};
\end{tikzpicture}
\\{\footnotesize (Le jeton en gris remplacé par les trois noirs.)}
\end{center}

On remarquera que les jetons sur la ligne ou la colonne $0$ ne peuvent
plus être retirés ou servir de quelque manière que ce soit : on pourra
dire que cette ligne et cette colonne $0$ sont la « défausse » des
jetons.  Le jeu se termine lorsque chacun des jetons est sur la ligne
ou la colonne $0$ (=dans la défausse), car il n'est alors plus
possible de jouer.  Les joueurs (Alice et Bob) jouent à tour de rôle
et le premier qui ne peut plus jouer a perdu.

\smallbreak

(0) Décrire brièvement le jeu complètement équivalent dans lequel il
n'y a pas de ligne ou de colonne $0$ (on fait démarrer la numérotation
à $1$), c'est-à-dire pas de défausse (les jetons disparaissent plutôt
qu'être défaussés) : quels sont les types de coups possibles à ce
jeu ?  (Distinguer selon que $\alpha'=0$ ou non, et selon que
$\beta'=0$ ou non.)  On se permettra dans la suite d'utiliser
librement l'une ou l'autre variante du jeu.

\begin{corrige}
Les jetons défaussés ne jouant aucun rôle dans le jeu, on peut les
ignorer et obtenir un jeu équivalent.  Les lignes et les colonnes sont
alors numérotés par des ordinaux non nuls, c'est-à-dire $\geq 1$.  Les
quatre types de coups possibles dans ce jeu, selon que $\alpha'$ et/ou
$\beta'$ vaut $0$, sont alors :
\begin{itemize}
\item simplement retirer un jeton de la case $(\alpha,\beta)$ [cas où
  $\alpha' = 0$ et $\beta' = 0$],
\item déplacer un jeton de la case $(\alpha,\beta)$ vers une case
  $(\alpha,\beta')$ plus à gauche (i.e., $\beta' < \beta$) dans la
  ligne [cas où $\alpha' = 0$ et $\beta' \neq 0$],
\item déplacer un jeton de la case $(\alpha,\beta)$ vers une case
  $(\alpha',\beta)$ plus haut (i.e., $\alpha' < \alpha$) dans la
  colonne [cas où $\alpha' \neq 0$ et $\beta' = 0$],
\item remplacer un jeton de la case $(\alpha,\beta)$ par un jeton sur
  une case $(\alpha,\beta')$ plus à gauche dans la ligne, un autre sur
  une case $(\alpha',\beta)$ plus haut dans la colonne, et un
  troisième sur la case $(\alpha',\beta')$ à l'intersection de la
  nouvelle ligne et de la nouvelle colonne, comme dans le jeu
  d'origine [cas où $\alpha' \neq 0$ et $\beta' \neq 0$].
\end{itemize}
\vskip-\baselineskip\vskip-\baselineskip
\end{corrige}

\smallbreak

(1) (a) Montrer qu'il existe une fonction $h(\alpha,\beta)$ de deux
ordinaux $\alpha,\beta$ et à valeurs ordinales qui soit strictement
croissante en chaque variable (c'est-à-dire que si $\alpha' < \alpha$
alors $h(\alpha',\beta) < h(\alpha,\beta)$ et que si $\beta' < \beta$
alors $h(\alpha,\beta') < h(\alpha,\beta)$).  Pour cela, on pourra,
comme on préfère, poser $h(\alpha,\beta) = \omega^{\max(\alpha,\beta)}
+ \omega^{\min(\alpha,\beta)}$ ou bien $h(\alpha,\beta) =
\alpha\boxplus\beta$ où $\alpha\boxplus\beta$ désigne l'ordinal dont
les chiffres de la forme normale de Cantor sont la somme des chiffres
correspondants de $\alpha$ et de $\beta$.\spaceout (b) En déduire que
le jeu considéré dans cet exercice termine toujours en temps fini.
(On pourra par exemple considérer la somme des $\omega^\gamma$ où
$\gamma$ parcourt, dans l'ordre décroissant, les valeurs
$h(\alpha,\beta)$ pour les cases $(\alpha,\beta)$ où il y a un jeton.)

\begin{corrige}
(a) Si on pose $h(\alpha,\beta) = \omega^{\max(\alpha,\beta)} +
  \omega^{\min(\alpha,\beta)}$ et si $\alpha' < \alpha$, alors soit
  $\max(\alpha',\beta) < \max(\alpha,\beta)$ soit $\max(\alpha',\beta)
  = \max(\alpha,\beta)$ et alors $\min(\alpha',\beta) <
  \min(\alpha,\beta)$ : dans les deux cas, $h(\alpha',\beta) <
  h(\alpha,\beta)$ par comparaison des formes normales de Cantor.
  Comme la fonction $h$ est symétrique en ses deux arguments, on a
  aussi $h(\alpha,\beta') < h(\alpha,\beta)$ si $\beta' < \beta$.

Si on préfère poser $h(\alpha,\beta) = \alpha\boxplus\beta$, et si
$\alpha' < \alpha$, le chiffre correspondant à la plus haute puissance
de $\omega$ qui diffère est strictement plus petit dans la forme
normale de Cantor de $\alpha'$ que celui de $\alpha$, et cette
affirmation est encore vraie lorsqu'on ajoute chiffre à chiffre la
forme normale de Cantor de $\beta$, donc on a bien $h(\alpha',\beta) <
h(\alpha,\beta)$.  Comme la fonction $h$ est symétrique en ses deux
arguments, on a aussi $h(\alpha,\beta') < h(\alpha,\beta)$ si $\beta'
< \beta$.

{\footnotesize
(\emph{Remarque :} En fait, la fonction $\boxplus$, aussi appelée
« somme naturelle » sur les ordinaux, est plus adaptée dans ce
contexte, parce qu'on peut se convaincre que $\alpha \boxplus \beta =
\sup^+ \big( \{\alpha'\boxplus\beta : \alpha' < \alpha\} \cup\,
\{\alpha\boxplus\beta' : \beta' < \beta\}\big)$, c'est-à-dire qu'elle
est justement la plus petite fonction strictement croissante en chacun
de ses arguments.  On comparera cette définition avec $\alpha \oplus
\beta = \mex \big( \{\alpha'\oplus\beta : \alpha' < \alpha\} \cup\,
\{\alpha\oplus\beta' : \beta' < \beta\}\big)$.  En revanche,
l'addition usuelle ne convient pas, parce que $\alpha'+\beta$ peut
être égal à $\alpha+\beta$ même si $\alpha'<\alpha$, par exemple
$1+\omega = 0+\omega$.)
\par}

(b) À une position du jeu ayant $s$ jetons sur les
cases $(\alpha_i,\beta_i)$ (pour $i=1,\ldots,s$) on peut associer
l'ordinal $\omega^{\gamma_1} + \cdots + \omega^{\gamma_s}$ où les
$\gamma_i := h(\alpha_i,\beta_i)$ ont été triés de façon à avoir
$\gamma_1 \geq \cdots \geq \gamma_s$ (donc, quitte à regrouper les
mêmes puissances, à avoir une forme normale de Cantor).  Un coup
consiste à remplacer un terme $\omega^{\gamma_i}$ par une somme de
$\omega^{\gamma'}$ pour des $\gamma' < \gamma_i$ vu que
$h(\alpha',\beta) < h(\alpha,\beta)$ et $h(\alpha,\beta') <
h(\alpha,\beta)$ et \textit{a fortiori} $h(\alpha',\beta') <
h(\alpha,\beta)$ si $\alpha'<\alpha$ et $\beta'<\beta$.  Ceci fait
donc décroître strictement l'ordinal en question, et en particulier,
la partie doit terminer en temps fini.
\end{corrige}

(2) Dans le cas particulier où il n'y a qu'une ligne de jetons
(numérotée $1$ ; ou bien deux lignes numérotées $0$ et $1$ si on garde
la défausse), expliquer pourquoi le jeu est simplement une
reformulation du jeu de nim.

\begin{corrige}
Un coup joué sur un jeton de la ligne $1$ consiste simplement soit à
le retirer soit à le déplacer vers une colonne plus à gauche ; on peut
identifier la position ayant $n_i$ jetons sur la case $(1,\alpha_i)$ à
un partie de nim ayant $n_i$ lignes avec $\alpha_i$ allumettes : le
coup consistant à déplacer un jeton de la case $\alpha_i$ vers la case
$\alpha_i' < \alpha_i$ peut se voir comme un coup de nim consistant à
diminuer le nombre d'allumettes de la ligne qui en a $\alpha_i$ pour
qu'il en reste $\alpha'_i$ ; et le fait de retirer un jeton sur la
case $(1,\alpha_i)$ comme le coup de nim consistant à retirer toutes
les allumettes de la ligne correspondante.  Les jeux sont donc
complètement équivalents.
\end{corrige}

\smallbreak

(3) Montrer que la valeur de Grundy d'un état quelconque du jeu vaut
\[
\bigoplus_{i=1}^s (\alpha_i\otimes\beta_i)
\]
où $(\alpha_1,\beta_1),\ldots,(\alpha_s,\beta_s)$ sont les cases où se
trouvent les $s$ jetons (répétées en cas de jetons multiples), où
$\oplus$ désigne la somme de nim, et où l'opération $\otimes$ sur les
ordinaux est définie inductivement par
\[
\alpha\otimes\beta := \mex\Big\{(\alpha'\otimes\beta)
\oplus (\alpha\otimes\beta') \oplus (\alpha'\otimes\beta')
: \alpha'<\alpha,\text{~et~}\beta'<\beta\Big\}
\tag{*}
\]

\begin{corrige}
Il n'y a aucune interaction entre les jetons dans le jeu.  En
particulier, jouer à une somme de nim de deux positions du jeu
considéré dans cet exercice revient au même que de jouer à la position
ayant la réunion de ces deux ensembles de jetons.  Si on note
$u_{\alpha,\beta}$ la position du jeu ayant un unique jeton sur la
case $(\alpha,\beta)$, on déduit de \ref{summary-of-grundy-theory} que
la valeur de Grundy d'un état quelconque du jeu vaut
$\bigoplus_{i=1}^s \gr(u_{\alpha_i,\beta_i})$.  Or
$\gr(u_{\alpha,\beta})$ est le plus petit ordinal qui n'est pas de la
forme $\gr(z)$ pour $z$ une position voisine sortante
de $u_{\alpha,\beta}$, et d'après ce qu'on vient de dire, on obtient
exactement $\gr(u_{\alpha,\beta}) = \mex\big\{\gr(u_{\alpha',\beta})
\oplus \gr(u_{\alpha,\beta'}) \oplus \gr(u_{\alpha',\beta'}) :
\alpha'<\alpha, \beta'<\beta\big\}$, c'est-à-dire bien
$\gr(u_{\alpha,\beta}) = \alpha\otimes\beta$ (même définition
inductive), et on a montré ce qui était demandé.
\end{corrige}

\smallbreak

(4) Calculer la valeur de $\alpha\otimes\beta$ pour $0\leq\alpha\leq
5$ et $0\leq\beta\leq 5$.  Pour accélérer les calculs ou bien pour les
confirmer, on pourra utiliser les résultats de
l'exercice \ref{inductions-on-nim-product} (il n'est pas nécessaire
d'avoir traité l'exercice en question).  On ne demande pas de
détailler les calculs, mais on recommande de les vérifier
soigneusement.

\begin{corrige}
En procédant inductivement, on trouve :
{\[
\begin{array}{c|cccccccccccccc}
\otimes&0&1&2&3&4&5\\\hline
0&0&0&0&0&0&0\\
1&0&1&2&3&4&5\\
2&0&2&3&1&8&10\\
3&0&3&1&2&12&15\\
4&0&4&8&12&6&2\\
5&0&5&10&15&2&7\\
\end{array}
\]}
(pour simplifier les calculs, on peut notamment utiliser la
commutativité, et le fait que $\alpha\otimes 3 = \alpha\otimes(2\oplus
1) = (\alpha\otimes 2) \oplus \alpha$ et de même $\alpha\otimes 5 =
\alpha\otimes(4\oplus 1) = (\alpha\otimes 4) \oplus \alpha$ ; si on
n'a pas l'habitude de calculer des sommes de nim, le mieux est sans
doute de tout écrire en binaire, quitte à reconvertir ensuite).
\end{corrige}

\smallbreak

(5) Si vous deviez jouer dans la position suivante (les lignes et
colonnes sont numérotées à partir de $1$, autrement dit la défausse
n'est pas figurée), quel coup feriez-vous ?

\begin{center}
\vskip-\baselineskip
\begin{tikzpicture}
\draw[step=.5cm,help lines] (0,-2.5) grid (2.5,0);
\begin{scope}[every node/.style={circle,fill,inner sep=1.2mm}]
\node at (0.25,-0.25) {};
\node at (0.75,-0.75) {};
\node at (1.25,-1.25) {};
\node at (1.75,-1.75) {};
\node at (2.25,-2.25) {};
\end{scope}
\node[anchor=east] at (0,-0.25) {$1$};
\node[anchor=east] at (0,-0.75) {$2$};
\node[anchor=east] at (0,-1.25) {$3$};
\node[anchor=east] at (0,-1.75) {$4$};
\node[anchor=east] at (0,-2.25) {$5$};
\node[anchor=south] at (0.25,0) {$1$};
\node[anchor=south] at (0.75,0) {$2$};
\node[anchor=south] at (1.25,0) {$3$};
\node[anchor=south] at (1.75,0) {$4$};
\node[anchor=south] at (2.25,0) {$5$};
\end{tikzpicture}
\end{center}

\begin{corrige}
La valeur de Grundy est $(1\otimes 1) \oplus (2\otimes 2) \oplus
(3\otimes 3) \oplus (4\otimes 4) \oplus (5\otimes 5) = 1 \oplus 3
\oplus 2 \oplus 6 \oplus 7 = 1$, et on veut l'annuler.  Plusieurs
coups sont possibles pour y arriver, par exemple : retirer le jeton en
$(1,1)$ (c'est-à-dire jouer $(\alpha,\beta,\alpha',\beta') =
(1,1,0,0)$) ; ou bien, remonter le jeton $(3,3)$ en $(1,3)$
(c'est-à-dire jouer $(\alpha,\beta,\alpha',\beta') = (3,3,1,0)$) ; ou
encore, remplacer le jeton $(5,5)$ par trois jetons en $(4,5)$,
$(5,4)$ et $(4,4)$ (c'est-à-dire jouer $(\alpha,\beta,\alpha',\beta')
= (5,5,4,4)$).
\end{corrige}


%
%
%

\exercice\label{inductions-on-nim-product}

On définit inductivement une opération $\alpha\otimes\beta$
(\defin{produit de nim}) de deux ordinaux $\alpha,\beta$ par
$\alpha\otimes\beta := \mex\{(\alpha'\otimes\beta) \oplus
(\alpha\otimes\beta') \oplus (\alpha'\otimes\beta') : \alpha'<\alpha,
\beta'<\beta\}$ (autrement dit, par la formule (*) de
l'exercice \ref{game-for-nim-product} ; il n'est pas nécessaire
d'avoir traité l'exercice en question, même s'il est possible de s'en
servir).  La notation $\oplus$ désigne la somme de nim.  On rappelle
par ailleurs que $\gamma = \mex S$ signifie que $\gamma\not\in S$ et
que tout ordinal $\gamma'<\gamma$ appartient à $S$.

(1) Montrer que $\otimes$ est commutative, c'est-à-dire que
$\beta\otimes\alpha = \alpha\otimes\beta$ quels que soient les
ordinaux $\alpha,\beta$.

\begin{corrige}
Par induction sur $\alpha$ et $\beta$, on prouve $\beta\otimes\alpha =
\alpha\otimes\beta$ en supposant que la même formule est vraie si l'un
de $\alpha$ ou $\beta$ est remplacé par un ordinal strictement plus
petit.  Or $\beta\otimes\alpha = \mex \{(\beta\otimes\alpha') \oplus
(\beta'\otimes\alpha) \oplus (\beta'\otimes\alpha') : \alpha'<\alpha,
\beta'<\beta\}$ (en utilisant la commutativité de $\oplus$), et par
hypothèse d'induction ceci vaut $\mex
\{(\alpha'\otimes\beta) \oplus (\alpha\otimes\beta') \oplus
(\alpha'\otimes\beta') : \alpha'<\alpha, \beta'<\beta\} =
\alpha\otimes\beta$.

Si on préfère, on peut aussi utiliser le jeu défini dans
l'exercice \ref{game-for-nim-product}, en remarquant qu'échanger les
coordonnées des cases de tous les jetons ne change rien au jeu (i.e.,
les règles sont symétriques ligne/colonne) donc ne change pas la
fonction de Grundy.
\end{corrige}

\smallbreak

(2) Montrer que $0$ est absorbant pour $\otimes$,
c'est-à-dire $\alpha\otimes 0 = 0$ pour tout ordinal $\alpha$.
Montrer que $1$ est neutre pour $\otimes$, soit $\alpha\otimes 1 =
\alpha$ pour tout ordinal $\alpha$.

\begin{corrige}
On a $\alpha \otimes 0 = \mex \varnothing = 0$ (puisqu'il n'existe pas
de $\beta'<0$).  Pour la seconde affirmation, par induction sur
$\alpha$, on prouve $\alpha \otimes 1 = \alpha$ : en effet, $\alpha
\otimes 1 = \mex \{(\alpha'\otimes 1) \oplus (\alpha\otimes 0) \oplus
(\alpha'\otimes 0): \alpha'<\alpha\}$, et en utilisant le fait que $0$
est aborbant pour $\otimes$ et neutre pour $\oplus$ et l'hypothèse
d'induction, ceci vaut $\mex \{\alpha': \alpha'<\alpha\} = \alpha$.

Si on préfère, on peut aussi utiliser le jeu défini dans
l'exercice \ref{game-for-nim-product}, en remarquant qu'un jeton sur
la ligne ou colonne $0$ est mort (défaussé), et que jouer sur la ligne
ou colonne $1$ revient à jouer au jeu de nim d'après la question (2)
de l'exercice \ref{game-for-nim-product}.
\end{corrige}

\smallbreak

(3) (a) Montrer que si $\alpha'\neq\alpha$ et $\beta'\neq\beta$, alors
$(\alpha'\otimes\beta) \oplus (\alpha\otimes\beta') \oplus
(\alpha'\otimes\beta') \neq \alpha\otimes\beta$.\spaceout (b) En
déduire que si $\alpha\otimes\beta = \alpha\otimes\beta'$ alors
$\alpha=0$ ou bien $\beta=\beta'$.

\begin{corrige}
(a) Grâce aux propriétés de la somme de nim ($\gamma \neq \gamma'$
  équivaut à $\gamma\oplus\gamma' \neq 0$), la condition qu'on veut
  montrer est équivalente à : $(\alpha\otimes\beta) \oplus
  (\alpha'\otimes\beta) \oplus (\alpha\otimes\beta') \oplus
  (\alpha'\otimes\beta') \neq 0$.  Sous cette forme, on voit qu'il y a
  symétrie entre $\alpha'$ et $\alpha$ et entre $\beta'$ et $\beta$ :
  on peut donc supposer $\alpha'<\alpha$ et $\beta'<\beta$, auquel cas
  la propriété est claire par la définition même de $\otimes$ comme
  un $\mex$.

(b) C'est le cas particulier du (a) lorsque $\alpha'=0$.
\end{corrige}

\smallbreak

(4) Montrer que $\otimes$ est distributive sur $\oplus$, c'est-à-dire
$\alpha \otimes (\beta\oplus\gamma) = (\alpha\otimes\beta) \oplus
(\alpha\otimes\gamma)$.  Pour cela, on pourra procéder par induction
et remarquer que pour montrer $\lambda = \mu$ il suffit de montrer que
(a) $\xi<\lambda$ implique $\xi\neq\mu$ et que (b) $\xi<\mu$ implique
$\xi\neq\lambda$.  (Et on rappelle que si $\xi < \mex S$ alors $\xi\in
S$.)

\begin{corrige}
Pour montrer $\lambda = \mu$ il suffit de montrer que
(a) $\xi<\lambda$ implique $\xi\neq\mu$ et que (b) $\xi<\mu$ implique
$\xi\neq\lambda$ : en effet, la contraposée de (a) est que
$\mu\geq\lambda$, et la contraposée de (b) est que $\lambda\geq\mu$.

Procédons par induction pour prouver $\alpha \otimes
(\beta\oplus\gamma) = (\alpha\otimes\beta) \oplus
(\alpha\otimes\gamma)$ : on peut supposer cette égalité connue si au
moins l'un des ordinaux $\alpha,\beta,\gamma$ est remplacé par un
strictement plus petit.

(a) Si $\xi < \alpha \otimes (\beta\oplus\gamma)$, alors on peut
écrire $\xi = (\alpha'\otimes(\beta\oplus\gamma)) \oplus
(\alpha\otimes\delta') \oplus (\alpha'\otimes\delta')$ où
$\alpha'<\alpha$ et $\delta' < \beta\oplus\gamma$.  La définition
inductive de $\oplus$ permet alors d'écrire soit $\delta' =
\beta'\oplus\gamma$ où $\beta'<\beta$, soit $\delta' =
\beta\oplus\gamma'$ où $\gamma'<\gamma$ : par symétrie, plaçons nous
sans perte de généralité dans le premier cas.  On a alors $\xi =
(\alpha'\otimes (\beta\oplus\gamma)) \oplus (\alpha\otimes
(\beta'\oplus\gamma)) \oplus (\alpha'\otimes (\beta'\oplus\gamma))$.
Par hypothèse d'induction, on peut développer les trois termes, et
quitte à simplifier les deux termes $\alpha'\otimes\gamma$ qui
apparaissent, on obtient $\xi = (\alpha'\otimes\beta) \oplus
(\alpha\otimes\beta') \oplus (\alpha'\otimes\beta') \oplus
(\alpha\otimes\gamma)$.  Puisque la somme $(\alpha'\otimes\beta)
\oplus (\alpha\otimes\beta') \oplus (\alpha'\otimes\beta')$ des trois
premiers termes est différente de $\alpha\otimes\beta$
(d'après (3)(a)), on en déduit que $\xi \neq (\alpha\otimes\beta)
\oplus (\alpha\otimes\gamma)$, ce qu'on voulait.

(b) Maintenant, si $\xi < (\alpha\otimes\beta) \oplus
(\alpha\otimes\gamma)$, on a soit $\xi = \delta' \oplus
(\alpha\otimes\gamma)$ avec $\delta' < \alpha\otimes\beta$, soit $\xi
= (\alpha\otimes\beta) \oplus \varepsilon'$ avec $\varepsilon' <
\alpha\otimes\gamma$ : par symétrie, plaçons nous sans perte de
généralité dans le premier cas.  On peut alors écrire $\delta' =
(\alpha'\otimes\beta) \oplus (\alpha\otimes\beta') \oplus
(\alpha'\otimes\beta')$ avec $\alpha'<\alpha$ et $\beta'<\beta$, donc
on a : $\xi = (\alpha'\otimes\beta) \oplus (\alpha\otimes\beta')
\oplus (\alpha'\otimes\beta') \oplus (\alpha\otimes\gamma)$.  Quitte à
faire apparaître deux termes $\alpha'\otimes\gamma$ qui s'annulent,
l'hypothèse d'induction permet de réécrire $\xi = (\alpha'\otimes
(\beta\oplus\gamma)) \oplus (\alpha\otimes (\beta'\oplus\gamma))
\oplus (\alpha'\otimes (\beta'\oplus\gamma))$.  Or $\beta'\oplus\gamma
\neq \beta\oplus\gamma$ : en utilisant (3)(a), on a $\xi \neq \alpha
\otimes (\beta\oplus\gamma)$, ce qu'on voulait.
\end{corrige}

\smallbreak

On \emph{admet} que $\otimes$ est associative, c'est-à-dire
$(\alpha\otimes\beta) \otimes \gamma = \alpha \otimes
(\beta\otimes\gamma)$ (ce n'est pas très difficile à prouver).

(5) On va enfin montrer que pour tout $\alpha > 0$ il existe un
$\alpha^*$ tel que $\alpha\otimes\alpha^* = 1$, c'est-à-dire, un
\emph{inverse} pour le produit de nim.  Pour cela, on suppose par
l'absurde le contraire, et on considère $\alpha$ le \emph{plus petit}
ordinal non nul qui n'a pas d'inverse, et on va arriver à une
contradiction.  Pour cela, appelons $\delta_0 = \sup^+ \big(\{\alpha\}
\cup \{\beta^* : \beta<\alpha\} \big)$ (où $\beta^*$ désigne l'inverse
de $\beta$, qu'on a supposé exister vu que $\beta<\alpha$) le plus
petit ordinal strictement supérieur à $\alpha$ et aux inverses des
ordinaux $<\alpha$, et par récurrence sur l'entier naturel $n$,
posons $\delta_{n+1} = \sup^+ \big(\{\beta_1 \oplus \beta_2 :
\beta_1,\beta_2<\delta_n\} \cup \{\beta_1 \otimes \beta_2 :
\beta_1,\beta_2<\delta_n\}\big)$ le
plus petit ordinal strictement supérieur à la somme ou au produit de
nim de deux ordinaux strictement plus petits que $\delta_n$ (on a bien
sûr $\delta_{n+1} \geq \delta_n$).  Soit enfin $\delta =
\lim_{n\to\omega} \delta_n = \sup\{\delta_n :
n\in\mathbb{N}\}$.\spaceout (a) Expliquer pourquoi si $\beta_1,\beta_2
< \delta$ alors $\beta_1\oplus\beta_2 < \delta$ et
$\beta_1\otimes\beta_2 < \delta$.\spaceout (b) Montrer que si $0 <
\alpha' < \alpha$ alors nécessairement $\alpha' \otimes \delta \geq
\delta$ (dans le cas contraire, considérer le produit de nim de $\alpha'
\otimes \delta$ par $(\alpha')^*$ et utiliser (a)).\spaceout (c) En
déduire que si $0 < \alpha' < \alpha$ et $\delta' < \delta$, alors
$(\alpha'\otimes\delta) \oplus (\alpha\otimes\delta') \oplus
(\alpha'\otimes\delta')$ est $\geq \delta$ (dans le cas contraire,
montrer que le premier terme serait $<\delta$).  En particulier, il
est $\neq 1$.\spaceout (d) Expliquer pourquoi si $\alpha' = 0$ alors
$(\alpha'\otimes\delta) \oplus (\alpha\otimes\delta') \oplus
(\alpha'\otimes\delta')$ est encore une fois $\neq 1$.\spaceout (e) En
déduire que $\alpha\otimes\delta = 1$ et conclure.

\begin{corrige}
(a) Si $\beta < \delta$, comme $\delta$ est la borne supérieure
  des $\delta_n$, il existe $n$ tel que $\beta < \delta_n$.  Ainsi, si
  $\beta_1,\beta_2 < \delta$, il existe $n$ tel que $\beta_1,\beta_2 <
  \delta_n$ (en prenant le maximum des deux $n$ obtenus), donc
  $\beta_1\oplus\beta_2 < \delta_{n+1}$ et $\beta_1\otimes\beta_2 <
  \delta_{n+1}$ d'après la définition de $\delta_{n+1}$, ce qui donne
  bien $\beta_1\oplus\beta_2 < \delta$ et $\beta_1\otimes\beta_2 <
  \delta$.

(b) Si $0 < \alpha' < \alpha$ et si $\alpha' \otimes \delta < \delta$,
  alors comme $(\alpha')^* < \delta_0 \leq \delta$, on a
  $(\alpha')^*\otimes (\alpha' \otimes \delta) < \delta$ d'après (a).
  Or $(\alpha')^* \otimes (\alpha' \otimes \delta) = \delta$ par
  associativité et par le fait que $(\alpha')^*$ est l'inverse
  de $\alpha'$ : contradiction.

(c) Si $0 < \alpha' < \alpha$ et $\delta' < \delta$, montrons que
  $\gamma := (\alpha'\otimes\delta) \oplus (\alpha\otimes\delta')
  \oplus (\alpha'\otimes\delta')$ est $\geq\delta$.  Dans le cas
  contraire, $\alpha'\otimes\delta$ serait la somme de nim
  de $\gamma$, qui est $<\delta$ par hypothèse, de
  $\alpha\otimes\delta'$ qui est $<\delta$ par (a), et de
  $\alpha'\otimes\delta'$ qui l'est aussi ; de nouveau par (a), on
  aurait $\alpha'\otimes\delta < \delta$, ce qui contredit (b).

(d) Si $\alpha'=0$ alors $(\alpha'\otimes\delta) \oplus
  (\alpha\otimes\delta') \oplus (\alpha'\otimes\delta') =
  \alpha\otimes\delta' \neq 1$ puisqu'on a supposé que $\alpha$
  n'avait pas d'inverse.

(e) Par définition, $\alpha\otimes\delta$ est le plus petit ordinal
  qui n'est pas de la forme $(\alpha'\otimes\delta) \oplus
  (\alpha\otimes\delta') \oplus (\alpha'\otimes\delta')$ pour
  $\alpha'<\alpha$ et $\delta'<\delta$.  Or
  on a montré en (c) et (d) que cette expression n'est jamais $1$, et
  en revanche elle peut être $0$ (pour $\alpha'=\delta'=0$).  Le plus
  petit ordinal qui n'est pas de cette forme est donc $1$ : mais on a
  alors prouvé $\alpha\otimes\delta = 1$, ce qui contredit l'hypothèse
  que $\alpha$ n'ait pas d'inverse.
\end{corrige}

\smallbreak

{\footnotesize
\emph{Remarque :} On a donc vu que les ordinaux, pour les opérations
$\oplus$ et $\otimes$, i.e., les « nimbres », forment donc un
\emph{corps} commutatif, corps dit de « caractéristique $2$ »
car $1\oplus 1 = 0$.  Un raisonnement assez semblable à celui fait
en (5) permettrait de montrer, en outre, que ce corps est
« algébriquement clos », c'est-à-dire que tout polynôme non constant y
a une racine.  Les entiers naturels strictement inférieurs
à $2^{2^r}$, quant à eux, forment le corps fini ayant ce nombre
d'éléments.)
\par}

%
%
%

\exercice\label{tree-hackenbush-exercise}

On s'intéresse dans cet exercice au jeu de \textbf{Hackenbush impartial
  en arbre}\index{Hackenbush}, défini comme suit.  L'état du jeu est représenté par un
arbre (fini, enraciné\footnote{C'est-à-dire que la racine fait partie
  de la donnée de l'arbre, ce qui est la convention la plus
  courante.}).  Deux joueurs alternent et chacun à son tour choisit
une arête de l'arbre et l'efface, ce qui fait automatiquement
disparaître du même coup tout le sous-arbre qui descendait de cette
arête (voir figure).  Le jeu se termine lorsque plus aucun coup n'est
possible (c'est-à-dire que l'arbre est réduit à sa seule racine),
auquel cas, selon la convention habituelle, le joueur qui ne peut plus
jouer a perdu.

\begin{center}
\begin{tikzpicture}[baseline=0]
\fill [gray!50!white] (-1.5,0) rectangle (1.5,-0.2);
\begin{scope}[every node/.style={circle,fill,inner sep=0.5mm}]
\node (P0) at (0,0) {};
\node (P1) at (0,1) {};
\node (P2) at (-1.5,2) {};
\node (P3) at (-2.0,3) {};
\node (P4) at (-1.0,3) {};
\node (P5) at (1.5,2) {};
\node (P6) at (0.75,3) {};
\node (P7) at (1.5,3) {};
\node (P8) at (2.25,3) {};
\node (P9) at (1.75,1) {};
\end{scope}
\begin{scope}[line width=1.5pt]
\draw (P0) -- (P1);
\draw (P1) -- (P2);
\draw (P1) -- (P5);
\draw (P2) -- (P3);
\draw (P2) -- (P4);
\draw (P5) -- (P6);
\draw (P5) -- (P7);
\draw (P5) -- (P8);
\draw (P0) -- (P9);
\end{scope}
\begin{scope}[line width=3pt,red]
\draw ($0.5*(P1) + 0.5*(P5) + (-0.2,-0.2)$) -- ($0.5*(P1) + 0.5*(P5) + (0.2,0.2)$);
\draw ($0.5*(P1) + 0.5*(P5) + (-0.2,0.2)$) -- ($0.5*(P1) + 0.5*(P5) + (0.2,-0.2)$);
\end{scope}
\end{tikzpicture}
devient
\begin{tikzpicture}[baseline=0]
\fill [gray!50!white] (-1.5,0) rectangle (1.5,-0.2);
\begin{scope}[every node/.style={circle,fill,inner sep=0.5mm}]
\node (P0) at (0,0) {};
\node (P1) at (0,1) {};
\node (P2) at (-1.5,2) {};
\node (P3) at (-2.0,3) {};
\node (P4) at (-1.0,3) {};
\node (P9) at (1.75,1) {};
\end{scope}
\begin{scope}[line width=1.5pt]
\draw (P0) -- (P1);
\draw (P1) -- (P2);
\draw (P2) -- (P3);
\draw (P2) -- (P4);
\draw (P0) -- (P9);
\end{scope}
\end{tikzpicture}
\end{center}

(1) Expliquer pourquoi une position de ce jeu peut être considérée
comme une somme de nim de différents jeux du même type.  Plus
exactement, soit $T$ un arbre de racine $x$, soient $y_1,\ldots,y_r$
les fils de $x$, soient $T_1,\ldots,T_r$ les sous-arbres ayant pour
racines $y_1,\ldots,y_r$ et soient $T'_1,\ldots,T'_r$ les arbres de
racine $x$ où $T'_i$ est formé de $x$ et de $T_i$ (avec une arête
entre $x$ et $y_i$) : expliquer pourquoi la position représentée par
l'arbre $T$ est la somme de nim de celles représentées par
$T'_1,\ldots,T'_r$.  Qu'en déduit-on sur la valeur de Grundy de la
position $T$ ?

\begin{corrige}
Il s'agit simplement d'observer que les différentes branches de
l'arbre n'interagissent pas du tout.  Jouer à la somme de nim des jeux
de Hackenbush représentés par les arbres $T'_1,\ldots,T'_r$ revient,
si on veut, à jouer à Hackenbush sur la réunion disjointe de ces
arbres, et comme la racine n'intervient pas, cela ne change rien au
jeu de réunir leurs racines en une seule, ce qui donne l'arbre $T$.

On en déduit que $\gr(T) = \gr(T'_1) \oplus \cdots \gr(T'_r)$ où
$\gr(T)$ désigne, par abus de notation, la valeur de Grundy de la
position de Hackenbush représentée par l'arbre $T$, et où $\oplus$
désigne la somme de nim des entiers naturels.
\end{corrige}

\smallbreak

\centerline{* * *}

Indépendamment de ce qui précède, on va considérer une nouvelle
opération sur les jeux : si $G$ est un jeu combinatoire impartial, vu
comme un graphe orienté (bien-fondé), on définit un jeu noté $*{:}G$
défini en ajoutant une unique position $0$ à $G$ comme on va
l'expliquer.  Pour chaque position $z$ de $G$ il y a une position
notée $*{:}z$ de $*{:}G$, et il y a une unique autre position,
notée $0$, dans $*{:}G$ ; pour chaque arête $z \to z'$ de $G$, il y a
une arête $*{:}z\, \to \, *{:}z'$ dans $*{:}G$, et il y a de plus une
arête $*{:}z\, \to 0$ dans $*{:}G$ pour chaque $z$ (en revanche, $0$
est un puits, c'est-à-dire qu'aucune arête n'en part) ; la position
initiale de $*{:}G$ est $*{:}z_0$ où $z_0$ est celle de $G$.  De façon
plus informelle, pour jouer au jeu $*{:}G$, chaque joueur peut soit
faire un coup normal ($*{:}z\, \to \, *{:}z'$) de $G$, soit appliquer
un coup « destruction totale » $*{:}z\, \to 0$ qui fait terminer
immédiatement le jeu (et celui qui l'applique a gagné\footnote{Ce jeu
  considéré tout seul n'est donc pas très amusant puisqu'on a toujours
  la possibilité de gagner instantanément.}).

\smallbreak

(2) Montrer par induction bien-fondée que si $G$ est un jeu
combinatoire impartial (bien-fondé) de valeur de Grundy $\alpha$,
alors $*{:}G$ a pour valeur de Grundy $1+\alpha$.

\begin{corrige}
Observons tout d'abord que la valeur de Grundy de la position $0$
de $*{:}G$ vaut $0$ puisque c'est un puits.  Montrons par induction
bien-fondée sur les sommets de $*{:}G$ que la valeur de Grundy de la
position $*{:}z$ vaut $1+\gr(z)$ où $\gr(z)$ désigne la valeur de
Grundy de la position $z$ dans $G$.  Les voisins sortants de $*{:}z$
sont $0$ et les $*{:}z'$ pour $z'$ voisin sortant de $z$ (dans $G$) ;
la définition de la valeur de Grundy assure donc que $\gr(*{:}z) =
\mex(\{0\} \cup \{\gr(*{:}z') : z'\in\outnb(z)\})$, c'est-à-dire que
la valeur de Grundy de $*{:}z$ est le plus petit ordinal qui n'est
ni $0$ ni un $\gr(*{:}z')$ pour $z'$ voisin sortant de $z$ ; mais par
hypothèse d'induction, on peut remplacer $\gr(*{:}z')$ par
$1+\gr(z')$, ce qui donne $\gr(*{:}z) = \mex(\{0\} \cup \{1+\gr(z') :
z'\in\outnb(z)\})$.  Montrons que ce $\mex$ vaut $1+\gr(z)$ : pour
cela, il suffit d'observer que (A) $1+\gr(z)$ ne vaut ni $0$ ni
$1+\gr(z')$, et (B) tout ordinal $<1+\gr(z)$ vaut soit $0$ soit
$1+\gr(z')$.  Or le (A) est clair car $1+\gr(z) > 0$ et que $\gr(z')
\neq \gr(z)$ assure $1+\gr(z') \neq 1+\gr(z)$, et le (B) est clair car
si $\beta < 1+\gr(z)$, alors soit $\beta=0$ soit $\beta = 1+\beta'$ où
$\beta' < \gr(z)$, auquel cas la définition de $\gr$ assure qu'il
existe $z'$ voisin sortant de $z$ tel que $\beta' = \gr(z')$.
\end{corrige}

\smallbreak

(3) On revient au jeu de Hackenbush impartial en arbre.  Soit $T$ un
arbre de racine $y$ et $T'$ l'arbre obtenu en ajoutant une nouvelle
racine $x$ à $T$, c'est-à-dire que les sommets de $T'$ sont ceux de
$T$ plus $x$, qui en est la racine, avec une arête entre $x$ et $y$.
Expliquer pourquoi le jeu de Hackenbush représenté par $T'$ s'obtient
par la construction « $*{:}$ » considérée en (2) à partir de celui
représenté par $T$.  Qu'en déduit-on sur la valeur de Grundy de la
position $T'$ par rapport à celle de $T$ ?

\begin{corrige}
Un coup dans le jeu de Hackenbush représenté par l'arbre $T'$ peut
consister soit à couper l'arbre à sa racine, c'est-à-dire couper
l'arête reliant $x$ et $y$, ce qui met fin au jeu immédiatement, soit
à jouer dans $T$ (i.e., couper une arête de celui-ci) ; c'est
précisément la définition qu'on a donnée de la
construction « $*{:}$ ».

On en déduit que $\gr(T') = 1 + \gr(T)$ (comme par ailleurs on a ici
affaire à des entiers naturels, le $+$ est commutatif, donc dans cette
question on peut aussi l'écrire $\gr(T') = \gr(T) + 1$).
\end{corrige}

\smallbreak

(4) Déduire des questions précédentes une méthode pour calculer la
valeur de Grundy d'une position quelconque au Hackenbush impartial en
arbre.

\begin{corrige}
Des questions précédentes, on déduit que la valeur de Grundy d'un
arbre $T$ au Hackenbush impartial se calcule comme la somme de nim des
$\gr(T'_i) = 1+\gr(T_i)$ où les $T_i$ sont les sous-arbres partant des
fils de la racine de $T$.

On peut donc calculer la valeur de Grundy d'un arbre en calculant
celle de ses sous-arbres enracinés aux différents nœuds, des feuilles
vers la racine : dès qu'on a calculé la valeur de Grundy de tous les
sous-arbres enracinés aux fils $y_1,\ldots,y_r$ d'un nœud $x$, on en
déduit celle du sous-arbre enraciné en leur père $x$ comme la somme de
nim des valeurs en question incrémentées de $1$.
\end{corrige}

\smallbreak

(5) Quelle est la valeur de Grundy de la position représentée
ci-dessous ?  (Il s'agit de la position utilisée en exemple plus
haut.)  Quel coup préconiseriez-vous dans cette situation ?

\begin{center}
\begin{tikzpicture}[baseline=0]
\fill [gray!50!white] (-1.5,0) rectangle (1.5,-0.2);
\begin{scope}[every node/.style={circle,fill,inner sep=0.5mm}]
\node (P0) at (0,0) {};
\node (P1) at (0,1) {};
\node (P2) at (-1.5,2) {};
\node (P3) at (-2.0,3) {};
\node (P4) at (-1.0,3) {};
\node (P5) at (1.5,2) {};
\node (P6) at (0.75,3) {};
\node (P7) at (1.5,3) {};
\node (P8) at (2.25,3) {};
\node (P9) at (1.75,1) {};
\end{scope}
\begin{scope}[line width=1.5pt]
\draw (P0) -- (P1);
\draw (P1) -- (P2);
\draw (P1) -- (P5);
\draw (P2) -- (P3);
\draw (P2) -- (P4);
\draw (P5) -- (P6);
\draw (P5) -- (P7);
\draw (P5) -- (P8);
\draw (P0) -- (P9);
\end{scope}
\end{tikzpicture}
\end{center}

\begin{corrige}
On trouve les valeurs de Grundy suivantes en notant à côté de chaque
nœud la valeur du sous-arbre enraciné en ce nœud :
\begin{center}
\begin{tikzpicture}[baseline=0]
\fill [gray!50!white] (-1.5,0) rectangle (1.5,-0.2);
\begin{scope}[every node/.style={circle,fill,inner sep=0.5mm}]
\node (P0) at (0,0) {};
\node (P1) at (0,1) {};
\node (P2) at (-1.5,2) {};
\node (P3) at (-2.0,3) {};
\node (P4) at (-1.0,3) {};
\node (P5) at (1.5,2) {};
\node (P6) at (0.75,3) {};
\node (P7) at (1.5,3) {};
\node (P8) at (2.25,3) {};
\node (P9) at (1.75,1) {};
\end{scope}
\begin{scope}[line width=1.5pt]
\draw (P0) -- (P1);
\draw (P1) -- (P2);
\draw (P1) -- (P5);
\draw (P2) -- (P3);
\draw (P2) -- (P4);
\draw (P5) -- (P6);
\draw (P5) -- (P7);
\draw (P5) -- (P8);
\draw (P0) -- (P9);
\end{scope}
\node[anchor=east] at (P2) {$0$};
\node[anchor=west] at (P5) {$1$};
\node[anchor=east] at (P1) {$3$};
\node[anchor=east] at (P0) {$5$};
\end{tikzpicture}
\end{center}

La valeur de Grundy recherchée est donc $5$.  En étudiant les
différentes possibilités, on trouve qu'un coup gagnant possible
consiste à retirer n'importe laquelle des arêtes les plus en haut sur
le dessin (dans tous les cas, le sommet étiqueté $3$ sur la figure
ci-dessus passe à $0$ et la racine de même).
\end{corrige}

\smallbreak

(La question qui suit est indépendante des questions précédentes et
concerne les jeux combinatoires \emph{partisans}.)

(6) On remarque que la construction $*{:}G$ définie avant la
question (2) peut se définir de façon identique lorsque $G$ est un jeu
partisan, en donnant à une arête $*{:}z\, \to \, *{:}z'$ la même
couleur que $z\to z'$, et à une arête $*{:}z\, \to 0$ la couleur verte
(ce qui signifie : à la fois bleue et rouge).  En décrivant une
stratégie, montrer que si $G \geq H$ on a aussi $*{:}G \geq *{:}H$, et
en déduire que si $G\doteq H$ alors $*{:}G \doteq *{:}H$ (où $\doteq$
désigne l'égalité au sens de Conway des jeux partisans).

\begin{corrige}
Tout d'abord, observons que $-(*{:}G) = *{:}(-G)$ puisque la
construction « $*{:}$ » est symétrique entre bleu et rouge.

La condition $G\geq H$ signifie que le joueur bleu (Blaise) possède
une stratégie gagnante au jeu $G-H = G+(-H)$ s'il joue en second.
Montrons qu'il en possède encore une à $(*{:}G) - (*{:}H) = (*{:}G) +
(*{:}(-H))$ en considérant comment il répond à un coup de son
adversaire (Roxane).  Si Roxane joue un coup de « destruction totale »
sur l'une des composantes $(*{:}G)$ ou $(*{:}(-H))$, Blaise réplique
sur l'autre et gagne.  Si Roxane joue un coup dans une des deux
composantes $G$ ou $-H$, Blaise répond selon la stratégie qu'il est
supposé posséder.  Dans tous les cas, Blaise peut répondre à tout coup
de Roxane, donc il gagne.  Ceci montre $*{:}G \geq *{:}H$.

Comme $G\doteq H$ signifie $G\geq H$ et $G\leq H$, on a bien $*{:}G
\geq *{:}H$ et $*{:}G \leq *{:}H$ d'après ce qu'on vient d evoir,
c'est-à-dire $*{:}G \doteq *{:}H$.  (La valeur d'un jeu partisan
$*{:}G$ est donc déterminée par la valeur de $G$.)
\end{corrige}



\setbox0=\vbox\bgroup
\subsection{Notions sur les combinatoires partisans à information parfaite}
\egroup



%
%
%

\printindex



%
%
%
\end{document}
